\documentclass[11pt]{amsart}
\usepackage[margin=1in]{geometry}
\usepackage[T1]{fontenc}
\usepackage{amsmath,amssymb,amsthm,mathtools}
\usepackage{enumitem}
\usepackage{needspace}
\usepackage{longtable}
\usepackage{microtype,mathrsfs}
\usepackage{xcolor}
\usepackage{mdframed}
\newmdenv[backgroundcolor=black!5,linecolor=black!65,linewidth=0.4pt,roundcorner=3pt,skipabove=8pt,skipbelow=8pt,innerleftmargin=6pt,innerrightmargin=6pt,innertopmargin=3pt,innerbottommargin=3pt]{mainresult}
\definecolor{citationblue}{RGB}{0,0,255}
\usepackage[colorlinks=true,linktoc=page,linkcolor=citationblue,citecolor=citationblue,urlcolor=black]{hyperref}
\raggedbottom
\numberwithin{equation}{section}
% Use the AMS address metadata in a title-page footnote, as requested by
% the journal, and suppress its default repetition at the end.
\makeatletter
\newcommand{\titlepageaddresses}{\@setaddresses\let\@addpunct\@gobble}
\newcommand{\cleartitlepageaddresses}{\global\let\addresses\@empty}
\makeatother
\pdfinfoomitdate=1
\pdftrailerid{}
\pdfsuppressptexinfo=15
\newtheorem{theorem}{Theorem}[section]
\newtheorem{proposition}[theorem]{Proposition}
\newtheorem{lemma}[theorem]{Lemma}
\newtheorem{corollary}[theorem]{Corollary}
\theoremstyle{definition}
\newtheorem{definition}[theorem]{Definition}
\theoremstyle{remark}
\newtheorem{remark}[theorem]{Remark}
\setlength{\emergencystretch}{1em}
\title[Uniform large-genus volume asymptotics]{Uniform large-genus asymptotics\\ of Weil--Petersson volumes}
\author{Zhiyi Li}
\address[Zhiyi Li]{Department of Modern Physics, University of Science and Technology of China, Hefei, Anhui 230026, China}
\address[Zhiyi Li]{CAS Key Laboratory of Theoretical Physics, Institute of Theoretical Physics, Chinese Academy of Sciences, Beijing 100190, China}
\author{Qihang Wang}
\thanks{Corresponding author: Qihang Wang. ORCID: \href{https://orcid.org/0009-0003-2504-8756}{0009-0003-2504-8756}.}
\address[Qihang Wang]{School of Mathematical Sciences, Peking University, Beijing 100871, China}
\email{\href{mailto:2501110049@stu.pku.edu.cn}{2501110049@stu.pku.edu.cn}}
\thanks{\titlepageaddresses}
\date{}
\subjclass[2020]{Primary 32G15; Secondary 14H10, 14H81, 41A60, 58J50, 60F05}
\keywords{Weil--Petersson volume, moduli space of curves, intersection number, large genus, local limit theorem}
\hypersetup{pdftitle={Uniform large-genus asymptotics of Weil-Petersson volumes},pdfauthor={Zhiyi Li and Qihang Wang}}
\begin{document}
\begin{abstract}
Let $V_{g,n}$ be the Weil--Petersson volume of the moduli space of hyperbolic surfaces of genus $g$ with $n$ cusps. The fixed-$n$ large-genus expansion of Mirzakhani and Zograf has correction terms polynomial in $n$ of degree $2i$ at order $g^{-i}$, and ceases to be uniform when $n$ is of order $\sqrt g$. We prove an asymptotic expansion in powers of $1/g$ that holds uniformly in all $n\ge0$, with coefficients analytic in a density parameter that ranges over a compact interval, including the limits $n/g\to0$ and $n/g\to\infty$. Its leading term is explicit, contains no undetermined constant and has relative error $O(1/g)$; truncation at an order proportional to $\sqrt g$ gives relative error $O(\mathrm{e}^{-a\sqrt g})$ for some $a>0$. For fixed $n$ this confirms Zograf's conjecture that the constant in the Mirzakhani--Zograf asymptotics is $1/\sqrt\pi$. We also determine, at every cusp density, the limits of the volume ratios $V_{g,n}/V_{g-1,n+2}$ and $V_{g,n+1}/(2gV_{g,n})$ and, in all degrees, the asymptotics of the ratios of adjacent coefficients of the one-boundary volume polynomial; this answers two questions of Mirzakhani. The proof deforms mixed $\kappa_1$--$\psi$ intersection numbers into pure $\psi$ intersection numbers along a family whose recursion has nonnegative coefficients, fixes the overall constant through Painlev\'e~I without assuming resurgence, and recovers individual volumes by a local limit theorem in the number of cusps. As an application, the spectral gap of a Weil--Petersson random surface with $n\gg\sqrt g$ cusps tends to zero in probability; this removes the restriction $n/g\to0$ from a theorem of Shen and Wu.
\end{abstract}
\maketitle
\cleartitlepageaddresses
\clearpage
\setcounter{tocdepth}{1}
\tableofcontents
\clearpage
\section{Introduction}

Let $V_{g,n}(L_1,\ldots,L_n)$ denote the Weil--Petersson volume of the moduli space of genus-$g$ hyperbolic surfaces with $n$ labeled boundary components of lengths $L_i$. We use the orbifold convention and put
\[
 V_{g,n}=V_{g,n}(0,\ldots,0),\qquad
 V_{g,n}(\mathbf L)=\frac{1}{(3g-3+n)!}
 \int_{\mathcal M_{g,n}(\mathbf L)}\omega_{\mathrm{WP}}^{3g-3+n}.
\]
We assume $2g-2+n>0$. A boundary component of length zero is a cusp. These volumes are computed by Mirzakhani's recursion~\cite{MirInv}, are generating functions for intersection numbers on $\overline{\mathcal M}_{g,n}$~\cite{MirJAMS}, and control the geometry of random hyperbolic surfaces of large genus~\cite{Mir}; see~\cite{Wright} for a survey, and~\cite{ChenMullane} for a survey comparing Weil--Petersson and Masur--Veech volumes.

Precise asymptotics are known in two limits. For fixed $n$ as $g\to\infty$, Mirzakhani obtained growth estimates and the limits of the basic volume ratios~\cite{Mir}, and Mirzakhani and Zograf obtained a complete asymptotic expansion~\cite{MZ}. For fixed $g$ as $n\to\infty$, the form of the asymptotics was predicted by Itzykson (see \cite[Remark~3]{Zog98}). In genus zero, where Zograf had given a recursion for the volumes~\cite{Zog93}, Kaufmann, Manin and Zagier proved this prediction and expressed the exponential growth rate through the first zero of $J_0$~\cite{KMZ}. Zograf then expressed the generating function of the volumes in every genus through the Kontsevich--Witten free energy at shifted times~\cite{Zog98}, and Manin and Zograf proved the asymptotics for every fixed genus~\cite[Theorem~6.1]{ManZ}.

These two limiting descriptions do not by themselves give a uniform approximation when both parameters grow. By \cite[Theorem~1.2]{MZ}, the coefficient of $g^{-i}$ in the fixed-$n$ expansion is a polynomial in $n$ of degree $2i$ with leading term $(-1)^in^{2i}/(i!\,(2\pi^2)^i)$. The expansion is thus effectively a series in $n^2/g$, whose leading terms sum formally to $\exp(-n^2/(2\pi^2g))$. Once $n$ is of order $\sqrt g$, all its terms are of comparable size, so no fixed truncation remains a uniform approximation.

\textbf{Main results.} We prove a single asymptotic expansion of $V_{g,n}$ as $g\to\infty$ that is uniform in $n\ge0$. Its leading term is explicit and contains no undetermined constant, and its relative remainder is uniform from bounded cusp counts to cusp counts growing arbitrarily fast with the genus. The main consequences are the following.
\begin{itemize}[leftmargin=*]
\item \emph{Absolute volumes} (Theorem~\ref{thm:intro-absolute} and Corollary~\ref{cor:constant}). The relative error is $O(1/g)$ uniformly in $n$, and $O(\mathrm e^{-a\sqrt g})$ after truncation at an order proportional to $\sqrt g$. For fixed $n$ this identifies the constant in the Mirzakhani--Zograf asymptotic formula~\cite[Theorem~1.2]{MZ} as $1/\sqrt\pi$, as conjectured by Zograf~\cite{Zog}.
\item \emph{Volume ratios} (Theorem~\ref{thm:volume}). The ratios $V_{g,n}/V_{g-1,n+2}$, $V_{g,n+1}/(2gV_{g,n})$ and $V_{g,n}^2/(V_{g,n-1}V_{g,n+1})$ have explicit limits at every cusp density, with rate $O(1/g)$. The first limit answers Mirzakhani's question on the asymptotics of $V_{g,n}/V_{g-1,n+2}$ when $n$ grows with $g$~\cite[Section~1.4]{Mir}. The third settles the remaining case of a question of Shen and Wu~\cite[Question~3]{SW}, who asked whether $V_{g,n}^2/(V_{g,n-1}V_{g,n+1})\to1$ as $g+n\to\infty$.
\item \emph{Coefficients and the diagonal chain} (Theorems~\ref{thm:coefficient} and~\ref{thm:chain}). We determine the ratios of adjacent coefficients of the one-boundary volume polynomial over all degrees, including degrees that grow with the genus; this answers Mirzakhani's second question~\cite[Section~1.4]{Mir}. We also determine, to logarithmic precision and including bounded genus, the whole chain $V_{G,0},V_{G-1,2},\ldots,V_{0,2G}$, about whose behavior she remarked in the same section that little was known.
\item \emph{Random surfaces with many cusps} (Subsection~\ref{subsec:intro-geometry}). Combined with the vanishing-gap argument of Shen and Wu~\cite{SW}, the volume ratios show that the spectral gap tends to zero in probability whenever $n/\sqrt g\to\infty$; the new range $n\asymp g$ completes the negative answer to their question whether such surfaces have a uniform positive spectral gap~\cite[Question~2]{SW}. The volume ratios also determine the statistics of short geodesics and of the systole at every density above $\sqrt g$, and the tight length spectrum for every cusp sequence. At the critical scale $n\asymp\sqrt g$ we identify the full primitive length spectrum; this last result needs only the zero-density case of the volume-ratio limits.
\end{itemize}

\textbf{Method.} The proof has three ingredients. A one-parameter deformation turns the alternating formula for mixed $\kappa_1$--$\psi$ intersection numbers into a recursion with nonnegative coefficients. The normalization is fixed at the pure-$\psi$ end of the deformation through Painlev\'e~I, with no resurgence assumption. Finally, a local limit theorem in the number of cusps, valid at both ends of the parameter range, passes from generating functions to individual volumes.

\subsection{The cusp-density parameter}\label{subsec:parameter}
All our formulas are expressed through a single real parameter. We first define it and then explain, heuristically, why it is the right parameter; the proofs come later.

\textbf{The genus-zero relation.} Let
\[
 X(y)=\sqrt y\,J_1(2\sqrt y)=y-\frac{y^2}{2}+\frac{y^3}{12}-\cdots,
 \qquad y_0=\frac{j_{0,1}^2}{4},\qquad x_c=X(y_0),
\]
where $j_{0,1}$ is the first positive zero of $J_0$. Since $X'(y)=J_0(2\sqrt y)$, the function $X$ increases from $0$ to $x_c$ on $[0,y_0]$; write $Y(x)$ for the inverse function. Consider the normalized cusp generating functions
\begin{equation}
 \Phi_g(x)=\sum_{n\ge0}\frac{V_{g,n}}{(2\pi^2)^{3g-3+n}}\,\frac{x^n}{n!},
 \label{eq:intro-Phi}
\end{equation}
with the convention that terms with $2g-2+n\le0$ are omitted. In genus zero, $\Phi_0''(x)=Y(x)$; this is the classical genus-zero computation (cf.~\cite{KMZ,ManZ}), and in our normalization it follows from the string equation (see Step~1 of the proof of Theorem~\ref{thm:pants-poisson}). We call $x$ the \emph{cusp weight}: it weights each cusp in~\eqref{eq:intro-Phi}. The cusp weight $x$ and the parameter $Y=Y(x)\in[0,y_0)$ determine each other, and we use whichever is more convenient.

\textbf{The parameter functions.} For $0\le Y<y_0$, the entire function
\[
 K_Y(z)=\int_0^1X'\!\Bigl(Y+\frac z2(1-s^2)\Bigr)\,\mathrm ds
\]
has only positive simple zeros (Section~\ref{the-positive-compact-family-of-kernels}); it is the kernel of the recursion in Section~\ref{exact-amplitudes-and-the-harmonic-model}. Write its first zero as $2\Delta(Y)$ and put $T(Y)=Y+\Delta(Y)$. The substitution $v=T-\Delta s^2$ shows that $T=T(Y)$ is the unique root in $(y_0,\pi^2/4]$ of
\begin{equation}
 \int_Y^T\frac{X'(v)}{\sqrt{T-v}}\,\mathrm{d}v=0.
 \label{eq:SP1}
\end{equation}
Finally define
\begin{equation}
 I(Y)=\int_Y^{T(Y)} X'(v)\sqrt{T(Y)-v}\,\mathrm{d}v,
 \qquad
 \lambda(Y)=\frac{2X(Y)\sqrt{\Delta(Y)}}{I(Y)}.
 \label{eq:SP2}
\end{equation}
These definitions will be motivated in the next paragraph. In Section~\ref{the-stationary-parameter} we show that $I>0$ and that $\lambda$ is a strictly increasing bijection from $[0,y_0)$ onto $[0,\infty)$. At $Y=0$ one has $\Delta=T=\pi^2/4$, $I=\pi/4$ and $\lambda=0$; as $Y\uparrow y_0$, $\Delta\to0$ and $\lambda\to\infty$.

\textbf{Why $\lambda$ is the cusp density.} Let $h=2g-1$ and $x=X(Y)$. We shall prove (Proposition~\ref{prop:prefactor}) that, uniformly for $0\le Y<y_0$,
\begin{equation}
 \Phi_g'(x)=\frac{\Gamma(h)\,\varpi(Y)}{\pi\,4^g\sqrt{\pi h}}\,
 \bigl(\sqrt2\,I(Y)\bigr)^{-h}\Bigl(1+O\Bigl(\frac{\log g}{g}\Bigr)\Bigr),
 \label{eq:intro-heuristic}
\end{equation}
where $\varpi$ is the explicit positive function in~\eqref{eq:intro-varpi} below. Thus the cusp weight changes the exponential growth only through the \emph{action} $I(Y)$.

Now regard the terms of the series
$\Phi_g'(x)=\sum_{k\ge0}V_{g,k+1}x^k/((2\pi^2)^{3g-2+k}k!)$ as the weights of a
probability distribution on the number $k$ of additional cusps. Its mean is
$x\,\mathrm d\log\Phi_g'/\mathrm dx$. Since $I'(Y)=-X'(Y)\sqrt{\Delta(Y)}$ and
$\mathrm dY/\mathrm dx=1/X'(Y)$, the leading term in~\eqref{eq:intro-heuristic} gives
\[
 x\frac{\mathrm d}{\mathrm dx}\log\bigl(\sqrt2\,I(Y)\bigr)^{-h}
 =h\,\frac{X(Y)\sqrt{\Delta(Y)}}{I(Y)}=\frac{h\lambda(Y)}2 .
\]
So, under the cusp weight $x$, the number of additional cusps is about
$h\lambda(Y)/2\approx g\lambda(Y)$. To recover a single volume $V_{g,n}$, one
chooses $Y$ so that the leading approximation to this mean equals $n-1$, that is,
$\lambda(Y)=2(n-1)/(2g-1)$, and estimates the probability of the value
$k=n-1$ by a local limit theorem. In particular $\lambda(Y)$ is, to leading
order, the cusp density $n/g$, and we use $Y$ and $\lambda$
interchangeably. This heuristic produces the leading term in
Theorem~\ref{thm:intro-absolute}. As $n/g$ increases from $0$ to $\infty$,
the parameter $Y$ increases from $0$ to $y_0$ and the cusp weight $x$ from
$0$ to $x_c$.

The difficulty is to make this heuristic rigorous with relative errors uniform in $n$. The proof of~\eqref{eq:intro-heuristic} uses a differential equation in $Y$ whose integration constant is fixed at the endpoint $Y=y_0$, and the passage from the generating function to its coefficients needs Fourier bounds that survive at both ends of the parameter interval. Subsection~\ref{subsec:strategy} describes these steps.

\textbf{A matrix-model interpretation.} The cusp weight is a defect weight in Jackiw--Teitelboim (JT) gravity, and $Y$ locates the edge of the corresponding spectral density. At the endpoint $Y=y_0$ the square-root edge becomes cubic. Subsection~\ref{subsec:jt-edge} gives the precise dictionary with the defect-deformed matrix model of Maxfield--Turiaci and Witten~\cite{MT,WittenDef,ET}. In this language the endpoint corresponds to the $(2,3)$ critical point, and this is the point at which our normalized recursion reduces to pure $\psi$ intersection numbers. The interpretation does not by itself control the remainders.

\subsection{Absolute volumes and Zograf's constant}
For $n\ge1$ put $h=2g-1$ and $k=n-1$, let $Y\in[0,y_0)$ be the solution of $\lambda(Y)=2k/h$, and put $x=X(Y)$. Define
\begin{equation}
 \widehat V_{g,n}=\frac{(2\pi^2)^{3g-3+n}\,\Gamma(h)}{\pi\,4^g\sqrt{\pi h}}\,
 \frac{\varpi(Y)}{(\sqrt2\,I(Y))^{h}}
 \left(\frac{k}{\mathrm e x}\right)^{k}
 \sqrt{\frac{k}{h\,x\Psi'(x)}},
 \label{eq:intro-leading}
\end{equation}
where $\Psi(X(y))=\lambda(y)/2$ for $0\le y<y_0$, and $\varpi(Y)>0$ is given by
\begin{equation}
 \varpi(Y)^2=\frac{I(Y)}{\Delta(Y)^{3/2}}
 \left(\int_0^1(1-s^2)\,\frac{X(v)}{v}\,\mathrm ds\right)^{-1},
 \qquad v=T(Y)-\Delta(Y)s^2;
 \label{eq:intro-varpi}
\end{equation}
here $X(v)/v$ is an entire function of $v$, equal to one at $v=0$. For $n=0$ we put $\widehat V_{g,0}=\Gamma(2g-2)(4\pi^2)^{2g-3}/\sqrt{\pi g}$.

The formula~\eqref{eq:intro-leading} is the heuristic of Subsection~\ref{subsec:parameter} made explicit. Its first two factors are $(2\pi^2)^{3g-3+n}$ times the approximation~\eqref{eq:intro-heuristic} of $\Phi_g'(x)$. The coefficient of $x^k$ is about $\Phi_g'(x)x^{-k}$ times the probability of the center of the distribution, and $k!$ times that probability is $(k/\mathrm e)^k\sqrt{k/(hx\Psi'(x))}$ to leading order; here $hx\Psi'(x)$ is the variance scale supplied by the leading action. We use the Poisson point mass $\mathrm e^{-k}k^k/k!$ at the center rather than its Gaussian approximation $(2\pi k)^{-1/2}$. This choice makes the same formula correct when the cusp count stays bounded, where the number of additional cusps is approximately Poisson rather than Gaussian. At $k=0$ the last two factors in~\eqref{eq:intro-leading} are interpreted as their limits, which are one; then $Y=0$, $\varpi(0)=1$ and $\widehat V_{g,1}=\Gamma(2g-1)(4\pi^2)^{2g-2}/\sqrt{\pi(g-1/2)}$. The case $n=0$ is reduced to $n=1$ by the dilaton identity~\eqref{eq:TL12} at $p=0$ (Section~\ref{sec:absolute}).

The correction coefficients $A_j(g,n)$ below, with $A_0=1$, are determined by finite recursive rules from the kernels together with one-dimensional integrals in $Y$; see \eqref{eq:CG3} and \eqref{eq:AV7}; the first is written out in \eqref{eq:CG12}. For $n\ge1$ they are analytic functions of the parameter $Y$, including both endpoints of its interval.

\begin{mainresult}
\begin{theorem}[Uniform absolute expansion]\label{thm:intro-absolute}
There are constants $C,B,c,g_0>0$ such that, for every $g\ge g_0$ and every integer $0\le K\le c\sqrt{2g-1}$,
\[
 \sup_{n\ge0}\left|\frac{V_{g,n}}{\widehat V_{g,n}}
       -\sum_{j=0}^K\frac{A_j(g,n)}{(2g-1)^j}\right|
 \le\frac{CB^{K+1}((K+1)!)^2}{(2g-1)^{K+1}},
\]
and the coefficients satisfy $|A_j(g,n)|\le CB^j(j!)^2$ for all $g\ge g_0$ and $n\ge0$. In particular,
\[
 \sup_{n\ge0}\left|\frac{V_{g,n}}{\widehat V_{g,n}}-1\right|=O(g^{-1}).
\]
Taking $K$ to be a sufficiently small fixed multiple of $\sqrt g$, rounded down to an integer, gives relative error $O(\mathrm{e}^{-a\sqrt g})$ for some $a>0$, uniformly in $n$.
\end{theorem}
\end{mainresult}

The uniformity in $n$ is the main point: Theorem~\ref{thm:intro-absolute} replaces the fixed-$n$ series in $n^2/g$ by an expansion in $1/g$ alone. The same estimate applies when $n$ is fixed, when $n$ is proportional to $g$, and when $n/g\to\infty$. In the first regime the saddle is near $x=0$; as $n/g\to\infty$ it tends to the radius of convergence $x_c$ of the fixed-genus generating functions of Manin and Zograf. In this way a single formula joins the two classical regimes. The order $g^{-1}$ is sharp for this choice of leading term (Theorem~\ref{thm:absolute}); including the first correction \eqref{eq:CG12}, or \eqref{eq:AV7} for $n=0$, gives error $O(g^{-2})$. The proof is completed in Section~\ref{sec:absolute}.

For fixed $n$, Mirzakhani and Zograf proved that $V_{g,n}\sim C(2g-3+n)!\,(4\pi^2)^{2g-3+n}/\sqrt g$ for a universal constant $C>0$~\cite[Theorem~1.2]{MZ}. On the basis of numerical data, Zograf had conjectured that $C=1/\sqrt\pi$~\cite{Zog}; see also \cite[Remark~1.3]{MZ}. The same value appears in physical predictions. Saad, Shenker and Stanford obtained the unmarked large-genus prediction from the JT matrix integral under a resurgence assumption~\cite[Section~5.6]{SSS}. Okuyama and Sakai extracted the constant $1/\sqrt\pi$, together with first corrections for the coefficients of $V_{g,1}(L)$ at each fixed degree, numerically from the KdV equation~\cite[Section~2.4, equation~(89)]{OS}. Kimura conjectured the corresponding asymptotics of $V_{g,n}(\mathbf L)$ from JT gravity~\cite[equation~(28)]{Kim}, and Eynard, Garcia-Failde, Gregori, Lewa\'nski and Schiappa derived it from the nonperturbative topological recursion, assuming that the generating series of Weil--Petersson volumes is resurgent~\cite[Section~6.1, equation~(6.7)]{EGGLS}.

Our argument uses no resurgence assumption. We fix the normalization at the endpoint $Y=y_0$ of the parameter interval, where the problem reduces to pure $\psi$ intersection numbers and only the single sequence $\langle\tau_2^{3g-3}\rangle_g$ enters. An identity of Itzykson and Zuber~\cite{IZ} expresses this sequence through the formal power-series solution of Painlev\'e~I, whose large-order asymptotics follow from the results of Joshi--Kitaev, Kapaev, and Garoufalidis, Its, Kapaev and Mari\~no~\cite{JK,Kap,GIKM}. The leading term and the first correction of this sequence also appear as a special case of the uniform asymptotics for $\psi$ intersection numbers announced by Guo, Yang and Zagier~\cite{GYZ}; we derive both from Painlev\'e~I in Subsections~\ref{sec:prefactor-endpoint} and~\ref{sec:endpoint-expansion}, and use~\cite{GYZ} only as an independent check. Comparison with the Mirzakhani--Zograf expansion gives the following corollary.
\begin{mainresult}
\begin{corollary}[Zograf's constant]\label{cor:constant}
For every fixed $n\ge0$, as $g\to\infty$,
\[
 V_{g,n}=\frac{(2g-3+n)!\,(4\pi^2)^{2g-3+n}}{\sqrt{\pi g}}
 \bigl(1+O_n(g^{-1})\bigr).
\]
\end{corollary}
\end{mainresult}
This is proved at the end of Section~\ref{sec:prefactor}. Together with the bracket estimate \eqref{eq:TL23}, it also gives the formula conjectured by Kimura~\cite[equation~(28)]{Kim}: for fixed $n\ge1$ and nonnegative lengths $L_i=o(\sqrt g)$,
\[
 V_{g,n}(\mathbf L)\sim\sqrt{\frac2\pi}\,(4\pi^2)^{2g-3+n}\,\Gamma\Bigl(2g+n-\frac52\Bigr)\prod_{i=1}^n\frac{\sinh(L_i/2)}{L_i/2}.
\]
The factors at $L_i=0$ are interpreted by continuity. A direct derivation is given in \eqref{eq:boundary-factor} below.

\textbf{Bounded geodesic boundaries.} Proposition~\ref{prop:bounded-boundaries} extends the leading comparison to a fixed number of nonzero boundary lengths in a compact set, uniformly over all cusp counts. Its explicit factor interpolates, as $n/g$ runs from zero to infinity, between the fixed-cusp factor $\sinh(L/2)/(L/2)$ of~\cite{MP} and the fixed-genus, large-cusp factor $I_0(j_{0,1}L/(2\pi))$ of~\cite{HTLargeN}. The interpolation is given in \eqref{eq:boundary-profile}; both the number and lengths of the nonzero boundaries remain bounded.

\subsection{Volume ratios}\label{subsec:intro-ratios}
For sufficiently large $g$, define
\[
 B_{g,n}=\frac{V_{g,n}}{V_{g-1,n+2}},\qquad
 C_{g,n}=\frac{V_{g,n+1}}{2gV_{g,n}},\qquad
 D_{g,n}=\frac{V_{g,n}^2}{V_{g,n-1}V_{g,n+1}}.
\]
The first two quantities are defined for $n\ge0$, and the third for $n\ge1$. For each fixed $n$, Mirzakhani proved that $B_{g,n}\to1$ and $C_{g,n}\to4\pi^2$~\cite[Theorem~3.5]{Mir}, and she asked for the asymptotics of $B_{g,n}$ when $n=n(g)\to\infty$~\cite[Section~1.4]{Mir}. An estimate of Mirzakhani and Zograf states that $|(2g-2+n)V_{g,n}/V_{g,n+1}-1/(4\pi^2)|\le cn/(2g-2+n)$ for stable $(g,n)$ with $n\ge1$~\cite[Lemma~5.1]{MZ}; together with the fixed-$n$ result, this gives $C_{g,n}\to4\pi^2$ whenever $n=o(g)$. When $n$ is comparable to or larger than $g$, the limits depend on the density $n/g$, through the functions
\begin{equation}
 \mathcal B(\lambda(Y))=\frac{\pi^2}{4\Delta(Y)},\qquad
 \mathcal C(\lambda(Y))=\frac{\lambda(Y)\pi^2}{X(Y)}.
 \label{eq:SP3}
\end{equation}
They are well defined because $\lambda$ is a bijection, and they take the continuous values $\mathcal B(0)=1$ and $\mathcal C(0)=4\pi^2$. Heuristically, $\mathcal C$ appears because $x\Phi_g''(x)/\Phi_g'(x)$ is approximately the mean number $n-1$ of additional cusps, and $\mathcal B$ because $\Phi_{g-1}'''(x)/\Phi_g'(x)$ is approximately $8\Delta(Y)$; see~\eqref{eq:TL30} and~\eqref{eq:TL31}.

\begin{mainresult}
\begin{theorem}[Uniform volume ratios]\label{thm:volume}
There are constants $C_0,g_0>0$ such that, for every $g\ge g_0$,
\[
 \sup_{n\ge0}\left|\frac{B_{g,n}}{\mathcal B(n/g)}-1\right|\le\frac{C_0}{g},
 \qquad
 \sup_{n\ge0}\left|\frac{C_{g,n}}{\mathcal C(n/g)}-1\right|\le\frac{C_0}{g},
\]
and
\[
 \sup_{n\ge1}|D_{g,n}-1|\le\frac{C_0}{g}.
\]
In particular, along every path $n/g\to\infty$ with $g\to\infty$,
\[
 B_{g,n}\sim\frac{\pi^2}{3\sqrt5\,j_{0,1}}\sqrt{n/g},\qquad
 C_{g,n}\sim\frac{\pi^2}{x_c}\frac ng.
\]
\end{theorem}
\end{mainresult}
Section~\ref{extracting-fixed-cusp-counts} proves the uniform limits directly from a local limit theorem. The inverse-genus error follows from the sharper coefficient estimates, as shown in Section~\ref{subsec:quantitative-volume-ratios}. Before this theorem, the limit of $C_{g,n}$ was known for $n=o(g)$, as described above, and the limit $B_{g,n}\to1$ for $n=o(\sqrt g)$, by the absolute estimate \cite[Theorem~1.8]{MZ}. Theorem~\ref{thm:volume} removes these restrictions and gives the rate $O(g^{-1})$ uniformly in $n$.

The statement on $D_{g,n}$ answers a question of Shen and Wu~\cite[Question~3]{SW}, which asks whether $D_{g,n}\to1$ as $g+n\to\infty$. As they observe, the answer is affirmative when $n=o(g)$, by \cite[Lemma~5.1]{MZ}, and when $g$ is fixed, by \cite[Theorem~6.1]{ManZ}; they also proved $1/2-\pi^2/20\le\liminf D_{g,n}\le\limsup D_{g,n}\le1$ for every cusp sequence~\cite[Proposition~36]{SW}. Theorem~\ref{thm:volume} supplies the remaining range, in which $g\to\infty$ and $n/g\not\to0$. Combined with the fixed-genus case, it gives $D_{g,n}\to1$ whenever $g+n\to\infty$.

\subsection{Adjacent coefficients}
Mirzakhani's second question~\cite[Section~1.4]{Mir} concerns adjacent coefficients when the degree grows with the genus. We write the one-boundary volume polynomial in the normalization of~\cite{Mir},
\[
 V_{g,1}(L)=\sum_{k=0}^{N}\frac{a_{g,k}}{(2k+1)!}L^{2k},
 \qquad N=3g-2.
\]
The normalized intersection bracket $b_{g,k}=[\tau_k]_{g,1}$ uses the expansion in $2L$, and hence $b_{g,k}=4^k a_{g,k}$ (Section~\ref{exact-amplitudes-and-the-harmonic-model}). For fixed $k$, its adjacent ratio tends to one~\cite[Theorem 3.5]{Mir}; the corresponding ratio $a_{g,k}/a_{g,k+1}$ tends to four. Thus the limit one stated in \cite[Section~1.4]{Mir} refers to the bracket normalization, which is also the one used in~\cite{LX}.

\begin{mainresult}
\begin{theorem}[All adjacent degrees]\label{thm:coefficient}
For $0\le k<N$, let $t_{g,k}\in(0,\pi)$ be determined by
\[
 \frac{2k+1}{2g-1}
 =\frac{t_{g,k}^{2}\sin t_{g,k}}
 {\sin t_{g,k}-t_{g,k}\cos t_{g,k}}.
\]
Then
\[
 \sup_{0\le k<3g-2}
 \left|\frac{a_{g,k}/a_{g,k+1}}{4\pi^2/t_{g,k}^{2}}-1\right|
 =O(g^{-1/8}).
\]
In particular, $a_{g,k}/a_{g,k+1}\to4$ for every $k=o(g)$, and, with $m=N-k\ge1$,
\[
 \frac{a_{g,N-m}}{a_{g,N-m+1}}
 \sim\frac{4\pi^2g}{5m}\qquad\text{if }m=o(g).
\]
\end{theorem}
\end{mainresult}
The density $(2k+1)/(2g-1)$ keeps the exact codimension $m$ in the saddle equation; this is why the last formula also covers bounded positive $m$. Section~\ref{uniform-one-boundary-coefficient-ratios} proves the uniform equivalent, and Appendix~\ref{sec:quantitative-coefficients} proves the stated error. On the range $k\ge\varepsilon g$, for each fixed $\varepsilon\in(0,3)$, the stronger error $O_\varepsilon(g^{-1/4})$ holds, again including bounded positive codimension.

For each fixed codimension, Liu and Xu obtained a complete expansion of the coefficients relative to the top one~\cite[Theorem~1.2]{LX}, a partial answer to Mirzakhani's question. The bounded-codimension case of Theorem~\ref{thm:coefficient} agrees with their result after the change of normalization. What is new is uniformity over all degrees, including degrees and codimensions that grow with the genus.

The trigonometric saddle equation is the stationary equation of an instanton prediction from physics; see Subsections~\ref{subsec:intro-instanton} and~\ref{subsec:instanton-model}.

\subsection{The diagonal chain}
Along the chain $V_{G,0},V_{G-1,2},\ldots,V_{0,2G}$ the Euler characteristic is constant, and genus is traded for pairs of cusps. Mirzakhani noted that the adjacent ratios at the two ends of this chain follow from the fixed-$n$ and fixed-$g$ asymptotics, but that little was known about the behavior of the whole sequence~\cite[Section~1.4]{Mir}. Let $Y(u)$ correspond to $\lambda=2u/(1-u)$ and define
\begin{equation}
 F(t)=\int_0^t\log\frac{4\Delta(Y(u))}{\pi^2}\,\mathrm{d}u
 \quad(0\le t<1),\qquad F(1)=-2\log(2x_c).
 \label{eq:WC1}
\end{equation}
The endpoint value agrees with the limit of the integral.
\begin{mainresult}
\begin{theorem}[A logarithmic rate for the full chain]\label{thm:chain}
Uniformly over stable $(g,n)$ with $2g+n\to\infty$,
\begin{equation}
 \log\frac{V_{g,n}}{(2g+n-3)!(4\pi^2)^{2g+n-3}}
 =\frac{2g+n}{2}F\!\left(\frac n{2g+n}\right)+o(2g+n).
 \label{eq:WC2}
\end{equation}
Consequently,
\begin{equation}
 \sup_{0\le j\le G}
 \left|\frac1G\log\frac{V_{G-j,2j}}{V_{G,0}}-F(j/G)\right|\to0.
 \label{eq:WC3}
\end{equation}
\end{theorem}
\end{mainresult}
The second statement determines each entry of the chain up to a factor $\exp(o(G))$, uniformly in its position and including the planar endpoint $V_{0,2G}$. Section~\ref{uniform-logarithmic-volume-rate-and-the-full-diagonal-chain} proves it by summing the logarithms of the cusp ratios $C_{g,j}$ and comparing with the fixed-genus asymptotics at the planar end. When the genus tends to infinity, Theorem~\ref{thm:intro-absolute} gives far more precise information; the point of Theorem~\ref{thm:chain} is that it also covers bounded genus.

\subsection{Short geodesics and small eigenvalues}\label{subsec:intro-geometry}
Volume ratios control the geometry of random surfaces through
Mirzakhani's integration formula~\cite{MirInv}, which expresses the expected number of
short curves of a given topological type through volumes of the surfaces
obtained by cutting along them. Let $\mathbb P^{\mathrm{WP}}_{g,n}$ be the
probability measure obtained by normalizing Weil--Petersson volume on
$\mathcal M_{g,n}$, and let $\Delta_X$ be the nonnegative Laplace--Beltrami
operator of a surface $X$. The short geometry changes at the scale
$n\asymp\sqrt g$: for $n=o(\sqrt g)$ the spectral gap is bounded below
with high probability~\cite[Theorem~3]{SW}, whereas above this scale it
closes.
\begin{mainresult}
\begin{corollary}[Vanishing spectral gap]\label{cor:spectral}
Suppose that $g\to\infty$ and that the integer cusp counts satisfy $n(g)/\sqrt g\to\infty$. Then, for every $\varepsilon>0$,
\[
 \mathbb P^{\mathrm{WP}}_{g,n(g)}
 \{X:\Delta_X\text{ has an eigenvalue in }(0,\varepsilon)\}\longrightarrow1.
\]
\end{corollary}
\end{mainresult}
This corollary is an application of Shen--Wu's vanishing-gap
argument~\cite[Theorem~1 and Proposition~35]{SW} with a stronger volume input.
Shen and Wu asked whether a generic surface in $\mathcal M_{g,n}$ has a
uniform positive spectral gap when $n$ grows significantly faster than
$\sqrt g$~\cite[Question~2]{SW}. They proved that it does not for
$\sqrt g\ll n\ll g$, noted that Zograf's deterministic
bound~\cite{ZogSpectral} gives the same conclusion for $n/g\to\infty$,
and expected that the restriction $n/g\to0$ in their theorem could be
removed, which would settle the question.
The remaining range is $n\asymp g$. As Shen and Wu point out, their
short-curve argument covers this range once the adjacent-volume limit
$D_{g,n-1}\to1$ is known, that is, once their question~\cite[Question~3]{SW} has an
affirmative answer there. Theorem~\ref{thm:volume} supplies this limit,
and Subsection~\ref{sec:spectral} applies their argument with this input.
Corollary~\ref{cor:spectral} therefore completes the negative answer
to~\cite[Question~2]{SW}.

\textbf{Short geodesics and the systole.} Set
$q_g=n(n-1)/(8g\mathcal C(n/g))\asymp n^2/(g+n)$. For small $L$, the
expected number of geodesics of length at most $L$ that cut off two cusps
is asymptotic to $q_gL^2$ (Proposition~\ref{prop:pants-concentration}), so
$q_g^{-1/2}$ is the natural length scale of the shortest curves. Under the
hypotheses of Corollary~\ref{cor:spectral}, and assuming in addition that
$n/g\to\lambda\in[0,\infty]$, Theorem~\ref{thm:pants-poisson} identifies
the Poisson process formed by all primitive closed geodesics on this
scale, together with the joint law of the systole and the number of cusps
it cuts off. With probability tending to one, the shortest geodesics
separate off a genus-zero subsurface that contains only cusps. At limiting
cusp density $\lambda$, the curves that cut off $k$ cusps have rescaled
intensity $2w_k(x_\lambda)t\,\mathrm{d}t$, where $x_\lambda=X(Y_\lambda)$
and the weights $w_k$ are given in \eqref{eq:PC1}. At zero density only
curves cutting off two cusps survive in the limit, whereas at positive
density every $k\ge2$ has positive intensity.

\textbf{Many small eigenvalues.} Counting disjoint short pairs of pants and
applying the localized min--max argument of Hide and
Thomas~\cite{HTSphere}, Corollary~\ref{cor:many-spectrum} gives a growing
number of eigenvalues tending to zero, with an explicit bound on the $j$th
eigenvalue whenever $j\to\infty$ and $j=o(q_g)$.

\textbf{The tight length spectrum.} The same estimates determine a different part of the length spectrum for every cusp sequence $n=n(g)\ge0$. A tight geodesic is shortest in its homotopy class after filling all cusps. In Theorem~\ref{thm:tight-spectrum}, the primitive tight lengths, multiplied by $\alpha(n/g)=\sqrt{2r(Y_{n/g})}/\pi$, converge to a Poisson process with intensity
\[
 \frac{\cosh t-1}{t}\,\mathrm dt.
\]
This extends the near-critical Boltzmann theorem of Budd and Lions~\cite{BL} to arbitrary prescribed cusp counts, using their tight integration formula and geometric moment argument. The scale satisfies $\alpha(0)=1$ and $\alpha(\lambda)\asymp\lambda^{-1/4}$ as $\lambda\to\infty$.

\textbf{The critical scale.} At $n/\sqrt g\to c\in(0,\infty)$, the full ordinary
primitive length spectrum has a different limit
(Theorem~\ref{thm:critical-cluster}). It is the sum of an independent
Poisson background of intensity $(\cosh t-1)\,\mathrm dt/t$ and
the complete primitive spectra of two-cusp pairs of pants. Their waist
lengths form a Poisson process of intensity
$c^2\sinh(s/2)\,\mathrm ds/(8\pi^2)$. Shen--Wu's short pants count
already exhibits this waist intensity~\cite[Lemmas~28 and~33]{SW};
the full limit retains the correlations between all lengths in each
pair of pants and is not Poisson. Besides the volume bounds
of~\cite{MZ,SW}, this critical limit uses only the zero-density
genus-ratio limit $V_{g,n}/V_{g-1,n+2}\to1$ from Theorem~\ref{thm:volume},
rather than deducing this two-sided limit from the volume bounds
of~\cite{MZ,SW}. Its additional content is the joint pants decomposition,
including the nonsimple curves.

\subsection{Strategy of the proof}\label{subsec:strategy}
The heuristic of Subsection~\ref{subsec:parameter} has two parts: the approximation~\eqref{eq:intro-heuristic} of $\Phi_g'$, uniform in the cusp weight, and the passage from $\Phi_g'$ to one of its coefficients. We describe how each part is proved.

Mirzakhani's volumes are intersection numbers of $\kappa_1$ and $\psi$ classes, and $\kappa_1$ can be expressed through pure $\psi$ intersections~\cite{KMZ,ManZ}. The resulting formula~\eqref{eq:H5} is an alternating sum over the number of additional $\psi$ insertions, and asymptotics for its terms do not control the cancellation. Our first step replaces it by a recursion with nonnegative coefficients.

\begin{enumerate}[label=(\roman*),leftmargin=*]
\item \emph{A positive deformation.}
Moving along the string flow by a time $Y$ replaces the alternating weights $-X^{(k)}(0)$ by the derivatives $-X^{(k)}(Y)$; see~\eqref{eq:H6}. At $Y=0$ the resulting \emph{shifted amplitudes} are the Weil--Petersson brackets. After a natural rescaling they become, at $Y=y_0$, pure $\psi$ intersection numbers with insertions of $\tau_2$. The Dijkgraaf--Verlinde--Verlinde equations turn into a recursion~\eqref{eq:H9} whose coefficients are the Taylor coefficients of $1/K_Y$, and these are nonnegative for every $Y$ (Section~\ref{the-positive-compact-family-of-kernels}). Comparing this recursion with an explicit solvable model, the \emph{harmonic model}~\eqref{eq:H12}, gives two-sided bounds that are uniform in $Y$ and in all degrees but lose a factor $\mathrm e^{O(\sqrt g)}$ (Proposition~\ref{prop:harmonic}). This step uses no information about volumes.

\item \emph{Ratios, a differential equation and the normalization.}
For a fixed number of distinguished markings, the dilaton equation expresses the $Y$-derivative of $\log\Phi_g^{(p)}(X(Y))$ through amplitudes with one more marking, and the KdV equation relates genus $g$ to genus $g-1$ with two more markings. Positivity makes the amplitudes decrease as a degree increases, and this upgrades the coarse bounds of step~(i) to relative errors $O(1/g)$ for both ratios (Proposition~\ref{prop:marked-ratios}). Integrating the resulting differential equation in $Y$ gives~\eqref{eq:intro-heuristic} up to a factor independent of $Y$. This factor is fixed at $Y=y_0$, where only $\langle\tau_2^{3g-3}\rangle_g$ enters and Painlev\'e~I supplies the asymptotics (Section~\ref{sec:prefactor}). Iterating the same equations gives all higher corrections, with remainders controlled as the order grows (Section~\ref{sec:higher-order}, with proofs in Appendix~\ref{app:marked}).

\item \emph{From generating functions to coefficients.}
An expansion of a generating function does not by itself approximate its coefficients. We show that $\Phi_g^{(p)}$ is $(1/X'(Y(x)))^{2g-2+p}$ times a power series with nonnegative coefficients (Lemma~\ref{lem:convolution}). This factorization gives Fourier bounds on the circle $|z|=x$, and hence a local limit theorem for the cusp count, uniformly from $x$ near $0$ to $x$ near $x_c$ (Proposition~\ref{prop:local-limit}). Near $x=0$ the cusp count is nearly Poisson. A change of variables of the type studied by Temme~\cite{Temme} then gives an expansion around the Poisson law that remains regular for bounded cusp counts (Section~\ref{sec:sharp-extraction}, with proofs in Appendix~\ref{app:extraction}). Together these estimates yield Theorem~\ref{thm:intro-absolute} (Section~\ref{sec:absolute}).
\end{enumerate}

\subsection{Relation to previous work}
We group earlier work into three threads: the geometry of moduli spaces
and curve counting, the asymptotics of volumes and intersection numbers,
and short geodesics and spectra of random hyperbolic surfaces. The last
subsection compares the present paper with contemporaneous work on
Zograf's constant.

\subsubsection{Moduli spaces, volumes and curve counting}
Weil--Petersson volumes play two roles in this paper: they are
generating functions for intersection numbers, and they compute
probabilities for random hyperbolic surfaces.
Penner developed a combinatorial method for computing these
volumes~\cite{Penner}.
Building on McShane's identity~\cite{McShane}, Mirzakhani derived
a recursion that computes $V_{g,n}(\mathbf L)$ as a polynomial in
the squared boundary lengths~\cite{MirInv}; its coefficients are
intersection numbers of tautological classes on
$\overline{\mathcal M}_{g,n}$~\cite{MirJAMS}, with the symplectic
normalization established by Wolpert~\cite{Wolpert}. Mulase and Safnuk
recast the recursion as Virasoro constraints and showed that the
generating function of mixed $\kappa_1$--$\psi$ intersections is a
one-parameter solution of the Korteweg--de Vries (KdV) hierarchy~\cite{MS};
Eynard and Orantin identified it with the loop equations of a matrix
model~\cite{EO}. The identities expressing mixed $\kappa_1$--$\psi$
intersections through pure $\psi$ intersections, which are the starting
point of Section~\ref{exact-amplitudes-and-the-harmonic-model}, are due
to Kaufmann, Manin and Zagier and to Manin and Zograf~\cite{KMZ,ManZ}.
Related exact formulas include the tight-volume recursions of Budd and
Zonneveld~\cite{BZ}, the KdV boundary-addition identities and
fixed-genus Bessel singularity studied by Johnson~\cite{Johnson}, and
the Virasoro-theoretic proof by Xiong and Zhou of the intersection-number
and cohomological-field-theory descriptions of topological recursion,
including variation of the spectral-curve
function~\cite[Theorem~1.1 and Section~5.7]{XZ}. These are exact
identities; obtaining uniform remainder estimates from them is a
separate analytic problem.

Mirzakhani's integration formula~\cite{MirInv}
expresses the average number of simple multicurves of a prescribed
type in terms of volumes of the surfaces obtained by cutting along them.
After division by $V_{g,n}$, these formulas become expectations under
Weil--Petersson probability, so relative volume estimates control
geometric probabilities. Mirzakhani used this connection to study random
surfaces at large genus~\cite{Mir}, and Mirzakhani and Petri used joint
curve counts to obtain Poisson length statistics for closed
surfaces~\cite{MP}.
A different curve-counting limit fixes the surface and lets the length
bound tend to infinity. Mirzakhani's counting theorem assigns frequencies
to simple-curve topological types in that limit~\cite{MirCount}.
Delecroix, Goujard, Zograf and Zorich related these frequencies to
intersection numbers and Masur--Veech volumes~\cite{DGZZ}; Ren
obtained the large-genus ratio of separating to nonseparating frequencies
with many cusps~\cite{Ren}. Those frequencies are topological quantities,
whereas the processes studied here count short curves on a random
surface as its topology varies.

Masur--Veech volumes are the flat counterpart of Weil--Petersson
volumes, and their large-genus theory has developed in parallel.
Eskin and Zorich conjectured the large-genus asymptotics of
Masur--Veech volumes of strata of Abelian differentials~\cite{EZ};
Aggarwal proved them for all strata~\cite{AggAbelian}, and Chen,
M\"oller, Sauvaget and Zagier expressed these volumes through
intersection numbers on the strata and evaluated their large-genus
limits~\cite{CMSZ}. Rafi and Zorich recently determined the large-$n$
asymptotics of Masur--Veech volumes of strata of quadratic differentials
with many simple poles at fixed genus~\cite{RZ}.

\subsubsection{Asymptotics of volumes and intersection numbers}
The fixed-genus regime was recalled at the beginning of the
introduction. Here we discuss large-genus estimates at fixed cusp
count, coefficients of growing degree, pure $\psi$ intersection numbers
and predictions from physics.

\textbf{Fixed cusp counts.}
Grushevsky~\cite{Gru} and, for closed surfaces, Schumacher and
Trapani~\cite{ST} showed that for fixed $n$ the volumes grow like
$(2g)!$ up to factors $c^g$; see the account in \cite[Section~1.4]{Mir}.
After Mirzakhani's estimates~\cite{Mir}, Mirzakhani and Zograf obtained
the expansion to every fixed order~\cite[Theorem~1.2]{MZ}. They showed that its coefficients
are polynomials in $n$ with coefficients in $\mathbb Q[\pi^{-2},\pi^2]$
and observed numerically that only $\mathbb Q[\pi^{-2}]$
occurs~\cite[Theorem~1.2 and Remark~1.3]{MZ}; their leading term also
holds when $n^2/g\to0$~\cite[Theorem~1.8]{MZ}, and their cusp-ratio
estimate has error $O(n/g)$ for $n\ge1$~\cite[Lemma~5.1]{MZ}. Their
first unmarked correction enters \eqref{eq:AV5}. He, Wu and Xue obtained
first-order volume-ratio expansions with error $O((1+n^2)/g^2)$ for
$n=o(\sqrt g)$~\cite[Theorem~4.11]{HWX}. Anantharaman and Monk expanded
the volume polynomials at large genus~\cite{AM}. Guillermit sharpened
the coefficient bounds and treated boundary lengths of order $\sqrt g$
for fixed $n$~\cite[Theorems~2, 4, 5 and Remark~3]{Guillermit}. For
fixed $n$, her order-$j$ bound $(C_nj)^j$ grows more slowly with $j$
than our bound involving $(j!)^2$; our constants, however, are uniform
over all cusp counts. Hide, Macera and Thomas gave bounds explicit in
the expansion order, cusp count and degrees~\cite[Theorem~4.2,
A1--A4]{HMT}. Their genus must exceed an order-dependent multiple of the
cusp count, and the errors grow polynomially in the cusp count or
degrees, so these bounds do not yield the relative remainder of
Theorem~\ref{thm:intro-absolute} at positive cusp density.

\textbf{Growing degrees.}
Wu and Xue established effective upper bounds for normalized
intersection numbers when some degrees are large compared with
$\sqrt{2g-2+n}$, together with bounds for volumes with large boundary
lengths~\cite{WXPrime}. These bounds control individual coefficients
and volumes, whereas the adjacent-ratio asymptotics of
Theorem~\ref{thm:coefficient} require relative control across
neighboring degrees. A further difficulty is that a bracket may itself
be small, so that an additive error polynomial in the degree gives no
relative control. We use the Mirzakhani--Zograf bracket-difference
estimates near the lowest degrees; the full degree range requires the
harmonic comparison and the log-concavity argument of
Section~\ref{uniform-one-boundary-coefficient-ratios}.

\textbf{Pure $\psi$ intersection numbers.}
Motivated by Masur--Veech volumes of quadratic differentials,
Delecroix, Goujard, Zograf and Zorich conjectured uniform large-genus asymptotics for pure $\psi$ intersections (Appendix~E of the arXiv version of~\cite{DGZZ}). Aggarwal proved them for $o(\sqrt g)$ insertions~\cite{Agg}, and Guo and Yang gave a different proof for at most $C\log g$ insertions, for every fixed $C$, based on an explicit formula for the $n$-point generating series~\cite[Theorem~1]{GY}. For a fixed number of insertions, Eynard, Garcia-Failde, Giacchetto, Gregori and Lewa\'nski obtained resurgent expansions to every fixed order, uniformly in the degrees~\cite[Theorem~3.13]{EGGGL}.

Our comparison needs only fixed positive upper and lower bounds for the normalized pure-$\psi$ intersections. Appendix~\ref{app:primitive-bounds} proves these bounds by a scalar induction, adapting the recursion estimates of Guo, Yang and Zagier~\cite{GYZ} and using the published two-point bound of Delecroix, Goujard, Zograf and Zorich~\cite{DGZZBounds}. The absolute constant is fixed by the single sequence $\langle\tau_2^{3g-3}\rangle_g$ through the Itzykson--Zuber identity and Painlev\'e~I. To match the expansions at the pure-$\psi$ endpoint when the truncation order grows with $g$, we use the exact identity~\cite[equation~(6.10)]{IZ}, the zero-instanton contour formula of Garoufalidis, Its, Kapaev and Mari\~no~\cite{GIKM}, and the Borel representation of Costin, Costin and Huang~\cite{CCH}. Alameddine, Marchal, Nikolaev and Orantin recently solved the direct monodromy problem for the Borel-summed zero-parameter solution of Painlev\'e~I and identified it with the tritronqu\'ee solutions~\cite{AMNO}; we do not need this. Borot, Eynard and Giacchetto use comparison with Painlev\'e~I amplitudes to bound topological-recursion correlators, uniformly in genus and number of markings~\cite[Main Theorem and Section~1.3]{BEG}. Their bounds concern the growth of the correlators themselves, whereas the factorial bounds in Theorem~\ref{thm:intro-absolute} concern the order of the inverse-genus expansion after division by an explicit leading term.

\textbf{Instanton predictions.}\label{subsec:intro-instanton}
Eynard, Garcia-Failde, Gregori, Lewa\'nski and Schiappa derived large-genus asymptotics for ordinary Weil--Petersson volumes under a resurgence assumption~\cite{EGGLS}. Griguolo, Papalini, Russo and Seminara predicted the coefficients of ordinary Weil--Petersson volumes from instanton contributions~\cite[Section~3.3]{GPRS}, Johnson and Rodrigues wrote down the corresponding one-boundary contour integral~\cite[equation~(4.22)]{JR}, and Hatsuda, Matsumoto and Okuyama constructed transseries solutions of the KdV equation that governs the genus expansion of the JT partition function, whose leading nonperturbative sector reproduces the instanton terms~\cite{HMO}. After a change of conventions, the denominator of this integral is our harmonic model at $Y=0$, and its stationary equation is the saddle equation of Theorem~\ref{thm:coefficient} (Subsection~\ref{subsec:instanton-model}). Unlike the leading factor $\sinh(L/2)/(L/2)$ at bounded boundary length, this integral retains the dependence of the coefficients on their degree up to the top degree. Theorem~\ref{thm:coefficient} establishes the resulting asymptotic formula for the true adjacent coefficient ratios, uniformly over all degrees, from exact intersection identities and estimates; no resurgent expansion is assumed.

\subsubsection{Short geodesics and spectra of random hyperbolic surfaces}
For closed surfaces, Mirzakhani established a positive spectral gap
with high probability~\cite[Section~4.5]{Mir}. Wu and Xue~\cite{WX} and,
independently, Lipnowski and Wright~\cite{LipWright} raised the lower
bound to $3/16-\varepsilon$. Anantharaman and Monk then obtained
$2/9-\varepsilon$, first in the original version of the first part of
their Friedman--Ramanujan series~\cite[Theorem~10.1]{AMPartOne} and now
in a separate article~\cite{AMTwoNinths}, and then
$1/4-\varepsilon$~\cite[Theorem~1.6]{AMGap};
Hide, Macera and Thomas gave a polynomial rate of approach to $1/4$
\cite{HMT}. These results concern closed surfaces sampled with
Weil--Petersson measure. In a different model, uniformly random
covers of a fixed surface, Magee, Naud and Puder obtained the relative
spectral gap $3/16-\varepsilon$ for covers of a compact
surface~\cite{MNP}. For covers of a fixed cusped surface, Hide and Magee
proved that no new eigenvalues appear below $1/4-\varepsilon$, and
deduced that there are closed surfaces of genus tending to infinity with
first eigenvalue tending to $1/4$~\cite{HideMagee}. Growing cusp counts introduce the separating
pants mechanism studied by Shen and Wu~\cite{SW}. See~\cite{MonkNaud,Petri}
for surveys of spectral gaps of random hyperbolic surfaces.

The transition occurs at the scale $\sqrt g$. Below it, Hide proved~\cite[Theorem~1.3]{Hide}
that for $n=O(g^\alpha)$ with $0\le\alpha<1/2$ there is, with
probability tending to one, no nonzero eigenvalue below
$1/4-((2\alpha+1)/4)^2-\varepsilon$. Shen and Wu proved a
uniform positive lower bound for the spectral gap throughout the range
$n=o(\sqrt g)$~\cite[Theorem~3]{SW}, and He, Wu and Xue improved the
bound for $n=O(g^\alpha)$ to $1/4-(6(1-\alpha))^{-2}-\varepsilon$,
which tends to $5/36-\varepsilon$ as $\alpha\to1/2$~\cite[Theorem~1.2]{HWX}. At
$n\asymp\sqrt g$, Shen and Wu's result~\cite[Theorem~2]{SW} describes the small-value
distribution of the Cheeger constant, not a limiting law for Laplace
eigenvalues. Above this scale the spectral gap vanishes in probability
(Corollary~\ref{cor:spectral}), which answers negatively the question
of Shen and Wu whether a uniform spectral gap persists when
$n/\sqrt g\to\infty$~\cite[Question~2]{SW}.
This is a statement about the Weil--Petersson distribution: Guo, Hua
and Shen construct surfaces with $n/g$ tending to any prescribed
positive limit and a uniform spectral gap~\cite{GHS}.

Poisson limits and small eigenvalues at fixed genus were obtained by Hide and Thomas~\cite{HTSphere,HTLargeN}; their recent deterministic linear count of eigenvalues below $1/4$~\cite{HTMany} concerns a fixed threshold rather than the vanishing thresholds of Corollary~\ref{cor:many-spectrum}. For closed surfaces the picture is different: Monk proved that, with probability tending to one, a Weil--Petersson random closed surface has $O(g(\log g)^{-3/4})$ eigenvalues below $1/4$~\cite[Theorem~4]{Monk}. In the range $\sqrt g\ll n\ll g$, the second-moment volume inputs for our short-curve and small-eigenvalue results already follow from~\cite{SW}; what the uniform volume theorem adds is a single intensity valid across all cusp densities as the genus tends to infinity.

Budd and Lions studied the tight length spectrum in a near-critical Boltzmann model, where the cusp count is random~\cite{BL}. They asked whether their limit persists for other cusp weights and under the fixed-count measures~\cite[Open Question~1.6]{BL}. Recovering a prescribed count requires a local coefficient estimate, beyond concentration of the random count. Theorem~\ref{thm:tight-spectrum} supplies the fixed-count extension: tight volume estimates and a positive Fourier bound allow coefficient extraction at every cusp density. Corollary~\ref{cor:tight-boltzmann} then treats every subcritical cusp weight, including arbitrary approach to the critical endpoint. The planar separating curves in Theorem~\ref{thm:pants-poisson} instead become contractible after filling the cusps, and lie outside the tight spectrum. In genus zero, Budd and Curien showed that the Weil--Petersson random sphere with $n+1$ cusps, with its metric rescaled by $n^{-1/4}$, converges to the Brownian sphere~\cite{BC}.

\subsubsection{Zograf's constant and contemporaneous work}
The first version of this paper was made public on Zenodo on 16 October 2026~\cite{LWZenodo}, at \url{https://doi.org/10.5281/zenodo.23414478}. All numbered results and sections cited in this paragraph refer to \emph{Zenodo v1}, not to the present numbering. That version already contained the constant $C=1/\sqrt\pi$ (Corollary~1.2 and Section~7), the uniform volume-ratio limits (Theorem~1.1), the diagonal chain (Theorem~1.3), the adjacent-coefficient asymptotics over all degrees with the rate $O(g^{-1/8})$ (Theorems~1.4 and~17.1), the absolute expansion uniform in the cusp count with growing-order remainder bounds (Theorems~14.1--14.2), the vanishing spectral gap and many small eigenvalues (Corollaries~1.5--1.6), and the short-curve processes and the systole (Theorem~19.1).

In contemporaneous work, Liu and Xu~\cite{LXConstant}, whose preprint was submitted to arXiv on 24 October 2026, also determine $C=1/\sqrt\pi$. Their deformation joins the cubic curve $y=z-4z^3$ to the sine curve $y=\sin(2\pi z)/(2\pi)$ while keeping the first positive zero fixed. For each fixed number $n\ge2$ of markings, their uniform-variation theorem~\cite[Theorem~1.1]{LXConstant} controls the logarithmic derivative of the normalized volumes uniformly in the curve parameter $t\in[0,1]$ by $O_n(g^{-1})$. Transport from the cubic endpoint determines the constant~\cite[Corollary~1.2]{LXConstant}, and their first-correction theorem~\cite[Theorem~1.3]{LXConstant} adds the $1/g$ term with an $O_n(g^{-2})$ remainder, uniformly in $t$.

Our deformation is instead the string flow associated with the cusp weight $x=X(Y)$. Its amplitudes with a fixed number $p$ of distinguished markings encode the derivatives $\Phi_g^{(p)}(x)$ of the cusp generating function. Both arguments use positivity, descendant recursions, the dilaton and KdV equations, and a Painlev\'e~I normalization, but along different deformation families. In the present paper, Section~\ref{sec:prefactor} already determines the constant by specializing to $p=1$ and $Y=0$, hence to $V_{g,1}$: uniformity in the total cusp count is not needed for this identification. Its role is the stronger conclusion of an expansion valid for every $n\ge0$, including $n$ growing with $g$. This requires the coefficient extraction in Sections~\ref{extracting-fixed-cusp-counts}--\ref{sec:sharp-extraction} and Appendix~\ref{app:extraction} of the present paper, beyond the estimates for a fixed number of distinguished markings.

\subsection{Organization and notation}\label{subsec:notation}
\textbf{Organization.}
The main text proves the volume-ratio and adjacent-coefficient limits and
the leading-order estimates. It states the higher-order expansions and
the tight and critical spectrum theorems with outlines of their proofs;
the supporting arguments and the proofs of these statements are collected
in Appendices~\ref{app:kernels}--\ref{app:critical}.
The paper has four parts.
\begin{itemize}[leftmargin=*]
\item \emph{The positive recursion}
(Sections~\ref{the-positive-compact-family-of-kernels}--\ref{uniform-comparison}).
Section~\ref{the-positive-compact-family-of-kernels} introduces the
kernels, Section~\ref{the-stationary-parameter} the density parameter,
Section~\ref{exact-amplitudes-and-the-harmonic-model} the shifted
amplitudes and their positive recursion, and
Section~\ref{uniform-comparison} compares the recursion with the harmonic
model. The special-function facts and the bounds for pure $\psi$
intersections are proved in Appendices~\ref{app:kernels}
and~\ref{app:primitive-bounds}.
\item \emph{Generating functions}
(Sections~\ref{ratios-with-a-fixed-number-of-distinguished-markings}--\ref{sec:higher-order}).
Section~\ref{ratios-with-a-fixed-number-of-distinguished-markings} proves
the cusp and genus ratios for a fixed number of markings,
Section~\ref{sec:prefactor} the leading term and Zograf's constant, and
Section~\ref{sec:higher-order} states the higher-order expansion, which is
proved in Appendix~\ref{app:marked}.
\item \emph{Individual volumes}
(Sections~\ref{extracting-fixed-cusp-counts}--\ref{sec:absolute}).
Section~\ref{extracting-fixed-cusp-counts} proves the local limit theorem
and the volume-ratio limits, Section~\ref{sec:sharp-extraction} states the
sharper extraction theorems proved in Appendix~\ref{app:extraction}, and
Section~\ref{sec:absolute} proves Theorem~\ref{thm:intro-absolute} and the
rate in Theorem~\ref{thm:volume}.
\item \emph{Applications}
(Sections~\ref{uniform-logarithmic-volume-rate-and-the-full-diagonal-chain}--\ref{sec:geometry}).
Section~\ref{uniform-logarithmic-volume-rate-and-the-full-diagonal-chain}
treats the diagonal chain,
Section~\ref{uniform-one-boundary-coefficient-ratios} the adjacent
coefficients, with the rate proved in
Appendix~\ref{sec:quantitative-coefficients}, and
Section~\ref{sec:geometry} the geometry of random surfaces, with the proofs
of the tight and critical spectrum theorems in Appendices~\ref{app:tight}
and~\ref{app:critical}.
\end{itemize}
The final section summarizes the picture, explains the matrix-model
meaning of the parameter and lists open questions.

\Needspace{16\baselineskip}
\textbf{Proof guide.}
The following table lists where the main results are proved and the principal
results and methods used in their proofs. References to Sections~\ref{the-positive-compact-family-of-kernels}--\ref{uniform-comparison} include the supporting Appendices~\ref{app:kernels} and~\ref{app:primitive-bounds}.
\begin{center}
\small
\renewcommand{\arraystretch}{1.15}
\begin{tabular}{@{}p{0.31\textwidth}p{0.17\textwidth}p{0.45\textwidth}@{}}
\hline
Result & Proof location & Main ingredients\\
\hline
Theorem~\ref{thm:intro-absolute} & Section~\ref{sec:absolute} & Theorems~\ref{thm:full-expansion} and~\ref{thm:canonical-expansion} (Appendices~\ref{app:marked} and~\ref{app:extraction})\\
Corollary~\ref{cor:constant} (Zograf's constant) & Section~\ref{sec:prefactor} & Sections~\ref{the-positive-compact-family-of-kernels}--\ref{ratios-with-a-fixed-number-of-distinguished-markings}; Painlev\'e~I at the pure-$\psi$ endpoint; \cite{MZ}\\
Theorem~\ref{thm:volume}, limits & Section~\ref{extracting-fixed-cusp-counts} & Sections~\ref{the-positive-compact-family-of-kernels}--\ref{ratios-with-a-fixed-number-of-distinguished-markings}; fixed-cusp limits~\cite{Mir}\\
Theorem~\ref{thm:volume}, rate $O(g^{-1})$ & Section~\ref{sec:absolute} & Theorem~\ref{thm:sharp-extraction} (Appendix~\ref{app:extraction}), which uses Theorem~\ref{thm:full-expansion} at a fixed order (Appendix~\ref{app:marked})\\
Theorem~\ref{thm:coefficient}, uniform equivalent & Section~\ref{uniform-one-boundary-coefficient-ratios} & Proposition~\ref{prop:harmonic}; log-concavity; \cite{MZ}\\
Theorem~\ref{thm:coefficient}, rate & Appendix~\ref{sec:quantitative-coefficients} & as above\\
Theorem~\ref{thm:chain} & Section~\ref{uniform-logarithmic-volume-rate-and-the-full-diagonal-chain} & limits in Theorem~\ref{thm:volume}; \cite{ManZ}\\
Corollaries~\ref{cor:spectral} and~\ref{cor:many-spectrum}; Theorem~\ref{thm:pants-poisson} & Section~\ref{sec:geometry} & limits in Theorem~\ref{thm:volume}; Sections~\ref{ratios-with-a-fixed-number-of-distinguished-markings}, \ref{sec:prefactor} and~\ref{extracting-fixed-cusp-counts}\\
Theorem~\ref{thm:tight-spectrum}, Corollary~\ref{cor:tight-boltzmann} & Appendix~\ref{app:tight} & Sections~\ref{ratios-with-a-fixed-number-of-distinguished-markings}, \ref{sec:prefactor} and~\ref{extracting-fixed-cusp-counts}; \cite{BL,BZ}\\
Theorem~\ref{thm:critical-cluster} & Appendix~\ref{app:critical} & volume bounds of~\cite{MZ,SW}; zero-density limits in Theorem~\ref{thm:volume}; Proposition~\ref{prop:bounded-boundaries}\\
\hline
\end{tabular}
\end{center}
Thus all limits, and all geometric applications except the two spectrum theorems proved in Appendices~\ref{app:tight} and~\ref{app:critical}, are proved in the main text using only Appendices~\ref{app:kernels} and~\ref{app:primitive-bounds}. Appendices~\ref{app:marked}--\ref{sec:quantitative-coefficients} are needed only for the inverse-genus rates and for the expansions to higher order. They can be read independently of one another: Appendix~\ref{app:marked} after Section~\ref{sec:higher-order}, Appendix~\ref{app:extraction} after Section~\ref{sec:sharp-extraction}, using Theorem~\ref{thm:full-expansion} only through its statement, and Appendix~\ref{sec:quantitative-coefficients} after Section~\ref{uniform-one-boundary-coefficient-ratios}.

\textbf{Notation.} The following table lists the principal symbols and where they are defined.
{\small
\renewcommand{\arraystretch}{1.15}
\begin{longtable}{@{}lll@{}}
\hline
Symbol & Meaning & Defined in\\
\hline
\endhead
$X$, $a=X'$, $y_0$, $x_c$ & genus-zero relation, $X'(y)=J_0(2\sqrt y)$ & Subsection~\ref{subsec:parameter}, Section~\ref{the-positive-compact-family-of-kernels}\\
$x=X(Y)$, $Y\in[0,y_0)$ & cusp weight and parameter & Subsection~\ref{subsec:parameter}\\
$\Phi_g$ & normalized cusp generating function & \eqref{eq:intro-Phi}\\
$K_Y$, $H_Y$ & positive kernels & \eqref{eq:P1}\\
$\bar K_Y$, $\bar H_Y$ & kernels rescaled by $r(Y)$ & \eqref{eq:P9}\\
$\theta_\ell$, $w_\ell$ & coefficients of $1/\bar K_Y$, increments & \eqref{eq:P11}\\
$\vartheta_\ell$ & coefficients of $1/K_Y$ & Lemma~\ref{lem:kernel-positivity}\\
$r(Y)=2\Delta(Y)$ & first positive zero of $K_Y$ & Sections~\ref{the-positive-compact-family-of-kernels}, \ref{the-stationary-parameter}\\
$T$, $\Delta=T-Y$ & stationary root and half-zero & \eqref{eq:SP1}, \eqref{eq:SP2}\\
$I$, $\lambda$ & action and cusp density & \eqref{eq:SP2}\\
$\Psi=\lambda/2$ & \begin{tabular}[t]{@{}l@{}}leading cusp-count scale\\per unit complexity\end{tabular} & Section~\ref{the-stationary-parameter}\\
$\mathscr D=x\,\mathrm d/\mathrm dx$ & logarithmic derivative operator & Section~\ref{the-stationary-parameter}\\
$h$ & complexity, $2g-2+n$ or $2g-2+p$ & see below\\
$m$ & deficit $3g-3+n-|\mathbf d|$ & Section~\ref{exact-amplitudes-and-the-harmonic-model}\\
$U^Y$, $\mathcal U$ & shifted amplitudes, normalized & \eqref{eq:H6}, \eqref{eq:H7}\\
$\mathcal M$, $A_m(h)$ & harmonic model & \eqref{eq:H12}\\
$\overline U$, $W_{g,p}$, $c_{g,p}$ & marked amplitudes and degree profiles & \eqref{eq:TL9}\\
$A$, $\Theta$, $M_1$, $\mathcal Z=\sqrt2\,I$ & kernel scalars and action & \eqref{eq:P9}, \eqref{eq:PF1}\\
$\mathcal J_{g,p}$ & leading term for $\Phi_g^{(p)}$ & \eqref{eq:PF3}\\
$\widehat V_{g,n}$, $\varpi$ & leading term for $V_{g,n}$ & \eqref{eq:intro-leading}, \eqref{eq:intro-varpi}\\
$B_{g,n}$, $C_{g,n}$, $D_{g,n}$ & volume ratios & Subsection~\ref{subsec:intro-ratios}\\
$\mathcal B$, $\mathcal C$ & volume-ratio profiles & Subsection~\ref{subsec:intro-ratios}\\
$C_g(\mathbf e)$ & normalized pure $\psi$ intersection number & \eqref{eq:H4}\\
$\mathfrak p_k=\mathrm e^{-k}k^k/k!$ & Poisson point mass at its mean & Section~\ref{sec:sharp-extraction}\\
\hline
\end{longtable}}

The complexity parameter $h$ is specified locally: we use $h=2g-2+n$ for volumes with $n$ boundary components and $h=2g-2+p$ for generating functions with $p$ distinguished markings. In particular, $h=2g-1$ in \eqref{eq:intro-leading}, where $p=1$ and the remaining cusps are counted by $k$. Positive constants $C$, $c$ may change from line to line, and a subscript, as in $O_p$, indicates the parameters on which an implied constant depends. A few letters have an additional local meaning, stated where it occurs: $D$ is the degree sum $|\mathbf d|$ in Sections~\ref{exact-amplitudes-and-the-harmonic-model} and~\ref{sec:prefactor} and in Appendix~\ref{app:marked}; $h_0=2g-2$ is the complexity of a closed surface (Appendices~\ref{app:marked}, \ref{app:tight} and~\ref{app:critical}); the harmonic coefficients $A_m(h)$ are unrelated to the kernel scalar $A(Y)$ and to the coefficients $A_j(g,n)$; the Laplacian $\Delta_X$ of a surface $X$ is unrelated to $\Delta(Y)$; and $T$ is a length window in Section~\ref{sec:geometry}.

\section{Positive kernels}\label{the-positive-compact-family-of-kernels}

Two entire functions of a variable $z$, depending on the parameter $Y$,
control everything that follows. The first, $K_Y$, arises in
Section~\ref{exact-amplitudes-and-the-harmonic-model} when the marking
singled out by the Dijkgraaf--Verlinde--Verlinde equation is merged with
the $\kappa_1$ insertions: these mergers act on the amplitudes as
$K_Y(E)$, where $E$ raises a degree by one, so the Taylor coefficients of
$1/K_Y$ become the coefficients of the recursion. The second, $H_Y$,
plays the same role for the dilaton equation and for the explicit
comparison model. The later arguments need three properties: all zeros
of $K_Y$ are positive and simple (Lemma~\ref{lem:kernel-zeros});
$1/K_Y$ and $1/H_Y$ have nonnegative Taylor coefficients, even as
functions of $Y$ (Lemma~\ref{lem:kernel-positivity}); and after
rescaling by the first zero of $K_Y$ both kernels form a compact family,
including the endpoint $Y=y_0$
(Subsection~\ref{compact-normalization-and-uniform-lattice-estimates}).
The last property is what makes all later constants uniform in $Y$.
The first two are special-function statements; their proofs are given
in Appendix~\ref{app:kernels}.

Recall from Subsection~\ref{subsec:parameter} that \[
 X(y)=\sum_{j\ge0}\frac{(-1)^j y^{j+1}}{j!(j+1)!}
      =\sqrt yJ_1(2\sqrt y),\qquad
 a(y)=X'(y)=J_0(2\sqrt y),
\] that \(y_0=j_{0,1}^2/4\) is the first positive zero of \(a\), and that
\(x_c=X(y_0)\). Thus \(a>0\) on \([0,y_0)\), and the inverse branch
\(Y(x)\), with \(Y(0)=0\), maps \([0,x_c)\) onto \([0,y_0)\).

We shall use \begin{equation}
\begin{split}
 K_Y(z)&=\int_0^1X'\bigl(Y+\tfrac z2(1-s^2)\bigr)\,\mathrm{d}s
       =\sum_{\ell\ge0}\frac{X^{(\ell+1)}(Y)}{(2\ell+1)!!}z^\ell,\\
 H_Y(z)&=3\int_0^1s^2X'\bigl(Y+\tfrac z2(1-s^2)\bigr)\,\mathrm{d}s
       =\sum_{\ell\ge0}\frac{3X^{(\ell+1)}(Y)}{(2\ell+3)!!}z^\ell.
 \end{split}                                                   \label{eq:P1}
\end{equation} Beta integration gives the series and hence \begin{equation}
K_Y=H_Y+\tfrac23z\partial_zH_Y.                    \label{eq:P2}
\end{equation} We write \(r(Y)\) for the first positive zero of \(K_Y\); it equals \(2\Delta(Y)\).

\subsection{Real zeros and coefficient positivity}\label{real-zeros-of-the-kernel}

\begin{lemma}[Real zeros]\label{lem:kernel-zeros}
On \([0,\pi^2/4]\), \(X(y)>0\) for \(y>0\), \(X''<0\), and \(X'''>0\).
Moreover \(y_0<3/2\); in particular \(X(Y)\le Y<3/2\) for
\(0\le Y\le y_0\). All zeros of \(a(y)\) are positive real; in particular
\(a\) has no zeros in \(|y|<y_0\). For \(0\le Y<y_0\), every zero of
\(K_Y\) is positive and simple, and \begin{equation}
K_Y(z)=a(Y)\prod_{j\ge1}(1-z/r_j(Y)),
 \qquad 0<r_1(Y)<r_2(Y)<\cdots.                     \label{eq:P5}
\end{equation}
\end{lemma}

At \(Y=0\) one has \(K_0(z)=\sin\sqrt{2z}/\sqrt{2z}\), and the lemma is
the classical description of the zeros of the sine function. For
\(Y>0\), integration along the characteristics of a first-order
equation for \(K_Y\) writes \(K_Y\) as \(\sin t\) minus a small positive
correction; concavity then locates the real zeros, and Rouch\'e's
theorem shows that there are no others.

Reality of the zeros of \(K_Y\) does not imply it for \(H_Y\); indeed
some higher zeros of \(H_Y\) can be nonreal. The required substitute is
the following stronger coefficient assertion.

\begin{lemma}[Coefficient positivity]\label{lem:kernel-positivity}
For \(0\le Y<y_0\), \begin{equation}
\log\frac{a(Y)}{H_Y(z)}=\sum_{m\ge1}t_m(Y)z^m,
 \qquad t_m^{(q)}(0)\ge0\quad(m\ge1,q\ge0),\quad t_m(0)>0, \label{eq:P6}
\end{equation} where each \(t_m\) is analytic on \(|Y|<y_0\) and its Taylor
series at zero converges to \(t_m(Y)\) there. Consequently \(H_Y^{-h}\)
has strictly positive Taylor coefficients for every real \(h>0\), and
\(-H_Y'/H_Y\) has nonnegative coefficients. Moreover every Taylor
coefficient in \(Y\) of \[
 \widehat\vartheta_\ell(Y):=a(Y)[z^\ell]K_Y(z)^{-1}
\] is nonnegative, and the same holds for \(a^{-1}\) and for
\(\vartheta_\ell=[z^\ell]K_Y^{-1}\).
\end{lemma}

The coefficients \(t_m\) satisfy a system of ordinary differential
equations in \(Y\) whose right-hand sides are polynomials with
nonnegative coefficients, and they are positive at \(Y=0\), where \(H_0\)
has only positive real zeros; induction on the order of differentiation
gives the lemma. Note that these are statements about Taylor coefficients
in the parameter \(Y\), which is stronger than positivity at each real
\(Y\); they supply the positive convolution factor of
Section~\ref{a-positive-convolution-factor}.

\subsection{Rescaled kernels and a local limit theorem}\label{compact-normalization-and-uniform-lattice-estimates}

Set \(r=r_1(Y)\), \begin{equation}
\bar K_Y(w)=K_Y(rw)/a(Y),\qquad
 \bar H_Y(w)=H_Y(rw)/a(Y),\qquad
 A(Y)=\bar H_Y(1)/3.                               \label{eq:P9}
\end{equation} The first zero \(r\) is simple and varies analytically. Integration
of \eqref{eq:P2} gives \[
 H_Y(z)=\frac3{2z^{3/2}}\int_0^zv^{1/2}K_Y(v)\,\mathrm{d}v,
\] so \(H_Y(z)>0\) on \([0,r]\). Since \(H_Y^{-1}\) has nonnegative
Taylor coefficients, its radius of convergence exceeds \(r\). Indeed, by
Pringsheim's theorem a power series with nonnegative coefficients is
singular at the positive point of its circle of convergence; if the
radius were at most \(r\), that point would be a zero of \(H_Y\) in
\((0,r]\). In particular \(H_Y\) has no complex zero of modulus at most
\(r\).

At the critical endpoint put \(b=-X''(y_0)>0\). The equation
\(K_Y(r)=0\) and Taylor expansion at \(Y=y_0,r=0\) give \begin{equation}
r(Y)\sim3a(Y)/b,\qquad
 \bar K_Y(w)\longrightarrow1-w,\qquad
 \bar H_Y(w)\longrightarrow1-3w/5.                 \label{eq:P10}
\end{equation} These convergences are locally uniform in complex \(w\). The family
therefore extends continuously to \(Y=y_0\). Compactness and
nonvanishing on the closed unit disk give one \(\eta>0\) such that all
\(\bar H_Y^{-1}\) are analytic and uniformly bounded on
\(|w|\le1+\eta\), including this endpoint. Likewise \(\bar K_Y/(1-w)\)
is nonvanishing on a common disk of radius greater than one: its zeros
are the higher positive roots from \eqref{eq:P5}, and \eqref{eq:P10} treats the endpoint.
Consequently, with \(\theta_{-1}=0\), \begin{equation}
1=\theta_0\le\theta_1\le\cdots\le C,
 \quad \theta_\ell=[w^\ell]\bar K_Y(w)^{-1},
 \quad w_\ell:=\theta_\ell-\theta_{\ell-1}\ge0,
 \quad\sup_Y\sum_{\ell\ge0}(1+\ell)^q w_\ell<\infty.  \label{eq:P11}
\end{equation} Indeed \(\sum w_\ell w^\ell=(1-w)/\bar K_Y(w)\) is the inverse
product over the higher roots. It has nonnegative coefficients and a
common analytic disk. Also \(A(Y)\) is bounded above and away from zero.

Finally, \begin{equation}
P_Y(w)=\frac{\bar H_Y(1)}{\bar H_Y(w)}              \label{eq:P12}
\end{equation} is the probability generating function of a nonnegative integer
variable. Its probabilities at zero and one have uniform positive lower
bounds, its exponential moments are uniformly bounded, and \eqref{eq:P2} at the
first root gives mean \(3/2\). Its variance \(\sigma_Y^2\) is bounded
above and away from zero. At \(Y=y_0\) the law is geometric, with
generating function \((2/5)/(1-3w/5)\), so its variance is \(15/4\).

For a sum \(S_h\) of \(h\) independent copies, Fourier inversion now
gives, uniformly in \(Y\) and in integers \(m\) with
\(|m-3h/2|\le L\sqrt h\), for each fixed \(L\), \begin{equation}
\Pr(S_h=m)=\frac1{\sqrt{2\pi h\sigma_Y^2}}
     \exp\!\left(-\frac{(m-3h/2)^2}{2h\sigma_Y^2}\right)
       +o(h^{-1/2}).                              \label{eq:P13}
\end{equation} Indeed, on a fixed neighborhood of zero the centered
characteristic logarithm is \(-\sigma_Y^2t^2/2+O(|t|^3)\), uniformly by
the exponential moments. Away from zero its modulus is uniformly less
than one, because the first two probabilities have uniform positive
lower bounds. Scaling \(t\) by \(h^{-1/2}\), these facts give an
integrable Gaussian bound and pointwise convergence; Fourier inversion
proves \eqref{eq:P13}. At a bounded displacement from the mean the error is
relative \(O(h^{-1})\), by expansion through fourth order and
cancellation of the odd cubic integral. Thus \[
 [w^m]\bar H_Y(w)^{-h}
   =\bar H_Y(1)^{-h}\Pr(S_h=m)
\] has a uniform central lower bound \(c_L\bar H_Y(1)^{-h}/\sqrt h\),
and its logarithm at a bounded displacement is
\(-h\log\bar H_Y(1)-\tfrac12\log h+O(1)\). These are the coefficient
estimates used below.

\section{The density parameter}\label{the-stationary-parameter}

This section establishes the properties of the functions $T$,
$\Delta=T-Y$, $I$ and $\lambda$ defined in \eqref{eq:SP1} and
\eqref{eq:SP2}. Their meaning is as follows. The first positive zero of
$K_Y$ is $r(Y)=2\Delta(Y)$; equation \eqref{eq:SP1} is the condition
$K_Y(2\Delta)=0$ written in the variable $v=T-\Delta s^2$. The integral
$I(Y)$ is the action that governs the exponential growth
in~\eqref{eq:intro-heuristic}, and $\lambda(Y)/2$ is the logarithmic
derivative of $I(Y)^{-1}$ with respect to $\log x$, that is, the leading
cusp-count scale per unit complexity. Lemma~\ref{lem:parameter} below
shows that $T$ is well defined and that $\lambda$ is an analytic
increasing bijection from $[0,y_0)$ onto $[0,\infty)$, so that the
profiles $\mathcal B,\mathcal C$ of \eqref{eq:SP3} are functions of the
density $\lambda\ge0$; it also records the behavior of all four
functions at the two ends of the interval. Its proof, which uses the
Sonine integral and elementary calculus, is given in
Appendix~\ref{app:parameter-proof}.

For the local-limit arguments we view $\lambda$ as a function of the cusp
weight and put \[
 \Psi(x)=\lambda(Y(x))/2,\qquad x=X(Y),\qquad
 \mathscr D=x\frac{\mathrm{d}}{\mathrm{d}x}.
\]

\begin{lemma}[The parameter]\label{lem:parameter}
\begin{enumerate}[label=\textup{(\roman*)},leftmargin=*]
\item The root \(T(Y)\) of \eqref{eq:SP1} exists, is unique and depends
analytically on \(Y\in[0,y_0)\). One has \(T(0)=\pi^2/4\),
\(-1\le T'<0\), \(T''>0\), and \(T\to y_0\) as \(Y\to y_0\). Moreover
\(I>0\).
\item Put \begin{equation}
J(Y)=\int_0^1X(T-\Delta s^2)\,\mathrm{d}s
     =X(Y)+\frac{I(Y)}{\sqrt\Delta},
 \qquad
 \beta(Y)=\frac{X(Y)}{J(Y)}=
           \frac{\lambda(Y)}{2+\lambda(Y)}.
 \label{eq:SP11}
\end{equation} Then \(\beta'>0\), and \(\lambda\) is an analytic strictly
increasing bijection from \([0,y_0)\) onto \([0,\infty)\). At \(Y=0\),
\(\Delta=T=\pi^2/4\), \(I(0)=\pi/4\), \(J(0)=1/2\), and
\(\lambda(Y)=4Y+O(Y^2)\).
\item Put \(\delta=y_0-Y\) and \(b=-X''(y_0)=x_c/y_0>0\). As
\(\delta\to0\), \begin{equation}
\Delta=\frac32\delta+O(\delta^2),\qquad
 I(Y)=\frac{8b}{45}\Delta^{5/2}(1+O(\Delta)),
 \label{eq:SP13}
\end{equation} \begin{equation}
\lambda\sim\frac{45x_c}{4b}\Delta^{-2}
       =\frac{45j_{0,1}^2}{16}\Delta^{-2},
 \label{eq:SP14}
\end{equation} and consequently \begin{equation}
\mathcal B(\lambda)\sim
 \frac{\pi^2}{3\sqrt5\,j_{0,1}}\sqrt\lambda,\qquad
 \mathcal C(\lambda)\sim\frac{\pi^2}{x_c}\lambda.
 \label{eq:SP15}
\end{equation} Also \(\Delta\asymp\sqrt{1-\beta}\).
\item The function \(\Psi\) is analytic and strictly increasing for
\(0<x<x_c\). With \(\varepsilon=1-x/x_c\), \begin{equation}
\begin{array}{ll}
 x\downarrow0:&
 \Psi=2x+O(x^2),\quad\mathscr D\Psi=2x+O(x^2),\quad
 \mathscr D^2\Psi=2x+O(x^2),\\[1mm]
 x\uparrow x_c:&
 \displaystyle\Psi=\frac5{4\varepsilon}(1+O(\sqrt\varepsilon)),\quad
 \mathscr D\Psi=\frac5{4\varepsilon^2}(1+O(\sqrt\varepsilon)),\quad
 \mathscr D^2\Psi=O(\varepsilon^{-3}).
 \end{array}
 \label{eq:SP16}
\end{equation} These expansions arise from convergent analytic expansions
in \(x\) or \(\sqrt\varepsilon\). In particular \(\Psi^2/(\mathscr D\Psi)\)
has finite endpoint limits \(0\) and \(5/4\), and is bounded on the full
interval.
\end{enumerate}
\end{lemma}

The estimate \(\Delta\asymp\sqrt{1-\beta}\) in~(iii) implies the
integrability at one of the logarithmic chain profile in
Section~\ref{uniform-logarithmic-volume-rate-and-the-full-diagonal-chain}.

\section{Intersection amplitudes and the comparison model}\label{exact-amplitudes-and-the-harmonic-model}

We deform the Weil--Petersson intersection brackets into pure $\psi$
intersection numbers and derive a positive recursion along the deformation.
We call the resulting functions shifted amplitudes. The harmonic model
introduced below is an explicit, solvable comparison for this recursion.
Section~\ref{uniform-comparison} compares the full recursion with that
model, uniformly in the degrees and in the deformation parameter.
This comparison is the common starting point for the volume and
coefficient asymptotics.

\subsection{Intersection identities and normalization}\label{exact-amplitudes-standard-intersection-identities-and-normalization}

We use the following notation for a stable pair \((g,n)\) and a degree
tuple \(\mathbf d\): \(h\) is the complexity, \(N\) the dimension of
\(\overline{\mathcal M}_{g,n}\), and the deficit \(m\) is the number of
\(\kappa_1\) factors needed to fill up the dimension. Write \[
h=2g-2+n,\qquad N=3g-3+n,\qquad
D=\sum_{i=1}^n d_i,\qquad m=N-D.
\] Throughout the comparison theorem \(n\ge1\), the pair is stable, and
\(0\le D\le N\). A tuple with a negative degree or an unstable pair has
amplitude zero.

Recall \(X\), \(a=X'>0\) on \([0,y_0)\) and the kernels \(K_Y\) and
\(H_Y\) of \eqref{eq:P1}, and write \(r_Y=r(Y)\) for the first positive
zero of \(K_Y\). The properties of these kernels used below were
established in Section~\ref{the-positive-compact-family-of-kernels}. Put \begin{equation}
\bar H_Y(z)=\frac{H_Y(r_Yz)}a,\qquad
\bar K_Y(z)=\frac{K_Y(r_Yz)}a
      =\bar H_Y(z)+\frac23z\bar H_Y'(z),\qquad
\theta_\ell(Y)=[z^\ell]\bar K_Y(z)^{-1}.
\label{eq:H2}
\end{equation} At \(Y=y_0\) these have the continuous definitions \begin{equation}
\bar H_{y_0}(z)=1-\frac35z,\qquad \bar K_{y_0}(z)=1-z,\qquad
\theta_\ell(y_0)=1.
\label{eq:H3}
\end{equation} In particular \(\theta_0=1\). These are the rescaled kernels and
coefficients of \eqref{eq:P9} and \eqref{eq:P11}, recalled here together
with their endpoint values.

We use the following facts from intersection theory. The ordinary string,
dilaton, and Dijkgraaf--Verlinde--Verlinde (DVV) identities hold for Witten's psi
intersections~\cite{Witten,Kontsevich,DVV}, with
\(\langle\tau_0^3\rangle_0=1\) and \(\langle\tau_1\rangle_1=1/24\). The normalized forms used here are written in Appendix~\ref{app:primitive-bounds}. For
pure $\psi$ intersection numbers we use the normalization \begin{equation}
C_g(\mathbf e)=
\frac{4^g\prod_i(2e_i+1)!!}
 {3^{\,2g-2+q}\Gamma(2g-2+q)}
\left\langle\prod_{i=1}^q\tau_{e_i}\right\rangle_g,
\qquad \sum_i e_i=3g-3+q.
\label{eq:H4}
\end{equation} For positive degrees we use only
\[
 \frac16\le C_g(\mathbf e)\le1.
\]
Appendix~\ref{app:primitive-bounds} proves these coarse bounds directly
from DVV and the published two-point estimate. The argument adapts the
recursion estimates of Guo, Yang and Zagier~\cite{GYZ}; their stronger
lower bound $C_g(3g-2)\le C_g(\mathbf e)$ is not needed here.

With \(m=3g-3+n-|\mathbf d|\), our mixed brackets are
\[
 [\tau_{d_1}\cdots\tau_{d_n}]_{g,n}
 =\prod_i(2d_i+1)!!\,4^{|\mathbf d|}\frac{(2\pi^2)^m}{m!}
 \int_{\overline{\mathcal M}_{g,n}}\kappa_1^m\prod_i\psi_i^{d_i}.
\]
Here \(\kappa_1\) is the logarithmic Mumford class, \(\psi_i\) is the cotangent class, and \([\omega_{\mathrm{WP}}]=2\pi^2\kappa_1\)~\cite{Wolpert}; see also~\cite[Section~2]{MZ}. Thus
\[
 V_{g,n}(2\mathbf L)=\sum_{|\mathbf d|\le3g-3+n}
 [\tau_{d_1}\cdots\tau_{d_n}]_{g,n}
 \prod_i\frac{L_i^{2d_i}}{(2d_i+1)!}.
\]
Angle brackets below denote unweighted pure-$\psi$ integrals. Intersections of the wrong dimension and negative-degree amplitudes are zero. Unstable factors in the positive recursion are zero; explicitly specified genus-zero initial derivatives are kept separate.

The mixed-kappa/psi transform~\cite{KMZ,LX} gives, at
\(Y=0\), \begin{equation}
\frac{[\tau_{d_1}\cdots\tau_{d_n}]_{g,n}}
 {4^D(2\pi^2)^m}
=\prod_i(2d_i+1)!!
\sum_{\substack{\ell_k\ge0,\ k\ge2\\
                 \sum_{k\ge2}(k-1)\ell_k=m}}
\left\langle\tau_{\mathbf d}\prod_{k\ge2}\tau_k^{\ell_k}\right\rangle_g
\prod_{k\ge2}\frac{((-1)^k/(k-1)!)^{\ell_k}}{\ell_k!}.
\label{eq:H5}
\end{equation}
Since \(X'(0)=1\) and \(-X^{(k)}(0)=(-1)^k/(k-1)!\), the weights in
\eqref{eq:H5} are the derivatives \(-X^{(k)}(0)\), \(k\ge2\). The shifted
amplitudes replace them by the derivatives \(-X^{(k)}(Y)\) at a point
\(Y\). For general \(Y<y_0\), define the shifted amplitudes by the finite sum
\begin{equation}
U^Y_{g,n}(\mathbf d)=\prod_i(2d_i+1)!!
\sum_{\substack{\ell_k\ge0,\ k\ge2\\
                 \sum(k-1)\ell_k=m}}
\left\langle\tau_{\mathbf d}\prod_{k\ge2}\tau_k^{\ell_k}\right\rangle_g
\frac{\prod_{k\ge2}(-X^{(k)}(Y))^{\ell_k}/\ell_k!}
 {a^{h+\sum\ell_k}}.
\label{eq:H6}
\end{equation} The denominator is obtained by summing degree-one background
insertions with time \(t_1=1-a\): the dilaton equation gives the factor
\(\sum_{j\ge0}(h+\sum\ell_k)_j t_1^j/j!=a^{-h-\sum\ell_k}\). Define the
normalized amplitudes \begin{equation}
\mathcal U_{g,n}(\mathbf d;Y)=a^h r_Y^m U^Y_{g,n}(\mathbf d).
\label{eq:H7}
\end{equation} They extend to \(Y=y_0\) through their positive recursion below.

When all degrees vanish, the shifted amplitudes are derivatives of the
cusp generating function. This is the familiar shift of the Kontsevich--Witten (KW) times
along the string flow~\cite{IZ,KMZ}; the identification of $\Phi_g$ with the
KW free energy at the shifted times $t_0=x$, $t_1=0$, $t_k=(-1)^k/(k-1)!$
is recorded in~\cite[Remark~1]{Zog98}, where it is attributed to Dubrovin, and the corresponding deformation in
$(t_0,t_1)$ is used in~\cite[Section~2.4]{OS}. We recall it in the present
normalization. Start with KW
times \(t_0=x,\ t_1=0,\ t_k=-X^{(k)}(0)\), \(k\ge2\). The finite formal
string flow sends them to \[
t_0'=x-X(Y),\qquad t_1'=1-X'(Y),\qquad
t_k'=-X^{(k)}(Y)\quad(k\ge2).
\] Its vector field is
\(\sum_{k\ge0}t_{k+1}'\partial_{t_k'}-\partial_{t_0'}\). The string
equation makes the genus-\(g\) free energy invariant for \(g\ge1\); its
genus-zero anomaly disappears after three \(t_0\) derivatives. Since
\(\partial_{t_0}\) commutes with this flow, setting \(x=X(Y)\) gives \begin{equation}
U^Y_{g,n}(0^n)=\Phi_g^{(n)}(x),\qquad
\Phi_g(x)=\sum_{n\ge0}\frac{V_{g,n}}{(2\pi^2)^{3g-3+n}}\frac{x^n}{n!},
\label{eq:H8}
\end{equation} whenever the amplitude is stable. Unmarked \(\Phi_g\) is used only
for \(g\ge2\). The unstable genus-zero derivatives are separate initial
terms and are not substituted into the positive recursion.

\subsection{The positive recursion}\label{deriving-the-exact-positive-recursion}

We single out the first marking and call it the \emph{pivot}. Write
\(\langle\!\langle\tau_{d_1}\cdots\tau_{d_q}\rangle\!\rangle_g
=\prod_i(2d_i+1)!!\,\langle\tau_{d_1}\cdots\tau_{d_q}\rangle_g\).
In this normalization the Dijkgraaf--Verlinde--Verlinde equation for the
pivot reads, with \((-1)!!=1\),
\begin{equation*}
\begin{split}
\langle\!\langle\tau_{d_1}\tau_{d_2}\cdots\tau_{d_q}\rangle\!\rangle_g
={}&\sum_{j=2}^q(2d_j+1)\,
 \langle\!\langle\tau_{d_1+d_j-1}\textstyle\prod_{i\ne1,j}\tau_{d_i}\rangle\!\rangle_g
 +\frac12\sum_{u+v=d_1-2}
 \langle\!\langle\tau_u\tau_v\tau_{d_2}\cdots\tau_{d_q}\rangle\!\rangle_{g-1}\\
&+\frac12\sum_{u+v=d_1-2}\ \sum_{\substack{g_1+g_2=g\\ I\sqcup J=\{2,\ldots,q\}}}
 \langle\!\langle\tau_u\tau_I\rangle\!\rangle_{g_1}
 \langle\!\langle\tau_v\tau_J\rangle\!\rangle_{g_2}.
\end{split}
\end{equation*}
For \(d_1=0\) and \(d_1=1\) this is the string and the dilaton equation.
In the amplitudes \eqref{eq:H6}, besides the \(n\) external markings, the
intersection numbers contain \emph{background} insertions: \(\tau_k\) with
weight \(t_k=-X^{(k)}(Y)\) for \(k\ge2\), and, after the dilaton sum that
produced the factor \(a^{-h-\sum\ell_k}\), also \(\tau_1\) with weight
\(t_1=1-a\). Apply the displayed equation with the pivot \(\tau_{d_1}\)
and sum over the background. Three kinds of terms appear.

\begin{enumerate}[label=(\alph*),leftmargin=*]
\item \emph{The pivot merges with an external marking $j\ge2$.} By the
first term above, this gives \((2d_j+1)\,U^Y_{g,n-1}(d_1+d_j-1,\ldots)\),
with the background unchanged.
\item \emph{The pivot merges with a background insertion $\tau_k$,
$k\ge1$.} Background insertions carry no double factorial, so the
factor \(2k+1\) of the first term is divided by \((2k+1)!!\) and becomes
\(1/(2k-1)!!\). Choosing which of the
\(\ell_k\) copies is used turns \(t_k^{\ell_k}/\ell_k!\) into
\(t_k\cdot t_k^{\ell_k-1}/(\ell_k-1)!\). The result is
\[
\frac{t_k}{(2k-1)!!}\,U^Y_{g,n}(d_1+k-1,d_2,\ldots,d_n).
\]
\item \emph{Contractions.} The two remaining terms give the nonseparating
and separating contributions with coefficient \(1/2\). In the separating
case the background splits between the two components with the binomial
factors of exponential generating functions, so the term is a product of
two amplitudes.
\end{enumerate}

Let \(E\) be the operator raising the pivot degree by one. The terms of
type~(b) form \((1-a)U^Y+\sum_{k\ge2}t_kE^{k-1}U^Y/(2k-1)!!\). Moving them
to the left-hand side and recalling \(t_k=-X^{(k)}(Y)\) gives
\[
\left(a+\sum_{k\ge2}\frac{X^{(k)}(Y)}{(2k-1)!!}E^{k-1}\right)U^Y
=K_Y(E)\,U^Y=\bigl(\text{terms of types (a) and (c)}\bigr),
\]
where we compared with the series~\eqref{eq:P1} for \(K_Y\). For a fixed
parent, raising the pivot degree more than \(m\) times leaves the
dimension range, so \(E^{m+1}U^Y=0\) and \(K_Y(E)\) has the finite inverse
\(\sum_{\ell\le m}[z^\ell]K_Y(z)^{-1}E^\ell\). This is how the Taylor
coefficients of \(1/K_Y\) enter. Components of the wrong dimension vanish,
and for external complexity \(h\ge2\) no exceptional initial tuple occurs.

After \eqref{eq:H7}, the exact recursion is \begin{equation}
\begin{aligned}
\mathcal U_{g,n}(\mathbf d)={}&
\sum_{j=2}^n(2d_j+1)\sum_{\ell=0}^m\theta_\ell
 \mathcal U_{g,n-1}(d_1+d_j+\ell-1,\mathbf d_{\widehat{1,j}})\\
&+\frac12\sum_{\ell=0}^m\theta_\ell
 \sum_{u+v=d_1+\ell-2}
 \mathcal U_{g-1,n+1}(u,v,\mathbf d_{\widehat1})\\
&+\frac12\sum_{\ell=0}^m\theta_\ell
 \sum_{u+v=d_1+\ell-2}
 \sum_{\substack{g_1+g_2=g\\I\sqcup J=\{2,\ldots,n\}}}
 \mathcal U_{g_1,|I|+1}(u,\mathbf d_I)
 \mathcal U_{g_2,|J|+1}(v,\mathbf d_J).
\end{aligned}\label{eq:H9}
\end{equation} Here $\mathbf d_{\widehat{1,j}}$ denotes the tuple $\mathbf d$ with the entries $d_1$ and $d_j$ removed, and $\mathbf d_{\widehat1}$ the tuple with $d_1$ removed. The indices $u$ and $v$ are the degrees at the two new markings corresponding to the branches of the node. The last split is ordered. We call the amplitudes on the right of \eqref{eq:H9} the \emph{children} of the amplitude on the left: a merging or nonseparating child has complexity $h-1$, and the two factors of a separating term have complexities summing to $h-1$. Its initial values, computed directly from
\eqref{eq:H6}, are \begin{equation}
\mathcal U_{0,3}(0,0,0)=1,\qquad
\mathcal U_{1,1}(1)=\frac18,\qquad
\mathcal U_{1,1}(0)=\frac{\theta_1}{8}.
\label{eq:H10}
\end{equation} Equations \eqref{eq:H9}--\eqref{eq:H10}, with negative degrees and unstable factors set
to zero, determine all amplitudes by increasing \(h\).

By \eqref{eq:P11}, \begin{equation}
1=\theta_0\le\theta_1\le\cdots\le C_\theta
\label{eq:H11}
\end{equation} uniformly in \(Y\in[0,y_0]\). Thus every term of \eqref{eq:H9} is nonnegative,
and the continuous endpoint \eqref{eq:H3} defines a positive limiting theory.

\subsection{The harmonic model}\label{harmonic-model-and-the-needed-uniform-coefficient-estimates}

We now introduce the explicit model with which the amplitudes will be
compared. It depends on the degrees only through the deficit \(m\), and it
is chosen so that the linear part of~\eqref{eq:H9} reproduces it exactly
on average; this is the mass identity~\eqref{eq:H13} below, and the
reason for the name: the model is harmonic for the linear part of the
recursion. At the
endpoint \(Y=y_0\) the model is the denominator in the normalization
\eqref{eq:H4} of pure \(\psi\) intersection numbers (see~\eqref{eq:PF13}),
so the comparison theorem extends the coarse bounds on pure \(\psi\)
intersections to the whole parameter interval.

Suppress \(Y\) temporarily and put \begin{equation}
A_m(h)=[z^m]\bar H(z)^{-h},\qquad
\mathcal M_{g,n}(\mathbf d)=3^h\Gamma(h)4^{-g}A_m(h),\qquad
R=\mathcal U/\mathcal M.
\label{eq:H12}
\end{equation} All \(A_m(h)\) are positive. Define \[
q_\ell=\theta_\ell\,\frac{A_{m-\ell}(h-1)}{A_m(h)},
\qquad S_0=\sum_{\ell=0}^m q_\ell,\qquad
L=2D+n=3h-2m.
\] For \(h\ge2\), the identity \(\bar K=\bar H+(2/3)z\bar H'\) gives exactly \begin{equation}
\sum_{\ell=0}^m q_\ell
\left(1-\frac{2(m-\ell)}{3(h-1)}\right)=1.
\label{eq:H13}
\end{equation} Indeed, the generating series of the left side, before division by
\(A_m(h)\), is \[
\bar K^{-1}\left(\bar H^{-(h-1)}
-\frac{2z}{3(h-1)}(\bar H^{-(h-1)})'\right)
=\bar K^{-1}\bar K\bar H^{-h}=\bar H^{-h}.
\]

Dividing \eqref{eq:H9} by the model turns its linear terms into a
weighted sum of the ratios \(R\) at the children. We check that these
weights have total mass one when the pivot degree \(d_1\) is positive;
this is the meaning of \eqref{eq:H13}. Assume \(g\ge2\); states of genus at
most one are terminal in the argument of Section~\ref{uniform-comparison}.
A merging child has type \((g,n-1)\) and a nonseparating child type
\((g-1,n+1)\); both have complexity \(h-1\) and deficit \(m-\ell\). Hence the
model ratio of a merging child to its parent is
\(A_{m-\ell}(h-1)/(3(h-1)A_m(h))\), and that of a nonseparating child is four
times larger, because of the factor \(4^{-g}\) in \eqref{eq:H12}. The
merging children carry the coefficients \((2d_j+1)\theta_\ell\), \(j\ge2\),
and there are \(d_1+\ell-1\) ordered nonseparating pairs, each with
coefficient \(\theta_\ell/2\). Since \(2D+n=3h-2m\), the total weight is
\[
 \sum_{\ell=0}^m q_\ell\,
 \frac{\sum_{j\ge2}(2d_j+1)+2(d_1+\ell-1)}{3(h-1)}
 =\sum_{\ell=0}^m q_\ell\,\frac{3(h-1)-2(m-\ell)}{3(h-1)}=1
\]
by \eqref{eq:H13}. Thus, for a positive pivot, the linear model weights form
a probability kernel \(P\), and positivity of the separating terms gives
\(R\ge PR\).

At an all-zero state, where every degree and hence the pivot vanishes,
the formal \(\ell=0\) count includes one
nonseparating pair with multiplicity \(-1\), and also includes \(n-1\)
invalid mergers. The true linear mass is instead \begin{equation}
s=1+\frac{3-n}{3(h-1)}
       \frac{A_m(h-1)}{A_m(h)}.
\label{eq:H14}
\end{equation} Since \(\bar H(0)=1\) and \(\bar H^{-1}\) has positive coefficients,
\(A_m(h-1)\le A_m(h)\). In the range \(g\ge2\), therefore, \begin{equation}
s\ge\frac23,\qquad
s\le1\ (n\ge3),\qquad s\le1+O(h^{-1})\ (n=1,2).
\label{eq:H15}
\end{equation} Thus for a zero pivot the linear terms need not have mass exactly one.

We now record the analytic information needed below. The function
\begin{equation}
P_Y(w)=\frac{\bar H_Y(1)}{\bar H_Y(w)}
\label{eq:H16}
\end{equation} is the probability generating function of
\eqref{eq:P12}, recalled here. By
Subsection~\ref{compact-normalization-and-uniform-lattice-estimates}, it has
mean \(3/2\) and, uniformly in the compact parameter interval, it is
analytic on a disk of radius \(1+\eta>1\); its constant and linear
coefficients are bounded below by positive constants; and its variance
is bounded above and away from zero. Let \(S_h\) be the sum of \(h\)
independent variables with probability generating function
\eqref{eq:H16}. The uniform local limit theorem \eqref{eq:P13}, proved
there by Fourier inversion, gives \begin{equation}
\Pr(S_h=m)\ge c_Bh^{-1/2}
\quad\text{when }|m-3h/2|\le B\sqrt h,
\label{eq:H17}
\end{equation} and at bounded displacement, where the odd cubic term
integrates to zero, \begin{equation}
\Pr(S_h=3h/2+d)
=\frac{\mathrm{e}^{-d^2/(2h\sigma_Y^2)}}{\sqrt{2\pi h\sigma_Y^2}}
  (1+O_B(h^{-1})),\qquad |d|\le B,
\label{eq:H18}
\end{equation} at integer arguments. The same Fourier argument applies after tilting at a
positive radius \(\rho\) in a compact subset of the analytic interval; here $\rho$ is a radius in the normalized $z$-coordinate. If
\(\mu_Y(\rho)=-\rho\bar H_Y'(\rho)/\bar H_Y(\rho)=m/h\), it gives \begin{equation}
A_m(h)=
\frac{\bar H_Y(\rho)^{-h}\rho^{-m}}{\sqrt{2\pi h\,\rho\mu_Y'(\rho)}}
 (1+O(h^{-1}))
\label{eq:H19}
\end{equation} uniformly on such compact tilted families. Near the mean, the tilted
rate function has bounded second derivative and vanishing value and
first derivative; hence, for a fixed sufficiently small \(c_0>0\), \begin{equation}
\Pr(S_h=3h/2-u)\ge ch^{-1/2}\mathrm{e}^{-Cu^2/h},
\qquad 0\le u\le c_0h.
\label{eq:H20}
\end{equation}

The coefficient/probability relation is \begin{equation}
A_m(h)=\bar H(1)^{-h}\Pr(S_h=m).
\label{eq:H21}
\end{equation} Consequently, when \(L\le\sqrt h\), \eqref{eq:H11}, \eqref{eq:H17}, and the sum of the
numerator probabilities give \(S_0\le C\sqrt h\). When \(L\ge6\), the
minimum factor in \eqref{eq:H13} is \((L-3)/(3(h-1))\), so \(S_0\le Ch/L\).
Combining the two estimates, \begin{equation}
S_0\le C\sqrt h,\qquad S_0\le Ch/L\quad(L\ge6).
\label{eq:H22}
\end{equation} All constants here are uniform in \(Y\), and the finitely many states
of bounded complexity are absorbed into them.

\section{Comparison with the harmonic model}\label{uniform-comparison}

\begin{proposition}[Uniform harmonic comparison]\label{prop:harmonic}
For every fixed \(n\ge1\), there is a constant \(C_n\) such that
\[
 -C_n\log^2(h+1)\le
 \log\frac{\mathcal U_{g,n}(\mathbf d;Y)}
 {3^h\Gamma(h)4^{-g}[z^m]\bar H_Y(z)^{-h}}
 \le C_n\sqrt h
\]
for every stable pair \((g,n)\), every \(\mathbf d\in\mathbb Z_{\ge0}^n\)
with \(|\mathbf d|\le3g-3+n\), and every \(Y\in[0,y_0]\). Here
\(h=2g-2+n\) and \(m=3g-3+n-|\mathbf d|\).
\end{proposition}
The idea is the following. Write \(R=\mathcal U/\mathcal M\) for the
ratio of a true amplitude to the model amplitude~\eqref{eq:H12}, and
divide the recursion~\eqref{eq:H9} by the model. Because the model
satisfies the linear part of the recursion with total mass one
(identity~\eqref{eq:H13}), the linear terms express \(R\) at a state as an
average of \(R\) over the states of complexity \(h-1\) reached in one step,
and the separating terms add a nonnegative amount. Thus \(R\) is bounded
below by its expectation along a random walk that lowers the complexity,
and it is bounded above as soon as the separating contribution is small.

The comparison of intersection recursions with random walks has an
important precedent in Aggarwal's proof of uniform asymptotics for pure
$\psi$ intersections~\cite{Agg}. Here the recursion and its comparison
model vary with $Y$; the estimates must also remain uniform as this
parameter approaches the pure-$\psi$ endpoint.

Throughout this section a \emph{state} is a stable pair \((g,n)\), \(n\ge1\),
with a degree tuple \(\mathbf d\in\mathbb Z_{\ge0}^n\), \(|\mathbf d|\le3g-3+n\),
and a parameter \(Y\in[0,y_0]\); constants are uniform in \(Y\) and in the
state unless stated otherwise. We write \(Z_0\) for the number of entries
of \(\mathbf d\) equal to zero. A state is \emph{terminal} if \(m=0\) or
\(g\le1\); the random walks below stop there. The proof has three steps.
\begin{itemize}[leftmargin=*]
\item Lemma~\ref{lem:terminal} bounds \(R\) at terminal states, using only
the string and dilaton equations and the coarse bounds of
Appendix~\ref{app:primitive-bounds}.
\item Lemma~\ref{lem:lower} proves the lower bound. The only defect of the
averaging is at entries of degree zero, where some linear terms are
missing; we show that the random walk does not create many such entries.
\item Lemmas~\ref{lem:weighted-linear}--\ref{lem:upper} prove the upper
bound. We divide \(R\) by a weight that penalizes entries of degree zero
and one, show that the weighted separating terms have relative size
\(O(h^{-1/2})\), and conclude by examining, at each complexity, the state
at which the weighted ratio is largest.
\end{itemize}

\subsection{Bounds at terminal states}\label{elementary-terminal-bounds-and-a-global-crude-estimate}

\begin{lemma}[Terminal bounds]\label{lem:terminal}
There are constants \(c_\psi,C_\psi,c,C>0\) with the following properties.
\begin{enumerate}[label=\textup{(\roman*)},leftmargin=*]
\item For every pure \(\psi\) tuple \(\mathbf e\) with \(Z_0\) zero entries,
\(c_\psi(2/3)^{Z_0}\le C_g(\mathbf e)\le C_\psi2^{Z_0}\), in the normalization
\eqref{eq:H4}.
\item At every terminal state, \(R\ge c\,\mathrm e^{-CZ_0}\).
\item At every state, \(R\le\mathrm e^{Ch}\).
\end{enumerate}
\end{lemma}

\begin{proof}
For~(i), the standard normalized string and dilaton identities imply \begin{equation}
c_\psi(2/3)^{Z_0}\le C_g(\mathbf e)\le C_\psi2^{Z_0},
\label{eq:H23}
\end{equation} where \(Z_0\) is the number of zero entries. To prove
this, first remove the entries equal to one by the normalized dilaton equation,
retaining the final genus-one tuple \((1)\). At a noninitial tuple with
a zero, let \(p_+\) denote the number of positive entries; normalized
string has total positive coefficient \[
\frac{2(3g-3+q)+p_+}{3(2g-3+q)}
\] between \(2/3\) and \(2\). Once the ones have been removed, each
string child has one fewer zero; any newly created one can again be
removed. With no zeros use the primitive bounds stated after \eqref{eq:H4}; the
remaining bases are \(C_0(0,0,0)=1/3\) and \(C_1(1)=1/6\). Induction
proves \eqref{eq:H23}, which is~(i).

For~(ii) we use a lower comparison that requires no asymptotic
information on volumes. Set all kernel coefficients in \eqref{eq:H9} equal to one. The
corresponding background of KW times has \(t_1=0,t_2=3,t_k=0\) for \(k\ge3\), so its amplitudes
are \begin{equation}
\mathcal U^{\rm PI}_{g,n}(\mathbf d)
=\frac{3^m}{m!}\prod_i(2d_i+1)!!
\left\langle\tau_{\mathbf d}\tau_2^m\right\rangle_g.
\label{eq:H24}
\end{equation} They have the same initial \((0,3)\) and degree-one \((1,1)\) values,
and the degree-zero \((1,1)\) value \(1/8\). By \eqref{eq:H11}, induction in the
positive recursion gives \(\mathcal U\ge\mathcal U^{\rm PI}\).

Uniform analyticity of \(\bar H_Y^{-1}\) at zero gives a fixed \(B\ge1\) such
that \([z^j]\bar H_Y^{-1}\le B^j\). Thus \[
A_m(h)\le B^m(h)_m/m!.
\] Dividing \eqref{eq:H24} by \eqref{eq:H12} and using \eqref{eq:H4} gives \begin{equation}
R\ge C_g(\mathbf d,2^m)
       \left(\frac3{5B}\right)^m.
\label{eq:H25}
\end{equation} If \(g\le1\), dimension and \(D\ge n-Z_0\) imply \(m\le Z_0\). Hence
\(R\ge c\mathrm{e}^{-CZ_0}\) in this terminal range. At \(m=0\), exactly
\(R=C_g(\mathbf d)\), so \eqref{eq:H23} applies directly. This proves~(ii).

Finally we prove the global bound~(iii), \begin{equation}
R\le \mathrm{e}^{Ch}.
\label{eq:H26}
\end{equation} Note that in the normalized finite sum \eqref{eq:H6}--\eqref{eq:H7} a
background insertion of degree \(k\) has weight \[
v_k(Y)=-r_Y^{k-1}X^{(k)}(Y)/a.
\] The corresponding absolute majorant is \[
B_Y(z)=\sum_{k\ge2}\frac{3|v_k(Y)|}{(2k+1)!!}z^{k-1}.
\] This family has a uniformly bounded analytic germ and vanishes at
zero. Indeed it is the coefficientwise absolute version of \(1-\bar H_Y\);
the normalized family in \eqref{eq:H2}--\eqref{eq:H3} is compact and analytic on a common
disk. Thus a fixed \(z_*>0\) makes \(B_Y(z_*)\le1/2\). Using the upper
bound in \eqref{eq:H23} on every primitive term of \eqref{eq:H6} gives \[
\mathcal U\le
C2^{Z_0}3^h\Gamma(h)4^{-g}[z^m](1-B_Y(z))^{-h}
\le C2^{Z_0}3^h\Gamma(h)4^{-g}z_*^{-m}2^h.
\] The linear coefficient of \(-\log \bar H_Y\) is uniformly positive.
Therefore \[
A_m(h)\ge (ch)^m/m!\ge \mathrm{e}^{-C'h},
\] because \(m\le3h/2\); the case \(m=0\) is immediate. Since
\(Z_0\le n\le h+2\), division proves \eqref{eq:H26}.
\end{proof}

\subsection{The lower bound: control of zero degrees}\label{lower-comparison-by-controlling-only-zeros}

\begin{lemma}[Lower bound]\label{lem:lower}
There are constants \(c,C>0\) such that, at every state,
\[
 R\ge c\exp\{-C(Z_0+\log(h+1))\log(h+1)\}.
\]
\end{lemma}

\begin{proof}
Since the separating terms of \eqref{eq:H9} are nonnegative, \(R\) is at
least the average of \(R\) over its linear children. Iterating, \(R\) is
bounded below by the expected terminal value along a random walk that
lowers the complexity by one at each step. Two losses occur: the terminal
bounds \eqref{eq:H23} and \eqref{eq:H25} lose a factor exponential in the
number of zero entries, and the averaging may have mass \(s<1\) at states
with all degrees zero, again with a loss controlled by the number of zeros. It
therefore suffices to show that the walk rarely creates many zero entries.

At a nonterminal state choose any positive pivot if one exists;
otherwise choose a zero. Stop only at \(m=0\) or \(g\le1\). Use the
linear model weights, divided by their mass \(s\), as transition
probabilities. For a positive pivot \(s=1\); at an all-zero state use
\eqref{eq:H14}--\eqref{eq:H15}.

Let \(Z\) be the current zero count. At a positive pivot, a merger with a marking of degree zero
lowers \(Z\) by one, except when the pivot has degree one and \(\ell=0\);
that merger leaves \(Z\) unchanged. A nonseparating ordered pair can increase \(Z\) by
at most two, and at each \(\ell\) there are at most two pairs containing
a zero. Each pair has twice the model weight of a merger with one zero
marking. Thus \begin{equation}
p_{\rm birth}\le4p_{\rm death}/Z\qquad(Z\ge1).
\label{eq:H27}
\end{equation} That exceptional term contributes neither a
birth nor a death. At an all-zero state, a birth occurs only at
\(\ell=2\), raises the count by one, and its weight is at most
\(2/(Z-1)\) times the mass of mergers with a zero marking. Those mergers have
\(\ell\ge1\) and lower the count by at least one. Hence \eqref{eq:H27} still
holds for \(Z\ge2\). Normalization by \(s\) preserves these comparisons.

For \(Z\ge24\), the conditional drift obeys \[
\frac{\mathbb E(2^{Z_{\rm next}})}{2^Z}-1
\le3p_{\rm birth}-\frac12p_{\rm death}\le0.
\] During each excursion above 24, optional stopping bounds the
probability of reaching \(M\) by \(2^{Z_{\rm start}-M}\). An excursion
entered from below starts at a value at most 26. Since complexity decreases by
one, there are at most \(h_{\rm initial}+1\) excursions. A union bound
therefore shows that, with probability at least \(1/2\), \begin{equation}
\max Z\le Z_{\rm initial}+C+C\log(h_{\rm initial}+1).
\label{eq:H28}
\end{equation} All stopping times are bounded, so optional stopping
applies without further integrability conditions.

Only all-zero steps can have \(s<1\). By \eqref{eq:H14}--\eqref{eq:H15}, \[
-\log\min(s,1)\le CZ/h.
\] On the event \eqref{eq:H28}, write its upper bound as \(M\). Summing over the
decreasing complexities costs at most \[
CM+CM\sum_{j=M+1}^{h_{\rm initial}}\frac1j
\le CM\log(h_{\rm initial}+1).
\] Iterated positivity gives \[
R_{\rm initial}\ge
\mathbb E\!\left[R_{\rm terminal}\prod_{\rm steps}s\right].
\] Here the initial zero count is \(Z_0\). The terminal bounds of
Lemma~\ref{lem:terminal}(ii), applied at terminal zero counts at most
\(M\), and the probability in \eqref{eq:H28} now prove \begin{equation}
R\ge c\exp\{-C(Z_0+\log(h+1))\log(h+1)\}.\label{eq:H29}
\end{equation}
No asymptotics for volumes enter at the terminal states.
\end{proof}

\subsection{The upper bound: a weight for degrees zero and one}\label{upper-comparison-the-clipped-inactive-weight}

For the upper bound the separating terms must be shown to be small, and
for general states they are not: a state with many entries of degree zero
or one can distribute these markings between the two components in many
ways. We therefore divide \(R\) by a weight that grows with the number of
such entries. It is chosen so that the linear terms still average
\(R/\mathcal W\) up to a factor \(1+O(h^{-1/2})\) (this is \eqref{eq:H31}),
while the weighted separating terms become small
(Subsection~\ref{the-weighted-separating-estimate}).

Call degrees zero or one inactive. Let \(\iota\) be their number, \(Z_0\)
the number of zeros, and \(R_{\rm act}\) the number of degrees at least
two. Fix a sufficiently small constant \(c_*>0\), and set \begin{equation}
\phi(\iota)=\prod_{j=1}^\iota\max(1,c_*\sqrt j),\qquad
\mathcal W(\mathbf d)=4^{Z_0}\phi(\iota).
\label{eq:H30}
\end{equation} We call the truncation \(\max(1,\cdot)\) clipping; it makes
\(\phi\) nondecreasing, since every factor is at least one. Up to fixed
multiplicative constants, \(\phi(\iota)\) equals \(c_*^\iota\sqrt{\iota!}\). Its monotonicity is
essential when an inactive pivot splits into two active degrees.

In the upper bound we choose the pivot as follows: an active entry when
one exists, an entry equal to one if no active entry exists, and a zero
only at an all-zero state. Again the walk stops only at terminal states.
The linear model weights \(p_{\rm child}\) of a state are the weights
that appear when \eqref{eq:H9} is divided by the model; their mass is
given by \eqref{eq:H13}--\eqref{eq:H15}, and may exceed one at an
all-zero state with one or two markings.

\begin{lemma}[Weighted linear terms]\label{lem:weighted-linear}
If \(c_*\) is sufficiently small, then at every nonterminal state
\begin{equation}
\sum_{\rm linear}p_{\rm child}
\frac{\mathcal W_{\rm child}}{\mathcal W}
\le1+CS_0/h\le1+Ch^{-1/2}.
\label{eq:H31}
\end{equation}
\end{lemma}

\begin{proof}
The second inequality is \eqref{eq:H22}. We consider the three possible
pivots in turn.

For an active pivot, an ordinary inactive merger reduces \(\iota\) by at
least one and does not increase \(Z_0\). For all sufficiently large
\(\iota\), its weight ratio is at most \(1/2\). The exception is a pivot of degree two
merged with a zero at \(\ell=0\): it turns that zero into a one, leaves
\(\iota\) fixed, and divides the weight by four. The favorable mass is
therefore at least a fixed positive multiple of \(\iota S_0/h\).

At each \(\ell\), at most four nonseparating pairs contain an inactive
degree. A pair creating two inactive degrees has weight ratio, for large
\(\iota\), at most \[
16c_*^2\sqrt{(\iota+1)(\iota+2)}.
\] Its total model mass per \(\ell\) is bounded by a fixed multiple of
\(q_\ell/h\). Choose \(c_*\) so small that the coefficient of \(\iota\) in
this upper bound is less than half the favorable merger coefficient.
Pairs creating only one inactive degree cost \(O(\sqrt \iota)\), and all
other pairs leave the weight unchanged. The remaining expression is
bounded above by \(CS_0/h\). Bounded \(\iota\), including the clipping
threshold, also costs only \(CS_0/h\), because there are at most four
exceptional pairs and boundedly many inactive mergers.

If the pivot has degree one and all markings are inactive, every merger
decreases \(\iota\), while a nonseparating step increases it by at most one.
Only four pairs per \(\ell\) can increase the weight; a pair with two
active new degrees decreases \(\iota\) and cannot increase the clipped
weight. The same argument applies, with \(\iota-1\) favorable merging
labels.

At an all-zero state, all valid mergers have \(\ell\ge1\) and decrease
both the zero count and the inactive weight. Increasing pairs have
\(\ell\ge2\) and increase \(\iota\) by at most one. Use \(S_0-q_0\) in both
birth and merger estimates. The unweighted mass correction \eqref{eq:H14} is
nonpositive for \(n\ge3\) and \(O(h^{-1})\) otherwise. This proves \eqref{eq:H31}
in every case.
\end{proof}

\subsection{Bounding the separating terms}\label{the-weighted-separating-estimate}

We now show that, after weighting, the separating terms of
\eqref{eq:H9} have total relative size \(O(h^{-1/2})\), uniformly in the
state. Splitting off a component of complexity \(h_1\le h_2\) costs the
factor \(\Gamma(h_1)\Gamma(h_2)/\Gamma(h)\), which is at most
\(\binom{h-3}{h_1-1}^{-1}\); the point is that this beats the number of
ways of distributing the marking labels, once the weight has absorbed the
inactive ones. For a separating term with node degrees \(u,v\) and kernel
index \(\ell\), write \(\mathcal M_1,\mathcal M_2\) and
\(\mathcal W_1,\mathcal W_2\) for the model amplitudes and weights of its two
factors.

\begin{lemma}[Weighted separating terms]\label{lem:weighted-separating}
Let \(\eta=1/12\). At every nonterminal state,
\[
 \sum_{\rm separating}\frac{\theta_\ell}{2}\,
 \frac{\mathcal M_1\mathcal M_2}{\mathcal M}\,
 \frac{\mathcal W_1\mathcal W_2}{\mathcal W}\,
 \mathrm e^{\eta\min(h_1,h_2)}
 \le CS_0h^{-3/2}(1+\iota)\le Ch^{-1/2},
\]
where the sum runs over all separating terms of \eqref{eq:H9}.
\end{lemma}

\begin{proof}
The degree-dimension constraint supplies \begin{equation}
4R_{\rm act}+n\le3h,\qquad R_{\rm act}\le3h/5.
\label{eq:H32}
\end{equation} Remove the pivot, leaving \(\iota_*\) inactive and \(R_*\) active labeled
markings, with \(\iota_*+R_*=n-1\), \(\iota_*\le \iota\), and \(R_*\le R_{\rm act}\). In
one separating component put \(u\) inactive and \(j\) active original
labels, and write \(s=u+j\). Each child receives one new marking.
Comparing \eqref{eq:H30} with \(4^{Z_0}c_*^\iota\sqrt{\iota!}\), uniformly over the two
new degrees, \begin{equation}
\frac{\mathcal W_1\mathcal W_2}{\mathcal W}
\le C\frac{\sqrt{(s+2)h}}{\sqrt{\binom{\iota_*}u}}.
\label{eq:H33}
\end{equation} Indeed, when the \(\iota_*\) inactive original labels split as \(u,\iota_*-u\), the factorial part of this estimate is bounded by
\[
 \frac{\sqrt{(u+1)!(\iota_*-u+1)!}}{\sqrt{\iota_*!}}
 =\frac{\sqrt{(u+1)(\iota_*-u+1)}}{\sqrt{\binom{\iota_*}u}}.
\]
Each child receives at most one inactive new marking. Since \(u\le s\) and \(\iota_*\le h+2\), the numerator is at most \(C\sqrt{(s+2)h}\). The powers of \(c_*\), the factor for zeros, and the clipping comparison change by bounded integer exponents. An inactive pivot supplies an additional factor \(\iota^{-1/2}\), which
was discarded. The exponential factors and clipping comparison
contribute only a fixed constant, since the total inactive and zero
counts change by bounded amounts.

Orient \(h_1\le h_2\), and put \[
h'=h-3,\qquad k=h_1-1\le h'/2,\qquad s=k+2-2g_1.
\] Then \begin{equation}
\frac{\Gamma(h_1)\Gamma(h_2)}{\Gamma(h)}
=\frac1{(h-1)(h-2)\binom{h'}k}.
\label{eq:H34}
\end{equation} For all sufficiently large \(h\), \eqref{eq:H32} gives \(R_*\le2h'/3\).

The required label sum is \[
\Lambda_s=\sum_{u+j=s}\binom{\iota_*}u^{1/2}\binom{R_*}j.
\] Cauchy--Schwarz and Vandermonde give \begin{equation}
\Lambda_s\le
\binom{\iota_*+R_*}s^{1/2}
\left(\sum_{j\le s}\binom{R_*}j\right)^{1/2}.
\label{eq:H35}
\end{equation} For \(k\ge1\), take \(\xi =k/(h'-k)\le1\). The two factors before square
roots are bounded by \(\xi ^{-s}(1+\xi )^{\iota_*+R_*}\) and \(\xi ^{-s}(1+\xi )^{R_*}\),
respectively. The binomial mass at its mean gives \[
\binom{h'}k\ge c(k+1)^{-1/2}\xi ^{-k}(1+\xi )^{h'}.
\] Thus, if \(g_1\ge1\), so \(s\le k\), \begin{equation}
\frac{\Lambda_s}{\binom{h'}k}
\le C\sqrt{k+1}\,\xi ^{k-s}(1+\xi )^{-h'/6}
\le C\sqrt{k+1}\mathrm{e}^{-k/6}.
\label{eq:H36}
\end{equation} Here \(\iota_*+R_*=n-1\le h'\), \(R_*\le2h'/3\), and \(h'\log(1+\xi )\ge k\). The
case \(k=0\) is immediate.

If \(g_1=0\), then \(s=k+2\). Validity of this component forces
\(2j\le s-2=u+j-2\), hence \(u\ge j+2\ge2\). Use \[
\binom{\iota_*}u=\frac{\iota_*(\iota_*-1)}{u(u-1)}\binom{\iota_*-2}{u-2}.
\] After taking the square root, discard the denominator and extract a
factor \(\iota_*\). The remaining sum has \(u-2+j=k\), so \eqref{eq:H35}--\eqref{eq:H36} apply
with \(\iota_*-2\) in place of \(\iota_*\). This use of the degree constraint is
what replaces a factorial penalty on active labels.

Combining these estimates with \eqref{eq:H33}--\eqref{eq:H34}, and summing the
polynomially many genus choices at fixed \(k\), proves, for example with
\(\eta=1/12\), \begin{equation}
\sum_{\rm genus/label\ splits}
\frac{\Gamma(h_1)\Gamma(h_2)}{\Gamma(h)}
\sup_{\rm valid\ node\ degrees}
\frac{\mathcal W_1\mathcal W_2}{\mathcal W}
\mathrm{e}^{\eta\min(h_1,h_2)}
\le C h^{-3/2}(1+\iota).\label{eq:H37}
\end{equation} An empty supremum contributes zero. Thus the bound is uniform in the
possible node degrees; their sum, with its model coefficients, is
handled next.

For a separating recursion term, exact model normalization gives \begin{equation}
\frac{\theta_\ell}{2}\frac{\mathcal M_1\mathcal M_2}{\mathcal M}
=\frac{\theta_\ell}{6}
\frac{\Gamma(h_1)\Gamma(h_2)}{\Gamma(h)}
\frac{A_{m_1}(h_1)A_{m_2}(h_2)}{A_m(h)}.
\label{eq:H38}
\end{equation} Since \(h_1+h_2=h-1\) and \(m_1+m_2=m-\ell\), positivity and
coefficient convolution bound the sum over node degrees by
\(A_{m-\ell}(h-1)\). More explicitly, for each genus/label partition
first apply the node-independent envelope \eqref{eq:H33}, then extend the allowed
nonnegative model-coefficient sum to the full convolution. The resulting
label envelope is exactly the one summed in \eqref{eq:H35}--\eqref{eq:H37}, so variable
degree supports cause no interchange issue. Consequently the entire
weighted separating mass, including the exponential factor
\(\mathrm e^{\eta\min(h_1,h_2)}\) of \eqref{eq:H37}, is at most \begin{equation}
CS_0h^{-3/2}(1+\iota)=O(h^{-1/2}).
\label{eq:H39}
\end{equation} For \(L\ge6\), use \(S_0\le Ch/L\) and \(\iota\le n\le L\). For
\(L\le5\), both \(\iota\) and \(L\) are bounded, and \(S_0\le C\sqrt h\)
gives the smaller \(O(h^{-1})\) bound.
\end{proof}

\subsection{The upper bound from successive maxima}\label{terminal-upper-bounds-and-the-two-record-arguments}

\begin{lemma}[Upper bound]\label{lem:upper}
There is a constant \(C\) such that \(R\le C\mathcal W\mathrm e^{C\sqrt h}\)
at every state.
\end{lemma}

\begin{proof}
By Lemmas~\ref{lem:weighted-linear} and~\ref{lem:weighted-separating}, the weighted ratio satisfies an
approximate maximum principle: at each nonterminal state, \(R/\mathcal W\)
is at most \(1+O(h^{-1/2})\) times the largest value at its linear
children, plus a separating contribution that is small. We apply this at
complexities where the maximum of \(R/\mathcal W\) reaches a record, first
to exclude exponential growth \eqref{eq:H40}, and then to obtain the bound
\(\mathrm e^{C\sqrt h}\) in \eqref{eq:H41}.

At \(m=0\), \eqref{eq:H23} and \eqref{eq:H30} give \(R/\mathcal W\le C\). At \(g\le1\),
one has \(D\le n\), so \(R_{\rm act}\le n/2\) and \(\iota\ge n/2\). Thus
\(\log\mathcal W\ge(n/4)\log n-O(n)\), which dominates \eqref{eq:H26}, since
\(h=n+O(1)\). Hence \(R/\mathcal W\) is bounded at every terminal state.

Let \[
M_h=\max_{\substack{Y\in[0,y_0]\\\text{valid states of complexity }h}}
\frac{R}{\mathcal W}.
\] The maximum exists: the tuple set is finite and the normalized
parameter interval is compact. Equation \eqref{eq:H26} gives
\(\gamma=\limsup h^{-1}\log M_h<\infty\). Suppose \(\gamma>0\), and
choose \(0<\gamma'<\gamma\), \(\delta>0\), so close that
\(\gamma+\delta-\gamma'<\eta\) from \eqref{eq:H37}. Call \(h\) a \emph{record} for a sequence \(a_h\) if \(a_h>a_j\) for all \(j<h\). There are arbitrarily large
records of \(M_h\mathrm{e}^{-\gamma'h}\), and their maximizing states are
nonterminal. At such a record, \[
M_j/M_h\le \mathrm{e}^{-\gamma'(h-j)}\quad(j<h),\qquad
M_j\le C_\delta \mathrm{e}^{(\gamma+\delta)j}.
\] The linear part, after division by the maximizing value, is bounded
by \((1+Ch^{-1/2})\mathrm{e}^{-\gamma'}\). For a separating term use the record
inequality on the larger child and the global bound on the smaller.
Their product divided by \(M_h\) is at most a fixed constant times
\(\mathrm{e}^{(\gamma+\delta-\gamma')h_{\rm small}}\). Equation \eqref{eq:H39}, which includes the factor
\(\mathrm e^{\eta\min(h_1,h_2)}\), makes the total separating contribution \(o(1)\). The recursion
would therefore give \(1\le \mathrm{e}^{-\gamma'}+o(1)\), a contradiction. Thus
\begin{equation}
\limsup h^{-1}\log M_h\le0.
\label{eq:H40}
\end{equation}

To obtain the quantitative envelope, fix \(0<\delta<\eta\), so \eqref{eq:H40}
gives \(M_j\le C_\delta \mathrm{e}^{\delta j}\). Consider record maxima of
\(M_h/\mathrm{e}^{A\sqrt h}\). If these were unbounded, their large maximizing
states would again be nonterminal. For the larger separating child, the
record bound supplies
\(M_{\rm large}/M_h\le \mathrm{e}^{A(\sqrt{h_{\rm large}}-\sqrt h)}\le1\); bound
the smaller child by \(C_\delta \mathrm{e}^{\delta h_{\rm small}}\). The
separating contribution is therefore \(C'h^{-1/2}\), with \(C'\)
independent of \(A\). The linear record ratio is \[
\mathrm{e}^{A(\sqrt{h-1}-\sqrt h)}
=1-\frac{A}{2\sqrt h}+O_A(h^{-1}).
\] Using \eqref{eq:H31}, the recursion would give \[
1\le1+\frac{C+C'-A/2}{\sqrt h}+O_A(h^{-1}).
\] Choosing \(A>2(C+C'+1)\) excludes unbounded records. We have proved
\begin{equation}
R/\mathcal W\le C \mathrm{e}^{C\sqrt h},
\label{eq:H41}
\end{equation} uniformly in the entire normalized parameter family.
\end{proof}

\subsection{Completion of the comparison}\label{common-comparison-theorem-and-the-original-coefficients}

\begin{proof}[Completion of the proof of Proposition~\ref{prop:harmonic}]
For fixed \(n\), the zero count and the inactive count are at most \(n\), so
\(\mathcal W\le4^n\phi(n)\) is bounded. Lemmas~\ref{lem:lower}
and~\ref{lem:upper} therefore give, for every fixed \(n\ge1\), \begin{equation}
-C_n\log^2(h+1)\le
\log\frac{\mathcal U_{g,n}(\mathbf d;Y)}
 {3^h\Gamma(h)4^{-g}[z^m]\bar H_Y(z)^{-h}}
\le C_n\sqrt h,\label{eq:H42}
\end{equation} uniformly for \(Y\in[0,y_0]\) and every degree tuple
in the range of Proposition~\ref{prop:harmonic}. This proves the proposition.
\end{proof}

In particular the uniform error is \(o(g)\).

We record the case \(Y=0\), which compares the Weil--Petersson brackets
themselves with an explicit model; it is the input for the coefficient
ratios in Section~\ref{uniform-one-boundary-coefficient-ratios}. At \(Y=0\), \(a=1\), \(r_0=\pi^2/2\), and \[
H_0(z)=\frac{3(\sin\sqrt{2z}-\sqrt{2z}\cos\sqrt{2z})}
 {(\sqrt{2z})^3}.
\] By \eqref{eq:H5} and \eqref{eq:H7}, \(\mathcal U=b/4^N\), where
\(b=[\tau_{d_1}\cdots\tau_{d_n}]_{g,n}\). Undoing this same scaling in
the model yields \begin{equation}
M^{\rm WP}_{g,n}(\mathbf d)
=4^D3^h\Gamma(h)4^{-g}(2\pi^2)^m
 [z^m]H_0(z)^{-h}.
\label{eq:H43}
\end{equation} Thus \eqref{eq:H42} applies unchanged to \(b/M^{\rm WP}\). For \(n=1\),
\(h=2g-1\), \(m=3g-2-k\), it gives exactly \begin{equation}
\mathrm{e}^{-C\log^2(g+1)}
\le
\frac{[\tau_k]_{g,1}}
 {4^k3^{2g-1}\Gamma(2g-1)4^{-g}(2\pi^2)^{3g-2-k}
 [z^{3g-2-k}]H_0(z)^{-(2g-1)}}
\le \mathrm{e}^{C\sqrt g},
\label{eq:H44}
\end{equation} uniformly over \(0\le k\le3g-2\).

\section{Ratios of marked generating functions}\label{ratios-with-a-fixed-number-of-distinguished-markings}

The comparison theorem controls logarithms, whereas the volume-ratio
problem requires relative errors tending to zero. We obtain that stronger
control here for generating functions with a fixed number of distinguished
markings. The two ratios correspond to adding a cusp and to replacing one
unit of genus by two cusps. Their errors are $O(1/g)$, uniformly in the
parameter. The passage to volumes with a prescribed cusp count will be
made in Section~\ref{extracting-fixed-cusp-counts}.

The mechanism is simple. Positivity of the recursion implies that an
amplitude decreases when one of its degrees is raised~\eqref{eq:TL15}.
The dilaton equation~\eqref{eq:TL12} writes the cusp ratio as a weighted
sum, over degrees \(d\ge1\), of the normalized amplitudes
\(c_{g,p+1}(d,0^p)\). These weights can have either sign; their absolute
moments are uniformly bounded by~\eqref{eq:TL13}. The KdV equation~\eqref{eq:TL14} writes the
genus ratio through \(c_{g,p}(1,0^{p-1})\) and separating terms. Both
ratios are therefore within \(O(1/g)\) of explicit functions of \(Y\) once
we know that \(c_{g,p}(\mathbf d)=1-O(|\mathbf d|^2/g)\). This last estimate,
\eqref{eq:TL23}, follows from the difference of the recursion at adjacent
degrees; the coarse bounds of Section~\ref{uniform-comparison} are used
only to show that the separating terms are small.

\subsection{Shifted amplitudes and normalization}\label{shifted-amplitudes-and-normalization}

Use \(p\) for the number of distinguished markings, and keep \(n\) for the number of cusps. Retain the kernels and shifted amplitudes of Sections~\ref{the-positive-compact-family-of-kernels} and~\ref{exact-amplitudes-and-the-harmonic-model}; write \(U_{g,p}(\mathbf d;Y)=U^Y_{g,p}(\mathbf d)\), and suppress \(Y\) when convenient. Put \(h=2g-2+p\) and \(m=3g-3+p-|\mathbf d|\). Define \begin{equation}
\overline U_{g,p}(\mathbf d)=a^h r^m U_{g,p}(\mathbf d),
 \qquad W_{g,p}=4^g\overline U_{g,p}(\mathbf0),
 \qquad c_{g,p}(\mathbf d)=
 \frac{\overline U_{g,p}(\mathbf d)}{\overline U_{g,p}(\mathbf0)}.
 \label{eq:TL9}
\end{equation} Thus $\overline U_{g,p}(\mathbf d)$ is the normalized amplitude $\mathcal U_{g,p}(\mathbf d;Y)$ of \eqref{eq:H7}; the bar indicates that the markings are now the distinguished ones. The barred amplitudes obey the exact positive recursion \eqref{eq:H9}.

We use two consequences of
Sections~\ref{the-positive-compact-family-of-kernels}--\ref{uniform-comparison}.

\begin{itemize}
\item
  The kernel bounds \eqref{eq:P11}: for each fixed \(k\), \begin{equation}
1=\theta_0\le\theta_1\le\cdots,
   \qquad w_\ell:=\theta_\ell-\theta_{\ell-1}\ge0,
   \quad\theta_{-1}=0,
   \quad\sup_Y\sum_{\ell\ge0}(1+\ell)^k w_\ell<\infty
   \label{eq:TL10}
\end{equation} Also \(A(Y)\) is bounded above and away from
  zero; see Subsection~\ref{compact-normalization-and-uniform-lattice-estimates}.
\item
  Uniform harmonic comparison: for each fixed positive \(p\), \begin{equation}
\log W_{g,p}=
   \log\Gamma(2g-2+p)-(2g-2+p)\log A(Y)+o(g),
   \label{eq:TL11}
\end{equation} uniformly for \(0\le Y<y_0\). The corresponding crude bound
  \(W_{g,p}\le C_p^{g+1}(2g+p)!\) holds uniformly, including bounded
  stable genera. Proposition~\ref{prop:harmonic}, together with the
  central local limit \eqref{eq:H17} for the model, proves
  \eqref{eq:TL11} without any input on volume ratios. Only its \(o(g)\)
  error is needed below, and only for a positive number of markings.
\end{itemize}

\subsection{Dilaton and KdV identities}\label{exact-dilaton-and-genus-exchange-identities}

The shifted weighted dilaton identity is \begin{equation}
h\overline U_{g,p}(\mathbf0)=
 \sum_{d\ge1}b_d(Y)\overline U_{g,p+1}(d,0^p),
 \qquad
 b_d(Y)=\frac{X^{(d)}(Y)r^{d-1}}{a(2d+1)!!}
 =\frac13[z^{d-1}]\bar H_Y(z).
 \label{eq:TL12}
\end{equation} It holds also at \(p=0\), \(g\ge2\). To track the shift, the
background insertion of degree one contributes \(1-a\); moving it and
the higher background insertions to the left leaves the coefficients
\(X^{(d)}/(2d+1)!!\). The factor in \eqref{eq:TL9} then produces exactly
\(r^{d-1}/a\) in \eqref{eq:TL12}. The compact analytic family gives \begin{equation}
\sum_{d\ge1}b_d=A(Y),\qquad
 \sup_Y\sum_{d\ge1}(1+d)^k|b_d(Y)|<\infty
 \label{eq:TL13}
\end{equation} for every fixed \(k\).

The genus ratio is governed instead by the KdV equation. Write $F$ for the Kontsevich--Witten free energy, with subscripts indicating differentiation with respect to the corresponding times. The integrated KdV identity \(F_{01}=F_{00}^2/2+F_{0000}/12\),
differentiated \(p-2\) times in the string variable, gives for \(p\ge2\)
\begin{equation}
\overline U_{g,p}(1,0^{p-1})
 =\frac14\overline U_{g-1,p+2}(\mathbf0)
 +\frac32\sum_{j=0}^g\sum_{s=0}^{p-2}\binom{p-2}s
 \overline U_{j,s+2}(\mathbf0)
 \overline U_{g-j,p-s}(\mathbf0).
 \label{eq:TL14}
\end{equation} The split is ordered. The factor three from the weighted degree-one
insertion turns the original KdV coefficients into \(1/4\) and \(3/2\).
In \eqref{eq:TL14}, every term has the same total powers of \(a\) and \(r\): the
two separating complexities sum to \(h\), and their deficits sum to
\(3g-4+p\). Thus no omitted parameter factor is present.

The key consequence of positivity, used throughout this section, is the
following monotonicity.

\begin{lemma}[Monotonicity in the degrees]\label{lem:monotone}
For every stable \((g,p)\), every degree tuple \(\mathbf d\), every \(i\) and
every \(Y\in[0,y_0]\),
\begin{equation}
\overline U_{g,p}(\mathbf d)\ge
 \overline U_{g,p}(\mathbf d+\mathbf e_i),
 \qquad 0\le c_{g,p}(\mathbf d)\le1.
 \label{eq:TL15}
\end{equation}
Amplitudes outside the dimension range are zero.
\end{lemma}

\begin{proof}
We subtract the recursion at adjacent degrees. With
$\mathbf d'=(d_2,\ldots,d_p)$ fixed, let
$\mathcal B_{\mathbf d'}(t)$ denote the merging and contraction sums
in \eqref{eq:H9}, with the kernel weight removed and $d_1+\ell=t$.
These sums are nonnegative and vanish for
$t>3g-3+p-|\mathbf d'|$ by dimension. If
$m=3g-3+p-|\mathbf d|\ge0$, shifting the kernel index gives
\[
 \begin{aligned}
 \overline U_{g,p}(\mathbf d)
   &=\sum_{\ell=0}^m\theta_\ell
                      \mathcal B_{\mathbf d'}(d_1+\ell),\\
 \overline U_{g,p}(\mathbf d+\mathbf e_1)
   &=\sum_{\ell=1}^m\theta_{\ell-1}
                      \mathcal B_{\mathbf d'}(d_1+\ell).
 \end{aligned}
\]
Their difference is therefore
$\sum_{\ell=0}^m w_\ell\mathcal B_{\mathbf d'}(d_1+\ell)\ge0$, which proves
\eqref{eq:TL15} for \(i=1\), and by symmetry for every \(i\). This includes
the dimension boundary by the zero convention. The difference identity is
used for $h\ge2$; the finite bases satisfy the inequality directly because
\(\theta_1\ge\theta_0=1\).
\end{proof}

From \eqref{eq:TL12}, \eqref{eq:TL13}, and \eqref{eq:TL15}, \begin{equation}
\frac{W_{g,p}}{W_{g,p+1}}\le\frac C{2g-2+p}
 \label{eq:TL16}
\end{equation} in all stable cases used below, with a constant independent of \(p\): the absolute sum in \eqref{eq:TL13} is uniform and the ratios in \eqref{eq:TL15} are at most one. Finite stable bases are included in the constant. Multiplication of \eqref{eq:TL14} by \(4^g\) and use of \eqref{eq:TL15} give \begin{equation}
W_{g-1,p+2}\le W_{g,p},\qquad p\ge2.
 \label{eq:TL17}
\end{equation} The factor \(4^g\) in \eqref{eq:TL9} is essential: it converts the coefficient
\(1/4\) into the comparison constant one.

\subsection{Separating terms}\label{fixed-marking-splitting-estimates}

The separating terms in the KdV identity \eqref{eq:TL14}, and in the
degree difference above, are controlled by the convolution sums \(S_{g,p}\)
and \(T_{g,p}\) below. They are small for two reasons. Splitting off a
component of bounded genus \(j\) removes \(2j+s\) units of complexity,
and the cusp insertions cost a factor \(O(1/g)\) each, by
\eqref{eq:TL16}--\eqref{eq:TL17}; and a split
into two components of comparable genus is exponentially small by the
binomial entropy in \eqref{eq:TL21}.

For fixed \(p\ge2\), define \[
 S_{g,p}=\frac1{W_{g,p}}
 \sum_{j=0}^g\sum_{s=0}^{p-1}\binom{p-1}s
 W_{j,s+1}W_{g-j,p-s},
\] \begin{equation}
T_{g,p}=\frac1{W_{g,p}}
 \sum_{j=0}^g\sum_{s=0}^{p-2}\binom{p-2}s
 W_{j,s+2}W_{g-j,p-s}.
 \label{eq:TL18}
\end{equation} We claim, uniformly in \(Y\), \begin{equation}
S_{g,p}=O_p(g^{-2}),\qquad T_{g,p}=O_p(g^{-1}),
 \label{eq:TL19}
\end{equation} with the positive-genus portion of \(T\) bounded by \(O_p(g^{-2})\).
In particular \(T_{g,2}=O(g^{-2})\) because its genus-zero terms are
unstable.

By symmetry orient the smaller genus as \(j\le g/2\). For
\(1\le j\le\varepsilon g\), insert \(2j\) cusps into the larger-genus
component using \eqref{eq:TL16}, apply \eqref{eq:TL17} \(j\) times, and insert the
remaining \(s\) cusps. This gives \begin{equation}
\frac{W_{g-j,p-s}}{W_{g,p}}
 \le (C_p/g)^{2j+s}.
 \label{eq:TL20}
\end{equation} If \(p-s=1\), instead insert \(2j+1\) cusps before raising the genus,
so every genus-comparison step ends with at least two markings;
afterward insert \(s-1\) cusps. This is possible because \(p\ge2\), and
the same exponent results. In the \(T\) convolution both relevant
marking counts are already at least two.

The crude bound following \eqref{eq:TL11} makes the smaller-component
contribution a fixed polynomial in \(j\) times \[
 (2j)!(C_p/g^2)^j g^{-s}.
\] Choose \(\varepsilon>0\) small enough that successive factorial terms
have a ratio bounded by a fixed number below one. Their positive-genus
sum is \(O_p(g^{-2})\), also with any fixed polynomial in \(j\). For
genus zero, stability gives \(s\ge2\) in \(S\) and \(s\ge1\) in \(T\);
direct use of \eqref{eq:TL16} gives the respective orders.

For \(\varepsilon g\le j\le g/2\), apply \eqref{eq:TL11} to both factors and the
parent. The marking counts belong to a fixed finite set. The powers of
\(A(Y)\) cancel up to bounded exponents, and Stirling's formula gives \begin{equation}
\log\frac{W_{j,s+c}W_{g-j,p-s}}{W_{g,p}}
 \le-2g\mathsf H(j/g)+o(g),\qquad c\in\{1,2\},
 \label{eq:TL21}
\end{equation} where \(\mathsf H(t)=-t\log t-(1-t)\log(1-t)\). Its positive minimum
on \([\varepsilon,1/2]\) makes the macroscopic sum exponentially small.
This proves \eqref{eq:TL19}. Every use of \eqref{eq:TL11} here has a positive fixed
marking count.

\subsection{Cusp and genus ratios}\label{uniform-fixed-marking-ratios}

For \(p\ge2\), the adjacent-degree difference obtained from \eqref{eq:H9} has a
merging prefactor \(O_p(1/g)\) by \eqref{eq:TL16}. Its nonseparating prefactor is
bounded by \begin{equation}
\frac{W_{g-1,p+1}}{W_{g,p}}
 \le\frac Cg\frac{W_{g-1,p+2}}{W_{g,p}}=O_p(g^{-1})
 \label{eq:TL22}
\end{equation} using \eqref{eq:TL16}--\eqref{eq:TL17}. To see all
three contributions, put $D=|\mathbf d|$,
$\Omega_0=\sum_{\ell\ge0}w_\ell$ and
$\Omega_1=\sum_{\ell\ge0}\ell w_\ell$. Bound each child amplitude
by its zero-degree value in the difference identity above. Since there
are at most $d_1+\ell+1$ new-degree pairs, it yields
\[
 \begin{aligned}
 c_{g,p}(\mathbf d)-c_{g,p}(\mathbf d+\mathbf e_1)
 \le{}&\Omega_0(2D+p-1)\frac{W_{g,p-1}}{W_{g,p}}\\
 &+\bigl((d_1+1)\Omega_0+\Omega_1\bigr)
   \left(2\frac{W_{g-1,p+1}}{W_{g,p}}+\frac12S_{g,p}\right).
 \end{aligned}
\]
The factor $2$ in the nonseparating term includes the factor $4$
from the genus normalization of $W$. The two kernel moments are
uniformly bounded by \eqref{eq:TL10}. Equations~\eqref{eq:TL16},
\eqref{eq:TL19} and~\eqref{eq:TL22} now give \[
 0\le c_{g,p}(\mathbf d)-c_{g,p}(\mathbf d+\mathbf e_i)
 \le\frac{C_p(|\mathbf d|+1)}g.
\] Summing degree increments proves \begin{equation}
0\le1-c_{g,p}(\mathbf d)
 \le C_p\frac{|\mathbf d|(|\mathbf d|+1)}g,
 \qquad p\ge2.
 \label{eq:TL23}
\end{equation} When the right side exceeds one the bound is trivial. It also
bounds the terms outside the dimension range when they are set to zero.

Substitute \eqref{eq:TL23} into \eqref{eq:TL12}. The errors sum absolutely by \eqref{eq:TL13},
giving \begin{equation}
(2g-2+p)\frac{W_{g,p}}{W_{g,p+1}}
 =A(Y)+O_p(g^{-1}),\qquad p\ge1.
 \label{eq:TL24}
\end{equation} Since \(A\) is uniformly positive, the error is relative
\(O_p(1/g)\).

To recover the one-marking profile without circularity, first use \eqref{eq:TL16}
twice, \eqref{eq:TL17} at parent marking count two, and the now two-sided ratio
\eqref{eq:TL24} at marking count one: \begin{equation}
\frac{W_{g-1,2}}{W_{g,1}}
 \le\frac C{g^2}\frac{W_{g-1,4}}{W_{g,2}}
 \frac{W_{g,2}}{W_{g,1}}=O(g^{-1}).
 \label{eq:TL25}
\end{equation} The separating terms at one marking have two positive genera.
Equation \eqref{eq:TL16}, the inequality \(1/[j(g-j)]\le C/g\), and
\(T_{g,2}=O(g^{-2})\) give \begin{equation}
\frac1{W_{g,1}}\sum_{j=1}^{g-1}W_{j,1}W_{g-j,1}
 \le\frac Cg\frac{W_{g,2}}{W_{g,1}}T_{g,2}=O(g^{-2}).
 \label{eq:TL26}
\end{equation} The same degree-difference argument now proves \eqref{eq:TL23} for \(p=1\).
Thus \eqref{eq:TL24} holds also for \(p=0\).

Finally \eqref{eq:TL14}, \eqref{eq:TL19}, and \eqref{eq:TL23} yield \begin{equation}
\frac{W_{g-1,p+2}}{W_{g,p}}=1+O_p(g^{-1}),\qquad p\ge2.
 \label{eq:TL27}
\end{equation} For \(p=0,1\), add two markings in numerator and denominator and use
\eqref{eq:TL24}; the two cusp products have the same complexity factors and
cancel. Hence \eqref{eq:TL27} holds for every fixed \(p\ge0\).

We now return to volumes. Recall the cusp generating function
\(\Phi_g\) of \eqref{eq:intro-Phi} and the string-shift identity
\eqref{eq:H8}: \(U_{g,p}(0^p;Y)=\Phi_g^{(p)}(x)\) with \(x=X(Y)\), whenever
the amplitude is stable. By \eqref{eq:TL9},
\(W_{g,p}=4^ga^hr^m\Phi_g^{(p)}(x)\) with \(m=3g-3+p\). Adding a marking
raises both \(h\) and \(m\) by one, so
\[
 \frac{x\Phi_g^{(p+1)}(x)}{\Phi_g^{(p)}(x)}
 =\frac{x}{ar}\,\frac{W_{g,p+1}}{W_{g,p}}
 =\frac{hx}{A\,a\,r}\bigl(1+O_p(g^{-1})\bigr)
\]
by \eqref{eq:TL24}. Replacing \((g,p)\) by \((g-1,p+2)\) leaves \(h\)
unchanged and lowers \(m\) by one, so
\(\Phi_{g-1}^{(p+2)}/\Phi_g^{(p)}=4r\,W_{g-1,p+2}/W_{g,p}\), and
\eqref{eq:TL27} applies. This proves, for every fixed positive \(p\),
\begin{equation}
x\frac{(\Phi_g^{(p)})'(x)}{\Phi_g^{(p)}(x)}
 =h\Psi(x)+O_p(\Psi(x)),\qquad h=2g-2+p,
 \label{eq:TL30}
\end{equation} \begin{equation}
\frac{\Phi_{g-1}^{(p+2)}(x)}{\Phi_g^{(p)}(x)}
 =4r(Y)(1+O_p(g^{-1})),
 \qquad
 \Psi(x)=\frac{x}{A(Y)a(Y)r(Y)}.
 \label{eq:TL31}
\end{equation} The same genus ratio holds at \(p=0\). The estimates are uniform
throughout \(0\le Y<y_0\); at \(x=0\), \eqref{eq:TL30} is interpreted by
continuity.

The scalar is exactly the stationary parameter: \begin{equation}
H_Y(r)=\frac{3I(Y)}{2\Delta^{3/2}},\qquad
 A(Y)a(Y)r(Y)=\frac{I(Y)}{\sqrt\Delta},\qquad
 \Psi(x)=\lambda(Y)/2.
 \label{eq:TL32}
\end{equation} The first equality follows by substituting \(v=T-\Delta s^2\) in
\(I\), and the other two follow by direct substitution. We summarize the
output of this section used in the extraction arguments.

\begin{proposition}[Marked cusp and genus ratios]\label{prop:marked-ratios}
For every fixed \(p\ge1\), uniformly for \(0\le Y<y_0\) and \(x=X(Y)\), the
cusp ratio of \(\Phi_g^{(p)}\) satisfies \eqref{eq:TL30} and the genus ratio
satisfies \eqref{eq:TL31}, with \(\Psi(x)=\lambda(Y)/2\). The genus ratio
holds also at \(p=0\).
\end{proposition}

\section{The leading term and its normalization}\label{sec:prefactor}

The marked ratios determine a differential equation for the leading
term, but leave a multiplicative constant. The equation comes from the
dilaton identity \eqref{eq:TL12}, which expresses the \(Y\)-derivative of
\(\log\Phi_g^{(p)}\) through the degree profiles \(c_{g,p+1}(d,0^p)\); to
obtain it with error \(O(h^{-2})\) we need the \(1/h\) term of these
profiles. We first compute this degree correction, then integrate it and determine the
constant from pure $\psi$ intersections. The resulting estimate is
Proposition~\ref{prop:prefactor}; its fixed-cusp specialization proves
Corollary~\ref{cor:constant}. Matching at a point that approaches the
critical endpoint polynomially in the genus keeps the integrated error
uniform on the whole interval.

Write $\bar K=\bar K_Y$ and $\bar H=\bar H_Y$, and introduce the kernel moments
\begin{equation}
 \Theta=\sum_{\ell\ge0}w_\ell=-\frac1{\bar K'(1)},\qquad
 M_1=\sum_{\ell\ge0}\ell w_\ell=\frac{\bar K''(1)}{2\bar K'(1)^2},\qquad
 \mathcal Z=Aar^{3/2}=\sqrt2 I.
 \label{eq:PF1}
\end{equation}
Primes on $\bar K,\bar H$ mean differentiation in their argument; derivatives of
parameter functions below are with respect to $Y$. All the moments in
\eqref{eq:PF1} are uniformly bounded on the closed normalized interval.
The variance $v_Y=\sigma_Y^2$ of the law in \eqref{eq:P12}, written $\sigma_Y^2$ in Section~\ref{the-positive-compact-family-of-kernels}, satisfies
\begin{equation}
 v_Y=\frac{15}{4}-\frac{\bar H''(1)}{\bar H(1)}=\frac1{2A\Theta}.
 \label{eq:PF2}
\end{equation}
Indeed $\bar H'(1)=-3\bar H(1)/2$ and
$\bar K'(1)=-5\bar H(1)/2+2\bar H''(1)/3$.

For fixed $p\ge1$, put $h=2g-2+p$ and define
\begin{equation}
 \mathcal J_{g,p}(Y)=
 \frac{\Gamma(h)4^{-g}}{\sqrt{\pi h}}
 \mathcal Z(Y)^{-h}r(Y)^{p/2}\sqrt{A(Y)\Theta(Y)}.
 \label{eq:PF3}
\end{equation}
\begin{proposition}[Absolute marked prefactor]\label{prop:prefactor}
For each fixed $p\ge1$, uniformly for $0\le Y<y_0$,
\begin{equation}
 \Phi_g^{(p)}(X(Y))=
 \frac{\mathcal J_{g,p}(Y)}{\pi}
 \left(1+O_p\!\left(\frac{\log(g+1)}g\right)\right).
 \label{eq:PF4}
\end{equation}
\end{proposition}

\subsection{The first degree correction}

We compute the first degree correction directly from the positive
recursion. The remainder must be uniform in $Y$ and in a growing degree
range, since an expansion at each fixed degree would not suffice in the
dilaton sum. Set $D=|\mathbf d|$ and $q=\#\{i:d_i>0\}$.

\begin{lemma}[First degree correction]\label{lem:first-degree}
For each fixed $p\ge1$, uniformly for $Y\in[0,y_0]$ and for degree tuples
with $D\le h^{1/4}$,
\begin{equation}
 \begin{aligned}
 c_{g,p}(\mathbf d)&=1-\frac{\delta_p(\mathbf d)}h
       +O_p\!\left(\frac{(D+1)^4}{h^2}\right),\\
 \delta_p(\mathbf d)&=A\left\{
 \Theta\bigl(D^2+(p-4)D\bigr)+2M_1D+
 \frac{q(q+5-2p)}2\right\}.
 \end{aligned}
 \label{eq:PF5}
\end{equation}
\end{lemma}

\begin{proof}
Subtract the exact recursion \eqref{eq:H9} at pivot degree $t=d_i$ and
$t+1$, shifting the second kernel index. The difference has coefficients
$w_\ell$. By \eqref{eq:TL24} and \eqref{eq:TL27}, its merging prefactor is
$A/h+O_p(h^{-2})$, and its nonseparating prefactor, including the factor
$1/2$ and the genus normalization, is $2A/h+O_p(h^{-2})$.
For $p=1$ there is no merging term. All child degrees can be replaced
by zero with error at most $C_p(D+\ell+1)^2/h$, by \eqref{eq:TL23}.
The coefficients and pair counts are $O_p(D+\ell+1)$, so the moments
of $w_\ell$ make the resulting error $O_p((D+1)^3/h^2)$.
The separating difference is $O_p((D+1)/h^2)$ by \eqref{eq:TL19};
at $p=1$ use \eqref{eq:TL26}. These bounds also cover dimension-invalid
children through the zero convention in \eqref{eq:TL23}.

Let $z$ count zeros among the other $p-1$ entries. A merger is admissible
when $t+d_j+\ell\ge1$; the nonseparating pair count is
$(t+\ell-1)_+$. Thus the leading adjacent difference is
\begin{equation}
 \frac Ah\left\{\Theta(2D+p-3)+2M_1+(2-z)\mathbf1_{t=0}\right\}.
 \label{eq:PF6}
\end{equation}
Here the zero-pivot term records both the missing zero-zero merger and
the truncated pair count. Since the parent deficit is of order $h$ in
the indicated degree range, extending the kernel sum to infinity costs
$O_p(h^{-2})$, using sufficiently high moments of $w_\ell$.
Summing increments from zero to $\mathbf d$ proves \eqref{eq:PF5}:
the exceptional increments at the creation of successive positive
entries are $3-p,4-p,\ldots,q+2-p$. The sum of the remainders is bounded
by $C_p(D+1)^4/h^2$.
\end{proof}

For $d>h^{1/4}$ in the dilaton identity, use $0\le c\le1$, the
quadratic growth of $\delta_p$, and sufficiently high uniform absolute
moments of $b_d$; moments of order twelve already bound the
tail of $1+d^2/h$ by $O_p(h^{-2})$. This justifies summing the expansion,
although it gives no useful pointwise remainder at large degree.

\subsection{The differential equation for the leading term}

The three signed dilaton moments are exact:
\[
 \sum b_d=A,\qquad \sum d b_d=-A/2,\qquad
 \sum d^2 b_d=\bar H''(1)/3-7A/2.
\]
Apply \eqref{eq:PF5} to $(d,0^p)$ in \eqref{eq:TL12}. Changing the child
complexity $h+1$ to $h$ costs $O_p(h^{-2})$. We obtain
\begin{equation}
 h\frac{W_{g,p}}{W_{g,p+1}}
 =A\left(1-\frac{\beta_p}h+O_p(h^{-2})\right),\qquad
 \beta_p=A\left\{\frac{7-2p}{4}\Theta-M_1+2-p\right\}-\frac12.
 \label{eq:PF7}
\end{equation}
This holds for each fixed $p\ge0$: at $p=0$ its only profile input
has one marking. Inverting and undoing the normalization gives
\begin{equation}
 \frac{\mathrm{d}}{\mathrm{d}Y}\log\Phi_g^{(p)}(X(Y))
 =\frac{h+\beta_p+O_p(h^{-1})}{A r}.
 \label{eq:PF8}
\end{equation}

The correction in \eqref{eq:PF8} has an explicit primitive. Differentiate
$K_Y(r)=0$, using $\partial_YK=2\partial_zK+(K-a)/z$, to get
\begin{equation}
 r'=-(\Theta+2),\qquad
 (\log\mathcal Z)'=-\frac1{Ar}.
 \label{eq:PF9}
\end{equation}
The second identity also follows from $I'=-a\sqrt\Delta$.
Differentiating the root derivative and the corresponding equation
for $H$ gives
\begin{equation}
 \begin{split}
 (\log\Theta)'&=\frac{a'}a+
                  \frac{2\Theta+1-2M_1}{r},\\
 (\log A)'&=\frac{3\Theta/2+3-1/A}{r}-\frac{a'}a.
 \end{split}
 \label{eq:PF10}
\end{equation}
For the first formula use $\Theta=-a/(r\,\partial_zK_Y(r))$ and
$\partial_Y\partial_zK_Y=2\partial_z^2K_Y+\partial_zK_Y/r+a/r^2$ at the root $z=r$.
Equations \eqref{eq:PF7} and \eqref{eq:PF10} yield
\begin{equation}
 \frac{\beta_p}{Ar}
 =\frac12(\log(A\Theta))'+\frac p2(\log r)',\qquad
 \left|\frac{\mathrm{d}}{\mathrm{d}Y}\log
 \frac{\Phi_g^{(p)}(X(Y))}{\mathcal J_{g,p}(Y)}\right|
 \le\frac{C_p}{h A r}.
 \label{eq:PF11}
\end{equation}
Only exact identities have been differentiated here; the error term in
\eqref{eq:PF8} will be integrated, never differentiated.

\subsection{Fixing the constant at the critical endpoint}\label{sec:prefactor-endpoint}

Here we need the value at one pure-$\psi$ tuple, not only two-sided
bounds. In the normalization \eqref{eq:H4},
\begin{equation}
 C_g(2^{3g-3})=\frac1\pi+O(g^{-1}).
 \label{eq:PF12a}
\end{equation}
Indeed, the identity of Itzykson and Zuber~\cite[equation~(6.10)]{IZ}
expresses $C_g(2^{3g-3})$ through the coefficients of the formal
power-series solution of Painlev\'e~I.
Their large-order asymptotics follow from the Stokes data of Joshi--Kitaev
and Kapaev~\cite{JK,Kap}; see~\cite[Section~2.3]{GIKM}.
Subsection~\ref{sec:endpoint-expansion} proves a quantitative version,
\eqref{eq:FE10}, without using the present section. Estimate
\eqref{eq:PF12a} is also the case $\mathbf e=(2^{3g-3})$ of the uniform
theorem of Guo, Yang and Zagier~\cite[Theorem~1]{GYZ}; no other case of
that theorem is used.
For $m=3g-3+p$, string followed by dilaton gives the exact reduction
\begin{equation}
 C_g(0^p,2^m)=
 \prod_{i=0}^{p-1}
 \frac{5(m-i)}{3(2g-3+p+m-2i)}\ C_g(2^{3g-3})
 =\frac1\pi+O_p(g^{-1}).
 \label{eq:PF12}
\end{equation}
Each string step chooses a degree-two marking with normalized
coefficient $5/[3(2g-3+p+m)]$; the resulting degree-one marking is
removed by normalized dilaton with coefficient one. Since $p$ is fixed,
the displayed finite product is $1+O_p(g^{-1})$, and \eqref{eq:PF12a}
gives the final equality.

At the endpoint, \eqref{eq:H24} gives exactly
\begin{equation}
 \frac{\overline U_{g,p}(0^p;y_0)}
      {\mathcal M_{g,p}(0^p;y_0)}=C_g(0^p,2^m).
 \label{eq:PF13}
\end{equation}
Indeed the numerator is $3^m\langle\tau_0^p\tau_2^m\rangle_g/m!$
and the denominator is
$3^{h+m}\Gamma(h+m)/(4^g5^m m!)$.

We now show that the endpoint value persists on a short interval
$0\le y_0-Y\le h^{-2}$. The reason is that the amplitudes are polynomials
with nonnegative coefficients in the kernel coefficients, and these are
within a factor $1+O(y_0-Y)$ of their endpoint value one.
Here $\theta_\ell=[z^\ell]\bar K_Y(z)^{-1}$ are the rescaled kernel coefficients. Put $\delta=y_0-Y$. The normalized kernels are analytic in $\delta$
near zero on a common disk of radius greater than one. To see the
removable parameter singularity, apply the implicit-function theorem
to $K_{y_0-\delta}(\delta\rho)/\delta$ at $(\delta,\rho)=(0,3)$;
its $\rho$ derivative is nonzero. Division by the persistent root
$1-z$ then gives $\bar K_Y(z)/(1-z)=1+O(\delta)$ on a smaller common disk.
Together with coefficient positivity this proves
\begin{equation}
 1\le\theta_\ell(Y)\le1+C\delta\quad\hbox{for all }\ell\ge0.
 \label{eq:PF14}
\end{equation}
By the positive recursion, $\overline U$ is a polynomial with nonnegative
coefficients in the variables $\theta_\ell$, homogeneous of degree $h$. At complexity one the bases have degree one, with $\theta_0=1$;
each recursion adds one factor, while child complexities sum to $h-1$.
Thus
\begin{equation}
 1\le\frac{\overline U_{g,p}(0^p;Y)}
                {\overline U_{g,p}(0^p;y_0)}
 \le(1+C\delta)^h.
 \label{eq:PF15}
\end{equation}

Next we compare $\mathcal J_{g,p}$ with the harmonic model at all
degrees zero, for every $Y$. Here $m=3g-3+p=3h/2-p/2$. Since
$\mathcal Z=Aar^{3/2}$, the powers of $r$ in $a^hr^m\mathcal J_{g,p}$ cancel
exactly. On the other hand, \eqref{eq:H12} and \eqref{eq:H21} give
$\mathcal M_{g,p}(0^p;Y)=3^h\Gamma(h)4^{-g}(3A)^{-h}\Pr(S_h=m)$, and the
central estimate \eqref{eq:H18} at displacement $-p/2$, with variance
$v_Y=1/(2A\Theta)$ from \eqref{eq:PF2}, gives
$\Pr(S_h=m)=\sqrt{A\Theta/(\pi h)}\,(1+O_p(h^{-1}))$. Hence
\begin{equation}
 a^h r^m\mathcal J_{g,p}(Y)
 =\Gamma(h)4^{-g}A^{-h}\sqrt{\frac{A\Theta}{\pi h}}
 =\mathcal M_{g,p}(0^p;Y)(1+O_p(h^{-1})),
 \label{eq:PF16}
\end{equation}
where the first equality is exact.
Its ratio to its endpoint value is $1+O_p(h\delta)$ when
$0\le\delta\le h^{-2}$, by analyticity and uniform positivity of
$A,\Theta$. Hence \eqref{eq:PF12}--\eqref{eq:PF16} prove
$\Phi_g^{(p)}(X(Y))/\mathcal J_{g,p}(Y)=1/\pi+O_p(h^{-1})$
throughout that interval.

\begin{proof}[Proof of Proposition~\ref{prop:prefactor}]
Take $Y_*=y_0-h^{-2}$. For $Y\le Y_*$ integrate \eqref{eq:PF11}:
by \eqref{eq:PF9}, its error is at most
\[
 \frac{C_p}{h}\log\frac{\mathcal Z(Y)}{\mathcal Z(Y_*)}
 \le C_p\frac{\log h}{h},
\]
since $\mathcal Z(0)>0$ and
$\mathcal Z(Y_*)\asymp h^{-5}$. The value at $Y_*$ was fixed above.
Exponentiation gives \eqref{eq:PF4} for $Y\le Y_*$; the direct endpoint
comparison gives the stronger $O_p(h^{-1})$ error for $Y\ge Y_*$.
\end{proof}

At $Y=0$, $r=\pi^2/2$, $A=1/\pi^2$, $\Theta=2$, and
$\mathcal Z=\pi/(2\sqrt2)$. Taking $p=1$ in \eqref{eq:PF4} gives
\begin{equation}
 V_{g,1}=\frac{\Gamma(2g-1)(4\pi^2)^{2g-2}}{\sqrt{\pi g}}
 \left(1+O\!\left(\frac{\log(g+1)}g\right)\right).
 \label{eq:PF17}
\end{equation}
\begin{proof}[Proof of Corollary~\ref{cor:constant}]
Comparing \eqref{eq:PF17} with the expansion of
Mirzakhani--Zograf~\cite[Theorem~1.2 and Remark~1.3]{MZ} at $n=1$
identifies their universal constant as $1/\sqrt\pi$. Since the constant
does not depend on $n$, their expansion then gives the asserted formula,
with relative error $O_n(g^{-1})$, for every fixed $n\ge0$.
\end{proof}

For completeness, we derive the boundary-length consequence stated in the
introduction. Fix $n\ge1$, put $x_i=L_i/2$, and set $S(x)=\sinh(x)/x$,
with $S(0)=1$. At $Y=0$, the normalization \eqref{eq:TL9} gives
$c_{g,n}(\mathbf d)=[\tau_{d_1}\cdots\tau_{d_n}]_{g,n}/V_{g,n}$.
Thus the volume polynomial, extended by zero outside its degree range, is
\[
 \frac{V_{g,n}(\mathbf L)}{V_{g,n}}
 =\sum_{\mathbf d\ge0}c_{g,n}(\mathbf d)
   \prod_{i=1}^n\frac{x_i^{2d_i}}{(2d_i+1)!}.
\]
Divide by $\prod_i S(x_i)$ and regard the product weights as the law of
independent nonnegative integers $d_i$. The identity
$\sum_{d\ge0}(2d)(2d+1)x^{2d}/(2d+1)!=x^2S(x)$ gives
$\mathbb E[d_i(d_i+1)]\le x_i^2/2$. Since
$D(D+1)\le n\sum_i d_i(d_i+1)$, where $D=\sum_i d_i$,
\eqref{eq:TL23}, including its one-marking case, yields
\begin{equation}
 0\le 1-\frac{V_{g,n}(\mathbf L)}{V_{g,n}\prod_i S(L_i/2)}
 \le \frac{C_n}{g}\sum_{i=1}^n L_i^2.
 \label{eq:boundary-factor}
\end{equation}
For $L_i=o(\sqrt g)$ the right side tends to zero. Corollary~\ref{cor:constant}
and Stirling's formula then give the asserted asymptotic formula.

\section{Higher-order corrections for marked generating functions}\label{sec:higher-order}

The estimate of Proposition~\ref{prop:prefactor} has relative error
$O_p(\log g/g)$. The logarithm comes from integrating the error
$O_p(1/(hAr))$ in~\eqref{eq:PF11} towards the critical endpoint, where
$Ar$ has a simple zero. This section states two refinements, still for a
fixed number $p$ of distinguished markings: an explicit first correction,
which also gives the leading term with relative error $O_p(g^{-1})$, and a
complete expansion whose truncation order may grow like $\sqrt g$. Both
results are proved in Appendix~\ref{app:marked}; the passage to a
prescribed cusp count is made in
Sections~\ref{extracting-fixed-cusp-counts} and~\ref{sec:sharp-extraction}.

\subsection{The first correction}\label{subsec:first-correction}
The mechanism is that of Section~\ref{sec:prefactor}, carried one order
further. Lemma~\ref{lem:second-profile} extends the degree expansion
\eqref{eq:PF5} of the normalized profiles by a term $e_p(\mathbf d)/h^2$,
uniformly for $|\mathbf d|\le h^{1/8}$. Inserted into the dilaton identity,
it adds a term $\gamma_p/h$ to the numerator of the differential
equation~\eqref{eq:PF8}; see~\eqref{eq:ST8}. The coefficient $\gamma_p$,
defined in~\eqref{eq:ST7}, is a finite expression in kernel moments
(see~\eqref{eq:ST11}), and it vanishes at the pure-$\psi$ endpoint
(see~\eqref{eq:ST12}). Its primitive against $1/(Ar)$ is therefore bounded,
although $Ar$ vanishes there. The integration constant is fixed at the
endpoint by the first correction to $C_g(2^{3g-3})$, which we derive from
Painlev\'e~I in Subsection~\ref{sec:endpoint-expansion}, together with
the first correction to an explicit negative-binomial coefficient.
Accordingly, define
\begin{equation}
 t_p=\frac{-2p^2+12p-11}{20},\qquad
 \mathcal L_p(Y)=t_p-\int_Y^{y_0}\frac{\gamma_p(v)}{A(v)r(v)}\,\mathrm{d}v.
 \label{eq:ST13}
\end{equation}

\begin{theorem}[Uniform marked first correction]\label{thm:first-correction}
For every fixed $p\ge1$, uniformly for $0\le Y<y_0$,
\begin{equation}
 \Phi_g^{(p)}(X(Y))=\frac{\mathcal J_{g,p}(Y)}\pi
 \left\{1+\frac{\mathcal L_p(Y)}h+
            O_p\!\left(\frac{\log(g+1)}{g^2}\right)\right\}.
 \label{eq:ST14}
\end{equation}
In particular the leading relative error is $O_p(g^{-1})$.
\end{theorem}

At $Y=0$ the coefficients of this expansion reduce to the first-order
formulas of Mirzakhani and Zograf~\cite[Theorems~1.4 and~3.2(i)]{MZ};
this check is recorded at the end of Subsection~\ref{app:second-degree}.

\subsection{All orders}\label{subsec:all-orders}
The same scheme continues to every order. The recursion is triangular:
each new degree coefficient of the profiles uses earlier cusp and genus
coefficients, and the dilaton and KdV equations then give the new cusp and
genus coefficients. The difficulty is to keep the constants under control
as the order grows.

\emph{Where the squared factorial comes from.} The first two profile coefficients are polynomials in the degrees of degrees two and four: $\delta_p$ in~\eqref{eq:PF5} is quadratic in $D=|\mathbf d|$, and $e_p$ in~\eqref{eq:ST3} is quartic. This pattern persists. Each new order is obtained by summing an adjacent difference along a coordinate path of length $D$, and each difference involves pair counts of size $D$ and one factor $1/h$. So the order-$j$ coefficient behaves like a polynomial of degree $2j$ in the degrees, of size about $(D+j)^{2j}$ up to a factor $B^j$. In the dilaton sum the degree $d$ is weighted by the geometrically decaying coefficients $b_d$, so the effective size of $D$ is of order $j$, and $(Mj)^{2j}$ is $(j!)^2$ up to a geometric factor. The bound $B^{K+1}((K+1)!)^2h^{-K-1}$ for the first omitted term is smallest at an order proportional to $\sqrt h$, and a sufficiently small fixed multiple of $\sqrt h$ gives an exponentially small remainder. This explains both the Gevrey-2 bounds and the truncation order in Theorem~\ref{thm:full-expansion}.

Appendix~\ref{app:marked} constructs, by the finite triangular rules
\eqref{eq:FE1}--\eqref{eq:FE6}, the formal coefficients of the profiles
and of the cusp and genus ratios, and from them functions $a_{p,j}(Y)$ and
\[
 \mathcal L_{p,j}(Y)=[t^j]\exp\!\Bigl(\sum_{i=1}^j a_{p,i}(Y)t^i\Bigr),
 \qquad \mathcal L_{p,1}=\mathcal L_p;
\]
see~\eqref{eq:FE11}. On one complex neighborhood of $[0,y_0]$ they are
holomorphic and satisfy
$\|a_{p,j}\|+\|\mathcal L_{p,j}\|\le C_p\mathfrak B_p^j(j!)^2$.

\begin{theorem}[Complete uniform marked expansion]\label{thm:full-expansion}
For each fixed $p\ge1$ there are $C_p,\mathfrak B_p,c_p,h_p>0$ such that,
for $h=2g-2+p\ge h_p$ and $0\le K\le c_p\sqrt h$,
\begin{equation}
 \sup_{0\le Y<y_0}\left|
 \frac{\pi\Phi_g^{(p)}(X(Y))}{\mathcal J_{g,p}(Y)}
       -\sum_{j=0}^K\frac{\mathcal L_{p,j}(Y)}{h^j}\right|
 \le\frac{C_p\mathfrak B_p^{K+1}((K+1)!)^2}{h^{K+1}},
 \qquad \mathcal L_{p,0}=1.
 \label{eq:FE12}
\end{equation}
The coefficients are specified by the finite rules and integrals of
Appendix~\ref{app:marked}, and have the Gevrey-2 norms stated above on
one complex neighborhood of $[0,y_0]$. The same real error holds for the
logarithm with coefficients $a_{p,j}$. A sufficiently small square-root
truncation has error $O_p(\mathrm{e}^{-a_p\sqrt h})$; no convergence of
the infinite series is required.
\end{theorem}

The proof has three steps. Proposition~\ref{prop:hierarchy} divides each
order-$j$ coefficient by $(D+p+Mj)^{2j}$. After this division the
coefficient rules become coefficientwise inequalities between power
series with nonnegative coefficients, and a system of analytic majorants
(Lemma~\ref{lem:majorant}) bounds the normalized coefficients and the
actual remainders, uniformly in the order. Next, at the pure-$\psi$
endpoint, the integration constants are computed from the
Itzykson--Zuber identity and the large-order behavior of Painlev\'e~I;
the Borel representation of Costin, Costin and Huang~\cite{CCH} gives
these constants with a Gevrey-1 remainder. Finally the logarithmic
derivative is integrated from a matching point at distance $h^{-K-2}$
from the endpoint, where the positive homogeneous structure of the
amplitudes controls the change from the endpoint value.

At $Y=0$ this is the established fixed-cusp expansion mechanism;
effective order-dependent versions appear in
Hide--Macera--Thomas~\cite[Theorem~4.2]{HMT}. The estimates here remain
uniform up to the pure-$\psi$ endpoint. Their dependence on the order is
what permits the growing-order truncation in the absolute-volume theorem.

\section{Extracting a prescribed cusp count}\label{extracting-fixed-cusp-counts}

The preceding estimates control generating functions, which sum over cusp counts. We now recover a single coefficient, and hence a volume with a prescribed number of cusps. A positive convolution factor supplies Fourier bounds that turn the real-axis estimates into a local limit theorem. Small cusp weight, the interior, and the critical endpoint require separate estimates, with uniform constants across the three ranges.

Subsection~\ref{canonical-coefficient-ratios-and-all-cusp-count-uniformity} proves the qualitative volume-ratio limits using the leading estimates from Sections~\ref{ratios-with-a-fixed-number-of-distinguished-markings}--\ref{sec:prefactor}. The final subsection extends the argument to bounded geodesic boundaries. The higher-order expansions are needed only for the sharper extraction estimates that follow.

\subsection{A positive convolution factor}\label{a-positive-convolution-factor}

\begin{lemma}[Positive convolution factor]\label{lem:convolution}
For each stable marked pair \((g,p)\) with \(p\ge1\), the marked generating
function has the exact factorization
\begin{equation}
\Phi_g^{(p)}(x)=V(x)^h R_{g,p}(x),\qquad
 V(x)=Y'(x)=1/a(Y(x)),\qquad h=2g-2+p,
 \label{eq:TL33}
\end{equation} where \(R_{g,p}\) has nonnegative Taylor coefficients and
converges for \(0\le x<x_c\).
\end{lemma}

The point of \eqref{eq:TL33} is the following. An approximation of \(\Phi_g^{(p)}\) on
the real axis alone says nothing about its individual coefficients. The
factorization shows that, under the cusp weight \(x\), the number of
additional cusps is distributed as the sum of \(h\) independent variables
with generating function \(V(xz)/V(x)\) and one further independent
nonnegative integer variable; this gives the Fourier bound \eqref{eq:TL35} below, which is what
a local limit theorem needs.

\begin{proof}
The proof uses positivity in the parameter: by
Lemma~\ref{lem:kernel-positivity}, the functions
\(\widehat\vartheta_\ell(Y)=a(Y)[z^\ell]K_Y(z)^{-1}\) and \(1/a(Y)\) have
nonnegative Taylor coefficients in \(Y\). Since this positivity in \(Y\) is
available for the unrescaled coefficients, we do not use the rescaling by
\(r(Y)\) of \eqref{eq:H7} here.

Undo the normalization in \eqref{eq:H9} and regard the unrescaled coefficients \(\vartheta_\ell=[z^\ell]K_Y(z)^{-1}\) as independent variables. In this convention the bases are \(U_{0,3}=\vartheta_0\), \(U_{1,1}(1)=\vartheta_0/8\), and \(U_{1,1}(0)=\vartheta_1/8\). Induction from these bases shows that every amplitude is a finite polynomial with nonnegative
coefficients in the variables \(\vartheta_\ell\); each monomial has total degree \(h\),
and weighted degree \(m\) when \(\vartheta_\ell\) is given weight \(\ell\). A
linear recursion term adds one factor \(\vartheta_\ell\) to a child of complexity \(h-1\); a
separating term adds one to children whose complexities sum to \(h-1\).
Thus factoring \(a^{-1}\) from each of the \(h\) factors \(\vartheta_\ell=a^{-1}\widehat\vartheta_\ell\) leaves a
positive polynomial in the \(\widehat\vartheta_\ell\).

Finally \(Y(x)\) has nonnegative coefficients and \(Y(0)=0\). To justify
the product used here, a nonzero zero \(y\) of \(X\) gives
\(u(s)=J_1(2\sqrt y\,s)\), with \(u(1)=0\) and regular behavior at zero.
Integration of \((s u')'+(4ys-1/s)u=0\) against its conjugate gives
\(4y\int_0^1s|u|^2\,\mathrm{d}s=\int_0^1(s|u'|^2+|u|^2/s)\,\mathrm{d}s>0\). Thus all
nonzero zeros are positive real. Since \(X\) has order one half and a
simple zero at zero with derivative one, its Hadamard product is
\(X(Y)=Y\prod_k(1-Y/\beta_k)\), \(\beta_k>0\). Lagrange inversion now
gives \[
 [x^k]Y(x)=\frac1k[Y^{k-1}](Y/X(Y))^k>0.
\] The inverse series has radius exactly \(x_c\): a smaller radius
would, by positivity and Pringsheim's theorem, give a positive real
singularity before \(x_c\), contradicting the analytic real inverse
there; the quadratic critical point at \(x_c\) supplies the opposite
inequality. Composition therefore preserves coefficient positivity and
converges for \(0\le x<x_c\), since \(Y(x)<y_0\). Finally, by
\eqref{eq:H8}, \(\Phi_g^{(p)}(x)=U^{Y}_{g,p}(0^p)\) with \(Y=Y(x)\), and the
factor \(a(Y)^{-h}\) removed above equals \(V(x)^h\). This proves
\eqref{eq:TL33}.
\end{proof}

For the coefficient probability law with generating function
\(\Phi_g^{(p)}(xz)/\Phi_g^{(p)}(x)\), positivity in \eqref{eq:TL33} gives \begin{equation}
\left|\frac{\Phi_g^{(p)}(x\mathrm{e}^{i\theta})}{\Phi_g^{(p)}(x)}\right|
 \le\left|\frac{V(x\mathrm{e}^{i\theta})}{V(x)}\right|^h.
 \label{eq:TL35}
\end{equation} This is the Fourier estimate that prevents a generating-function
ratio from concealing different local coefficient behavior.

\subsection{Mean, variance and a central limit theorem}\label{mean-variance-and-a-real-central-limit-theorem}

We first compute the mean and variance of the cusp count under the cusp
weight \(x\), and show that it is asymptotically Gaussian; the next
subsection upgrades this to a local limit theorem. Recall
\(\mathscr D=x\,\mathrm{d}/\mathrm{d}x\) and the function \(\Psi\) of
Lemma~\ref{lem:parameter}. For convenience we repeat its endpoint expansions
\eqref{eq:SP16}, together with that of \(V=1/a(Y(x))\), which follows from
\(a(Y)=1-Y+O(Y^2)\) and \(Y(x)=x+O(x^2)\): \begin{equation}
\begin{aligned}
 x\downarrow0:\quad& V(x)=1+x+O(x^2),\quad \Psi(x)=2x+O(x^2),\\
 &\mathscr D\Psi(x)=2x+O(x^2),\quad \mathscr D^2\Psi(x)=2x+O(x^2);\\[1mm]
 x\uparrow x_c:\quad&\Psi(x)=\frac5{4\varepsilon}(1+O(\sqrt\varepsilon)),\\
 &\mathscr D\Psi(x)=\frac5{4\varepsilon^2}(1+O(\sqrt\varepsilon)),\quad
 \mathscr D^2\Psi(x)=O(\varepsilon^{-3}),
 \qquad\varepsilon=1-x/x_c.
 \end{aligned}
 \label{eq:TL36}
\end{equation} On compact subintervals of \((0,x_c)\), all needed
derivatives are bounded and \(\mathscr D\Psi\) is bounded away from
zero.

For \(F=\Phi_g^{(p)}\) at a positive cusp weight \(x\), let \begin{equation}
\mathbb P_x(N=k)=\frac{[z^k]F(z)x^k}{F(x)},\qquad
 \mu=h\Psi(x),\qquad \sigma^2=h\mathscr D\Psi(x).
 \label{eq:TL37}
\end{equation} The last two quantities are model center and variance, not claims
about their exact finite-genus values. Suppose \(h\to\infty\) and
\(\mu\to\infty\). For every fixed real \(u\), integrating \eqref{eq:TL30} gives
\begin{equation}
\log\mathbb E_x \mathrm{e}^{u(N-\mu)/\sigma}
 =\int_0^{u/\sigma}
 \{h\Psi(x\mathrm{e}^s)-h\Psi(x)+O_p(\Psi(x\mathrm{e}^s))\}\,\mathrm{d}s.
 \label{eq:TL38}
\end{equation} The interval of integration stays inside \((0,x_c)\). Expanding
\(h\Psi(x\mathrm e^s)-h\Psi(x)=sh\mathscr D\Psi(x)+O(s^2h\mathscr D^2\Psi)\),
the first part of the integral is \(u^2/2\) plus a Taylor remainder, and the
second part is an integrated error; both are \(o(1)\).
More precisely, in the critical regime its Taylor remainder and
integrated error are \(O_u(h^{-1/2})\); on an interior compact interval
they are also \(O_u(h^{-1/2})\). At \(x\to0\), \(\sigma^2\sim\mu\), the
Taylor remainder is \(O_u(\mu^{-1/2})\) and the integrated error is
\(O_u(x/\sqrt\mu)=O_u(\sqrt\mu/h)=o(1)\). These assertions follow
directly from \eqref{eq:TL36}; no estimate for the derivative of the
error in \eqref{eq:TL30} is required.

Thus the real moment generating functions converge to \(\exp(u^2/2)\)
along every such sequence. Exponential Markov bounds give tightness;
bounds at slightly larger real arguments give uniform integrability
along weakly convergent subsequences. Their limiting moment generating
functions identify the standard Gaussian. Therefore \begin{equation}
(N-\mu)/\sigma\Longrightarrow\mathcal N(0,1).
 \label{eq:TL39}
\end{equation}

\subsection{Fourier bounds and the local limit theorem}\label{uniform-fourier-bounds-and-the-lattice-local-limit}

A central limit theorem alone does not determine ratios of individual
coefficients. We now prove the corresponding local limit theorem.

\begin{proposition}[Uniform local limit theorem]\label{prop:local-limit}
Fix \(p\ge1\) and use the notation \eqref{eq:TL37}. Along every sequence
with \(h\to\infty\), \(\mu\to\infty\) and \(x\in(0,x_c)\),
\begin{equation}
\mathbb P_x(N=k)=\frac{1+o(1)}{\sqrt{2\pi}\sigma}
 \mathrm{e}^{-u^2/2}\qquad\text{whenever }(k-\mu)/\sigma\longrightarrow u.
 \label{eq:TL47}
\end{equation}
No relation between \(h\) and the distance from \(x\) to \(x_c\) is required.
\end{proposition}

\begin{proof}
Passing to a subsequence, we may assume that \(x\) tends to zero, to an
interior point, or to \(x_c\). In each case we bound the characteristic
function of \((N-\mu)/\sigma\) by an integrable function, uniformly along
the sequence; Fourier inversion and the pointwise convergence implied by
\eqref{eq:TL39} then give \eqref{eq:TL47}.

On a compact subinterval of \((0,x_c)\), the coefficient probabilities
at zero and one of \(V(xz)/V(x)\) have positive uniform lower bounds. If
these probabilities are \(p_0,p_1\), then \begin{equation}
\left|\sum_jp_j\mathrm{e}^{ij\theta}\right|^2
 \le1-2p_0p_1(1-\cos\theta).
 \label{eq:TL40}
\end{equation} Equations \eqref{eq:TL35} and \eqref{eq:TL40} give an upper bound
\(\exp(-c h\theta^2)\) on \([-\pi,\pi]\). Since \(\sigma\) is comparable
to \(\sqrt h\), this is integrable Gaussian domination on the rescaled
Fourier axis.

At \(x\to0\), one has \(p_0\asymp1\), \(p_1\asymp x\). The same argument
gives \(\exp(-c h x\theta^2)\). Since \(hx\asymp\mu\) and
\(\sigma^2\sim\mu\to\infty\), it again gives integrable domination after
rescaling.

Near the critical point a different bound is needed, because the law
with generating function \(V(xz)/V(x)\) then has variance of order
\(\varepsilon^{-2}\) and its probabilities at zero and one tend to zero. Quadratic
analytic inversion at \(x_c\) gives \begin{equation}
V(z)=A_0(1-z/x_c)^{-1/2}
 \Xi\bigl(\sqrt{1-z/x_c}\bigr),\qquad
 A_0>0,\quad \Xi(0)=1,
 \label{eq:TL41}
\end{equation} where \(\Xi\) is analytic and nonzero near zero and has real
coefficients. Put \(\Lambda(w)=\log \Xi(\sqrt w)\) and
\(w(\theta)=1-(1-\varepsilon)\mathrm{e}^{i\theta}\). Then \begin{equation}
\log\left|\frac{V(x\mathrm{e}^{i\theta})}{V(x)}\right|
 =-\frac14\log\left(1+
 \frac{2(1-\varepsilon)(1-\cos\theta)}{\varepsilon^2}\right)
 +\Re \Lambda(w(\theta))-\Lambda(\varepsilon).
 \label{eq:TL42}
\end{equation} There are fixed \(\delta,\varepsilon_0>0\) such that \begin{equation}
\left|\frac{V(x\mathrm{e}^{i\theta})}{V(x)}\right|
 \le(1+\theta^2/(4\varepsilon^2))^{-1/8}
 \quad(|\theta|\le\delta,\ 0<\varepsilon\le\varepsilon_0).
 \label{eq:TL43}
\end{equation} Indeed, for \(|\theta|\le\varepsilon\), the first derivative at zero
of the real part of \(\Lambda(w(\theta))\) vanishes, while its second
derivative is \(O(\varepsilon^{-3/2})\). Its change is therefore
\(O(\theta^2\varepsilon^{-3/2})\), absorbed by half of the leading
logarithm in \eqref{eq:TL42} when \(\varepsilon\) is small. For
\(\varepsilon\le|\theta|\le\delta\), its change is bounded by
\(C\sqrt\delta\), while the logarithm is at least \(\log(5/4)\). Choose
\(\delta\) small enough for the same absorption. We used the elementary
bound \(2(1-\varepsilon)(1-\cos\theta)\ge\theta^2/4\) on the small arc.

On the outer arc \(\delta\le|\theta|\le\pi\), \begin{equation}
|V(x\mathrm{e}^{i\theta})/V(x)|\le C_\delta\sqrt\varepsilon.
 \label{eq:TL44}
\end{equation} To see that \(V\) is analytic there, note that nonnegative coefficients and
\(Y(x)\to y_0\) give absolute convergence of \(Y\) on \(|x|\le x_c\).
Its first two coefficients are positive, so strict triangle inequality
gives \(|Y(x_c \mathrm{e}^{i\theta})|<y_0\) away from \(\theta=0\). By Lemma~\ref{lem:kernel-zeros}, \(X'\) has no zero in \(|Y|<y_0\). The inverse-function theorem
extends \(Y,V\) across the outer arc, where \(V\) is uniformly bounded.
Since \eqref{eq:TL41} gives \(V(x)\asymp\varepsilon^{-1/2}\), \eqref{eq:TL44} follows.

On the critical Fourier scale \(\theta=\varepsilon t/\sqrt h\), \eqref{eq:TL35}
and \eqref{eq:TL43} are bounded by \begin{equation}
(1+t^2/(4h))^{-h/8}\le(1+t^2/64)^{-2},\qquad h\ge16.
 \label{eq:TL45}
\end{equation} The inequality follows because \((1+a/h)^h\) increases with \(h\).
The right side is integrable. The outer contribution after
multiplication by \(\sqrt h/\varepsilon\) is at most \begin{equation}
2\pi\sqrt h\,C_\delta^h\varepsilon^{h/2-1}\longrightarrow0.
 \label{eq:TL46}
\end{equation} This convergence is uniform for arbitrarily small \(\varepsilon\):
fix sufficiently small \(\varepsilon_0\) with
\(C_\delta\sqrt{\varepsilon_0}<1\), and bound \eqref{eq:TL46} by
\(2\pi\varepsilon_0^{-1}\sqrt h(C_\delta\sqrt{\varepsilon_0})^h\) for
\(h\ge2\).

These integrable bounds, in the three regimes, complete the proof.
\end{proof}

The same Fourier bounds give an upper bound for every point mass, not
only near the center. We record it for later use in
Section~\ref{sec:geometry} and Appendix~\ref{app:tight}.

\begin{lemma}[Uniform point-mass bound]\label{lem:point-mass}
For every $c_0>0$ there is a constant $C$ with the following property.
Let $x\in(0,x_c)$, let $h,h'$ be positive integers with $h'\ge\max(16,c_0h)$,
and let $G(z)=V(z)^{h'}R(z)$, where $R$ has nonnegative Taylor
coefficients, converges for $|z|<x_c$ and satisfies $R(x)>0$. Then,
with $\sigma^2=h\mathscr D\Psi(x)$,
\[
 \sup_{j\ge0}\frac{[z^j]G(z)\,x^j}{G(x)}\le\frac C\sigma.
\]
\end{lemma}

\begin{proof}
By Fourier inversion and positivity of the coefficients of $R$,
\[
 \frac{[z^j]G(z)\,x^j}{G(x)}
 \le\frac1{2\pi}\int_{-\pi}^{\pi}
       \left|\frac{V(x\mathrm e^{i\theta})}{V(x)}\right|^{h'}\mathrm d\theta.
\]
On a compact subinterval of $(0,x_c)$, \eqref{eq:TL40} bounds the
integrand by $\mathrm e^{-ch'\theta^2}$, so the integral is
$O(h^{-1/2})$, while $\sigma\asymp\sqrt h$. Near zero the same argument
with $p_1\asymp x$ gives $\mathrm e^{-ch'x\theta^2}$ and the bound
$O((hx)^{-1/2})$, while $\sigma^2\asymp hx$. Near $x_c$, put
$\theta=\varepsilon t/\sqrt{h'}$ on $|\theta|\le\delta$; \eqref{eq:TL43}
and the argument of \eqref{eq:TL45} bound this part by
$(\varepsilon/\sqrt{h'})\int(1+t^2/64)^{-2}\,\mathrm dt$. On the outer arc
\eqref{eq:TL44} gives at most
$2\pi(C_\delta\sqrt\varepsilon)^{h'}\le C\varepsilon h^{-1/2}$ once
$C_\delta\sqrt{\varepsilon_0}<1$. Since $\sigma\asymp\sqrt h/\varepsilon$
there, all three regimes give $O(1/\sigma)$.
\end{proof}

\subsection{Ratios at a prescribed number of cusps}\label{canonical-coefficient-ratios-and-all-cusp-count-uniformity}

We now prove the limits in Theorem~\ref{thm:volume}. Each volume ratio is
a ratio of coefficients, at the same index, of two marked generating
functions. We evaluate both at the same cusp weight, chosen so that the
index is, up to a negligible shift, the model center of both. The local limit theorem \eqref{eq:TL47} shows
that the two local probabilities are asymptotically equal, so the
coefficient ratio is asymptotic to the ratio of the generating functions
themselves, which Section~\ref{ratios-with-a-fixed-number-of-distinguished-markings}
determines.

\begin{proof}[Uniform convergence in Theorem~\ref{thm:volume}]
First suppose \(g,n\to\infty\). Put \begin{equation}
k=n-1,\qquad h=2g-1,
 \qquad h\Psi(x)=k.
 \label{eq:TL48}
\end{equation} The last equation has a unique solution in \((0,x_c)\) by \eqref{eq:TL32} and
the parameter bijection. The two families \(\Phi_g'\) and
\(\Phi_{g-1}'''\) have exactly the same complexity \(h\), model center,
and model variance. Thus their two local probabilities in \eqref{eq:TL47}, at the
common coefficient index \(k\), have ratio tending to one. Equations
\eqref{eq:TL31} and \eqref{eq:TL47} give \begin{equation}
\frac{[x^k]\Phi_{g-1}'''(x)}{[x^k]\Phi_g'(x)}
 =4r(Y)(1+o(1)).
 \label{eq:TL49}
\end{equation} By the definition \eqref{eq:intro-Phi}, the coefficients are \begin{equation}
[x^k]\Phi_g'(x)=
 \frac{V_{g,n}}{(2\pi^2)^{3g-3+n}k!},\qquad
 [x^k]\Phi_{g-1}'''(x)=
 \frac{V_{g-1,n+2}}{(2\pi^2)^{3g-4+n}k!}.
 \label{eq:TL50}
\end{equation} Consequently \begin{equation}
B_{g,n}=\frac{\pi^2}{2r(Y)}(1+o(1))
 =\mathcal B\left(\frac{2(n-1)}{2g-1}\right)(1+o(1)).
 \label{eq:TL51}
\end{equation} The factorial \(k!\) cancels, and the powers of \(2\pi^2\) differ by
exactly one; this explains the factor in \eqref{eq:TL51}.

For the cusp ratio, compare \(\Phi_g''\) with \(\Phi_g'\) at the same
index \(k\) and the same cusp weight \eqref{eq:TL48}. Their complexities are \(h+1\)
and \(h\). Their centers differ by \(\Psi(x)\), and their variances have
ratio \((h+1)/h\). Moreover \begin{equation}
\frac{\Psi(x)}{\sqrt{h\mathscr D\Psi(x)}}=O(h^{-1/2})
 \label{eq:TL52}
\end{equation} uniformly in \(x\), because \(\Psi^2/(\mathscr D\Psi)\) is continuous
in the interior and has finite endpoint limits zero and \(5/4\),
respectively. Thus the local probability ratio again tends to one.
Undoing \eqref{eq:TL24}, \begin{equation}
C_{g,n}=\frac{\pi^2}{g}
 \frac{[x^k]\Phi_g''}{[x^k]\Phi_g'}
 =\frac{\pi^2}{g}\frac{h}{Aar}(1+o(1))
 =\frac{\pi^2(n-1)}{gx}(1+o(1)).
 \label{eq:TL53}
\end{equation} Here, as in \eqref{eq:TL50}, the coefficients \([x^k]\Phi_g''\) and
\([x^k]\Phi_g'\) carry the normalizing powers \((2\pi^2)^{3g-2+n}\) and
\((2\pi^2)^{3g-3+n}\); division by \(2g\) gives the factor \(\pi^2/g\). With
\(\lambda_{\mathrm{eff}}=2(n-1)/(2g-1)\), the last expression equals \begin{equation}
\left(1-\frac1{2g}\right)\mathcal C(\lambda_{\rm eff})(1+o(1)).
 \label{eq:TL54}
\end{equation}

One may replace \(\lambda_{\mathrm{eff}}\) by \(n/g\) in both \eqref{eq:TL51} and
\eqref{eq:TL54}. At a zero limiting slope, use the continuous values
\(\mathcal B(0)=1\) and \(\mathcal C(0)=4\pi^2\) of \eqref{eq:SP3}; at a
finite positive slope use continuity; and at an infinite slope use the
endpoint equivalents \eqref{eq:SP15} of Lemma~\ref{lem:parameter}.
Since \(\lambda_{\mathrm{eff}}/(n/g)\to1\) whenever \(n\to\infty\),
these give the desired relative replacement in all cases.

Finally the exact identity \begin{equation}
D_{g,n}=C_{g,n-1}/C_{g,n}
 \label{eq:TL56}
\end{equation} and the adjacent continuity of the profile \(\mathcal C\) give
\(D_{g,n}\to1\) for every growing-cusp sequence. For bounded cusp
counts, the classical fixed-cusp theorems give \(B_{g,n}\to1\),
\(C_{g,n}\to4\pi^2\), and hence \(D_{g,n}\to1\); these are the only
bounded-cusp inputs used here.

If any supremum in Theorem~\ref{thm:volume} did not tend to zero, select a violating
sequence of cusp counts. A subsequence either has a fixed bounded cusp
count or tends to infinity. The preceding conclusions contradict a fixed
violation in both cases. This proves the uniform convergence assertions.
Section~\ref{subsec:quantitative-volume-ratios} supplies the stated
inverse-genus rate using the sharper extraction theorem.

\end{proof}

\subsection{Bounded geodesic boundaries}\label{subsec:bounded-boundaries}

The positive factor also permits a fixed number of cusps to be replaced
by geodesic boundaries. The resulting factor depends on the cusp density
and joins the familiar large-genus and large-cusp boundary profiles.
For $0\le Y\le y_0$ and $L\ge0$, put
\begin{equation}
 \mathfrak b_Y(L)=\int_0^1
 I_0\!\left(\frac L\pi\sqrt{T(Y)-\Delta(Y)s^2}\right)\,\mathrm ds,
 \label{eq:boundary-profile}
\end{equation}
where $I_0$ is the modified Bessel function of order zero, $T=Y+\Delta$
and $r=2\Delta$. At $Y=y_0$ use $T=y_0$ and $\Delta=0$.
Writing $Y_\lambda$ for the inverse of $\lambda(Y)$, we have
$Y_0=0$ and $Y_\infty=y_0$.

\begin{proposition}[Boundary lengths at arbitrary cusp density]
\label{prop:bounded-boundaries}
Fix an integer $p\ge1$ and $L_0<\infty$. As $g\to\infty$,
\begin{equation}
 \sup_{\substack{n\ge p\\\mathbf L\in[0,L_0]^p}}
 \left|
 \frac{V_{g,n}(L_1,\ldots,L_p,0^{n-p})}
 {V_{g,n}\prod_{i=1}^p\mathfrak b_{Y_{n/g}}(L_i)}-1
 \right|\longrightarrow0.
 \label{eq:boundary-uniform}
\end{equation}
The endpoint factors are
\begin{equation}
 \mathfrak b_0(L)=\frac{\sinh(L/2)}{L/2},\qquad
 \mathfrak b_{y_0}(L)=I_0\!\left(\frac{j_{0,1}L}{2\pi}\right),
 \label{eq:boundary-endpoints}
\end{equation}
with value one at $L=0$.
\end{proposition}

The first factor occurs in the fixed-cusp estimate of Mirzakhani and
Petri~\cite[Proposition~3.1]{MP}; the second is the fixed-genus,
large-cusp factor of Hide and Thomas~\cite[Theorem~1.5]{HTLargeN}.
Here the genus tends to infinity and the cusp count is unrestricted.
The number and lengths of the nonzero boundaries remain bounded.
For comparison, Budd and Lions obtain a boundary factor for tight volumes
with a random cusp count in a near-critical regime
\cite[Proposition~4.2]{BL}. Exact transformations between ordinary and
tight boundary generating functions were established in~\cite{BZ}.

\begin{proof}
Put $h=2g-2+p$, $F(x)=\Phi_g^{(p)}(x)$ and
\[
 G(x;\mathbf L)=\sum_{k\ge0}
 \frac{V_{g,p+k}(\mathbf L,0^k)}{(2\pi^2)^{3g-3+p+k}}
 \frac{x^k}{k!}.
\]
We first compare $G$ with $F$ on the real axis, and then show that the
same comparison holds for their coefficients at the saddle.

\emph{The boundary insertion.}
Use the string flow of \eqref{eq:H8}, holding its parameter $Y$ fixed
while differentiating the KW times. Its linear part gives
\[
 \partial_{t_d}=\sum_{b=0}^d\frac{Y^{d-b}}{(d-b)!}\partial_{t_b'}.
\]
For $g\ge2$ the free energy is invariant under the flow. Apply the
displayed identity to each external insertion before setting $x=X(Y)$.
In the volume normalization of \eqref{eq:H5}, an external degree $d$
has coefficient $(L^2/(4\pi^2))^d/d!$ times the unweighted derivative.
Consequently, with $t_i=L_i^2/(4\pi^2)$, the exact identity is
\begin{equation}
 \frac{G(X(Y);\mathbf L)}{F(X(Y))}
 =\sum_{\mathbf a,\mathbf b\ge0}c_{g,p}(\mathbf b;Y)
 \prod_{i=1}^p
 \frac{t_i^{a_i+b_i}Y^{a_i}r^{b_i}}
 {a_i!(a_i+b_i)!(2b_i+1)!!}.
 \label{eq:boundary-insertion}
\end{equation}
Indeed, \eqref{eq:TL9} gives
$U^Y_{g,p}(\mathbf b)/U^Y_{g,p}(\mathbf0)=r^{|\mathbf b|}c_{g,p}(\mathbf b;Y)$.
Terms with invalid $\mathbf b$ are zero; there is no bound on $\mathbf a$.

The identity
\[
 \int_0^1(1-s^2)^b\,\mathrm ds=\frac{2^b b!}{(2b+1)!!}
\]
and the series for $I_0$ show that replacing $c_{g,p}$ by one in
\eqref{eq:boundary-insertion} gives
$\prod_i\mathfrak b_Y(L_i)$. The same double series with an extra
factor $|\mathbf b|(|\mathbf b|+1)$ is bounded uniformly for
$Y\in[0,y_0]$ and $\mathbf L\in[0,L_0]^p$: the inequality
$(a+b)!\ge a!b!$ supplies a summable factorial majorant for every fixed
polynomial moment. Thus \eqref{eq:TL23}, including its one-marking case,
gives
\begin{equation}
 \frac{G(X(Y);\mathbf L)}{F(X(Y))}
 =\prod_i\mathfrak b_Y(L_i)\bigl(1+O_{p,L_0}(g^{-1})\bigr).
 \label{eq:boundary-real}
\end{equation}
All terms in \eqref{eq:boundary-insertion} are nonnegative. The
majorant also justifies convergence of the formal string-flow identity
for $0\le x<x_c$. The factors $\mathfrak b_Y(L)$ are at least one;
they and their first $Y$ derivatives are uniformly bounded on the
compact set under consideration.

\emph{Extraction at a prescribed cusp count.}
The proof of \eqref{eq:TL33} applies to every external degree tuple:
$a(Y)^hU^Y_{g,p}(\mathbf b)$ has nonnegative Taylor coefficients in $Y$.
Before division by $F$, the insertion identity is a sum of these
amplitudes with nonnegative powers of $Y$ and nonnegative constants.
Since $Y(x)$ also has nonnegative coefficients, it follows that
\[
 G(x;\mathbf L)=V(x)^h R_{g,p,\mathbf L}(x),\qquad
 [x^j]R_{g,p,\mathbf L}(x)\ge0.
\]
Convergence follows from \eqref{eq:boundary-real}. Hence
\begin{equation}
 \left|\frac{G(x\mathrm e^{i\theta};\mathbf L)}{G(x;\mathbf L)}\right|
 \le\left|\frac{V(x\mathrm e^{i\theta})}{V(x)}\right|^h.
 \label{eq:boundary-fourier}
\end{equation}

Suppose first that $k=n-p\to\infty$, and choose $x$ by
$h\Psi(x)=k$. Put $\sigma^2=h\mathscr D\Psi(x)$. For fixed real $u$,
the ratio of the centered moment generating functions associated with
$G$ and $F$ is
\[
 \frac{(G/F)(x\mathrm e^{u/\sigma};\mathbf L)}
 {(G/F)(x;\mathbf L)}.
\]
It tends to one, uniformly in $\mathbf L$. To see this, the endpoint
expansions \eqref{eq:TL36} and the quadratic inverse at $x_c$ show that
\[
 \sup_{0<x<x_c}\frac{xY'(x)}{\sqrt{\mathscr D\Psi(x)}}<\infty.
\]
On the tilted interval $x\mathrm e^{u/\sigma}$ the model variance is
comparable to its value at $x$, and the interval remains subcritical.
Therefore $Y(x\mathrm e^{u/\sigma})-Y(x)=O_u(h^{-1/2})$.
Now use \eqref{eq:boundary-real} and the bounded derivative of
$\mathfrak b_Y$.

The real Gaussian limit \eqref{eq:TL39} thus holds for $G$ with the same
center $k$ and scale $\sigma$. Its Fourier bound
\eqref{eq:boundary-fourier} is identical to the one used in
\eqref{eq:TL40}--\eqref{eq:TL46}. Fourier inversion consequently gives
the local limit \eqref{eq:TL47} for $G$ as well. At the common index
$k$, both local probabilities are asymptotic to
$(\sqrt{2\pi}\sigma)^{-1}$. Their ratio tends to one, and hence
\[
 \frac{[x^k]G(x;\mathbf L)}{[x^k]F(x)}
 =\prod_i\mathfrak b_{Y}(L_i)(1+o(1)),\qquad
 \lambda(Y)=\frac{2(n-p)}{2g-2+p}.
\]
All errors are uniform on the fixed length box, including along
sequences tending to either endpoint of the parameter interval.

For bounded $k$, pass to a subsequence on which $k$ is fixed.
Equation~\eqref{eq:TL23} at $Y=0$, with $p+k$ markings, gives the
fixed-cusp boundary estimate \eqref{eq:boundary-factor}. This yields the
same conclusion since the saddle then tends to zero.

Finally, replacing $2(n-p)/(2g-2+p)$ by $n/g$ changes the corresponding
$Y$ by $o(1)$ uniformly in $n\ge p$. Indeed, along a subsequence on
which $n/g$ tends to zero, a positive finite limit, or infinity, both
parameters have the same limit in $[0,y_0]$. Continuity of
$\mathfrak b_Y(L)$, uniformly on the length box, proves
\eqref{eq:boundary-uniform}. At $Y=0$, $T=\Delta=\pi^2/4$;
expanding $I_0$ and integrating gives the first formula in
\eqref{eq:boundary-endpoints}. At $Y=y_0$, the integrand is constant,
which gives the second.
\end{proof}


\section{Extraction with an inverse-genus error}
\label{sec:sharp-extraction}

The local limit theorem \eqref{eq:TL47} has no rate. Its quantitative form, Proposition~\ref{prop:quantitative-local} in Appendix~\ref{app:extraction}, gives only a power rate in $\min(h,\mu)$, and its Gaussian approximation degenerates when the coefficient index stays bounded. To obtain an error of order $g^{-1}$ uniformly in the index, we retain an exact Poisson factor rather than its Gaussian approximation. Theorem~\ref{thm:sharp-extraction} supplies the estimate used for the leading absolute-volume formula in Section~\ref{sec:absolute}, and Theorem~\ref{thm:canonical-expansion} extends it to all orders. Both theorems are proved in Appendix~\ref{app:extraction}.

\emph{Why a Poisson factor.} Near zero cusp weight the leading term of $\Phi_g^{(p)}$ behaves like $\mathrm e^{hS(x)}$ with $S(x)=2x+O(x^2)$; see~\eqref{eq:CE1} below. At a positive weight $x$, the normalized terms $\mathrm e^{-2hx}(2hx)^k/k!$ are exactly Poisson probabilities with mean $2hx$. Hence, when the coefficient index $k$ stays bounded, the cusp count under the saddle weight is approximately Poisson, and a Gaussian approximation of its point mass makes a relative error of order $1/k$, which does not tend to zero. If the leading term is written instead with the exact Poisson mass $\mathfrak p_k=\mathrm e^{-k}k^k/k!$ at the mean, the ratio of the true coefficient to this model is $1+O(1/h)$ for every $k$. For large $k$ one has $\sqrt{2\pi k}\,\mathfrak p_k\to1$, and the Gaussian formula is recovered.

Fix $p\ge1$, put $F_h=\Phi_g^{(p)}$, $h=2g-2+p$, and write
\begin{equation}
 \mathcal J_{g,p}(x)=\mathcal J_{g,p}(0)\mathrm{e}^{hS(x)}B_p(x),\qquad
 S=\log\frac{\mathcal Z(0)}{\mathcal Z(Y(x))},\qquad
 B_p=\frac{r^{p/2}\sqrt{A\Theta}}{r(0)^{p/2}\sqrt{A(0)\Theta(0)}}.
 \label{eq:CE1}
\end{equation}
We view $\mathcal J$ as a function of $x=X(Y)$. Then
$\mathscr DS=\Psi$ and $\mathscr D\log B_p=\beta_p\Psi$.
For each integer $k\ge1$, choose $x$ by $h\Psi(x)=k$, and put
$v=\mathscr D\Psi$, $\sigma^2=hv$, $M=\min(h,k)$.
Let $\mathfrak p_k=\mathrm{e}^{-k}k^k/k!$, with $\mathfrak p_0=1$.

\begin{theorem}[Uniform Poisson-corrected extraction]\label{thm:sharp-extraction}
For fixed $p\ge1$, uniformly over integers $k\ge0$,
\begin{equation}
 [z^k]F_h(z)=\frac{\mathcal J_{g,p}(x)}\pi x^{-k}\mathfrak p_k
       \sqrt{\frac{k}{h\mathscr D\Psi(x)}}\,(1+O_p(h^{-1})).
 \label{eq:CE2}
\end{equation}
At $k=0$ take $x=0$ and interpret both $k$-dependent factors by
continuity as one.
\end{theorem}

The proof, given in Subsections~\ref{app:local-complex}--\ref{app:small-indices}, has three parts: transfer a sufficiently long marked expansion, namely Theorem~\ref{thm:full-expansion} at a fixed finite order, to small complex neighborhoods of the real saddle point; identify the Poisson correction; and estimate the small indices directly. The argument uses no estimate for a growing number of markings and never differentiates a remainder; away from zero, the model is continued analytically only to small neighborhoods of the real axis.

\subsection{The full expansion at a prescribed cusp count}\label{sec:canonical-expansion}

Theorem~\ref{thm:sharp-extraction} recovers the leading term with an inverse-genus error. We now extract the whole marked expansion, keeping the remainder uniform in both the cusp count and the truncation order. The result is Theorem~\ref{thm:canonical-expansion}, which provides the remaining input for the full absolute-volume expansion in Section~\ref{sec:absolute}.

\emph{The idea.} Away from zero cusp weight, including near the critical point, the coefficient integral has a nondegenerate saddle point, and the Gaussian method of steepest descent gives an expansion in $1/h$. Near zero this fails for bounded $k$: the saddle point of $\mathrm e^{hS(z)}z^{-k}$ approaches the singularity of $z^{-k}$ at the origin. Following Temme, we change variables so that the phase $S(z)-(k/h)\log z$ becomes exactly $u-(k/h)\log u$. The coefficient is then a Poisson coefficient $[u^k]\mathrm e^{hu}f(u)$ with an analytic amplitude $f$. Expanding $f$ about $u=k/h$ produces the signed central polynomials $T_m(k)$ of the coefficient operator in~\eqref{eq:CG8}, divided by powers of $h$. These are not the central moments of a probability distribution: for example, $T_2(k)=-k$. The resulting expansion in $1/h$ is uniform in $k$, including $k=0$. The two expansions are identified on an overlap of their ranges.

The proof, given in Subsections~\ref{app:phase-change}--\ref{app:canonical-proof}, uses a Poisson integral near zero cusp weight and a Gaussian integral away from zero, including near criticality. A change of phase identifies the two coefficient formulas on their overlap. Near zero we follow Temme's method for uniform Laplace and loop integrals~\cite[Sections~2.3 and~3.7]{Temme}.

Retain $h=2g-2+p$, $k=h\Psi(x)$, $v=\mathscr D\Psi(x)$ and the
functions $S,B_p$ in \eqref{eq:CE1}. Regard $\mathcal L_{p,j}$
as functions of $x$ via $Y(x)$. Write
\begin{equation}
 \mathcal Q_{g,p,k}=
 \frac{\mathcal J_{g,p}(x)}\pi x^{-k}\mathfrak p_k
                   \sqrt{\frac{k}{hv}}.
 \label{eq:CG1}
\end{equation}
The conventions at $k=0$ are those of Theorem~\ref{thm:sharp-extraction}.


Put $\mathcal L_{p,0}=1$. Let $\mathfrak G_w$ be the linear functional
on polynomials which maps $w^{2a}$ to $(-1)^a(2a-1)!!$ and odd
powers to zero. Define the formal even series
\begin{equation}
\begin{split}
 \mathcal N_p(\tau,x)=\mathfrak G_w\Bigg[
 &\exp\!\left\{
 \sum_{m\ge3}\frac{\tau^{m-2}\mathscr D^mS}{m!v^{m/2}}w^m
 +\sum_{m\ge1}\frac{\tau^m\mathscr D^m\log B_p}{m!v^{m/2}}w^m
 \right\}\\[-2pt]
 &\hspace{20pt}\cdot\sum_{s\ge0}\tau^{2s}
       \mathcal L_{p,s}(x\mathrm{e}^{\tau w/\sqrt v})\Bigg].
\end{split}
 \label{eq:CG2}
\end{equation}
Formal composition in the last line means Taylor expansion at $x$.
At every order only finitely many derivatives are used. If
$\sqrt{2\pi k}\,\mathfrak p_k\sim\sum_{j\ge0}s_jk^{-j}$, $s_0=1$, set
\begin{equation}
 \mathcal A_{p,j}(Y)=[\tau^{2j}]
       \frac{\mathcal N_p(\tau,X(Y))}
            {\displaystyle\sum_{i\ge0}s_i\Psi(X(Y))^{-i}\tau^{2i}}.
 \label{eq:CG3}
\end{equation}
This initially defines the functions for $0<Y<y_0$.

\begin{theorem}[Uniform expansion at a prescribed cusp count]\label{thm:canonical-expansion}
For every fixed $p\ge1$ there are constants $C_p,\mathfrak B_p,c_p,h_p>0$
such that, for $h=2g-2+p\ge h_p$, every integer $k\ge0$ and
$0\le K\le c_p\sqrt h$,
\begin{equation}
 \left|\frac{[z^k]\Phi_g^{(p)}}{\mathcal Q_{g,p,k}}
       -\sum_{j=0}^K\frac{\mathcal A_{p,j}(Y)}{h^j}\right|
 \le\frac{C_p\mathfrak B_p^{K+1}((K+1)!)^2}{h^{K+1}},
 \qquad h\Psi(X(Y))=k,
 \label{eq:CG4}
\end{equation}
where $\mathcal A_{p,0}=1$.
The coefficients are analytic on one complex neighborhood $\Omega_p$
of $[0,y_0]$, with
$\|\mathcal A_{p,j}\|_{\Omega_p}\le C_p\mathfrak B_p^j(j!)^2$, and
$\mathcal A_{p,j}(0)=\mathcal L_{p,j}(0)$.
For a sufficiently small fixed $c'_p>0$, truncation at
$K=\lfloor c'_p\sqrt h\rfloor$ has relative error
$O_p(\mathrm{e}^{-a_p\sqrt h})$, uniformly in $k$, for some $a_p>0$.
The marking count is fixed. The infinite series need not converge.
\end{theorem}

The first coefficient $\mathcal A_{p,1}$ is given explicitly in \eqref{eq:CG12}, and its value at the critical endpoint in \eqref{eq:CG13}. Retaining it improves the relative remainder to $O_p(h^{-2})$ over every coefficient index.


\section{Absolute volumes at every cusp count}\label{sec:absolute}

We now assemble the preceding estimates to prove Theorem~\ref{thm:intro-absolute}. The extraction theorem, Theorem~\ref{thm:sharp-extraction}, gives the explicit leading term with relative error $O(g^{-1})$, and also the quantitative volume ratios. The expansion of Theorem~\ref{thm:canonical-expansion} then gives all higher corrections and their uniform remainder. We treat positive cusp counts first and recover $n=0$ from the exact dilaton relation.

For $g\ge2$ and $n\ge1$, set $h=2g-1$, $k=n-1$, and choose $x=X(Y)$
by $h\Psi(x)=k$. At $k=0$ take $x=Y=0$. Define
\begin{equation}
 \begin{split}
 \widehat V_{g,n}&=
 (2\pi^2)^{3g-3+n}\frac{\mathcal J_{g,1}(Y)}{\pi}
 \left(\frac{k}{\mathrm{e} x}\right)^k
 \sqrt{\frac{k}{h\mathscr D\Psi(x)}}\quad(n\ge1),\\
 \widehat V_{g,0}&=
 \frac{\Gamma(2g-2)(4\pi^2)^{2g-3}}{\sqrt{\pi g}}.
 \end{split}
 \label{eq:AV1}
\end{equation}
At $k=0$, both $k$-dependent factors on the first line are defined by
their limits, which are one. All functions in this formula were defined
in \eqref{eq:TL32}, \eqref{eq:PF1}, and \eqref{eq:PF3}.
Since $r=2\Delta$, $Aar=I/\sqrt\Delta$ by \eqref{eq:TL32}, and
$\Theta=-1/\bar K_Y'(1)=-a/(r\,\partial_zK_Y(r))$ by \eqref{eq:PF1}, one has
$r^{1/2}\sqrt{A\Theta}=\varpi(Y)$ and $\mathcal Z=\sqrt2\,I$, with
$\varpi$ as in the introduction. Hence \eqref{eq:AV1} is the formula
\eqref{eq:intro-leading}, with $\mathscr D\Psi(x)=x\Psi'(x)$.

\begin{theorem}[The leading term]\label{thm:absolute}
There is a constant $C>0$ such that, for all sufficiently large $g$,
\begin{equation}
 \sup_{n\ge0}\left|\frac{V_{g,n}}{\widehat V_{g,n}}-1\right|
       \le\frac Cg.
 \label{eq:AV2}
\end{equation}
The order $g^{-1}$ is sharp for the leading model \eqref{eq:AV1}:
\begin{equation}
 \liminf_{g\to\infty}g\sup_{n\ge0}
       \left|\frac{V_{g,n}}{\widehat V_{g,n}}-1\right|
       \ge\frac7{12}-\frac{17}{6\pi^2}>0.
 \label{eq:AV3}
\end{equation}
\end{theorem}
\begin{proof}
For $n\ge1$, apply Theorem~\ref{thm:sharp-extraction} with $p=1$.
The coefficient identity \eqref{eq:TL50} gives
\begin{equation}
 V_{g,n}=(2\pi^2)^{3g-3+n}(n-1)!x^{-k}
     \frac{\mathcal J_{g,1}(Y)}\pi\,
       \mathfrak p_k\sqrt{\frac{k}{h\mathscr D\Psi(x)}}
       (1+O(g^{-1})),
 \label{eq:AV4}
\end{equation}
uniformly, including $k=0$ by continuity.
Since $k!\mathfrak p_k=(k/\mathrm{e})^k$, this is precisely the first line of
\eqref{eq:AV1} with uniform relative error $O(g^{-1})$.
For $n=0$, the constant was independently identified in
Section~\ref{sec:prefactor}. The fixed-cusp expansion of
Mirzakhani--Zograf~\cite[Theorem~1.2 and Remark~1.3]{MZ} yields
\begin{equation}
 \frac{V_{g,0}}{\widehat V_{g,0}}
 =1+\frac1g\left(\frac7{12}-\frac{17}{6\pi^2}\right)+O(g^{-2}).
 \label{eq:AV5}
\end{equation}
This closes \eqref{eq:AV2} and proves \eqref{eq:AV3}.
The latter concerns the uncorrected model chosen here, rather than an
obstruction to adding further asymptotic corrections.
\end{proof}

The model \eqref{eq:AV1} is explicit and contains no unknown constant.
Its action $I(Y)$ and the spectral-edge parameter $Y$ have the
matrix-model interpretation recalled in Subsection~\ref{subsec:jt-edge};
what is proved here is the uniform relative comparison with the true
volumes. The normalization is fixed at the pure-$\psi$ endpoint
(Section~\ref{sec:prefactor}); the fixed-cusp expansion of~\cite{MZ}
enters only afterwards, for $n=0$ in the proof above.

\subsection{An inverse-genus error for the volume ratios}
\label{subsec:quantitative-volume-ratios}

The sharper extraction estimate also gives a uniform rate in the
volume-ratio theorem. The two genus-exchange generating functions have
the same saddle. Adding a cusp changes the complexity by one; the
following argument controls the resulting saddle displacement even
near the endpoints.

\begin{proof}[Completion of the proof of Theorem~\ref{thm:volume}]
For $n\ge1$ put $h=2g-1$, $k=n-1$, and $h\Psi(x)=k$.
Apply Theorem~\ref{thm:sharp-extraction} at $(g,p)=(g,1)$ and
$(g-1,3)$. Their Poisson and variance factors cancel exactly, and
\eqref{eq:PF3} gives $\mathcal J_{g-1,3}/\mathcal J_{g,1}=4r$.
Consequently \eqref{eq:TL50} yields
\[
 B_{g,n}=\mathcal B(2k/h)(1+O(h^{-1})).
\]

For the cusp ratio assume first that $k\ge1$, and let $x'$ satisfy
$(h+1)\Psi(x')=k$. Write $q=\log x$, $q'=\log x'$ and
$v=\mathscr D\Psi$. The convergent endpoint expansions give
\[
 \sup_{0<x<x_c}\left\{
 \frac{\Psi(x)^2}{v(x)}+
 \left|\frac{\Psi(x)}{v(x)}\mathscr D\log v(x)\right|
 +|\beta_2(x)|\right\}<\infty.
\]
Indeed, near zero $\Psi\sim v\sim2x$; near criticality
$\Psi\sim5/(4\varepsilon)$ and $v\sim5/(4\varepsilon^2)$.
The function $\beta_2$ is bounded by \eqref{eq:PF7}.
Since $\Psi(x')/\Psi(x)=h/(h+1)$, integration with respect to
$\log\Psi$ shows that
\[
 \frac{v(x')}{v(x)}=1+O(h^{-1}),\qquad
 |q'-q|\le\frac{C\Psi(x)}{h v(x)}.
\]
For $L(q)=\log\mathcal J_{g,2}(\mathrm{e}^q)-kq$, the exact model identities give
\[
 L'(q)=(h+1)\Psi(\mathrm{e}^q)-k+\beta_2(\mathrm{e}^q)\Psi(\mathrm{e}^q).
\]
On $[q',q]$ its absolute value is at most $C\Psi(x)$, so
$|L(q')-L(q)|\le C\Psi(x)^2/(hv(x))=O(h^{-1})$.
Thus the saddle displacement and the variance factor in
\eqref{eq:CE2} each change the coefficient ratio by $1+O(h^{-1})$.
At the common point $x$, \eqref{eq:PF3} gives exactly
\[
 \frac{\mathcal J_{g,2}(x)}{\mathcal J_{g,1}(x)}
 =\frac{h}{Aar}\sqrt{\frac{h}{h+1}}.
\]
Using the coefficient normalization in \eqref{eq:TL53}, we obtain
\[
 C_{g,n}=\frac{h}{2g}\mathcal C(2k/h)(1+O(h^{-1})).
\]
At $k=0$ both saddles are zero. The conventions in
Theorem~\ref{thm:sharp-extraction} give the same two ratio formulas
directly, without logarithmic coordinates.

The positive profiles are analytic at zero, and their expansions at
infinity in \eqref{eq:SP13}--\eqref{eq:SP15} are convergent in the
inverse square root of the density. Hence
\[
 \sup_{\lambda\ge0}(1+\lambda)
 \bigl(|(\log\mathcal B)'(\lambda)|+
       |(\log\mathcal C)'(\lambda)|\bigr)<\infty.
\]
For $\lambda_*=2(n-1)/(2g-1)$ and $\lambda=n/g$,
$|\lambda_* -\lambda|/(1+\min(\lambda_*,\lambda))=O(g^{-1})$
uniformly in $n\ge1$. Integrating the preceding bound replaces
$\lambda_*$ by $\lambda$ with relative error $O(g^{-1})$.
For $n=0$, the fixed-cusp expansion of~\cite{MZ} at the finitely
many cusp counts $0,1,2$ gives the required rate directly.
Finally $D_{g,n}=C_{g,n-1}/C_{g,n}$ and the same logarithmic
derivative bound prove $\sup_{n\ge1}|D_{g,n}-1|=O(g^{-1})$.
\end{proof}

\subsection{Proof of Theorem~\ref{thm:intro-absolute}}

For $n\ge1$ the expansion is obtained by applying the extraction
theorem, Theorem~\ref{thm:canonical-expansion}, with $p=1$, keeping the
same leading term. The unmarked case $n=0$ is recovered from $n=1$ by the
exact dilaton relation, so its higher coefficients need no separate
asymptotic normalization. For $h=2g-1$, the scalar $\mathsf C_0$ of \eqref{eq:FE1} obeys
\begin{equation}
 V_{g,0}=\frac{V_{g,1}\mathsf C_0(h-1;0)}{4\pi^2(h-1)},\qquad
 \frac{\widehat V_{g,1}}{4\pi^2(h-1)\widehat V_{g,0}}
       =\sqrt{1+\frac1h}.
 \label{eq:AV6}
\end{equation}
Indeed \eqref{eq:TL12} holds at $p=0$; at $Y=0$, undo its
normalizations with $r=\pi^2/2$, $A=1/\pi^2$.
The scalar expansion of $\mathsf C_0$ follows directly from the last
line of \eqref{eq:FE5}, using the proved one-marking profile and its
summable remainder. This step does not require an unmarked harmonic
comparison. The second identity follows from \eqref{eq:AV1} exactly.

The coefficients in Theorem~\ref{thm:intro-absolute} are defined as
follows. For $n=0$, define constants $a_{0,j}$ by the finite-order formal
rule
\begin{equation}
 1+\sum_{j\ge1}a_{0,j}t^j
 =(1+t)^{1/2}\mathsf C_0\!\left(\frac t{1-t};0\right)
       \left(1+\sum_{j\ge1}\mathcal L_{1,j}(0)t^j\right).
 \label{eq:AV7}
\end{equation}
Here the second argument denotes $Y=0$, and $t=1/h$; thus the
cusp scalar uses complexity $h-1$. Put $A_j(g,0)=a_{0,j}$, and, for
$n\ge1$, put $A_j(g,n)=\mathcal A_{1,j}(Y_{g,n})$ with the coefficients
\eqref{eq:CG3} of Theorem~\ref{thm:canonical-expansion}, where
$(2g-1)\Psi(X(Y_{g,n}))=n-1$.

\begin{proof}[Proof of Theorem~\ref{thm:intro-absolute}]
For $n\ge1$, take $p=1,k=n-1$ in
Theorem~\ref{thm:canonical-expansion}. The coefficient identity \eqref{eq:TL50}
and $k!\mathfrak p_k=(k/\mathrm{e})^k$ identify $[z^k]\Phi_g^{(1)}/\mathcal Q_{g,1,k}$
with $V_{g,n}/\widehat V_{g,n}$, including $k=0$.
For $n=0$, the dilaton scalar in \eqref{eq:AV6} has a one-marking
child of complexity $h$, so exactly
\[
 \mathsf C_0(h-1;0)=\frac1{A(0)}
       \sum_{d\ge1}b_d(0)c_{g,1}(d;0).
\]
Apply \eqref{eq:FE7} at this child with $L=K+1$.
Choose its fixed $M$ large enough for the absolute exponential
$b_d$ moment. With $S=1+ML$,
\[
 \sum_d\frac{|b_d(0)|}{A(0)}(S+d)^{2L}
 \le S^{2L}\sum_d\frac{|b_d(0)|}{A(0)}\mathrm{e}^{2d/M}
 \le C D^L(L!)^2.
\]
This gives a Gevrey-2 remainder for the whole scalar sum, including
invalid degrees. Its coefficient series is precisely
$\mathsf C_0(t/(1-t);0)$: the child substitution in \eqref{eq:FE5}
becomes $(t/(1-t))/(1+t/(1-t))=t$.

The other two exact factors in \eqref{eq:AV6} are
$\sqrt{1+1/h}$, with a convergent expansion, and the $n=1$ ratio,
with the marked estimate at zero. All three exact factors are bounded.
The finite-product identity used in Proposition~\ref{prop:hierarchy}
contains one remainder per term; its factorial weights obey
$(j!)^2((L-j)!)^2\le(L!)^2$. Applied twice, it gives the same
bound for their product and exactly the coefficients in \eqref{eq:AV7}.
Take common constants with the $n\ge1$ case. The coefficient bound
$|A_j(g,n)|\le CB^j(j!)^2$ follows from those of
Theorem~\ref{thm:canonical-expansion} and Proposition~\ref{prop:hierarchy}.
Finally, for $K=\lfloor c'\sqrt{2g-1}\rfloor$ with $c'\le c$ chosen after
the final base $B$ so that $B(c')^2<1$, the inequality
$B^L(L!)^2h^{-L}\le(BL^2/h)^L$ gives the exponentially small error.
\end{proof}

In particular, the explicit first correction \eqref{eq:CG12}
gives a relative error $O(g^{-2})$ uniformly over all cusp counts.
At the unmarked endpoint $a_{0,1}=7/6-17/(3\pi^2)$, consistently
with \eqref{eq:AV5} after $h=2g-1$.
The fixed-cusp existence of an all-order expansion is the theorem
of Mirzakhani--Zograf~\cite{MZ}; here the coefficients and remainders
remain controlled as the cusp count crosses every density scale.

\section{The full diagonal chain}\label{uniform-logarithmic-volume-rate-and-the-full-diagonal-chain}

The absolute expansion assumes $g\to\infty$. To reach the bounded-genus
end of the diagonal chain, we instead sum the logarithms of the cusp
ratios of Theorem~\ref{thm:volume} and compare with the fixed-genus
asymptotics of Manin and Zograf. This proves Theorem~\ref{thm:chain}, with
$F$ as in \eqref{eq:WC1}. The result has only logarithmic precision,
although the cusp ratios are known with relative errors: summing $n$ of
their logarithms accumulates an error $n\varepsilon_g$ with
$\varepsilon_g\to0$.

We use the functions \(T,\Delta,I,J\) and the profile \(\mathcal C\)
of Section~\ref{the-stationary-parameter}, and the variable
\(t=\lambda/(2+\lambda)\in[0,1)\); thus \(Y(u)\) in \eqref{eq:WC1} is the
parameter with \(\lambda=2u/(1-u)\).

\begin{proof}[Proof of Theorem~\ref{thm:chain}]
We first give a closed form for \eqref{eq:WC1}. From \eqref{eq:SP11}
and differentiation under the integral, \[
 I'=-X'(Y)\sqrt\Delta,
 \qquad J'=-\frac{I\Delta'}{2\Delta^{3/2}}.
\] Writing \(t=\beta(Y)=X/J\), these identities show that \begin{equation}
F(t)=-2\log(2J(Y))
      -(1-t)\log\frac{4\Delta(Y)}{\pi^2}
 \label{eq:WC4}
\end{equation} has derivative \(\log(4\Delta/\pi^2)\). Its value at zero is zero,
because \(J(0)=1/2\) and \(\Delta(0)=\pi^2/4\); hence it is precisely
\eqref{eq:WC1}. As \(t\to1\), \(J\to x_c\) and \((1-t)\log\Delta\to0\), proving
its stated continuous endpoint.

Set \[
 Z_{g,n}=\frac{V_{g,n}}{(2g+n-3)!(4\pi^2)^{2g+n-3}},\qquad
 f(\lambda)=\log\frac{2\mathcal C(\lambda)}{4\pi^2(2+\lambda)}
            =-\log(2J(Y)).
\] Since \(J'>0\), \(f\) decreases from zero to \(-\log(2x_c)\) and has
finite total variation. Equation \eqref{eq:WC4} gives the exact identity \begin{equation}
\int_0^\lambda f(v)\,\mathrm{d}v
   =\left(1+\frac\lambda2\right)
       F\left(\frac\lambda{2+\lambda}\right).
 \label{eq:WC5}
\end{equation}

For \(g\ge2\), the exact adjacent-cusp relation is \[
 \log Z_{g,n}=\log Z_{g,0}
  +\sum_{j=0}^{n-1}
      \log\frac{2g C_{g,j}}{4\pi^2(2g+j-2)}.
\] Let \(\varepsilon_g\to0\) bound the logarithmic error in the uniform
\(C\) profile. Replacing \(C\) by \(\mathcal C\), and replacing
\(2g+j-2\) by \(2g+j\), gives total error
\(O(n\varepsilon_g+1+\log(1+n/g))\). The bounded variation of \(f\) gives
the uniform Riemann-sum estimate \[
 \sum_{j=0}^{n-1}f(j/g)
       =g\int_0^{n/g}f(v)\,\mathrm{d}v+O(1),
\] with a constant independent of \(n\). The fixed-cusp closed-volume
asymptotic gives \(\log Z_{g,0}=o(g)\). Using \eqref{eq:WC5}, we obtain \begin{equation}
 \log Z_{g,n}
  =\frac{2g+n}{2}F\left(\frac{n}{2g+n}\right)+o(2g+n),
 \label{eq:WC2a}
\end{equation} uniformly for all \(n\) as \(g\to\infty\).

The same statement holds uniformly over \emph{all} stable pairs with
\(2g+n\to\infty\). To check the remaining case, keep \(g\) fixed and use
the fixed-genus leading cusp asymptotic~\cite[Section~6]{ManZ}, which gives \[
 \log Z_{g,n}=-n\log(2x_c)+O_g(\log(n+1)).
\] Meanwhile \(F(1)-F(t)=O((1-t)|\log(1-t)|)\), by \eqref{eq:SP13}, so the right
side of \eqref{eq:WC2} differs from the same leading term by only
\(O_g(\log(n+1))\). Every sequence of stable pairs has a subsequence on
which the genus is fixed or tends to infinity. This proves the full
uniformity claimed for \eqref{eq:WC2}.

Along \((g,n)=(G-j,2j)\), the value \(2g+n=2G\) and the factorial
normalization are constant. Subtracting the closed value
\(\log Z_{G,0}=o(G)\) from \eqref{eq:WC2} proves \eqref{eq:WC3}, including \(j=0\) and
\(j=G\). This full-chain statement includes bounded genera, whereas the
absolute equivalent in Theorem~\ref{thm:absolute} is uniform in cusp count
as the genus tends to infinity.\qedhere
\end{proof}

\section{Uniform one-boundary coefficient ratios}\label{uniform-one-boundary-coefficient-ratios}

The coefficient problem uses the same harmonic comparison, but not the
higher-order marked expansion. Log-concavity converts the logarithmic
comparison into adjacent ratios in the interior. Near the top degree we
use an exact intersection number; near degree zero we use bracket
difference estimates. We first prove the uniform equivalent and then,
in Appendix~\ref{sec:quantitative-coefficients}, quantify its error.

Put \(N=3g-2\) and retain the defining normalization \[
 V_{g,1}(L)=\sum_{k=0}^{N}\frac{a_{g,k}}{(2k+1)!}L^{2k},
 \qquad b_{g,k}=[\tau_k]_{g,1}=4^k a_{g,k}.
\] For every valid adjacent index \(0\le k<N\), define
\(t_{g,k}\in(0,\pi)\) by \begin{equation}
\alpha_{g,k}:=\frac{2k+1}{2g-1}
 =\frac{t_{g,k}^{2}\sin t_{g,k}}
        {\sin t_{g,k}-t_{g,k}\cos t_{g,k}}.
 \label{eq:C1}
\end{equation} The assertion to prove is \begin{equation}
\sup_{0\le k<3g-2}
 \left|\frac{a_{g,k}/a_{g,k+1}}{4\pi^2/t_{g,k}^{2}}-1\right|
 \longrightarrow0.
 \label{eq:C2}
\end{equation} Equivalently, the comparison for \(b_{g,k}/b_{g,k+1}\) is
\(\pi^2/t_{g,k}^2\). The ratio at \(k=N\) is undefined and is not
included. In the boundary-length normalization of \(V_{g,1}(L)\), the bottom limit is
four, while the top adjacent ratios diverge.

\subsection{Coefficient identities and further inputs}\label{exact-dictionary-and-the-inputs-needed-here}

We first prove the uniform equivalent \eqref{eq:C2}; Appendix~\ref{sec:quantitative-coefficients} supplies the quantitative error in Theorem~\ref{thm:coefficient}. Let \[
 J_k=\int_{\overline{\mathcal M}_{g,1}}
       \kappa_1^{N-k}\psi_1^k,
 \qquad m=N-k,\quad h=2g-1.
\] The exact coefficient identities are \begin{equation}
a_{g,k}=(2k+1)!!\frac{(2\pi^2)^m}{m!}J_k,
 \qquad
 \frac{a_{g,k}}{a_{g,k+1}}
 =\frac{2\pi^2}{(2k+3)(N-k)}\frac{J_k}{J_{k+1}}.
 \label{eq:C3}
\end{equation} Note that \(J_k\) is a mixed \(\kappa_1\)--\(\psi_1\)
intersection number, not a pure \(\psi\) intersection number.

We use the following consequences of Section~\ref{uniform-comparison}.
Define \[
 H(z)=\frac{3(\sin t-t\cos t)}{t^3},\quad t=\sqrt{2z},
 \qquad \widetilde A_m(h)=[z^m]H(z)^{-h},
\] with the removable value \(H(0)=1\), and \begin{equation}
M_{g,k}=4^k\frac{3^h\Gamma(h)}{4^g}
                  (2\pi^2)^m \widetilde A_m(h).
 \label{eq:C4}
\end{equation} Thus \(H=H_0\) is the kernel of Section~\ref{the-positive-compact-family-of-kernels}
at \(Y=0\). The tilde indicates that \(\widetilde A_m(h)\) is a coefficient of the
unscaled kernel \(H_0\); thus \(\widetilde A_m(h)=r_0^{-m}A_m(h)\), with
\(A_m(h)\) from \eqref{eq:H12} at \(Y=0\) and \(r_0=\pi^2/2\).
There is an absolute \(C\) such that, for every sufficiently large
\(g\), uniformly in \(0\le k\le N\), \begin{equation}
\exp[-C\log^2(g+1)]
 \le\frac{b_{g,k}}{M_{g,k}}
 \le\exp(C\sqrt g).
 \label{eq:C5}
\end{equation} This is the case of one marking of
\eqref{eq:H42}--\eqref{eq:H44}. At \(Y=0\), the unscaled amplitude is
\(U^0=b/[4^k(2\pi^2)^m]\); undoing that scaling gives exactly
\eqref{eq:C4}, without an additional genus-dependent factor. In particular \begin{equation}
\sup_{0\le k\le N}\frac1g
       \left|\log\frac{b_{g,k}}{M_{g,k}}\right|\longrightarrow0.
 \label{eq:C6}
\end{equation} The lower estimate in \eqref{eq:C5} will also control the bottom endpoint.

We also use the following three facts, which are independent of the
harmonic comparison.

\begin{enumerate}[label=\arabic*.]
\item
  The cotangent classes are nef on the compactification. Indeed, the
  nef logarithmic relative dualizing class in~\cite[Theorem~0.4]{Keel}
  is the cotangent class at the additional mark under the identification
  of the universal curve with the moduli space with one additional mark.
  The logarithmic
  kappa class is nef as well: for the forgetful universal curve
  \(\pi:\overline{\mathcal M}_{g,n+1}\to\overline{\mathcal M}_{g,n}\),
  one has \(\kappa_1=\pi_*(\psi_{n+1}^2)\)~\cite{AC}. Its intersection against any
  complete curve is the square of the pulled-back nef cotangent class on
  the pulled-back universal surface, and hence is nonnegative. This
  argument can be carried out after a finite cover supporting a family
  and on its projective coarse surface; passage to the coarse space
  multiplies all mixed intersection numbers by the same positive factor
  and does not affect the inequality.
  Thus the nef rational divisor classes on the projective coarse moduli
  space satisfy the Khovanskii--Teissier inequalities~\cite[Section~1.6]{Laz}; see also~\cite[Theorem~2.15]{BFJ}. This gives \begin{equation}
J_k^2\ge J_{k-1}J_{k+1},\qquad 1\le k<N.
  \label{eq:C7}
\end{equation} Thus \(q_k:=J_k/J_{k+1}\) is nondecreasing. No log-concavity of the
  normalized bracket sequence is assumed.
\item
  Bracket monotonicity and the fixed-marking absolute difference estimate of Mirzakhani--Zograf~\cite[Theorem~3.2(ii)]{MZ} give, with \(c_k=b_{g,k}/V_{g,1}\), \begin{equation}
c_0=1,\qquad 0\le c_k-c_{k+1}\le C\frac{k+1}{g},
  \qquad 0\le k<N.
  \label{eq:C8}
\end{equation} For \(k\ge1\), their explicit bound is
  \(b_{g,k}-b_{g,k+1}\le16((k+1)V_{g,0}+kV_{g-1,2})\); the fixed-marking
  volume ratios turn it into \eqref{eq:C8}. The \(k=0\) case is the fixed-degree
  estimate \(1-c_1=O(1/g)\). Summing \eqref{eq:C8} gives \(1-c_k\le Ck(k+1)/g\).
  These are absolute estimates, so they cannot be divided by an
  arbitrarily small \(c_k\).
\item
  The fixed-marking leading volume asymptotic gives \begin{equation}
V_{g,1}\asymp g^{-1/2}\Gamma(2g-1)
                      (4\pi^2)^{2g-2},
  \label{eq:C9}
\end{equation} and the one-point pure-$\psi$ identity is \(J_N=1/(24^g g!)\). Only
  the upper bound in \eqref{eq:C9} is needed below. The exact value of the top
  ratio is derived below from the ordinary DVV and string equations.
\end{enumerate}

\subsection{The saddle equation for the degree}\label{the-positive-saddle-and-its-effective-density}

The reciprocal of \(H\) has strictly positive Taylor coefficients and is
analytic on a neighborhood of \([0,\pi^2/2]\). These facts also follow
directly from \(2zH''+5H'+H=0\): writing \(-H'/H=\sum_{j\ge0}r_jz^j\)
gives \(r_0=1/5\) and \[
 (2j+5)r_j=2\sum_{u+v=j-1}r_ur_v>0\quad(j\ge1).
\] Hence \(-\log H\) and \(H^{-1}\) have positive nonconstant
coefficients. There is no positive real zero of \(H\) through
\(\pi^2/2\), because \(\sin t-t\cos t>0\) for \(0<t\le\pi\); by
Pringsheim's theorem this excludes a smaller complex singularity of
\(H^{-1}\).

For \(0<z\le\pi^2/2\), tilt these coefficients to a probability law with
generating function \(H(z)/H(zw)\). Its mean is \begin{equation}
\mu(z)=-\frac{zH'(z)}{H(z)}
 =\frac12\left(3-
     \frac{t^2\sin t}{\sin t-t\cos t}\right),
 \qquad t=\sqrt{2z}.
 \label{eq:C10}
\end{equation} Its variance is \(z\mu'(z)>0\). Thus the right side of \eqref{eq:C1}, as a
function of \(t\), decreases strictly from three to zero on \((0,\pi)\),
proving existence and uniqueness in \eqref{eq:C1}. Notice the exact dimension
identity \begin{equation}
\frac mh=\frac{3-\alpha_{g,k}}2,
 \qquad \alpha_{g,k}=3-\frac{2m}{h}.
 \label{eq:C11}
\end{equation}

For \(m/h\) in a compact subset of \((0,3/2)\), there is a unique saddle
\(z\) with \(\mu(z)=m/h\). Fourier inversion for the tilted iid sum,
whose laws have uniform exponential moments and positive probabilities
at zero and one, gives \begin{equation}
\log \widetilde A_m(h)=-h\log H(z)-m\log z+O(\log h).
 \label{eq:C12}
\end{equation} The error is uniform on such compact subsets.

\subsection{Interior degrees: exponential comparison and concavity}\label{interior-degrees-exponential-comparison-plus-concavity}

Suppose \(k/g\to\alpha\in(0,3)\) and let \(z=t^2/2\) solve
\(\mu(z)=(3-\alpha)/2\). Interpolate linearly on the degree mesh the
values \begin{equation}
f_g(k/g)=\frac1g\log
               \frac{J_k}{J_Ng^{2(N-k)}}.
 \label{eq:C13}
\end{equation} By \eqref{eq:C7}, these functions are concave. Equations \eqref{eq:C3}--\eqref{eq:C6}, \eqref{eq:C12},
and Stirling's formula give their pointwise limit on \((0,3)\): \begin{equation}
\begin{aligned}
 f(\alpha)={}&-2\log H(z)-(3-\alpha)\log z
 +(3-\alpha)\log(3-\alpha)\\
 &+3\log6-\alpha\log(2\alpha)-2(3-\alpha).
 \end{aligned}
 \label{eq:C14}
\end{equation} Explicitly, the model for the intersection number \(J_k\) is
\(J_k^{\mathrm{mod}}=3^h\Gamma(h)4^{-g}m!\widetilde A_m(h)/(2k+1)!!\). Its quotient
by the top model value is exactly \(m! \widetilde A_m(h)(2N+1)!!/(2k+1)!!\); the
two occurrences of \((N-k)\log g\) in its Stirling expansion cancel the
affine normalization in \eqref{eq:C13}. The logarithmic errors in \eqref{eq:C6}, including
the top one, vanish after division by \(g\).

Differentiation of \eqref{eq:C14} gives \begin{equation}
f'(\alpha)=\log\frac{z}{2\alpha(3-\alpha)}.
 \label{eq:C15}
\end{equation} The terms containing \(dz/d\alpha\) cancel by the saddle equation.
Concavity now upgrades the exponential-scale information to an adjacent
ratio. For fixed small \(\varepsilon>0\), the slope on \([k/g,(k+1)/g]\)
is at most the secant slope on
\([\alpha-2\varepsilon,\alpha-\varepsilon]\) and at least the one on
\([\alpha+\varepsilon,\alpha+2\varepsilon]\), for all sufficiently large
\(g\). These four endpoints are fixed, so pointwise convergence
suffices. First let \(g\) tend to infinity and then let \(\varepsilon\)
decrease to zero. Since \eqref{eq:C14} is differentiable, the neighboring slope
converges to \eqref{eq:C15}. Consequently \begin{equation}
\frac{J_k}{g^2J_{k+1}}\longrightarrow
            \frac{2\alpha(3-\alpha)}z.
 \label{eq:C16}
\end{equation} Using the exact identity \eqref{eq:C3}, \begin{equation}
\frac{a_{g,k}}{a_{g,k+1}}\longrightarrow
      \frac{2\pi^2}{z}=\frac{4\pi^2}{t^2}.
 \label{eq:C17}
\end{equation} The passage from macroscopic secants to adjacent ratios uses the
concavity of $\log J_k$.

\subsection{The top endpoint, including bounded codimension}\label{the-top-endpoint-including-bounded-codimension}

The one-kappa pushforward identity gives
\(J_{N-1}=\langle\tau_2\tau_{3g-3}\rangle_g\). Applying DVV at the
degree-two insertion and then the string equation twice gives \[
 15J_{N-1}=
 \{(6g-3)(6g-5)+12g\}J_N.
\] The nonseparating term is
\((1/2)\langle\tau_0^2\tau_{3g-3}\rangle_{g-1}=12gJ_N\); separating
terms vanish by dimension. At \(g=1\), the same identity follows
directly from \(\langle\tau_0^3\rangle_0=1\). Therefore \begin{equation}
q_{N-1}=\frac{12g^2-12g+5}{5},
 \qquad
 \frac{a_{g,N-1}}{a_{g,N}}
 =\frac{2\pi^2(12g^2-12g+5)}{5(6g-3)}.
 \label{eq:C18}
\end{equation}

Let \(m=N-k\ge1\) and \(m=o(g)\). For each fixed \(\alpha<3\),
eventually \(q_{\lfloor\alpha g\rfloor}\le q_{N-m}\le q_{N-1}\) by \eqref{eq:C7}.
As \(\alpha\) increases to three, the limit on the right side of \eqref{eq:C16}
tends to \(12/5\), because \begin{equation}
\frac{t^2\sin t}{\sin t-t\cos t}
       =3-\frac{t^2}{5}+O(t^4)\quad(t\to0).
 \label{eq:C19}
\end{equation} The exact top value \eqref{eq:C18} gives the same limit. Hence \begin{equation}
\frac{q_{N-m}}{g^2}\longrightarrow\frac{12}{5},
 \qquad
 \frac{a_{g,N-m}}{a_{g,N-m+1}}
       \sim\frac{4\pi^2g}{5m}\quad(m=o(g),\ m\ge1).
 \label{eq:C20}
\end{equation} In particular the statement includes every fixed positive
codimension.
For fixed $m$, this leading term agrees, after the length normalization
above, with the fixed-codimension expansion of
Liu--Xu~\cite[Theorem~1.2]{LX}. The same monotonicity argument also covers
unbounded $m=o(g)$ and joins this endpoint to the interior limit.

\subsection{The lowest degrees}\label{the-bottom-endpoint-without-a-relative-error-assumption}

Near degree zero the interior argument gives no information: the
comparison \eqref{eq:C5} determines \(\log b_{g,k}\) only up to an error
\(o(g)\), whereas for \(k=o(g)\) the quantities to be determined vary on a
smaller scale. We combine two other inputs. For degrees up to a power of \(\log g\), the absolute estimate
\eqref{eq:C8} of Mirzakhani and Zograf gives the ratios directly. For
larger \(k=o(g)\), we argue by contradiction: by the monotonicity of the
ratios \(q_k\), a ratio bounded away from one would persist over a range
of degrees comparable to \(k\), and would force \(b_{g,k}\) to drop below
the lower bound \eqref{eq:C23}.

Set \(z_0=\pi^2/2\), and let \(S_h\) be the sum of \(h\) iid variables
with generating function \(H(z_0)/H(z_0w)\). Since \(H(z_0)=3/\pi^2\),
\eqref{eq:C4} becomes the exact identity \begin{equation}
M_{g,k}=\pi^2\Gamma(h)(4\pi^2)^{h-1}
                     \Pr(S_h=N-k).
 \label{eq:C21}
\end{equation} The summands have mean \(3/2\), nonzero variance, span one, and
exponential moments. Uniform exponential tilting in a fixed neighborhood
of that mean gives, for \(0\le k\le c_0g\) with one fixed small
\(c_0>0\), \begin{equation}
\Pr(S_h=N-k)\ge c g^{-1/2}
            \exp[-C(k+1)^2/g].
 \label{eq:C22}
\end{equation} Indeed the tilted local central limit estimate gives the factor
\(g^{-1/2}\), while the rate function has a quadratic zero at the mean;
the displacement here is exactly \(k+1/2\). Combining \eqref{eq:C5}, \eqref{eq:C9}, and
\eqref{eq:C21}--\eqref{eq:C22} yields \begin{equation}
0\le\log\frac{V_{g,1}}{b_{g,k}}
 \le C\log^2(g+1)+C\frac{(k+1)^2}{g}.
 \label{eq:C23}
\end{equation}

Write \[
 \rho_k=\frac{b_{g,k}}{b_{g,k+1}}\ge1,
 \qquad \omega_k=(2k+3)(N-k),
 \qquad \rho_k=\frac{\pi^2q_k}{2\omega_k}.
\] Suppose \(k=o(g)\) but, along a subsequence, \(\rho_k\ge1+\eta\) for
a fixed \(\eta>0\). First consider \(k>\log^3(g+1)\); then \eqref{eq:C23} is
\(o(k)\) at every index comparable to \(k\). Put
\(\delta=\eta/[2(2+\eta)]\). For \(k\le j\le\lfloor(1+\delta)k\rfloor\),
\begin{equation}
\frac{\omega_j}{\omega_k}\le\frac{2j+3}{2k+3}\le1+\delta,
 \qquad
 \rho_j\ge\rho_k\frac{\omega_k}{\omega_j}
            \ge1+\eta/2.
 \label{eq:C24}
\end{equation} These inequalities follow from the monotonicity of the raw ratios $q_j$.
Multiplying these inequalities forces
\(\log(V_{g,1}/b_{g,\lfloor(1+\delta)k\rfloor+1})\ge c_\eta k\),
contradicting \eqref{eq:C23}.

For the remaining indices \(k\le\log^3(g+1)\), summing the absolute
estimate \eqref{eq:C8} gives \(c_{k+1}=1-o(1)\) uniformly, so division in \eqref{eq:C8} is
legitimate. In fact there is a fixed small \(\delta_0>0\) for which \begin{equation}
\rho_k=1+O((k+1)/g)
       \quad(0\le k\le\delta_0\sqrt g).
 \label{eq:C25}
\end{equation} This proves \begin{equation}
\frac{a_{g,k}}{a_{g,k+1}}\longrightarrow4
                  \quad\text{whenever }k=o(g).
 \label{eq:C26}
\end{equation} No positive lower bound on \(c_k\) was assumed outside the small
square-root window.

\subsection{Combining the degree ranges}\label{uniformity-and-the-corrected-top-density}

\begin{proof}[Proof of the uniform equivalent \eqref{eq:C2}]
Take any sequence of valid indices and a subsequence on which \(k/g\)
converges in \([0,3]\). An interior limit is covered by \eqref{eq:C17}, and a
zero limit by \eqref{eq:C26}, while the comparison in \eqref{eq:C1} also tends to four. If
the limit is three, then \(m=N-k=o(g)\) and \eqref{eq:C20} applies. Equations
\eqref{eq:C11} and \eqref{eq:C19} give \[
 t_{g,k}^2\sim\frac{10m}{2g-1},\qquad
 \frac{4\pi^2}{t_{g,k}^2}\sim\frac{4\pi^2g}{5m}.
\] Thus the true and the proposed ratios are asymptotically equal on that
subsequence too. A sequence violating the supremum statement \eqref{eq:C2} would
contradict one of these three cases.
\end{proof}

The effective density \((2k+1)/(2g-1)\) in \eqref{eq:C1} is essential near the
top: it retains the exact codimension \(m\) in
\(3-\alpha_{g,k}=2m/(2g-1)\). Replacing it by \(k/g\) before taking the
endpoint limit would change the fixed-codimension denominator.

\subsection{Relation to the instanton prediction}\label{subsec:instanton-model}
The denominator \(H\) in \eqref{eq:C4} is the normalized ordinary Weil--Petersson instanton action. To see the correspondence, choose the positive sign convention
\[
 A_{\rm inst}(w)=\frac{\sin(2\pi w)-2\pi w\cos(2\pi w)}{4\pi^3}
     =\frac23 w^3 H(2\pi^2w^2).
\]
With $x=w^2$, this is the effective potential of~\cite[equation~(3.15)]{EGGLS}, whose values at its stationary points are the instanton actions. The ordinary-volume polynomial prediction in~\cite[Section~3.3]{GPRS} and the one-boundary term in~\cite[equation~(4.22)]{JR} contain, up to an overall factor independent of the coefficient degree, the contour expression
\[
 \frac{\Gamma(h)}{L}\oint
       \frac{\mathrm{e}^{Lw}}{w A_{\rm inst}(w)^h}\,\mathrm{d}w,\qquad h=2g-1.
\]
The coefficient of \(L^{2k}\), multiplied by \((2k+1)!\), is proportional to
\[
 \Gamma(h)\operatorname*{Res}_{w=0} w^{2k}A_{\rm inst}(w)^{-h}\,\mathrm{d}w.
\]
Since \(3h-2k-1=2m\), expansion at zero gives the exact identity
\[
 \Gamma(h)\operatorname*{Res}_{w=0} w^{2k}A_{\rm inst}(w)^{-h}\,\mathrm{d}w
 =\Gamma(h)(3/2)^h(2\pi^2)^m \widetilde A_m(h)
 =2\,4^{-k}M_{g,k}.
\]
Thus the model for the coefficients \(a_{g,k}\) has precisely the normalization \(4^{-k}M_{g,k}\), up to the degree-independent overall factor in the prediction. Moreover, the exponent in the \(\mathrm{d}w/w\) representation has its stationary point at
\[
 \frac{2k+1}{h}=\frac{wA_{\rm inst}'(w)}{A_{\rm inst}(w)}
 =\frac{t^2\sin t}{\sin t-t\cos t},\qquad t=2\pi w.
\]
Thus both the comparison model and the saddle equation already appear in the instanton calculus. The estimates \eqref{eq:C5} and the concavity argument above justify the asymptotic formula for adjacent ratios of the true mixed intersection numbers, including both endpoint regimes; that proof does not use the contour expression.

\section{Short curves and small eigenvalues}\label{sec:geometry}

The analytic estimates enter geometry through ratios of volumes of cut
surfaces, by Mirzakhani's integration formula~\cite{MirInv}. We first apply the
adjacent-volume limit to Shen--Wu's vanishing-gap argument~\cite{SW}
(Subsection~\ref{sec:spectral}). We then determine the short-curve process
and the systole above the scale $\sqrt g$
(Subsection~\ref{sec:pants-spectrum}) and deduce many small eigenvalues
(Subsection~\ref{sec:pants-eigenvalues}). The last two subsections state
the limit of the tight length spectrum at every cusp density and of the
full primitive spectrum at the critical scale $n\asymp\sqrt g$; their
proofs are given in Appendices~\ref{app:tight} and~\ref{app:critical}. The
critical result uses volume bounds of Mirzakhani--Zograf and Shen--Wu
and only the zero-density case of the volume-ratio limits, not the
uniform expansion.

\subsection{The first positive eigenvalue}\label{sec:spectral}

The simplest spectral consequence uses only the discrete cusp ratio
$D_{g,n-1}$. A short curve enclosing two cusps bounds a pair of pants
of fixed area $2\pi$, joined to the rest of the surface through a short
boundary. This gives a small Cheeger ratio and hence a small Rayleigh
quotient. Shen--Wu's second-moment argument makes such a curve probable.
Shen and Wu asked for the adjacent-volume limit needed to extend
this argument beyond $n=o(g)$~\cite[Question~3]{SW}. Theorem~\ref{thm:volume}
provides that limit along every sequence with $g\to\infty$.

\begin{proof}[Proof of Corollary~\ref{cor:spectral}]
Put $n=n(g)$ and
\[
 L_g=\left(\frac{\sqrt g}{n}\right)^{1/2}.
\]
Let $N_g(X)$ count the simple closed geodesics of length at most $L_g$ whose complement has components of topological types $(0,3)$ and $(g,n-1)$. Proposition~35 of~\cite{SW}, in the notation of its arXiv version~4, states under the sole assumption $n/\sqrt g\to\infty$ that
\[
 \liminf_{g\to\infty}\mathbb P^{\mathrm{WP}}_{g,n}\{N_g\ge1\}
 \ge \liminf_{g\to\infty}
 \frac{V_{g,n-1}^{2}}{V_{g,n}V_{g,n-2}}.
\]
The ratio on the right is $D_{g,n-1}$, which tends to one by Theorem~\ref{thm:volume}. Thus $\mathbb P^{\mathrm{WP}}_{g,n}\{N_g\ge1\}\to1$. The same conclusion also follows within this paper: $N_g$ is the count $N_{L_g}$ of Subsection~\ref{sec:pants-spectrum}, with $q_g$ as in \eqref{eq:PS1}, and since $L_g\to0$ and $q_gL_g^2\asymp n\sqrt g/(g+n)\to\infty$, Proposition~\ref{prop:pants-concentration} gives $N_g/(q_gL_g^2)\to1$ in probability. Its proof is the same second-moment argument, with $D_{g,n-1}\to1$ as the volume input.

Write $\mathfrak h(X)$ for the Cheeger constant. On the event $N_g\ge1$, a counted geodesic bounds a hyperbolic pair of pants of area $2\pi$; the other component has area $2\pi(2g+n-3)\ge2\pi$ for all sufficiently large $g$. This separating curve is an admissible Cheeger cut, so
\[
 \mathfrak h(X)\le \frac{L_g}{2\pi}.
\]
Let $\mathcal R(X)$ be the infimum of the Rayleigh quotient over nonzero $H^1(X)$ functions orthogonal to constants. Buser's inequality~\cite[Theorem~7.1]{Buser}, in the form recorded in~\cite[Theorem~11]{SW} for complete noncompact hyperbolic surfaces, gives $\mathcal R(X)\le c\mathfrak h(X)$ with an absolute constant $c$. Since $L_g\to0$, this upper bound is eventually smaller than both $\varepsilon$ and $1/4$. The essential spectrum starts at $1/4$; the min--max principle, as recalled in~\cite[Theorem~7]{SW}, therefore yields a positive discrete eigenvalue $\lambda_1(X)=\mathcal R(X)<\varepsilon$ on $N_g\ge1$, for all sufficiently large $g$. This proves the asserted probability limit.
\end{proof}

The use of $D_{g,n-1}$ above is the only change to the final step of Shen--Wu's proof: under their condition $n=o(g)$ this limit follows from \cite[Lemma~5.1]{MZ}, as they note. The short-curve estimate, the Cheeger cut, and the passage to a discrete eigenvalue are their geometric and spectral inputs.

\subsection{Short-curve processes and the systole}
\label{sec:pants-spectrum}

We now determine the limiting process of short curves and the systole.
The moment estimates proved here will also be used in
Subsection~\ref{sec:pants-eigenvalues} to count many disjoint short pairs of pants.
Throughout this and the next subsection,
$g\to\infty$, $n/\sqrt g\to\infty$, and $X$ has law
$\mathbb P^{\mathrm{WP}}_{g,n}$. Let $\mathcal P_2(X)$ be the set of
unoriented simple closed geodesics bounding a pair of pants with two
cusps. Define the deterministic intensity parameter
\begin{equation}
 q_g=\frac{n(n-1)}{8g\mathcal C(n/g)}.
 \label{eq:PS1}
\end{equation}
The profile endpoints in Theorem~\ref{thm:volume} imply
$q_g\asymp n^2/(g+n)\to\infty$. In particular
$q_g\sim n^2/(32\pi^2g)$ when $n/g\to0$, and
$q_g\sim x_cn/(8\pi^2)$ when $n/g\to\infty$.

We write $\operatorname{PPP}(\nu)$ for a Poisson point process with intensity measure $\nu$.

\begin{theorem}[Short-curve processes and the systole]\label{thm:pants-poisson}
\emph{(i) Curves cutting off two cusps.} The random locally finite counting measures
\begin{equation}
 \sum_{\gamma\in\mathcal P_2(X)}
       \delta_{\sqrt{q_g}\ell_\gamma}
 \ \Longrightarrow\ \operatorname{PPP}(2t\,\mathrm{d}t)
 \quad\hbox{on }[0,\infty)
 \label{eq:PS2}
\end{equation}
converge in distribution for the vague topology. Thus counts in finitely
many disjoint bounded intervals $[a_i,b_i]$ converge jointly to independent
Poisson variables with means $b_i^2-a_i^2$. If $\ell_{\min}^{(2)}$ is the
shortest length in $\mathcal P_2(X)$, with value infinity when this set is
empty, then
\[
 \mathbb P^{\mathrm{WP}}_{g,n}
 \{\sqrt{q_g}\ell_{\min}^{(2)}>t\}\longrightarrow \mathrm{e}^{-t^2}
 \qquad(t\ge0).
\]
\smallskip
\noindent\emph{(ii) All planar cusp types.} Let $\mathcal P_k(X)$ consist of simple geodesics whose
genus-zero side contains exactly $k$ cusps. If $n/g\to\lambda\in[0,\infty]$,
put $x_\lambda=X(Y_\lambda)$, where $\lambda(Y_\lambda)=\lambda$, with
$x_0=0$ and $x_\infty=x_c$. Define
\begin{equation}
 w_k(x)=\frac{2V_{0,k+1}}{k!}
             \left(\frac{x}{2\pi^2}\right)^{k-2},\qquad
 W(x)=\sum_{k\ge2}w_k(x)=\frac{2\Phi_0'(x)}{x^2}.
 \label{eq:PC1}
\end{equation}
Here $w_2(0)=W(0)=1$, $w_k(0)=0$ for $k>2$, and $W(x_c)<\infty$.
Then
\begin{equation}
 \sum_{k\ge2}\ \sum_{\gamma\in\mathcal P_k(X)}
       \delta_{(\sqrt{q_g}\ell_\gamma,k)}
 \ \Longrightarrow\
 \operatorname{PPP}\!\left(2t\,\mathrm{d}t\otimes
                    \sum_{k\ge2}w_k(x_\lambda)\delta_k\right).
 \label{eq:PC2}
\end{equation}
Convergence is for the vague topology on the product with the discrete
type space. The projection to lengths also converges, to
$\operatorname{PPP}(2W(x_\lambda)t\,\mathrm{d}t)$.

\noindent\emph{(iii) Primitive geodesics and the systole.} This is also the limit of
the process of all unoriented primitive closed geodesics, at the same
rescaling. Writing $\operatorname{sys}(X)$ for the shortest closed
geodesic length, we have
\[
 \mathbb P^{\mathrm{WP}}_{g,n}
 \{\sqrt{q_g}\operatorname{sys}(X)>t\}
 \longrightarrow \mathrm{e}^{-W(x_\lambda)t^2}\qquad(t\ge0).
\]
With probability tending to one every systolic curve is planar
separating. Select one by any measurable tie rule, and let $K_g$ be
its planar-side cusp count, with value zero if it is not planar. Then
$\bigl(\sqrt{q_g}\operatorname{sys}(X),K_g\bigr)$ converges jointly
to a pair $(Z,K)$ with
\begin{equation}
 \begin{aligned}
 \mathbb P\{Z\in \mathrm{d}t,K=k\}
 &=2t\,w_k(x_\lambda)\mathrm{e}^{-W(x_\lambda)t^2}\,\mathrm{d}t,\\
 \mathbb P\{K=k\}&=\frac{w_k(x_\lambda)}{W(x_\lambda)}
 \qquad(t>0,\ k\ge2).
 \end{aligned}
 \label{eq:FS1}
\end{equation}
\end{theorem}

The method of factorial moments for curves separating off a prescribed
number of cusps has a direct precedent in Hide--Thomas~\cite[Theorem~1.7]{HTSphere}
for punctured spheres, which also treats summed types and nested
configurations~\cite[Remark~1.9 and Proposition~3.6]{HTSphere}. Here the uniform volume ratios
give the entire family of type intensities for diverging genus at every
cusp density above the square-root scale. These planar curves
become contractible after the cusps are filled; they are excluded from
the tight spectrum of~\cite[Definition~2.1]{BL}. The all-primitive
process of Mirzakhani--Petri~\cite[Theorem~4.1]{MP} concerns closed
surfaces at fixed physical lengths. Here the full process is observed
on a shrinking scale; a tight-spectrum limit at its longer natural
scale is a separate question.

\emph{Factorial moments.} For sufficiently small $L$, the collar lemma forces all curves counted by
$N_L=\#\{\gamma\in\mathcal P_2(X):\ell_\gamma\le L\}$ to be disjoint.
Their pants interiors are also disjoint. Otherwise one pants boundary
would be an interior essential geodesic of the other pants, hence
boundary-parallel, contradicting uniqueness of geodesic representatives.
Thus their pairs of cusp labels are disjoint, and their complement has
type $(g,n-r)$ for $r$ pants. We have $g\ge2$ eventually, so the
exceptional double-pants decomposition of a four-punctured sphere does
not arise.

Mirzakhani's integration formula~\cite{MirInv,Mir}, with the same conventions as
\cite[Lemma~29]{SW}, gives, for any fixed $r$ and short length windows
$I_1,\ldots,I_r$, the ordered distinct-curve integral
\begin{equation}
 \mathbb E\!\sum_{\substack{\gamma_1,\ldots,\gamma_r\in\mathcal P_2(X)\\
                            \text{distinct}}}
       \prod_{j=1}^r\mathbf1_{\ell_{\gamma_j}\in I_j}
 =\frac{(n)_{2r}}{2^r V_{g,n}}
    \int_{I_1\times\cdots\times I_r}
       V_{g,n-r}(x_1,\ldots,x_r,0^{n-2r})
       \prod_{j=1}^r x_j\,\mathrm{d}x_j.
 \label{eq:PS3}
\end{equation}
Here $(n)_{2r}=n(n-1)\cdots(n-2r+1)$. The factor $2^{-r}$ removes
the order of the two labels within each pair. The tuple of curves is
ordered, so there is no $r!$ divisor. Each pants volume is $V_{0,3}=1$.
Mirzakhani's boundary estimate~\cite[Lemma~3.2]{Mir}, in the form
of~\cite[Lemma~14]{SW},
\begin{equation}
 1\le\frac{V_{g,m}(\mathbf x)}{V_{g,m}}
       \le\exp\!\left(\frac12\sum_jx_j\right)
 \label{eq:PS4}
\end{equation}
is uniform in the topological type. For fixed $r$, it replaces the
boundary volume in \eqref{eq:PS3} by $V_{g,n-r}(1+O_r(L))$ when
all $I_j\subset[0,L]$.

Theorem~\ref{thm:volume} also implies, for every fixed $r$,
\begin{equation}
 \frac{V_{g,n-r}}{V_{g,n}}
  =(2g\mathcal C(n/g))^{-r}(1+o(1)).
 \label{eq:PS5}
\end{equation}
Indeed this is a product of $r$ adjacent cusp ratios. Replacing $(n-j)/g$
by $n/g$ changes their profile by $1+O_r(n^{-1})$, because
$\lambda\partial_\lambda\log\mathcal C(\lambda)$ is bounded: this follows
from analyticity in the interior and the two endpoint expansions.

\begin{proof}[Proof of Theorem~\ref{thm:pants-poisson}]
We prove the planar process by bounding the tail over cusp types and
computing the mixed factorial moments for any fixed collection of types.
We then exclude the other topological types on the same length scale;
this identifies the full primitive process and the systole. Write $P=2\pi^2$, $h=2g-1$, $m=n-1$, and $F_h=\Phi_g'$.
Choose $x\in(0,x_c)$ by $h\Psi(x)=m$ and put $q_g^*=nx/(4P)$.
The effective density is $2(n-1)/(2g-1)$. The boundedness of
$\lambda\partial_\lambda\log X(Y_\lambda)$, including its endpoint limits, implies
$q_g^*/q_g\to1$ and $q_g^*\to\infty$.

\emph{Step 1: a uniform tail bound over cusp types.}
We first show that curves cutting off a large number of cusps contribute
negligibly, uniformly along the sequence. The input is the
coefficient-ratio bound \eqref{eq:PC3}, which compares volumes whose cusp
counts differ by $j$.
By the positive convolution \eqref{eq:TL33} and
Lemma~\ref{lem:point-mass}, every point mass is at most $C/\sigma$,
where $\sigma^2=h\mathscr D\Psi(x)$. The local limit
\eqref{eq:TL47} gives a lower bound $c/\sigma$ at $m$.
Consequently, for all sufficiently large $h,m$,
\begin{equation}
 \frac{[z^{m-j}]F_h(z)}{[z^m]F_h(z)}
 =x^j\frac{\mathbb P_x(N=m-j)}{\mathbb P_x(N=m)}
 \le Cx^j\qquad(0\le j\le m).
 \label{eq:PC3}
\end{equation}
Only the denominator uses the local limit; the numerator may be at any
index, including zero.

Let $N_{k,L}$ count curves of type $k$ with length at most $L$, and put
$a_k=V_{0,k+1}/(k!P^{k-2})$. Mirzakhani integration gives
\begin{equation}
 \begin{aligned}
 \mathbb E N_{k,L}=\frac{\binom nk}{V_{g,n}}
 \int_0^L\ell\,V_{0,k+1}(\ell,0^k)
                  V_{g,n-k+1}(\ell,0^{n-k})\,\mathrm{d}\ell,\\
 \binom nk\frac{V_{0,k+1}V_{g,n-k+1}}{V_{g,n}}
 =\frac nP a_k\frac{[z^{n-k}]F_h}{[z^{n-1}]F_h}.
 \end{aligned}
 \label{eq:PC4}
\end{equation}
The two sides have different genera, so there is no side-exchange
factor. The coefficient identity follows from \eqref{eq:TL50}; it
holds up to $k=n$. Bound \eqref{eq:PS4} changes the integrand's two
volumes by a factor between one and $\mathrm{e}^L$, independently of $k$.
Using \eqref{eq:PC3} with $j=k-1$ therefore gives, for each fixed $T$,
\begin{equation}
 \mathbb E\sum_{k\ge K}N_{k,T/\sqrt{q_g^*}}
 \le CT^2\sum_{k\ge K}a_kx_c^{k-2}\longrightarrow0
 \quad(K\to\infty),
 \label{eq:PC5}
\end{equation}
uniformly along the sequence. To justify summability directly, the
genus-zero string solution gives $\Phi_0'''(x)=1/a(Y(x))=Y'(x)$.
Since $V_{0,3}=1$, its zero initial conditions give
$\Phi_0''=Y$ and $\Phi_0'(x)=\int_0^xY(u)\,\mathrm{d}u$. This is finite at
$x_c$; positivity gives convergence of the coefficient series there.
Its initial term $x^2/2$ also gives $W(0)=1$.

\emph{Step 2: fixed types and nested configurations.}
Next we compute the mixed factorial moments for finitely many fixed cusp
types. The main contribution comes from configurations whose planar
sides are pairwise disjoint, which we call \emph{stars}; nested
configurations are shown to be negligible.
Fix $r$ ordered distinct curves of prescribed types $k_1,\ldots,k_r$,
and write $K=\sum k_i$. The collar lemma makes their boundaries
disjoint. For a star, the ordered
cusp-label multiplicity is $(n)_K/\prod_i k_i!$, and the complement
has type $(g,n-K+r)$. There is no $r!$ divisor. The boundary error is
$1+O_r(L)$, as each cut occurs twice in \eqref{eq:PS4}.
For fixed $d$, \eqref{eq:PS5} and $q_g^*/q_g\to1$ give
$V_{g,n-d}/V_{g,n}=(x/(Pn))^d(1+o(1))$.
Thus for windows $[a_i,b_i]/\sqrt{q_g^*}$ the star contribution tends to
\begin{equation}
 \prod_{i=1}^r
       w_{k_i}(x_\lambda)(b_i^2-a_i^2).
 \label{eq:PC6}
\end{equation}
The positive expression has a relative error before taking the limit,
even when some limiting weights vanish.

Filling cusps turns each planar side into a disk. Disjoint boundaries
therefore have nested or disjoint disks; their cut graph is a tree with
one genus-$g$ root and $r$ planar vertices. Let $d$ be the root degree,
$a_v$ the number of cusps at a planar vertex and $d_v$ its degree;
put $a=\sum_v a_v$. Prescribed types bound $a$ by $K$.
There are finitely many such trees with ordered edges and cusp
allocations, depending only on the fixed types. Classification of
surfaces gives the corresponding mapping-class orbits with bounded
symmetry factors; assigning cusp labels costs at most $Cn^a$.
The root volume ratio is $O((x/(Pn))^{a-d})$, the planar volumes
are fixed, and the $r$ length integrals cost $O((nx)^{-r})$.
Each cut type contributes at most
\begin{equation}
 Cn^{d-r}x^{a-d-r}=O(n^{-1})\quad\hbox{if it is not a star}.
 \label{eq:PC7}
\end{equation}
Indeed stability gives $a_v+d_v\ge3$ and
$\sum_vd_v=2r-d$, hence $a\ge r+d$. The power of $x$ is
nonnegative, so no loss occurs as $x\to0$. A nonstar has $d\le r-1$;
if $d=r$, every planar vertex is a leaf and the configuration is a star.
This accounts for every fixed-type mixed factorial moment.

The limits \eqref{eq:PC6} are the mixed factorial moments of independent
Poisson counts. Higher fixed moments give uniform integrability, and
the product Poisson law is moment determinate. First moments give
tightness on compact sets in the marked space. This proves
\eqref{eq:PC2}. Its unmarked projection requires more than vague
convergence: truncate the types, pass to the finite-type limit, then
let the cutoff grow using \eqref{eq:PC5} and the convergent limiting
series. This proves the stated unmarked process. Finally replace
$q_g^*$ by $q_g$. The type-two restriction gives \eqref{eq:PS2}
without assuming a limit for $n/g$, by subsequences and compactness of
$[0,\infty]$. Its zero-count probabilities give the pants minimum law.

\emph{Step 3: the remaining curves and the systole.}
It remains to show that, on the scale $q_g^{-1/2}$, there are no short
curves other than planar separating ones: nonseparating curves, and
separating curves with positive genus on both sides, are rare.
Nonseparating simple curves of length at most $L\le L_0$ have expected
count at most $C\mathrm{e}^LL^2V_{g-1,n+2}/V_{g,n}\le CL^2$, by integration,
\eqref{eq:PS4}, and the uniform lower bound for $B_{g,n}$.
To sum every positive-genus separating cut, put
$S_g(z)=\sum_{i=1}^{g-1}\Phi_i'(z)\Phi_{g-i}'(z)$. Exact positivity
\eqref{eq:TL33} includes genus one, so the child powers add to $h-1$:
\begin{equation}
 S_g(z)=V(z)^{h-1}Q_g(z),\qquad Q_g\ge0
 \ \hbox{coefficientwise},\qquad
 \frac{S_g(x)}{F_h(x)}\le\frac{Ca(Y)r(Y)^2}{g^2}.
 \label{eq:FS2}
\end{equation}
The last bound is \eqref{eq:TL26}: undoing \eqref{eq:TL9} leaves
exactly $ar^2$, as the child complexities sum to $2g-2$ and their
deficits to $3g-4$. In particular bounded child genera are retained.
Lemma~\ref{lem:point-mass}, applied to the positive factor $V^{h-1}$,
gives $\sup_j\mathbb P_x^{S}(N=j)\le C/\sigma$. Against the parent lower
bound $\mathbb P_x^{F}(N=n-1)\ge c/\sigma$, the standard deviations
cancel. The coefficient identity \eqref{eq:TL50} therefore yields
\begin{equation}
 \begin{aligned}
 \mathcal A_{g,n}
 &:=\sum_{i=1}^{g-1}\sum_{k=0}^n\binom nk
       \frac{V_{i,k+1}V_{g-i,n-k+1}}{V_{g,n}}
   =\frac1P\frac{[z^n]S_g(z)}{[z^n]\Phi_g(z)}\\
 &=\frac{n}{Px}\frac{S_g(x)}{F_h(x)}
       \frac{\mathbb P_x^S(N=n)}{\mathbb P_x^F(N=n-1)}
   \le\frac{Cnar^2}{xg^2}\le\frac Cg.
 \end{aligned}
 \label{eq:FS3}
\end{equation}
Indeed $(n-1)/x=h/(Aar)$, while $A$ is bounded below and $r$ above.
Thus no small-density factor is lost. Each separating orbit with two
positive genera occurs in this ordered sum. Its integration symmetry
and one-handle factors are at most one; applying \eqref{eq:PS4} to
both components bounds its total expected count by
$C\mathrm{e}^LL^2\mathcal A_{g,n}\le CL^2/g$. This includes all cusp partitions.

For fixed $T$, take $L=T/\sqrt{q_g}$. Markov's inequality shows that
with probability $1-O(T^2/q_g)$ every simple geodesic in this window
is planar. Primitive geodesics shorter than $2\operatorname{arcsinh}1$
are simple, by the standard collar theorem used in
\cite[Section~1.1]{HTSphere}. Since $q_g\to\infty$, the full primitive
process and planar process agree on each bounded rescaled window
with probability tending to one. This transfers the process limit
and its zero-count law to the systole.

Finally restrict the first point to $[\eta,T]$ and finitely many cusp
types. The first-point map is continuous when its first length is
unique and avoids the endpoints, as happens almost surely in the
limiting Poisson process. Let the type cutoff grow using
\eqref{eq:PC5}, then $\eta\downarrow0$ and $T\to\infty$.
The omitted length probabilities tend to zero by the zero-count law;
$W\ge1$. Nonplanar competitors are excluded on each fixed window.
This also makes any tie rule immaterial and proves planarity of every
systolic curve with probability tending to one. The first-arrival
density of the limiting marked process is \eqref{eq:FS1}.
\end{proof}

\subsection{Many small eigenvalues}\label{sec:pants-eigenvalues}

\begin{proposition}[Many short pants]\label{prop:pants-concentration}
If $L=L_g\to0$ and $q_gL_g^2\to\infty$, then
\begin{equation}
 \frac{N_{L_g}}{q_gL_g^2}\longrightarrow1
 \quad\hbox{in probability}.
 \label{eq:PS6}
\end{equation}
Consequently, for each fixed $\varepsilon\in(0,1)$, with probability
tending to one there are at least
$\lfloor(1-\varepsilon)q_gL_g^2\rfloor-1$ positive discrete eigenvalues
in $(0,6L_g]$, counted with multiplicity.
\end{proposition}
\begin{proof}
The first two cases of \eqref{eq:PS3} give
\[
 m_L:=\mathbb E N_L=\frac{(n)_2L^2}{4}
       \frac{V_{g,n-1}}{V_{g,n}}(1+O(L))\sim q_gL^2,
\]
and
\[
 \frac{\mathbb E(N_L)_2}{m_L^2}
 =\frac{(n-2)(n-3)}{n(n-1)}D_{g,n-1}^{-1}(1+O(L))\longrightarrow1.
\]
Since $m_L\to\infty$, the relative variance tends to zero, proving
\eqref{eq:PS6}. This second-moment argument does not use the point-process
limit at a window growing with $g$.

Hide--Thomas~\cite[Lemma~5.3]{HTSphere} prove that $N$ such pants
boundaries of length at most $L$ give $N-1$ positive eigenvalues below
$6L$, provided $6L<1/4$. For the reader's convenience we recall the
variational argument.
Inside each pants choose a function equal to inward distance $\rho$ on
the half-collar $0\le\rho\le1$, equal to one beyond it, and zero outside
the pants. For all sufficiently small boundary lengths $\ell$, the
collar lemma supplies this embedded half-collar. The function lies in
$H^1(X)$, and
\[
 \int|\nabla f|^2=\ell\sinh1,\qquad
 \int|f|^2\ge2\pi-\ell\sinh1.
\]
Different pants have disjoint supports up to their boundaries, so the
span of these functions has dimension $N$ and Rayleigh quotient at most
$6L$ for small $L$. Its intersection with the orthogonal complement of
constants has dimension at least $N-1$. Since the essential spectrum
starts at $1/4$, the form version of min--max gives the claimed positive
discrete eigenvalues. Apply \eqref{eq:PS6} and let $g$ grow.
\end{proof}

\begin{corollary}[A growing number of small eigenvalues]\label{cor:many-spectrum}
Suppose $g\to\infty$ and $n/\sqrt g\to\infty$. If $j=j(g)\to\infty$
and $j=o(q_g)$, then
\[
 \mathbb P^{\mathrm{WP}}_{g,n}
 \{\lambda_j(X)\text{ exists and }\lambda_j(X)\le12\sqrt{j/q_g}\}
 \longrightarrow1,
\]
where positive discrete eigenvalues are ordered with multiplicity.
\end{corollary}

\begin{proof}[Proof of Corollary~\ref{cor:many-spectrum}]
For $j\to\infty$ with $j=o(q_g)$, choose $L_g=2\sqrt{j/q_g}$.
Then $L_g\to0$, $q_gL_g^2=4j$, and \eqref{eq:PS6} implies
$N_{L_g}\ge j+1$ with probability tending to one. The preceding
variational estimate gives $\lambda_j\le6L_g=12\sqrt{j/q_g}$.
\end{proof}

The spectral mechanism is established in the fixed-genus work of
Hide--Thomas~\cite[Theorem~1.10 and Remark~5.5]{HTSphere}; their later
bordered-volume asymptotics give linear counts below each fixed positive
threshold at fixed genus~\cite[Theorem~1.1]{HTLargeN}.
Their deterministic result~\cite[Theorem~1.3]{HTMany}, in the July 2026
revision, gives a linear number below $1/4$ whenever $g<an$ for fixed
$a>0$. Our threshold tends to zero and the genus grows. In the range
$\sqrt g\ll n\ll g$, the corresponding second-moment volume inputs
already follow from Shen--Wu~\cite{SW}. The extension supplied here is
the density-dependent intensity and removal of the upper restriction on
$n/g$. No sharp eigenvalue asymptotic, matching lower bound, or
Poisson eigenvalue law is asserted.


\subsection{The tight length spectrum}\label{subsec:tight-spectrum}

We now consider curves that remain noncontractible after filling all
cusps. The corresponding tight Weil--Petersson volumes were introduced
by Budd~\cite{Budd}. Following Budd and Lions~\cite{BL}, a closed geodesic is
\emph{tight} if it is the unique shortest representative of its free
homotopy class on the filled surface. The distinction matters: several
ordinary geodesics can belong to the same class after the cusps are
filled. The tight integration formula and the geometric argument used
in the proof are due to~\cite{BL}; the new input is a volume estimate that
allows arbitrary prescribed cusp counts.

Let $Y_\lambda$ satisfy $\lambda(Y_\lambda)=\lambda$,
as in Proposition~\ref{prop:bounded-boundaries}, and set
\begin{equation}
 \alpha(\lambda)=\frac{\sqrt{2r(Y_\lambda)}}\pi,
 \qquad \nu(\mathrm dt)=\frac{\cosh t-1}{t}\,\mathrm dt.
 \label{eq:tight-scale}
\end{equation}
Here $\alpha(0)=1$, and the density of $\nu$ is defined to be zero
at the origin.

\begin{theorem}[Tight lengths at arbitrary cusp density]
\label{thm:tight-spectrum}
Let $n=n(g)\ge0$ be any integer sequence, and let $X$ have law
$\mathbb P^{\mathrm{WP}}_{g,n}$. As $g\to\infty$, the point process
\begin{equation}
 \sum_{\substack{\gamma\text{ primitive, unoriented}\\
                  \gamma\text{ tight on }X}}
 \delta_{\alpha(n/g)\ell_X(\gamma)}
 \ \Longrightarrow\ \operatorname{PPP}(\nu)
 \label{eq:tight-process}
\end{equation}
converges in distribution for the vague topology on $[0,\infty)$.
For every fixed $L<\infty$, with probability tending to one all the
curves counted in $[0,L]$ are simple, pairwise disjoint, and have
connected complement after filling the cusps.
\end{theorem}

In particular, no limit of $n/g$ is required. At large density,
the endpoint expansion of $r$ gives
\[
 \alpha(\lambda)\sim
 \frac{\sqrt{3j_{0,1}\sqrt5}}\pi\lambda^{-1/4}.
\]
This agrees with the scale in~\cite[Corollary~1.3]{BL}.
Budd and Lions asked whether their limit persists for other cusp
weights and under the fixed-count measures
$\mathbb P^{\mathrm{WP}}_{g,n}$~\cite[Open Question~1.6]{BL}.
Theorem~\ref{thm:tight-spectrum} treats the fixed-count part of this
question; Corollary~\ref{cor:tight-boltzmann} below treats the
remaining cusp-weight regimes. The planar curves in
Theorem~\ref{thm:pants-poisson} are contractible after filling the
cusps and therefore do not occur in \eqref{eq:tight-process}.

The proof follows the moment method of Budd and Lions. Their tight
integration formula~\cite[Proposition~3.2]{BL} expresses factorial moments
of tight curve counts through tight volumes of the cut surfaces. We first
write the generating functions of tight volumes through the shifted
amplitudes of Section~\ref{exact-amplitudes-and-the-harmonic-model}
(Lemma~\ref{lem:tight-gf}); then extract the coefficient at the
prescribed cusp count, as in Section~\ref{extracting-fixed-cusp-counts}
(Lemma~\ref{lem:tight-connected}); then show that configurations with
disconnected complement are negligible (Lemma~\ref{lem:tight-separating});
and finally identify the Poisson limit (Subsection~\ref{subsec:tight-proof}).
These steps are carried out in Appendix~\ref{app:tight}.

\emph{Random cusp counts.}
For $0\le\mu<\mu_{\rm cr}=x_c/(2\pi^2)$, the Boltzmann measure
$\mathbb P^{\mathrm{WP}}_{g,\mu}$ first chooses the cusp count $N$
with probabilities proportional to $\mu^nV_{g,n}/n!$, then samples
$\mathbb P^{\mathrm{WP}}_{g,n}$ conditionally on $N=n$.
Its probability generating function is
$\Phi_g(2\pi^2\mu z)/\Phi_g(2\pi^2\mu)$, as in~\cite{BL}.

\begin{corollary}[Arbitrary subcritical cusp weights]
\label{cor:tight-boltzmann}
For any sequence $0\le\mu_g<\mu_{\rm cr}$, put
\[
 a_g=\frac{\sqrt{2r(Y(2\pi^2\mu_g))}}\pi.
\]
Under $\mathbb P^{\mathrm{WP}}_{g,\mu_g}$, the primitive tight
lengths multiplied by $a_g$ converge to $\operatorname{PPP}(\nu)$.
In particular, if $\mu_g\to\mu<\mu_{\rm cr}$, the unscaled lengths have
limiting intensity $(\cosh(a_\mu t)-1)\,\mathrm dt/t$, where
$a_\mu=\sqrt{2r(Y(2\pi^2\mu))}/\pi$. If $\mu_g\to\mu_{\rm cr}$, no condition on
the speed of convergence is required.
\end{corollary}
\noindent The proof is given in Subsection~\ref{subsec:tight-boltzmann-proof}.


\subsection{The full length spectrum at the critical cusp scale}
\label{sec:critical-cluster}

The tight spectrum omits curves that become contractible when the cusps
are filled. At $n\asymp\sqrt g$, these curves contribute on the same
length scale as the nonseparating curves. Their contribution is not
Poisson: a two-cusp pair of pants contributes its whole length spectrum,
not only its boundary.

For $s>0$, let $P_s$ be the hyperbolic pair of pants with two cusps and
one geodesic boundary of length $s$. Write $K_s$ for its primitive
unoriented length counting measure, including the boundary once and
excluding parabolic classes. Equal lengths are counted with their
geometric multiplicities. Let
\begin{equation}
 \mu_c(\mathrm ds)=\frac{c^2}{8\pi^2}\sinh(s/2)\,\mathrm ds,
 \qquad \nu(\mathrm dt)=\frac{\cosh t-1}{t}\,\mathrm dt.
 \label{eq:cluster-intensities}
\end{equation}

\begin{theorem}[Critical primitive length spectrum]
\label{thm:critical-cluster}
Suppose $g\to\infty$ and $n/\sqrt g\to c\in(0,\infty)$, and let
$X$ have law $\mathbb P^{\mathrm{WP}}_{g,n}$. Then
\begin{equation}
 \sum_{\substack{\gamma\text{ primitive}\\\text{unoriented on }X}}
       \delta_{\ell_X(\gamma)}
 \ \Longrightarrow\ \Xi_c:=\mathcal N+\sum_{s\in\mathcal Q_c}K_s,
 \label{eq:critical-cluster}
\end{equation}
as locally finite counting measures on $[0,\infty)$ with the vague
topology. Here $\mathcal N$ and $\mathcal Q_c$ are independent Poisson
point processes of intensities $\nu$ and $\mu_c$. In particular, for
every nonnegative continuous compactly supported function $f$,
\begin{equation}
 \mathbb E\mathrm e^{-\langle f,\Xi_c\rangle}
 =\exp\left\{-\int(1-\mathrm e^{-f(t)})\,\nu(\mathrm dt)
 -\int(1-\mathrm e^{-\langle f,K_s\rangle})\,\mu_c(\mathrm ds)\right\}.
 \label{eq:cluster-laplace}
\end{equation}
\end{theorem}

Shen and Wu obtained the factorial-moment asymptotics for short
two-cusp pants and the limiting probability that such a curve occurs
at this cusp scale~\cite[Lemmas~28 and~33]{SW}. The corresponding Poisson mean is
$\mu_c([0,L])=c^2(\cosh(L/2)-1)/(4\pi^2)$ on their stated short
window. Theorem~\ref{thm:critical-cluster} concerns the full primitive
process on every fixed length window, including the correlated lengths
inside each pair of pants. The background intensity and the geometric
moment method come from Mirzakhani and Petri~\cite{MP}. The critical
regime needs the volume bounds of Mirzakhani--Zograf and
Shen--Wu~\cite{MZ,SW} together with the zero-density case of
Theorem~\ref{thm:volume} and Proposition~\ref{prop:bounded-boundaries}:
at $n\asymp\sqrt g$ the earlier estimates bound the genus ratio
$V_{g-1,n+2}/V_{g,n}$ only from above. The full uniform expansion is not
used.

The nonsimple part distinguishes this regime from closed surfaces.
He, Shen, Wu and Xue proved that the nonsimple systole of a
Weil--Petersson random closed surface is asymptotic to $\log g$,
with a more precise concentration about $\log g-\log\log g$~\cite[Theorem~1.1]{HSWX}.
Thus nonsimple curves leave every
fixed length window in that model. At the critical cusp scale here,
two-cusp pants instead retain nonsimple curves of bounded length.

\emph{Outline of the proof.} Short multicurves are organized by their cut
graphs. Using only the volume bounds of~\cite{MZ,SW},
Lemma~\ref{lem:critical-cuts} shows that in the limit a cut graph
contributes only if it has one component carrying the genus, possibly
with nonseparating loops, and two-cusp pairs of pants attached to it as
leaves. A bounded-multiplicity argument (Lemmas~\ref{lem:primitive-count}
and~\ref{lem:critical-complex}) reduces intersecting simple curves and
nonsimple curves to such cut graphs; hence, with probability tending to
one, every short nonsimple curve lies in a two-cusp pair of pants. The mixed
factorial moments of nonseparating lengths and pants waists converge to those
of independent Poisson processes (Lemma~\ref{lem:critical-skeleton}, where the
zero-density volume ratios enter), and attaching to each waist the full
spectrum of its pair of pants gives $\Xi_c$. The details are given in
Appendix~\ref{app:critical}.

The limit is not Poisson. If $m_L(s)=K_s([0,L])$, then
Lemma~\ref{lem:primitive-count} and the support restriction $s\le2L$
give finite moments and
\begin{equation}
 \operatorname{Var}\Xi_c([0,L])-\mathbb E\Xi_c([0,L])
 =\int m_L(s)(m_L(s)-1)\,\mu_c(\mathrm ds).
 \label{eq:cluster-overdispersion}
\end{equation}
Besides its waist, $P_s$ contains the primitive figure-eight around
its two cusps, of length $2\operatorname{arccosh}(2+\cosh(s/2))$
by \cite[Theorem~1.3]{Rivin}. For $L>2\operatorname{arccosh}3$,
both lengths lie below $L$ on an interval of positive $\mu_c$-measure.
The right side of \eqref{eq:cluster-overdispersion} is therefore
strictly positive. The theorem identifies these correlations in the
limit; no convergence of the full-spectrum second moments is asserted.


\section{Concluding remarks and open questions}

The results above show that one real parameter, determined by the cusp
density, governs the large-genus behavior of $V_{g,n}$ in every regime.
As $n/g$ runs from zero to infinity, the saddle $X(Y_{n/g})$ moves from
the fixed-cusp regime of Mirzakhani and Zograf to the Bessel singularity
$x_c$ of the fixed-genus theory of Manin and Zograf. The volume ratios,
the boundary factor $\mathfrak b_Y$, the short-curve intensities $w_k$
and the tight-length scale $\alpha$ interpolate in the same way.

The coefficients of the normalized volume expansion are analytic in
$Y$ on a neighborhood of the compact interval $[0,y_0]$. This is the
precise sense in which the two endpoint regimes fit into one expansion.
It avoids the loss of uniformity of the fixed-$n$ expansion at
$n\asymp\sqrt g$, while retaining an error bound independent of $n$.

The short geometry nevertheless changes at the scale $\sqrt g$.
For $n=o(\sqrt g)$, Shen and Wu proved a uniform spectral
gap~\cite[Theorem~3]{SW}. At $n/\sqrt g\to c\in(0,\infty)$, the
primitive length spectrum is a Poisson background together with
independent clusters formed by two-cusp pairs of pants
(Theorem~\ref{thm:critical-cluster}). Above this scale, short curves
bounding genus-zero subsurfaces containing cusps have the statistics of
Theorem~\ref{thm:pants-poisson}, and the spectral gap closes
(Corollary~\ref{cor:spectral}).

The proof rests on three features that may be useful beyond
Weil--Petersson volumes. Along the string flow, the alternating
$\kappa_1$--$\psi$ formula becomes a recursion with nonnegative
coefficients, so that comparison and monotonicity arguments hold
uniformly in the deformation parameter. The deformation ends at pure
$\psi$ intersection numbers, where Painlev\'e~I fixes the integration
constants without any resurgence assumption. Finally, the positive
convolution factor of Lemma~\ref{lem:convolution} turns estimates for
generating functions into a local limit theorem that holds uniformly
from bounded cusp counts up to the critical cusp weight. The geometric
consequences combine these volume estimates with the integration
formulas and short-curve arguments credited in
Section~\ref{sec:geometry}.

\subsection{Cusp weights and the JT spectral edge}\label{subsec:jt-edge}

Saad, Shenker and Stanford related the genus expansion of
Jackiw--Teitelboim (JT) gravity to a matrix integral, with
Weil--Petersson volumes as its geometric input~\cite{SSS}.
The parameter $Y$ describes a spectral edge in its defect deformation,
introduced by Maxfield--Turiaci and by Witten~\cite{MT,WittenDef}.
We record the normalization to relate the
uniform estimates to that model. A defect of opening angle $2\pi\alpha$
is a cusp at $\alpha=0$~\cite[Section~2]{ET}.
Let $\lambda_{\mathrm{def}}$ be its weight, distinct from the cusp
density $\lambda=n/g$, and use the convention $2\gamma=1$ of
Maxfield--Turiaci~\cite{MT}. Their string equation at $\alpha=0$
and its branch connected to the undeformed edge are
\[
 \frac{\sqrt{E_0}}{2\pi}I_1(2\pi\sqrt{E_0})
       +\lambda_{\mathrm{def}}=0,\qquad
 E_0=-\frac{Y}{\pi^2},\qquad x=2\pi^2\lambda_{\mathrm{def}}.
\]
The Bessel expressions are interpreted by their entire power series.
Substitution gives exactly $X(Y)=x$. Moreover, the definition
\eqref{eq:H8} gives, for $g\ge2$,
\[
 \sum_{n\ge0}V_{g,n}\frac{\lambda_{\mathrm{def}}^n}{n!}
       =(2\pi^2)^{3g-3}\Phi_g(x).
\]
Thus a fixed cusp count is a coefficient of the defect-weighted
amplitude, whereas $Y$ parametrizes its spectral edge.

The genus-zero density, with its overall topological factor removed,
can be written using our kernel. Specializing the Abel integral in
\cite[equation~(3.50)]{MT} to $\alpha=0$ gives
\[
 \rho_0(E)=\frac1{4\pi}\int_{E_0}^{E}
        \frac{I_0(2\pi\sqrt u)}{\sqrt{E-u}}\,\mathrm du
   =\frac{\sqrt q}{2\pi}K_Y(-2\pi^2q),\qquad q=E-E_0\ge0.
\]
For the second equality set $u=E_0+q(1-s^2)$ and use
\eqref{eq:P1}. At $Y=0$ this is the JT density
$\sinh(2\pi\sqrt E)/(4\pi^2)$. At $Y=y_0$, put
$b=-X''(y_0)>0$. Since $a(y_0)=0$, expansion of $K_{y_0}$ gives
\[
 \rho_0(E_0+q)=\frac{\pi b}{3}q^{3/2}+O(q^{5/2}).
\]
The square-root edge has become a cubic critical edge, the local
geometry associated with the $(2,3)$ critical point. This identifies
the physical meaning of the same endpoint at which the normalized
intersection recursion reduces to pure $\psi$ intersections.

The exponent in the large-genus action agrees with this critical
geometry. With $\delta=y_0-Y$, the endpoint
expansions of Section~\ref{the-stationary-parameter} give
\[
 x_c-x=\frac b2\delta^2(1+O(\delta)),\qquad
 \mathcal Z=\frac{2\sqrt3\,b}{5}\delta^{5/2}(1+O(\delta))
       =\frac{2^{9/4}\sqrt3}{5b^{1/4}}
          (x_c-x)^{5/4}\bigl(1+O(\sqrt{x_c-x})\bigr).
\]
The marked expansion in Theorem~\ref{thm:full-expansion} uses the
full action $\mathcal Z(Y)$ and remains uniform as $Y$ approaches
$y_0$. Replacing the action by its leading critical power would instead
introduce a multiplicative factor $\exp(O(g\sqrt{x_c-x}))$ for fixed
marking count. That replacement is not uniform in an arbitrary joint
limit. In formal genus counting, the corresponding balance is between
the string coupling $\hbar$ and $(x_c-x)^{5/4}$; the genus index $g$
is not the inverse string coupling. A double-scaled sum over all genera
requires a separate summation prescription and estimates beyond the
large-genus coefficient asymptotics proved here.


\subsection{Open questions}
We conclude with some questions that the present methods leave open.

\emph{Actual Gevrey-1 remainders.} The most direct analytic strengthening would replace the squared factorial in Theorem~\ref{thm:intro-absolute} by a single factorial, with constants uniform in the cusp count and truncation order $K\le cg$. A sufficiently small linear truncation would then give error $O(\mathrm e^{-ag})$. The pure-$\psi$ endpoint already has this order of remainder (Subsection~\ref{sec:endpoint-expansion}), but transporting it across the whole parameter interval requires sharper estimates for the actual recursion. In the present proof the squared factorial comes from the degree paths in the profile recursion (Subsection~\ref{subsec:all-orders}); bounds on the formal coefficients alone would not settle the question, since the actual remainders must be controlled as well.

\emph{Rate for adjacent ratios.} Theorem~\ref{thm:coefficient} has the
uniform error $O(g^{-1/8})$, and $O_\varepsilon(g^{-1/4})$ for
$k\ge\varepsilon g$; neither exponent is claimed to be optimal
(Appendix~\ref{sec:quantitative-coefficients}). What is the optimal
uniform rate?

\emph{Other backgrounds.} The derivation of the recursion
\eqref{eq:H9} uses only the string, dilaton and
Dijkgraaf--Verlinde--Verlinde equations, and applies to other
backgrounds of Kontsevich--Witten times; the higher Weil--Petersson
volumes of Kaufmann, Manin and Zagier arise in this way~\cite{KMZ,ManZ}.
For which backgrounds does the deformed recursion have nonnegative
coefficients, and does the uniform analysis of this paper then extend?

\emph{Several boundary components.} Can one obtain relative asymptotics
for mixed intersection numbers when both the number of nonzero degrees
and their sum grow with the genus, uniformly over their allocation among
the markings? Proposition~\ref{prop:bounded-boundaries} treats a fixed
number of bounded nonzero lengths; it does not provide this control of
individual mixed intersections.

\emph{Geometry and spectrum.} Theorem~\ref{thm:tight-spectrum} determines
the tight spectrum at every cusp density, while
Theorem~\ref{thm:critical-cluster} identifies the pants clusters in the
ordinary spectrum when $n/\sqrt g$ has a positive finite limit.
What is the joint distribution
of ordinary geodesics in the same homotopy class after filling the
cusps, when $n$ is proportional to $g$? Our spectral results are upper bounds
for small eigenvalues; their limiting distribution would require
additional geometric input.

\appendix
\section{Proofs for the kernels and the parameter}\label{app:kernels}

This appendix proves Lemmas~\ref{lem:kernel-zeros}, \ref{lem:kernel-positivity}
and~\ref{lem:parameter}. Notation is that of
Sections~\ref{the-positive-compact-family-of-kernels}
and~\ref{the-stationary-parameter}.

\subsection{Proof of Lemma~\ref{lem:kernel-zeros}}\label{app:kernel-zeros}

The elementary identities \(X''(y)=-X(y)/y\) and
\(X'''(y)=\sum_{j\ge0}(-1)^jy^j/(j!(j+2)!)\) will also be useful. On
\([0,\pi^2/4]\), \(X(y)>0\) for \(y>0\), \(X''<0\), and \(X'''>0\). For
the first assertion we use the alternating-series
lower bound \[
 X(y)/y\ge1-y/2+y^2/12-y^3/144>0\quad(0\le y\le5/2);
\] the cubic on the right is decreasing and is positive at \(5/2\). The
other assertions follow from the displayed identities and
\(X'''(y)\ge1/2-y/6>0\). Similarly, the partial sum of the series for
\(a(3/2)\) through \(j=4\) is an upper bound by the alternating-series
test, and it is negative; this proves \(y_0<3/2\). In particular
\(X(Y)\le Y<3/2\).

All zeros of \(a(y)\) are positive real. Indeed, if \(a(y)=0\), the
function \(u(s)=J_0(2\sqrt y\,s)\) solves \((su')'+4ysu=0\), is regular
at zero, and vanishes at one. Integration against \(\bar u\) gives \[
 4y\int_0^1s|u|^2\,\mathrm{d}s=\int_0^1s|u'|^2\,\mathrm{d}s>0.
\] The series is entire of order \(1/2\), so its Hadamard product has
genus zero and no nonconstant exponential factor. In particular it has
no zeros in \(|y|<y_0\).

We next prove that every zero of \(K_Y\) is positive and simple. At
\(Y=0\), \[
 K_0(z)=\frac{\sin\sqrt{2z}}{\sqrt{2z}}.
\] Coefficient comparison in \eqref{eq:P1} gives \[
 \partial_YK_Y=2\partial_zK_Y+(K_Y-a(Y))/z.
\] Integration along the characteristics \(z+2Y=\text{constant}\) yields
\begin{equation}
\sqrt{t^2-4Y}\ K_Y((t^2-4Y)/2)=\sin t-R_Y(t),\qquad
 R_Y(t)=2\int_0^Y\frac{a(u)}{\sqrt{t^2-4u}}\,\mathrm{d}u.       \label{eq:P3}
\end{equation} For \(0<Y<y_0\) and real \(t>2\sqrt Y\), \(R_Y>0\), \(R_Y'<0\), \(R_Y''>0\), and \[
 R_Y(t)\le\frac{2X(Y)}{\sqrt{t^2-4Y}}.
\] At \(Y=0\), \(R_0\) vanishes identically, as in the sinc formula above. For \(Y>0\), \(\sin t-R_Y(t)\) is strictly concave on each positive
sine interval. On the first interval \((2\sqrt Y,\pi)\), it tends to
zero at the left endpoint, is positive immediately afterwards by
\(K_Y(0)=a(Y)>0\), and is negative at \(\pi\); hence it has exactly one
simple interior zero. On \((2j\pi,(2j+1)\pi)\), \(j\ge1\), it is
negative at both endpoints and positive at the midpoint, since \[
 R_Y((2j+\tfrac12)\pi)
 <\frac3{\sqrt{(5\pi/2)^2-6}}<1.
\] Strict concavity gives exactly two simple zeros there. Negative sine
intervals contain no zero. For real \(z\le0\), every argument of \(X'\)
in \eqref{eq:P1} is less than \(y_0\), where \(X'>0\); for negative arguments
this follows from the positive modified-Bessel series. Thus \(K_Y\) has
no nonpositive real zero.

To exclude nonreal zeros, set \(G_Y(t)=K_Y((t^2-4Y)/2)\). It is entire
and even. Formula \eqref{eq:P3}, with the square root asymptotic to \(t\), holds
outside the segment \([-2\sqrt Y,2\sqrt Y]\). On \(|t|=(N+1/2)\pi\), \begin{equation}
\frac{G_Y(t)}{\sin t/t}
 =\frac t{\sqrt{t^2-4Y}}
    \left(1-\frac{R_Y(t)}{\sin t}\right)=1+O(N^{-1}).   \label{eq:P4}
\end{equation} The error is uniform: \(R_Y=O(N^{-1})\), and \(|\sin t|\) is bounded
below on these circles. For the latter assertion use
\(|\sin t|\ge\sinh1\) when \(|\Im t|\ge1\); otherwise \(|\Re t|\)
differs from \((N+1/2)\pi\) by \(O(N^{-1})\). Rouché's theorem gives
exactly \(2N\) zeros of \(G_Y\) inside each sufficiently large circle.
The real zeros already counted give exactly \(N\) positive and \(N\)
negative zeros there: for even \(N\) the circle ends at the midpoint of
a positive sine interval, after its first zero, and for odd \(N\) it
ends in the subsequent negative interval. They exhaust the count,
proving the assertion. Consequently the product formula \eqref{eq:P5}
holds. This is again a genus-zero product, since \eqref{eq:P1} has order \(1/2\).

\subsection{Proof of Lemma~\ref{lem:kernel-positivity}}\label{app:kernel-positivity}

We first prove \eqref{eq:P6}. By \eqref{eq:P1}, \[
 \partial_YH_Y=2\partial_zH_Y+3(H_Y-a(Y))/z.
\] Writing \(T=\log(a/H_Y)\) and \(\gamma=-a'/a\), this becomes \[
 \partial_YT=-\gamma+2\partial_zT+3(\mathrm{e}^T-1)/z.
\] The constant identity is \(\gamma=5t_1\); for \(m\ge1\), \begin{equation}
t_m'=(2m+5)t_{m+1}
  +3\sum_{\ell=2}^{m+1}\frac1{\ell!}
       \sum_{\substack{i_1+\cdots+i_\ell=m+1\\ i_j\ge1}}
          t_{i_1}\cdots t_{i_\ell}.                 \label{eq:P7}
\end{equation} At \(Y=0\), \[
 H_0(z)=3\frac{\sin\sqrt{2z}-\sqrt{2z}\cos\sqrt{2z}}
                    {(\sqrt{2z})^3}.
\] Its zeros are positive real. For a direct proof, a nonzero root \(t\)
of \(\sin t-t\cos t\) gives the regular function \(u(s)=j_1(ts)\),
satisfying \((s^2u')'+(t^2s^2-2)u=0\) and \(u(1)=0\). Integration
against \(\bar u\) shows \[
 t^2\int_0^1s^2|u|^2\,\mathrm{d}s
   =\int_0^1(s^2|u'|^2+2|u|^2)\,\mathrm{d}s>0.
\] Here \(j_1(w)=\sin w/w^2-\cos w/w\) is regular at zero. The root is
simple because \((\sin t-t\cos t)'=t\sin t\) does not vanish at a
nonzero root. The order-\(1/2\) product therefore gives
\(t_m(0)=m^{-1}\sum_j\rho_j^{-m}>0\).

Induct on the derivative order \(q\), simultaneously for all indices
\(m\), in \eqref{eq:P7}. Every differentiated term is a product of derivatives of
order at most \(q\) with nonnegative coefficients, so the induction is
valid; each fixed equation has only finitely many terms. For each \(m\),
\(t_m(Y)\) is a polynomial in the ratios \(X^{(j)}(Y)/a(Y)\), so it is
analytic on \(|Y|<y_0\), and its Taylor series converges to
\(t_m(Y)\) there. This proves \eqref{eq:P6}, and gives strictly positive
coefficients of \(H_Y^{-h}\) for every real \(h>0\), and nonnegative
coefficients of \(-H_Y'/H_Y\).

The same argument applied to \(S=\log(a/K_Y)=\sum_{m\ge1}s_m(Y)z^m\)
gives \begin{equation}
s_m'=(2m+3)s_{m+1}
  +\sum_{\ell=2}^{m+1}\frac1{\ell!}
       \sum_{\substack{i_1+\cdots+i_\ell=m+1\\ i_j\ge1}}
           s_{i_1}\cdots s_{i_\ell}.               \label{eq:P8}
\end{equation} The initial sinc product and the same analytic argument imply that
every Taylor coefficient in \(Y\) of \[
 \widehat\vartheta_\ell(Y):=a(Y)[z^\ell]K_Y(z)^{-1}
\] is nonnegative. The product for \(a\) implies the same for \(a^{-1}\)
and for \(\vartheta_\ell=[z^\ell]K_Y^{-1}\). This completes the proof.

\subsection{Proof of Lemma~\ref{lem:parameter}}\label{app:parameter-proof}

We use from Lemma~\ref{lem:kernel-zeros} that $X>0$,
$X''<0$ and $X'''>0$ on the relevant interval, and that $0<y_0<3/2$.

Normalize the left side of \eqref{eq:SP1} by \(2\sqrt{T-Y}\). Its root equation
becomes \begin{equation}
Q(Y,T)=\int_0^1X'(T(1-s^2)+Ys^2)\,\mathrm{d}s=0.
 \label{eq:SP8}
\end{equation} It is strictly decreasing in \(T\). At \(T=y_0\) it is positive when
\(Y<y_0\). The series and the beta integral give the Sonine identity \begin{equation}
\int_0^T\frac{X'(v)}{\sqrt{T-v}}\,\mathrm{d}v
       =\sin(2\sqrt T).
 \label{eq:SP9}
\end{equation} At \(T_0=\pi^2/4\), the full integral is zero, and removing \([0,Y]\)
makes it negative when \(Y>0\). This proves existence and uniqueness in
\eqref{eq:SP1}, including \(T(0)=T_0\). The derivative \(Q_T\) never vanishes
there, so the implicit function theorem gives analyticity. Also
\(T'<0\), and \(T\to y_0\) as \(Y\to y_0\).

Twice differentiating \eqref{eq:SP8}, with \(v=T(1-s^2)+Ys^2\), gives \[
 T''(Y)=
 -\frac{\int_0^1X'''(v)[T'(1-s^2)+s^2]^2\,\mathrm{d}s}
        {\int_0^1X''(v)(1-s^2)\,\mathrm{d}s}>0.
\] At \((Y,T)=(0,T_0)\), \eqref{eq:SP9} gives \(Q_T=Q_Y=-1/(2T_0)\), and hence
\(T'(0)=-1\). Consequently \(-1\le T'<0\).

The integral \(I\) is positive. Indeed, subtracting \((T-y_0)\) times
\eqref{eq:SP1} from its definition gives \begin{equation}
I(Y)=\int_Y^T
 \frac{X'(v)(y_0-v)}{\sqrt{T-v}}\,\mathrm{d}v>0.
 \label{eq:SP10}
\end{equation} Introduce \(J\) and \(\beta\) as in \eqref{eq:SP11}.
The equality for \(J\) is integration by parts, and the last equality
for \(\beta\) follows from the definition of \(\lambda\). If
\(Q_2=\int_0^1s^2X'(T-\Delta s^2)\,\mathrm{d}s\), then \(Q_2>0\): the function
\(X'(T-\Delta s^2)\) is strictly increasing in \(s\) and has mean zero
by \eqref{eq:SP8}. Moreover \(Q_2\le X'(Y)/3\). Differentiating \(J\) gives \begin{equation}
J'=(1-T')Q_2\le\frac23X'(Y),
 \qquad
 \beta'=\frac{X'J-XJ'}{J^2}\ge\frac{X'}{3J}>0.
 \label{eq:SP12}
\end{equation} At zero, \eqref{eq:SP9} gives \(I(0)=\pi/4\), \(J(0)=1/2\), and
\(\lambda(Y)=4Y+O(Y^2)\). At the other endpoint \(J\to x_c\) and
\(X(Y)\to x_c\), so \(\beta\to1\). This proves the asserted analytic
bijection and the values in \eqref{eq:SP3} at zero.

The endpoint constants are also explicit. Put \(\delta=y_0-Y\) and
\(b=-X''(y_0)=x_c/y_0>0\). The implicit equation at \((y_0,y_0)\) has
\(T'=-1/2\); this gives the first formula in \eqref{eq:SP13}. For the
second formula, substitute
\(X'(T-\Delta s^2)=b\Delta(s^2-1/3)+O(\Delta^2)\) into
\(I=2\Delta^{3/2}\int_0^1s^2X'(T-\Delta s^2)\mathrm{d}s\). Together with
\(X(Y)\to x_c\), this proves \eqref{eq:SP14}, and \eqref{eq:SP15} follows
from \eqref{eq:SP3}. Since \(1-\beta=2/(2+\lambda)\), \eqref{eq:SP14}
also gives \(\Delta\asymp\sqrt{1-\beta}\).

By part~(ii), \(\Psi\) is analytic and strictly increasing for
\(0<x<x_c\). The first line of \eqref{eq:SP16} follows from
\(\lambda(Y)=4Y+O(Y^2)\) and \(x=X(Y)=Y+O(Y^2)\). For the second line, quadratic inversion gives
\(x_c-x=(b/2)\delta^2(1+O(\delta))\). Combining this with \eqref{eq:SP13} and
\eqref{eq:SP14} gives the constant \(5/4\). All expansions here arise from
convergent analytic expansions in \(x\) or \(\sqrt\varepsilon\), so
their asserted derivatives follow by termwise differentiation. In
particular \(\Psi^2/(\mathscr D\Psi)\) has finite endpoint limits \(0\)
and \(5/4\), and is bounded on the full interval.


\section{Uniform bounds for positive intersection numbers}
\label{app:primitive-bounds}

We give the coarse bounds used in Section~\ref{uniform-comparison}.
They are consequences of the stronger bounds in~\cite[Lemmas~2--6]{GYZ}.
The proof below uses their separating-term estimate, which we reproduce,
and replaces the two-variable majorant by a scalar one. The only
intersection estimates imported into the proof are the one-point formula
and the two-point bound of Delecroix, Goujard, Zograf and
Zorich~\cite[equation~(6)]{DGZZBounds}.

\begin{lemma}\label{lem:primitive-bounds}
For $g\ge1$ and positive integers $d_1,\ldots,d_n$ with
$\sum d_i=3g-3+n$, the normalization in~\eqref{eq:H4} satisfies
\[
 \frac16\le C_g(d_1,\ldots,d_n)\le1.
\]
\end{lemma}

Write $C(\mathbf d)=C_g(\mathbf d)$ and $h=2g-2+n$ when the
genus is determined by the degrees. The normalized string and dilaton
identities are
\begin{equation}
 C(0,\mathbf d)=\sum_{d_j>0}\frac{2d_j+1}{3(Z-1)}
                    C(\mathbf d-\mathbf e_j),\qquad
 C(1,\mathbf d)=C(\mathbf d),
 \label{eq:primitive-string}
\end{equation}
where $Z$ is the complexity of $(0,\mathbf d)$; the terminal value
$C_1(1)=1/6$ is kept separately. For a pivot $d_1\ge2$, DVV gives
\begin{align}
 C(\mathbf d)&=\sum_{j=2}^n\frac{2d_j+1}{3(h-1)}
 C(d_2,\ldots,d_j+d_1-1,\ldots,d_n)\notag\\
 &\quad+\frac{2}{3(h-1)}\sum_{a+b=d_1-2}
 C_{g-1}(a,b,d_2,\ldots,d_n)\notag\\
 &\quad+\sum_{\substack{a+b=d_1-2\\I\sqcup J=\{2,\ldots,n\}}}
 \frac{\Gamma(h_1)\Gamma(h_2)}{6\Gamma(h)}
 C(a,\mathbf d_I)C(b,\mathbf d_J).
 \label{eq:primitive-dvv}
\end{align}
Here $h_1,h_2$ are the complexities of the two factors, and only stable,
dimension-compatible terms are included. The weights in the first two
lines sum to one.

\begin{proof}[Proof of Lemma~\ref{lem:primitive-bounds}]
\emph{The lower bound.}
We prove the same lower bound for positive-genus tuples with at most two
zero entries, by induction on complexity. For a tuple with $z\in\{1,2\}$
zeros, the total string weight is $(3h-z)/(3h-3)\ge1$.
String lowers complexity without increasing the number of zeros.
For a positive tuple, remove the ones by dilaton. Unless the tuple
reduces to the genus-one base, its genus is at least two. Its merger
children are positive and its handle children have positive genus and
at most two zeros. Dropping the nonnegative separating terms in
\eqref{eq:primitive-dvv} therefore proves the lower bound.

\emph{The upper majorant and zero entries.}
Define the nondecreasing function
\[
 F(h)=\begin{cases}1/3,&1\le h\le7,\\
                   1-4/(h-1),&h\ge7.
       \end{cases}
\]
We will prove $C(\mathbf d)\le F(h)$ for positive tuples.
Its increment for $h\ge8$ is $4/((h-1)(h-2))$.
The corresponding bound for a tuple with one zero is $C\le1$:
remove all ones, and let $Z$ be the remaining complexity. String bounds
the result by
\[
 \left(1+\frac{2}{3(Z-1)}\right)F(Z-1).
\]
For $Z\le8$ this is at most $5/9$; for $Z\ge9$, substitution of
$F(Z-1)=1-4/(Z-2)$ gives a bound below one. This observation uses only
positive tuples of smaller complexity and so is available in the induction.

\emph{The separating contribution.}
Suppose all $d_i\ge2$. Both separating genera are positive by dimension.
For a fixed subset $I$ of size $m$ and genus $g_1$, the degree $a$ is
determined. The total separating weight is thus at most
\[
 \sum_{m=0}^{n-1}\binom{n-1}{m}
 \sum_{\substack{g_1+g_2=g\\g_1,g_2\ge1}}
 \frac{(2g_1+m-2)!(2g_2+n-m-3)!}{6(h-1)!}.
\]
As in~\cite[Lemma~3]{GYZ}, the terms with at least one genus equal to
one, allowing double counting, are bounded by
\[
 \frac{1}{3(h-1)(h-2)}
 \sum_{m=0}^{n-1}\prod_{j=1}^m\frac{n-j}{h-2-j}
 \le\frac{3}{2(h-1)(h-2)}.
\]
Indeed $n\le3h/5$ and $n\le h-2$. The first two factors are at most
one, and each subsequent factor is at most $3/5$, giving a sum at most
$2+\sum_{r\ge0}(3/5)^r=9/2$.
For $g_1,g_2\ge2$, the inequality
$\binom ab\binom cd\le\binom{a+c}{b+d}$ gives the bound
\[
 \frac{n(g-3)}{6(h-1)(h-2)\binom{2g-4}{2}}
 \le\frac{3}{8(h-1)(h-2)}.
\]
This part vanishes for $g\le3$; otherwise use $n\le3g-3$.
The total weight is therefore at most $2/((h-1)(h-2))$.
Each factor is positive or has one zero, so the same bound applies to
the full separating contribution once the smaller-complexity bounds hold.

\emph{Initial cases.}
Put $a_g=C_g(3g-2)$ and $b_g=C_g(0,3g-1)$.
The one-point formula and the published two-point bound give
\[
 C_g(d_1,d_2)\le b_g=\frac{6g-1}{6g-3}a_g<\frac1\pi<\frac8{25}.
\]
For completeness, $b_{g+1}/b_g=(6g+5)(6g+1)(6g-3)/
(108g(g+1)(2g-1))>1$, and Stirling's formula gives $b_g\to1/\pi$.
All positive genus-two tuples reduce by dilaton to $(4)$, $(2,3)$ or
$(2,2,2)$. String gives
\[
 C_1(0,2)=\frac5{18},\quad C_1(0,0,3)=\frac{35}{108},\quad
 C_1(0,0,2,2)=\frac{25}{81}.
\]
Using these values in~\eqref{eq:primitive-dvv} yields
\[
 C_2(4)=\frac{35}{144},\quad C_2(2,3)=\frac{1015}{3888},\quad
 C_2(2,2,2)=\frac{175}{648}<\frac3{11}.
\]
The other two values are also below $3/11$. Hence a genus-two tuple
with one zero is at most $(11/9)(3/11)=1/3$ after removing ones;
its remaining complexity is at least four. The genus-one one-zero
case reduces to $(0,2)$ and has the same bound.

For $h\le7$, the only additional positive cores with at least three
markings are the genus-three tuples $(2,2,5)$, $(2,3,4)$ and $(3,3,3)$.
For the first two, two string steps bound the handle term by
$(16/15)(11/9)(3/11)=16/45$. Each separating factor has genus one
or two and is at most $1/3$. Thus
\[
 C\le\frac89\frac8{25}+\frac19\frac{16}{45}
              +\frac19\frac2{30}=\frac{671}{2025}<\frac13.
\]
For $(3,3,3)$, both handle terms reduce to $(0,3,3)$ in genus two, so
\[
 C\le\frac79\frac8{25}+\frac29\frac13
              +\frac19\frac2{30}=\frac{223}{675}<\frac13.
\]
Together with dilaton, the one- and two-point bounds and the genus-one
base, this proves every initial case $h\le7$.

\emph{Induction step.}
Let $h\ge8$. After removing ones, assume $n\ge3$ and choose a
minimal degree $d=d_1\ge2$. If $d\ge3$, all handle children except
the two endpoint terms have positive degrees. Each endpoint has one
zero and at most one degree-one entry. Removing that entry if necessary
and applying string bounds it by
$(1+2/(3(h-3)))F(h-1)$. The total endpoint weight is $4/(3(h-1))$.
Since $F\le1$, the excess over $F(h-1)$ is at most
\[
 \frac{8}{9(h-1)(h-3)}+\frac2{(h-1)(h-2)}
 \le\frac{46}{15(h-1)(h-2)}<F(h)-F(h-1).
\]
If $d=2$, the handle has two zeros. Two string steps, allowing an
intervening dilaton step, give the factor
\[
 R_h=\frac{(3h-5)(3h-10)}{(3h-6)(3h-12)}.
\]
The resulting positive tuples have complexity at most $h-3$.
The handle weight is $2/(3(h-1))$, and
\[
 \frac{2(R_h-1)}{3(h-1)}
 =\frac{2(9h-22)}{27(h-4)(h-1)(h-2)}
 \le\frac{25}{27(h-1)(h-2)}.
\]
Adding separation gives $79/(27(h-1)(h-2))$, again below the increment
of $F$. This closes the induction and proves the upper bound.
\end{proof}

\section{The marked expansion beyond the leading term}\label{app:marked}

This appendix proves Theorems~\ref{thm:first-correction}
and~\ref{thm:full-expansion}. Subsections~\ref{app:second-degree}--\ref{app:critical-matching}
compute one further term in the recursion. The coefficient to be integrated
vanishes at the endpoint, so its primitive stays bounded. This gives both
an $O_p(g^{-1})$ relative error for the leading term and its first explicit
correction. The number $p$ of distinguished markings remains fixed.
Subsections~\ref{sec:full-expansion} and~\ref{sec:endpoint-expansion}
then construct the full series and control its remainder as the
truncation order grows.

\subsection{The second degree coefficient}\label{app:second-degree}

We first extend the degree expansion \eqref{eq:PF5} of the normalized
profiles by one order (Lemma~\ref{lem:second-profile}). The argument is
the adjacent-degree subtraction used for \eqref{eq:PF5} in
Section~\ref{sec:prefactor}, now
keeping the $h^{-2}$ terms of the prefactors and the two smallest
separating configurations.

We retain $\delta_p$ and $\beta_p$ from \eqref{eq:PF5} and
\eqref{eq:PF7}. The kernel coefficients $\theta_\ell$ are those of \eqref{eq:H2}. Set
\begin{equation}
 \alpha_p=A\{(p-3)\Theta+2M_1+2p-3\},\qquad
 \mathcal T_j=\sum_{\ell\ge\max(0,j)}w_\ell.
 \label{eq:ST1}
\end{equation}
For a pivot $i$, put $t=d_i$ and let $z$ count zero entries among the
other $p-1$ markings. Let $\mathbf d^{ij,\ell}$ be the merger child,
with new degree $t+d_j+\ell-1$, and let $\mathbf d^{i;u,v}$ be the
nonseparating child with the pivot replaced by $u,v$. Define
\begin{equation}
\begin{split}
 B_{p,i}(\mathbf d)={}&
 A\sum_{j\ne i}(2d_j+1)
 \sum_{\ell\ge\max(0,1-t-d_j)}w_\ell
 [1-\beta_{p-1}-\delta_{p-1}(\mathbf d^{ij,\ell})]\\
 &+2A\sum_{\ell\ge\max(0,2-t)}w_\ell
 \sum_{\substack{u+v=t+\ell-2\\u,v\ge0}}
 [1-\beta_{p+1}-\alpha_p-\delta_{p+1}(\mathbf d^{i;u,v})]\\
 &+A^2\left\{\left[\binom z2+\frac{\theta_1}2\right]
                    \mathcal T_{2-t}+\frac12\mathcal T_{3-t}\right\}.
\end{split}
 \label{eq:ST2}
\end{equation}
The first line is absent at $p=1$. Starting at zero along a fixed
coordinate-increment path, define $e_p$ by
\begin{equation}
 e_p(\mathbf0)=0,\qquad
 e_p(\mathbf d+\mathbf e_i)-e_p(\mathbf d)=-B_{p,i}(\mathbf d).
 \label{eq:ST3}
\end{equation}
The formulas converge absolutely by \eqref{eq:TL10}; they give
$B_{p,i}=O_p((D+1)^3)$ and $e_p=O_p((D+1)^4)$.

\begin{lemma}[Second profile correction]\label{lem:second-profile}
For fixed $p\ge1$, uniformly in $Y$ and $D\le h^{1/8}$,
\begin{equation}
 c_{g,p}(\mathbf d)=1-\frac{\delta_p(\mathbf d)}h
       +\frac{e_p(\mathbf d)}{h^2}
       +O_p\!\left(\frac{(D+1)^6}{h^3}\right).
 \label{eq:ST4}
\end{equation}
The function $e_p$ is independent of the chosen path.
\end{lemma}
\begin{proof}
We first need the linear term in the genus exchange. For $p\ge2$,
the two genus-zero three-marking contributions to \eqref{eq:TL18}
give $T_{g,p}=2(p-2)A/h+O_p(h^{-2})$. Every other genus-zero term
loses at least two cusp factors, and the positive-genus portion is
$O_p(h^{-2})$ by \eqref{eq:TL19}. Thus \eqref{eq:TL14} and
\eqref{eq:PF5} imply
\begin{equation}
 \frac{W_{g-1,p+2}}{W_{g,p}}=1-\frac{\alpha_p}h+O_p(h^{-2}).
 \label{eq:ST5}
\end{equation}
Indeed the coefficient is
$\delta_p(\mathbf e_1)+3(p-2)A=\alpha_p$.
For $p=0,1$, descend from $p+1$ using \eqref{eq:PF7}; the cusp ratio
changes its first coefficient by
$\beta_{p+2}-\beta_p=-A(\Theta+2)$, while
$\alpha_{p+1}-\alpha_p=A(\Theta+2)$.
Consequently the two adjacent-difference prefactors are
\begin{equation}
\begin{split}
 \frac{W_{g,p-1}}{W_{g,p}}
 &=\frac Ah\left(1+\frac{1-\beta_{p-1}}h\right)+O_p(h^{-3}),\\
 2\frac{W_{g-1,p+1}}{W_{g,p}}
 &=\frac{2A}h\left(1+\frac{1-\beta_{p+1}-\alpha_p}h\right)
     +O_p(h^{-3}).
\end{split}
 \label{eq:ST6}
\end{equation}
The first line is used only at $p\ge2$. Both children have complexity
$h-1$, which accounts for the two constants $+1$.

We next retain the leading separating terms of the adjacent difference.
If the smaller component has genus $j$ and contains $s$ of the other
external labels, the complementary component has complexity $h-(2j+s)$.
Stability forces $2j+s\ge2$.
Equality occurs only for $(j,s)=(0,2)$ and $(1,0)$. The first component
is a three-marked sphere, so its three degrees must be zero; it gives
$\binom z2\mathcal T_{2-t}$. The second is a one-marked torus, whose
$W$ amplitudes are $\theta_1/2$ and $1/2$ at degrees zero and one.
It gives $(\theta_1/2)\mathcal T_{2-t}+(1/2)\mathcal T_{3-t}$.
The two orientations cancel the factor $1/2$ in the recursion.
In both cases the large all-zero component divided by the parent is
$A^2/h^2+O_p(h^{-3})$, by fixed iterations of \eqref{eq:PF7} and
\eqref{eq:ST5}. Its degrees can be replaced by zero at cost
$O_p((D+\ell+1)^2/h)$.

After removing these two types the all-zero separating mass is
$O_p(h^{-3})$. To verify this sharper bound, repeat the proof of
\eqref{eq:TL20}: the remaining genus-zero terms have $s\ge3$, the
remaining genus-one terms have $s\ge1$, and $j\ge2$ starts at
$h^{-4}$. The same factorial-geometric series sums these bounds,
with any fixed polynomial in $j$; \eqref{eq:TL21} handles the
macroscopic splits. At $p=1$, insert $2j+1$ cusps in the larger
component before raising its genus to marking count two, then use
$W_{g,2}/W_{g,1}=O(g)$. This bounds its ratio by
$O((C/g)^{2j})$. After the two genus-one orientations are removed,
the factorial series starts at $j=2$ and is $O(h^{-4})$.
This argument uses the marking-independent constant in \eqref{eq:TL16}.
Pair counts and the moments of $w_\ell$ give an adjacent-difference
error $O_p((D+1)^3/h^3)$ for the separating part.

Now substitute \eqref{eq:PF5} and \eqref{eq:ST6} into the exact
adjacent subtraction of \eqref{eq:H9}. The two retained separating
types give the last line of \eqref{eq:ST2}. The result is
\[
 c(\mathbf d)-c(\mathbf d+\mathbf e_i)
 =\frac{\delta_p(\mathbf d+\mathbf e_i)-\delta_p(\mathbf d)}h
  +\frac{B_{p,i}(\mathbf d)}{h^2}
  +O_p((D+1)^5/h^3).
\]
To justify the substitution, cut the kernel index $\ell$ at $h^{1/8}$. Below this cutoff, child
degrees are at most $2h^{1/8}+O_p(1)$, inside the range of
\eqref{eq:PF5} for all sufficiently large $h$. Its remainder, multiplied
by the linear prefactor and degree count, is
$O_p((D+\ell+1)^5/h^3)$. Above the cutoff use $0\le c\le1$ and
polynomial growth of the displayed coefficients. Since the uniform
moments \eqref{eq:TL10} of $w_\ell$ are finite to every order, these
tails are $O_p(h^{-3})$, even after summing at most $D$ increments. The dimension cutoff is farther away and is
covered by the same argument. Summing increments proves \eqref{eq:ST4}.
For a fixed tuple, expansions along two paths must have the same
$h^{-2}$ coefficient as $g\to\infty$, proving path independence.
\end{proof}

\emph{Comparison with the fixed-cusp expansion.}
At $Y=0$ one has $A=1/\pi^2$, $\Theta=2$ and $M_1=3/2$, and the
normalization \eqref{eq:TL9} gives $W_{g,p}/W_{g,p+1}=4V_{g,p}/V_{g,p+1}$
and $W_{g-1,p+2}/W_{g,p}=V_{g-1,p+2}/V_{g,p}$. With $h=2g-2+p$,
equations \eqref{eq:PF7} and \eqref{eq:ST5} become
\[
 \frac{V_{g,p+1}}{8\pi^2gV_{g,p}}
 =1+\frac{(\frac12-\pi^{-2})p-\frac54+2\pi^{-2}}{g}+O_p(g^{-2}),
 \qquad
 \frac{V_{g-1,p+2}}{V_{g,p}}=1+\frac{3-2p}{\pi^2g}+O_p(g^{-2}),
\]
and \eqref{eq:PF5} gives, for fixed $k\ge1$,
\[
 \frac{[\tau_k\tau_0^{p-1}]_{g,p}}{V_{g,p}}
 =1-\frac{k^2+(p-\frac52)k-\frac p2+\frac32}{\pi^2g}+O_p(g^{-2}).
\]
These are the first-order formulas of Mirzakhani and
Zograf~\cite[Theorems~1.4 and~3.2(i)]{MZ}, recovered here as the
$Y=0$ case of the parameter-dependent coefficients.

\subsection{The coefficient in the differential equation}\label{app:transport-coefficient}

Inserting Lemma~\ref{lem:second-profile} into the dilaton identity gives
the next coefficient $\gamma_p$ of the differential equation~\eqref{eq:PF8}.
We then evaluate it with finitely many kernel moments and check that it
vanishes at the pure-$\psi$ endpoint.

Apply \eqref{eq:ST4} to the dilaton child of complexity $h+1$ and set
\begin{equation}
 E_p=\sum_{d\ge1}b_d e_{p+1}(d,0^p),\qquad
 \gamma_p=\beta_p^2-\beta_p-E_p/A.
 \label{eq:ST7}
\end{equation}
The moments of $b_d$ justify this sum and yield
\begin{equation}
 \frac{\mathrm{d}}{\mathrm{d}Y}\log\Phi_g^{(p)}(X(Y))
 =\frac{h+\beta_p+\gamma_p/h+O_p(h^{-2})}{Ar}.
 \label{eq:ST8}
\end{equation}
Indeed, replacing $h+1$ by $h$ in the child profile gives
$1-\delta/h+(\delta+e)/h^2$.
Thus the dilaton ratio is
$A\{1-\beta_p/h+(\beta_p+E_p/A)/h^2+O_p(h^{-3})\}$,
and inversion proves \eqref{eq:ST8}. To justify the sum, truncate the
dilaton degree at $h^{1/8}$; the tail is bounded using
$1+d^2/h+d^4/h^2$ and a sufficiently high absolute moment of $b_d$
from \eqref{eq:TL13}.

The coefficient in \eqref{eq:ST7} can be evaluated using finitely many
kernel moments, as follows.
Put $P=p+1$, and write $B_P(t)=B_{P,1}(t,0^{P-1})$. With
$k=t+\ell-1$ and $s=t+\ell-2$, \eqref{eq:ST2} becomes
\[
\begin{split}
 R_P(k)&=1-\beta_{P-1}
       -A\{\Theta(k^2+(P-5)k)+2M_1k+(4-P)\mathbf1_{k\ge1}\},\\
 S_P(s)&=(s+1)(1-\beta_{P+1}-\alpha_P)\\
 &\quad-A\{(s+1)(\Theta(s^2+(P-3)s)+2M_1s)
                   +(5-2P)s-1+\mathbf1_{s=0}\}.
\end{split}
\]
Then
\begin{equation}
\begin{split}
 B_P(t)={}&A(P-1)\sum_{\ell\ge\max(0,1-t)}w_\ell R_P(t+\ell-1)\\
 &+2A\sum_{\ell\ge\max(0,2-t)}w_\ell S_P(t+\ell-2)\\
 &+A^2\left\{\left[\binom{P-1}2+\frac{\theta_1}2\right]
                     \mathcal T_{2-t}+\frac12\mathcal T_{3-t}\right\}.
\end{split}
 \label{eq:ST9}
\end{equation}
The zero-degree correction in $S_P$ follows by summing the
$q$ term of $\delta_{P+1}$: for $s\ge1$ the two endpoint pairs have
$q=1$ and the $s-1$ interior pairs have $q=2$.

For $t\ge3$, remove both lower cutoffs, replace
$\mathbf1_{k\ge1}$ by one, $\mathbf1_{s=0}$ by zero, and both tails
by $\Theta$. This gives a cubic polynomial
$B_\infty(t)=a_3t^3+a_2t^2+a_1t+a_0$.
Its coefficients use only $M_j=\sum\ell^jw_\ell$, $0\le j\le3$.
Put $\Delta_j=B_P(j)-B_\infty(j)$ for $j=0,1,2$, and define
\begin{equation}
 \begin{split}
 \Sigma(d)&=\frac{a_3}4d^2(d-1)^2+
       \frac{a_2}6d(d-1)(2d-1)+\frac{a_1}2d(d-1)+a_0d,\\
 e_P(d,0^{P-1})&=-\Sigma(d)-\sum_{j=0}^2\Delta_j\mathbf1_{d\ge j+1}.
 \end{split}
 \label{eq:ST10}
\end{equation}
The three exceptional values use the same moments and $w_0,w_1,w_2$;
there is no remaining infinite recursion. Set
\[
 L_j=\frac13[(z\partial_z+1)^j\bar H_Y(z)]_{z=1}\quad(1\le j\le4),
 \qquad b_1=\frac13,\quad b_2=-\frac{\theta_1}5,
\]
and
\[
 \Sigma_b=\frac{a_3}4(L_4-2L_3+L_2)
   +\frac{a_2}6(2L_3-3L_2+L_1)
   +\frac{a_1}2(L_2-L_1)+a_0L_1.
\]
Then the finite formula for \eqref{eq:ST7} is
\begin{equation}
 \gamma_p=\beta_p^2-\beta_p+
 \frac{\Sigma_b+\Delta_0A+\Delta_1(A-b_1)+\Delta_2(A-b_1-b_2)}A.
 \label{eq:ST11}
\end{equation}
All ingredients extend analytically to $Y=y_0$, by the removable
normalized-kernel extension used in \eqref{eq:PF14}.

At that endpoint $A=2/15$, $\Theta=1$, $M_j=0$ for $j\ge1$,
$w_0=\theta_1=1$, and $\beta_p=-p/5$. Equation~\eqref{eq:ST9}
gives $B_P(0)=0$ and $B_P(1)=A(P-1)(P+4)/5$.
Only $b_1=1/3$ and $b_2=-1/5$ remain, so
$E_p/A=p(p+5)/25=\beta_p^2-\beta_p$. In particular,
\begin{equation}
 \gamma_p(y_0)=0,\qquad \gamma_p(Y)=O_p(y_0-Y),\qquad
 \gamma_p/(Ar)=O_p(1).
 \label{eq:ST12}
\end{equation}
This agrees with an exact computation: string and dilaton at the pure-$\psi$
endpoint give $W_{g,p+1}/W_{g,p}=3(h+m)=(15h-3p)/2$,
so that $A\,W_{g,p+1}/(hW_{g,p})=1-p/(5h)$ exactly, with $A=2/15$;
there is no term of order $h^{-2}$, and $\gamma_p=0$.

\subsection{Critical matching to the next order}\label{app:critical-matching}

Recall $t_p$ and $\mathcal L_p$ from \eqref{eq:ST13}.
By \eqref{eq:ST9}--\eqref{eq:ST12}, $\mathcal L_p$ is an explicit
bounded function, given by finitely many kernel derivatives and a
convergent one-dimensional integral.

\begin{proof}[Proof of Theorem~\ref{thm:first-correction}]
The sharper critical input is
$C_g(2^{3g-3})=\pi^{-1}(1-2/(9g)+O(g^{-2}))$.
It is derived in Subsection~\ref{sec:endpoint-expansion}, as the
coefficient $d_1=-2/9$ in \eqref{eq:FE10}, from the Itzykson--Zuber
identity, the large-order relation of Garoufalidis--Its--Kapaev--Mari\~no
and the first one-instanton coefficient $v_1=-1/8$; that derivation is
independent of the transport argument of
Subsections~\ref{app:second-degree}--\ref{app:critical-matching}, and the value
agrees with \cite[equation~(25)]{GYZ}. A leading estimate alone would not identify
this correction. Expanding the finite product \eqref{eq:PF12} gives
\[
 \pi C_g(0^p,2^m)=
 1+\frac{p(5-p)/15-4/9}{h}+O_p(h^{-2}).
\]
The critical harmonic coefficient also has a first correction. Its
summand law is geometric with mean $3/2$, variance $v=15/4$,
third cumulant $15$, and fourth cumulant $705/8$. At displacement
$d=-p/2$ from $3h/2$, expansion of its exact negative-binomial
coefficient by Stirling's formula gives
\[
 \frac{\mathcal M_{g,p}(0^p;y_0)}
 {\Gamma(h)4^{-g}A_*^{-h}/\sqrt{2\pi hv}}
 =1+\frac{-p^2/30+4p/15-19/180}{h}+O_p(h^{-2}),
 \qquad A_*=2/15.
\]
For example the coefficient before substituting $d=-p/2$ is
$-d^2/(2v)-15d/(2v^2)+(705/8)/(8v^2)-5(15)^2/(24v^3)$.
Combining these two expansions with \eqref{eq:PF13} proves
\[
 \lim_{Y\uparrow y_0}
 \frac{\pi\Phi_g^{(p)}(X(Y))}{\mathcal J_{g,p}(Y)}
 =1+\frac{t_p}{h}+O_p(h^{-2}).
\]
Take $Y_*=y_0-h^{-3}$. The positive homogeneous-polynomial comparison
\eqref{eq:PF15} and the analytic normalized factor in \eqref{eq:PF16}
change this ratio by $O_p(h(y_0-Y))=O_p(h^{-2})$ for $Y\ge Y_*$.
Moreover $\mathcal L_p(Y)=t_p+O_p(y_0-Y)$ there.
For $Y\le Y_*$, subtract the exact logarithmic derivative of
$\mathcal J_{g,p}$ in \eqref{eq:PF11} from \eqref{eq:ST8}, and integrate:
\[
 \log\frac{\pi\Phi_g^{(p)}(X(Y))}{\mathcal J_{g,p}(Y)}
 =\frac{\mathcal L_p(Y)}h+
 O_p\!\left(h^{-2}+\frac1{h^2}
             \log\frac{\mathcal Z(Y)}{\mathcal Z(Y_*)}\right).
\]
Here \eqref{eq:ST12} controls the terminal integral, and
$\mathcal Z(Y_*)\asymp h^{-15/2}$. The remainder is
$O_p(\log h/h^2)$ uniformly. Boundedness of $\mathcal L_p$ permits
exponentiation, proving \eqref{eq:ST14} in both regions.
\end{proof}

\subsection{The formal hierarchy}\label{sec:full-expansion}

We now construct the full series and control its remainder as the truncation order grows. These are still estimates for generating functions: the passage to a prescribed cusp count is made in Subsection~\ref{sec:canonical-expansion} and Appendix~\ref{app:extraction}.

The recursion is triangular. Each new degree coefficient uses earlier cusp and genus coefficients; the dilaton and KdV equations then give the new cusp and genus coefficients. Proposition~\ref{prop:hierarchy} bounds the coefficients and the actual remainders. Matching at the pure-$\psi$ endpoint in Subsection~\ref{sec:endpoint-expansion} completes the proof of Theorem~\ref{thm:full-expansion}.

The proof makes the heuristic of Section~\ref{sec:higher-order} precise. It divides each order-$j$ coefficient by $S_j^{2j}$, where $S_j=D+p+Mj$ is the weight described there. After this division, the coefficient rules become coefficientwise inequalities between power series with nonnegative coefficients, and an analytic majorant (Lemma~\ref{lem:majorant}) bounds the normalized coefficients by $CB^j$. The true remainders are then controlled by a second, linear majorant system with the same normalization.


Use a formal variable $t$ for inverse complexity, and put
\begin{equation}
 \mathsf P_p(\mathbf d;t)=\sum_{j\ge0}e_{p,j}(\mathbf d)t^j,
 \quad \mathsf C_p(h)=\frac hA\frac{W_{g,p}}{W_{g,p+1}},
 \quad \mathsf G_p(h)=\frac{W_{g-1,p+2}}{W_{g,p}}.
 \label{eq:FE1}
\end{equation}
The formal series for the last two quantities will also be denoted by
$\mathsf C_p(t),\mathsf G_p(t)$. Their constant terms, and $e_{p,0}$,
are one. Here $e_{p,1}=-\delta_p$ and $e_{p,2}=e_p$ from
\eqref{eq:PF5} and \eqref{eq:ST3}. All series operations below are
truncated at a fixed finite order; convergence of a formal series is
never asserted. A complexity change from $h$ to $h-c$ means replacing
$t$ by $t/(1-ct)$.

For a bounded component of genus $j$ containing $s$ other external
labels, set $r_*=2j+s$. Thus $r_*$ is the decrease in complexity
from the parent to the complementary component. Its all-zero ratio has the exact
factorization
\begin{equation}
\begin{split}
 \frac{W_{g-j,p-s}}{W_{g,p}}
 ={}&\prod_{a=0}^{r_*-1}
   \frac{A\mathsf C_{p-s+a}(h-r_*+a)}{h-r_*+a}
      \prod_{b=0}^{j-1}\mathsf G_{p+2b}(h),\\
 \mathsf R_{p;j,s}(t)
 ={}&(At)^{r_*}\prod_{a=0}^{r_*-1}
 \frac{\mathsf C_{p-s+a}(t/(1-(r_*-a)t))}{1-(r_*-a)t}
       \prod_{b=0}^{j-1}\mathsf G_{p+2b}(t).
\end{split}
 \label{eq:FE2}
\end{equation}
The genus steps have unchanged complexity; the cusp steps remove
$r_*$ markings at genus $g-j$. All their marking counts are positive
in the applications below. To obtain $\mathsf R$ through order $t^K$
requires scalar coefficients only through order $K-r_*$.

The coefficients are computed as follows. Use the child tuples of
\eqref{eq:ST2}, and set $\mathcal W_{j,q}(\mathbf d)=4^j\overline
U_{j,q}(\mathbf d)$ for the exact bounded-complexity amplitudes.
At each order $K\ge1$, form
\begin{equation}
\begin{split}
 \mathsf A_{p,i}={}&\frac{At}{1-t}\mathsf C_{p-1}\!\left(\frac t{1-t}\right)
 \sum_{j\ne i}(2d_j+1)
 \sum_{\ell\ge\max(0,1-d_i-d_j)}w_\ell
 \mathsf P_{p-1}\!\left(\mathbf d^{ij,\ell};\frac t{1-t}\right),\\
 \mathsf N_{p,i}={}&\frac{2At}{1-t}\mathsf C_{p+1}\!\left(\frac t{1-t}\right)
 \mathsf G_p(t)\sum_{\ell\ge\max(0,2-d_i)}w_\ell
 \sum_{u+v=d_i+\ell-2}\mathsf P_{p+1}\!\left(\mathbf d^{i;u,v};\frac t{1-t}\right).
\end{split}
 \label{eq:FE3}
\end{equation}
Nonnegative node degrees are understood, and $\mathsf A$ is absent
at $p=1$. The retained separating contribution is
\begin{equation}
 \mathsf S_{p,i}=
 \sum_{\substack{j\ge0,\ I\subset\{1,\ldots,p\}\setminus\{i\}\\
        s=|I|,\ (j,s+1)\text{ stable}\\2\le2j+s\le K}}
 \mathsf R_{p;j,s}(t)\sum_{\ell\ge0}w_\ell
 \sum_{u+v=d_i+\ell-2}
 \mathcal W_{j,s+1}(u,\mathbf d_I)
 \mathsf P_{p-s}\!\left(v,\mathbf d_J;\frac t{1-(2j+s)t}\right),
 \label{eq:FE4}
\end{equation}
where $J$ is the remaining set of labels. Only bounded $j$ occur;
for large $g$ these orient the smaller component uniquely. The two
orientations cancel the separating factor $1/2$. Define
\begin{equation}
 \begin{split}
 e_{p,K}(\mathbf0)&=0,\\
 e_{p,K}(\mathbf d+\mathbf e_i)-e_{p,K}(\mathbf d)
   &=-[t^K](\mathsf A_{p,i}+\mathsf N_{p,i}+\mathsf S_{p,i}),\\
 \mathsf C_p(t)&=\frac1A\sum_{d\ge1}b_d
       \mathsf P_{p+1}\!\left(d,0^p;\frac t{1+t}\right).
 \end{split}
 \label{eq:FE5}
\end{equation}
Finally, for $p\ge2$, KdV gives
\begin{equation}
\begin{split}
 \mathsf G_p(t)={}&\mathsf P_p(\mathbf e_1;t)
 -3\sum_{\substack{0\le s\le p-2,\ j\ge0\\
             (j,s+2)\text{ stable}\\1\le2j+s\le K}}
  \binom{p-2}s\mathcal W_{j,s+2}(\mathbf0)\mathsf R_{p;j,s}(t),\\
 \mathsf G_p(h)&=\mathsf G_{p+1}(h+1)
                   \mathsf C_{p+2}(h)/\mathsf C_p(h).
\end{split}
 \label{eq:FE6}
\end{equation}
The last line descends from $p=2$ to $p=1$, and optionally to $p=0$.
It is an exact identity; its shift $h+1$ must be retained.

The coefficient and remainder bounds below are obtained by comparing
coefficientwise inequalities with an analytic system. We isolate the
elementary principle used for this comparison.

\begin{lemma}[Analytic majorants]\label{lem:majorant}
Let $\Phi=(\Phi_1,\ldots,\Phi_m)$ be power series in
$(t,v_1,\ldots,v_m)$ with nonnegative coefficients, convergent near the
origin, such that $\Phi(0,0)=0$ and the matrix $\partial_v\Phi(0,0)$ is
strictly lower triangular. Then $U=\Phi(t,U)$ has a unique formal
solution with $U(0)=0$. Its coefficients are nonnegative, and it
converges near $t=0$. Let $e_{i,j}\in[0,\infty]$, for $1\le i\le m$ and
$j\ge1$, and put $E_i=\sum_{j\ge1}e_{i,j}t^j$. If
$e_{i,j}\le[t^j]\Phi_i(t,E)$ for all $i,j$, then
$e_{i,j}\le[t^j]U_i$. In particular every $e_{i,j}$ is finite, and
$e_{i,j}\le C\rho^{-j}$ with $C,\rho>0$ depending only on $\Phi$.
\end{lemma}
\begin{proof}
Order the pairs $(j,i)$ lexicographically. Since $E$ has no constant
term, $[t^j]\Phi_i(t,E)$ is a polynomial with nonnegative coefficients
in the numbers $e_{k,l}$. A monomial $t^av^\beta$ with $a\ge1$ or
$|\beta|\ge2$ contributes only coefficients with $l<j$, since each
factor has order at least one. The remaining terms are
$\partial_{v_k}\Phi_i(0,0)\,e_{k,j}$, and they occur only for $k<i$.
Thus $[t^j]\Phi_i(t,E)$ involves only pairs preceding $(j,i)$. The
equations $U=\Phi(t,U)$ therefore determine the coefficients of $U$
recursively and uniquely, and these are nonnegative. The same induction
and the monotonicity of polynomials with nonnegative coefficients give
$e_{i,j}\le[t^j]U_i$. The derivative of $(t,v)\mapsto v-\Phi(t,v)$
with respect to $v$ at the origin is $I-\partial_v\Phi(0,0)$, which is
invertible. The analytic implicit-function theorem gives a solution
holomorphic near $t=0$ and vanishing at $t=0$; its Taylor series is the
formal solution. Cauchy's estimate on a small circle gives the final
bound.
\end{proof}

\begin{proposition}[Uniform recursive expansions]\label{prop:hierarchy}
There are $C,B,M>0$ and one complex neighborhood $\Omega$ of
$[0,y_0]$ such that the preceding rules determine path-independent
holomorphic coefficients with
\[
 \|e_{p,j}(\mathbf d)\|_\Omega\le CB^j(D+p+Mj)^{2j},\qquad
 \|\mathsf C_{p,j}\|_\Omega+\|\mathsf G_{p,j}\|_\Omega
       \le CB^j(p+Mj)^{2j}.
\]
For every stable state $p\ge1$, $h=2g-2+p\ge1$, every $L\ge1$,
and every nonnegative degree tuple, uniformly for $0\le Y\le y_0$,
\begin{equation}
 \begin{split}
 \left|c_{g,p}(\mathbf d)-\sum_{j<L}e_{p,j}(\mathbf d)h^{-j}\right|
       &\le CB^L(D+p+ML)^{2L}h^{-L},\\
 \left|\mathsf C_p(h)-\sum_{j<L}\mathsf C_{p,j}h^{-j}\right|
 +\left|\mathsf G_p(h)-\sum_{j<L}\mathsf G_{p,j}h^{-j}\right|
       &\le CB^L(p+ML)^{2L}h^{-L}.
 \end{split}
 \label{eq:FE7}
\end{equation}
Here $D=\sum d_i$, $\mathsf C_{p,j}=[t^j]\mathsf C_p(t)$ and
$\mathsf G_{p,j}=[t^j]\mathsf G_p(t)$; invalid-degree profiles and invalid genus numerators
are zero. For fixed $p,j$, the rules use finitely many kernel moments,
low kernel coefficients and exact bounded-complexity amplitudes.
\end{proposition}

The proof of Proposition~\ref{prop:hierarchy} occupies
Subsections~\ref{app:weights}--\ref{app:remainder-bounds}. At each order the
construction runs in the sequence profile, cusp ratio, genus ratio at
$p\ge2$, and descent to $p=1$; every profile on the right of
\eqref{eq:FE3}--\eqref{eq:FE4} has smaller order, and the dilaton and KdV
rules use the current profile only after it has been constructed.
Subsection~\ref{app:weights} fixes the analytic neighborhood and the
weights $S_j$ and records four elementary bounds.
Lemma~\ref{lem:coefficient-bounds} bounds the formal coefficients,
uniformly in the order, by a system of analytic majorants; this is the
first assertion of the proposition. Lemmas~\ref{lem:tails}
and~\ref{lem:first-order} control the omitted topological types and give
first-order estimates for the true quantities. Finally,
Subsection~\ref{app:remainder-bounds} proves the remainder bounds
\eqref{eq:FE7} by induction on the order, using a linear majorant system
for the normalized errors.

\subsection{Weights and elementary bounds}\label{app:weights}
Shrink one complex neighborhood $\Omega$ of $[0,y_0]$ so that the
normalized kernels $\bar K_Y,\bar H_Y$ and their nonvanishing quotients
from Section~\ref{the-positive-compact-family-of-kernels} are holomorphic
and bounded on a common disk $|w|\le1+2\eta$, $\eta>0$, for every
$Y\in\Omega$, and so that $A(Y)$ is bounded above and away from zero on
$\Omega$. Since $\sum_\ell w_\ell w^\ell=(1-w)/\bar K_Y(w)$ and
$b_d=\frac13[w^{d-1}]\bar H_Y(w)$, Cauchy's estimate on the circle
$|w|=1+\eta$ gives constants $C_0$ and $\rho=1+\eta>1$ with
\begin{equation}
 |w_\ell(Y)|\le C_0\rho^{-\ell},\qquad |b_d(Y)|\le C_0\rho^{-d}
 \qquad(Y\in\Omega,\ \ell,d\ge0).
 \label{eq:FE-geom}
\end{equation}
The absolute-series argument in the proof of Lemma~\ref{lem:terminal}(iii) also works on $\Omega$:
the background majorant $B_Y$ of
Section~\ref{elementary-terminal-bounds-and-a-global-crude-estimate}
has a uniformly bounded germ, so one small radius $z_*$ makes
$B_Y(z_*)\le1/2$ for all $Y\in\Omega$. A stable component of complexity
$r$ has at most $r+2$ markings and deficit $m\le3r/2$, so the primitive
bound of Lemma~\ref{lem:terminal}(i) and this majorant give
\begin{equation}
 |\mathcal W_{j,q}(\mathbf d)|\le C_0^{\,r}\,\Gamma(r)\qquad(r=2j-2+q\ge1),
 \label{eq:FE-small}
\end{equation}
uniformly in the markings, the degrees and $Y\in\Omega$; enlarging
$C_0$ covers the components of complexity one.

Fix an integer $M\ge2$ with $\mathrm{e}^{2/M}<\rho$. For a profile at a
state with $p$ markings and degree sum $D$, and for a scalar with $p$
markings, the order-$j$ weights are
\begin{equation}
 S_j=D+p+Mj\qquad\text{and}\qquad S_j=p+Mj
 \label{eq:FE-weights}
\end{equation}
respectively, so that $S_j\ge Mj\ge j$ and $S_j\ge p$. The following four
bounds are used repeatedly.

\begin{enumerate}[label=(\alph*),leftmargin=*]
\item\emph{Child weights.} Let $i\le j$. The merger child
$\mathbf d^{ij',\ell}$ has $p-1$ markings and degree sum $D+\ell-1$, the
nonseparating child $\mathbf d^{i;u,v}$ has $p+1$ markings and degree
sum $D+\ell-2$, and the large child $(v,\mathbf d_J)$ of a separating
term has $p-s$ markings and degree sum at most $D+\ell-2$. In all three
cases the order-$i$ weight of the child is at most $S_j+\ell$. A scalar
factor occurring at loss $r$ has at most $p+r$ markings; if it enters at
order $i$ with $i+r\le j$, its weight is at most $p+r+Mi\le p+Mj\le S_j$,
because $r\le j-i\le M(j-i)$. The dilaton child $(d,0^p)$ of a scalar
with $p$ markings has weight at most $S_j+d+1$, and the children in the
KdV and descent rules have weights at most $S_j+2$.
\item\emph{Moment bound.} For $0\le k\le2j$ and $a\ge0$,
$(S_j+a)^k\le S_j^k(1+a/S_j)^{2j}\le S_j^k\mathrm{e}^{2a/M}$, since
$S_j\ge Mj$. Hence, with
$C_1=C_0\sum_{n\ge0}\rho^{-n}\mathrm{e}^{2(n+1)/M}<\infty$,
\begin{equation}
 \sum_{\ell\ge0}|w_\ell|(S_j+\ell+1)^k\le C_1S_j^k,\qquad
 \sum_{d\ge1}|b_d|(S_j+d+1)^k\le C_1S_j^k\qquad(0\le k\le2j).
 \label{eq:FE-moment}
\end{equation}
\item\emph{Shifts.} For integers $a\ge1$, $q\ge0$ and $0\le c\le c_0$,
\[
 [t^q](1-ct)^{-a}=\frac{a(a+1)\cdots(a+q-1)c^q}{q!}
 \le\frac{((a+q)c_0)^q}{q!},
\]
with the empty product equal to one. For a product
$\prod_k(1-c_kt)^{-a_k}$, taking absolute values if necessary, the
coefficient is bounded by that of $(1-c_0t)^{-a}$, where
$a=\sum_k a_k$ and $c_0\ge\max_k|c_k|$. In an order-$j$ term, a
prefactor of order $r$ and coefficients of total order $i$ leave shift
degree $q=j-r-i$. The total denominator exponent is at most $r+i$,
so $a+q\le j$; also $c_0\le\max(1,r)\le j$. Thus the coefficient is
at most $j^{2q}/q!\le S_j^{2q}/q!$. The case $a=0$ contributes only
$q=0$. After normalization all shifts are therefore majorized by
$\mathrm{e}^{t}$.
\item\emph{Counts.} At a profile state, $\sum_{j'\ne i}(2d_{j'}+1)\le
2D+p\le2S_j$, the number of pairs $(u,v)$ with $u+v=d_i+\ell-2$ is at
most $D+\ell\le S_j+\ell$, and a coordinate path from $\mathbf0$ to
$\mathbf d$ has $D\le S_j$ steps. In a separating term of loss
$r=2j_0+s$ the number of label subsets is $\binom{p-1}s\le p^s/s!$,
there are at most $r+1\le2^r$ pairs $(j_0,s)$, and the small component
has complexity $r-1$ in the profile rule and $r$ in the KdV rule, so
\eqref{eq:FE-small} gives
\begin{equation}
 \binom{p-1}{s}\,|\mathcal W|\le C_0^{\,r}\frac{p^s\Gamma(r)}{s!}
 \le C_0^{\,r}(p+r)^r\le C_0^{\,r}S_j^{\,r}.
 \label{eq:FE-count}
\end{equation}
Indeed $\Gamma(r)/s!\le r^{r-1-s}$ when $s\le r-1$, and $\Gamma(r)/r!\le1$
when $s=r$; and $p+r\le p+Mj=S_j$ because $r\le j$.
\end{enumerate}

\subsection{Bounds for the formal coefficients}\label{app:coefficient-bounds}

\begin{lemma}[Coefficient bounds]\label{lem:coefficient-bounds}
The rules \eqref{eq:FE3}--\eqref{eq:FE6} determine coefficients
$e_{p,j}(\mathbf d)$, $\mathsf C_{p,j}$ and $\mathsf G_{p,j}$ that are
holomorphic on $\Omega$, and there are $C,B_0$ such that, on $\Omega$, for
all $p\ge1$, all degree tuples and all $j\ge0$,
\[
 |e_{p,j}(\mathbf d)|\le CB_0^{\,j}S_j^{2j},\qquad
 |\mathsf C_{p,j}|+|\mathsf G_{p,j}|\le CB_0^{\,j}S_j^{2j},
\]
with the weights \eqref{eq:FE-weights}.
\end{lemma}

\begin{proof}
Define the normalized coefficient sequences
\[
 \varepsilon_j=\sup\frac{|e_{p,j}(\mathbf d)|}{S_j^{2j}},\quad
 \gamma_j=\sup\frac{|\mathsf C_{p,j}|}{S_j^{2j}},\quad
 \eta_j=\sup_{p\ge2}\frac{|\mathsf G_{p,j}|}{S_j^{2j}},\quad
 \zeta_j=\sup\frac{|\mathsf G_{1,j}|}{S_j^{2j}},
\]
the suprema being taken over $Y\in\Omega$, all marking counts $p\ge1$ and
all degree tuples; $\varepsilon_0=\gamma_0=\eta_0=\zeta_0=1$. Write
$\varepsilon(t)=\sum_j\varepsilon_jt^j$, and similarly $\gamma,\eta,\zeta$,
and put $\bar\eta=\max(\eta,\zeta)$ coefficientwise, which bounds every
genus factor. Each coefficient of order $j$ is a finite combination of
coefficients of lower order, summed against $w_\ell$ or $b_d$; by
\eqref{eq:FE-geom} and the polynomial growth in $\ell$ or $d$ of the
lower-order weights these sums converge absolutely and uniformly on
$\Omega$, so each coefficient is holomorphic on $\Omega$ and each
supremum is finite, by induction on $j$. We bound the order-$j$
coefficients term by term.

\textbf{Profile.} Fix a state $(\mathbf d,p)$ and a direction $i$. By
\eqref{eq:FE3}, $[t^j]\mathsf A_{p,i}$ is the sum over $j'\ne i$,
$\ell\ge0$ and $i_1+i_2+q=j-1$ of
\[
 A(2d_{j'}+1)\,w_\ell\;\mathsf C_{p-1,i_1}\;
 e_{p-1,i_2}(\mathbf d^{ij',\ell})\;[t^q](1-t)^{-(1+i_1+i_2)}.
\]
By (a) of Subsection~\ref{app:weights} the two coefficients are at most
$\gamma_{i_1}S_j^{2i_1}$ and $\varepsilon_{i_2}(S_j+\ell)^{2i_2}$; by
(b) the sum over $\ell$ is at most $C_1S_j^{2(i_1+i_2)}$; by (c) the
shift coefficient is at most $S_j^{2q}/q!$; and by (d) the merger count is
at most $2S_j$. Hence
\[
 |[t^j]\mathsf A_{p,i}|\le2|A|C_1\,S_j^{\,2j-1}\,
 [t^{j-1}]\bigl(\mathrm{e}^{t}\gamma(t)\varepsilon(t)\bigr).
\]
The term $\mathsf N_{p,i}$ has in addition one genus coefficient and the
pair count $S_j+\ell$, which \eqref{eq:FE-moment} absorbs with
$k=2(i_1+i_2+i_3)+1\le2j$; it is at most
$2|A|C_1S_j^{2j-1}[t^{j-1}](\mathrm{e}^{t}\gamma\bar\eta\varepsilon)$.
A separating term of loss $r\ge2$ in \eqref{eq:FE4} carries the
prefactor $(At)^r$, $r$ cusp coefficients, $j_0\le r/2$ genus
coefficients, one profile coefficient, the count \eqref{eq:FE-count}, the
pair count and the shifts. Since $\gamma$ and $\bar\eta$ have constant
term one and nonnegative coefficients,
$\gamma^r\bar\eta^{j_0}\preceq\gamma^{2r}\bar\eta^{2r}$, and the term is
at most
$(2|A|C_0)^rC_1S_j^{\,r+1}S_j^{\,2(j-r)}[t^{j-r}](\mathrm{e}^t\gamma^{2r}\bar\eta^{2r}\varepsilon)$,
where $S_j^{r+1}\le S_j^{2r-1}$ because $r\ge2$. Summing the increments
along a coordinate path of at most $S_j$ steps supplies one more factor
$S_j$. After division by $S_j^{2j}$ we obtain, for a constant $b_1$
depending only on $\sup_\Omega|A|$, $C_0$ and $C_1$,
\begin{equation}
 \varepsilon_j\le\Bigl[b_1\mathrm{e}^{t}\varepsilon(t)
 \Bigl(t\gamma\bar\eta+\sum_{r\ge2}(b_1t)^r\gamma^{2r}\bar\eta^{2r}\Bigr)\Bigr]_j
 \qquad(j\ge1).
 \label{eq:FE-maj-E}
\end{equation}

\textbf{Cusp ratio.} By \eqref{eq:FE5}, $[t^j]\mathsf C_p$ is the sum over
$d\ge1$ and $i+q=j$ of $A^{-1}b_d\,e_{p+1,i}(d,0^p)\,[t^q](1+t)^{-i}$.
The child weight is at most $S_j+d+1$ by (a), so \eqref{eq:FE-moment}
and (c) give
\begin{equation}
 \gamma_j\le\bigl[b_2\,\mathrm{e}^{t}\varepsilon(t)\bigr]_j\qquad(j\ge1),
 \qquad b_2=C_1/\inf_\Omega|A|.
 \label{eq:FE-maj-C}
\end{equation}

\textbf{Genus ratio at $p\ge2$.} In \eqref{eq:FE6} the term
$\mathsf P_p(\mathbf e_1;t)$ has weight at most $S_j+1$ and contributes
at most $\mathrm{e}^{2/M}\varepsilon_j$. A separating term of loss $r\ge1$
carries $3\binom{p-2}s|\mathcal W_{j_0,s+2}(\mathbf0)|\le3C_0^rS_j^r\le3C_0^rS_j^{2r}$
by \eqref{eq:FE-count}, the prefactor $(At)^r$, at most $2r$ scalar
coefficients and the shifts. Hence
\begin{equation}
 \eta_j\le\Bigl[b_3\varepsilon(t)+b_3\mathrm{e}^{t}
 \sum_{r\ge1}(b_3t)^r\gamma^{2r}\bar\eta^{2r}\Bigr]_j\qquad(j\ge1),
 \label{eq:FE-maj-H}
\end{equation}
for a constant $b_3$ depending only on $M$, $\sup_\Omega|A|$ and $C_0$.

\textbf{Descent.} The identity
$\mathsf G_1(h)=\mathsf G_2(h+1)\mathsf C_3(h)/\mathsf C_1(h)$ is a
product of a genus series with the shift $t/(1+t)$, a cusp series and
the formal inverse of a cusp series; the three weights are at most
$S_j+2$. If $|\mathsf C_{1,i}|\le\gamma_iS_j^{2i}$ for $i\le j$, then the
coefficients of $1/\mathsf C_1$ through order $j$ are bounded by
$S_j^{2i}$ times those of $1/(2-\gamma(t))=\sum_{k\ge0}(\gamma-1)^k$.
Therefore
\begin{equation}
 \zeta_j\le\Bigl[b_4\,\mathrm{e}^{t}\frac{\eta(t)\gamma(t)}{2-\gamma(t)}\Bigr]_j
 \qquad(j\ge1),\qquad b_4=\mathrm{e}^{12/M}.
 \label{eq:FE-maj-G}
\end{equation}

\textbf{The majorant system.} Put $b=\max(1,b_1,b_2,b_3,b_4)$ and let
$E_*,C_*,H_*,G_*$ be the formal solution, with constant terms one, of
\begin{equation}
\begin{aligned}
 E_*&=1+b\,\mathrm{e}^{t}E_*\Bigl(tC_*G_*+\frac{(btC_*^2G_*^2)^2}{1-btC_*^2G_*^2}\Bigr),&
 C_*&=1+b(\mathrm{e}^{t}E_*-1),\\
 H_*&=1+b(E_*-1)+b\,\mathrm{e}^{t}\frac{btC_*^2G_*^2}{1-btC_*^2G_*^2},&
 G_*&=1+b\Bigl(\mathrm{e}^{t}\frac{H_*C_*}{2-C_*}-1\Bigr).
\end{aligned}
\label{eq:FE-system}
\end{equation}
Here $E_*$ bounds the profile coefficients after division by their
weights, $C_*$ the cusp coefficients, $H_*$ the genus coefficients at
$p\ge2$ and $G_*$ the descent. Since $b\ge1$, $C_*/(2-C_*)\succeq1$ and
$\mathrm{e}^t\succeq1$, the last equation gives $G_*\succeq H_*$; so
$G_*$ majorizes $\bar\eta$ as soon as $H_*,G_*$ majorize $\eta,\zeta$.
Writing $E_*=1+v_1$, $C_*=1+v_2$, $H_*=1+v_3$, $G_*=1+v_4$, the
right-hand sides of \eqref{eq:FE-system} are power series in $(t,v)$ with
nonnegative coefficients, convergent near the origin and vanishing there.
At $t=0$ the first equation has no linear term, the second and third are
linear in $v_1$ only, and the fourth depends only on $v_2,v_3$; so the
Jacobian at the origin is strictly lower triangular, and
Lemma~\ref{lem:majorant} shows that the formal solution of
\eqref{eq:FE-system} converges near $t=0$. The inequalities
\eqref{eq:FE-maj-E}--\eqref{eq:FE-maj-G} now give
$\varepsilon_j\le[E_*]_j$, $\gamma_j\le[C_*]_j$, $\eta_j\le[H_*]_j$ and
$\zeta_j\le[G_*]_j$ for all $j$, by induction on $j$ in the stage order:
at order $j$ the right-hand sides involve $\bar\eta$ only through orders
below $j$, or through $\eta_j$ in the descent, and
$\max(\eta_i,\zeta_i)\le[G_*]_i$ as soon as $\eta_i\le[H_*]_i$ and
$\zeta_i\le[G_*]_i$. Cauchy's estimate for the convergent solution gives
$\varepsilon_j,\gamma_j,\eta_j,\zeta_j\le CB_0^{\,j}$ for all $j$, with
$C,B_0$ depending only on the system.
\end{proof}

\subsection{Tails over topological types and first-order estimates}\label{app:tails}

We next bound the tails of the sums over topological types, before
truncating any expansion. In the sums $S_{g,p}$ and $T_{g,p}$ of
\eqref{eq:TL18}, orient the smaller-complexity component. In $S_{g,p}$ of \eqref{eq:TL18} its loss is its complexity
plus one; in $T_{g,p}$ its loss is its complexity. Write
$S_{g,p}^{[L]}$ and $T_{g,p}^{[L]}$ for the parts of these sums in which
the loss is at least $L$; for $p=1$, $S_{g,1}$ is the one-marking
separating sum of \eqref{eq:TL26}.

\begin{lemma}[Tails over topological types]\label{lem:tails}
Fix $M_0>0$. There are $C_{M_0},B_{M_0}$ such that, for all sufficiently
large $h$ and all $1\le p\le M_0\sqrt h$,
\begin{equation}
 \begin{aligned}
 S_{g,p}^{[L]}&\le C_{M_0}\left(\frac{B_{M_0}(p+L)}h\right)^L,
                  &&2\le L\le\sqrt h,\\
 T_{g,p}^{[L]}&\le C_{M_0}\left(\frac{B_{M_0}(p+L)}h\right)^L,
                  &&1\le L\le\sqrt h,
 \end{aligned}
 \label{eq:FE13}
\end{equation}
the second bound being used only for $p\ge2$. The ranges of $L$ start at
the stability minima two and one respectively.
\end{lemma}

\begin{proof}
Note first that, before specializing to fixed markings,
Lemmas~\ref{lem:lower} and~\ref{lem:upper}, with zero count and marked count
both $p$, give, using the common local limit \eqref{eq:H17},
\[
 \log W_{g,p}=\log\Gamma(h)-h\log A
 +O_{M_0}(\sqrt h\log(h+1)+\log^2(h+1)).
\]
Indeed the harmonic probability has displacement $-p/2=O(\sqrt h)$
and lies between $c_{M_0}/\sqrt h$ and one; the terminal factors cost
$O(p\log(p+1)+(p+\log(h+1))\log(h+1))$.
For losses between $\varepsilon h$ and half the parent complexity,
Stirling entropy dominates this error. Both children have comparable
complexity and $O(\sqrt h)$ markings. The $2^p$ subsets cost only
$\mathrm{e}^{O(\sqrt h)}$, so their total is $O(\mathrm{e}^{-c_\varepsilon h})$.
For $r\le\varepsilon h$, \eqref{eq:TL16}--\eqref{eq:TL17} give the
complementary ratio at most $(C/h)^r$. If its marking count is one,
insert $2j_0+1$ cusps before raising genus $j_0$ times, then the
remaining $s-1$ cusps; this uses $p\ge2$ and avoids a one-marking
genus comparison. The same subset/Gamma bound gives
$a_r=[C(p+r)/h]^r$, after absorbing the number of types into $C^r$.
Their consecutive ratio satisfies
\[
 \frac{a_{r+1}}{a_r}
 \le \frac{eC(p+r+1)}h
 \le eC\left(\frac{M_0}{\sqrt h}+\varepsilon+\frac1h\right).
\]
Choose $\varepsilon$ first, and then a lower bound on $h$ depending on
$M_0$, so this is at most $1/2$. Thus the sum over $L\le r\le\varepsilon h$
is at most $2a_L$. The macroscopic tail also fits \eqref{eq:FE13}:
for $1\le L\le\sqrt h$,
\[
 \left(\frac{p+L}{h}\right)^L
 \ge \exp\{-\sqrt h\log h\},
\]
which eventually exceeds $\mathrm{e}^{-c_\varepsilon h}$. Enlarging the constants
covers the bounded remaining complexities and proves the asserted bound.
At $p=1$, compare to $T_{g,2}$: two uses of \eqref{eq:TL16} and the
proved fixed-one-marking cusp ratio bound multiply its summand by
at most $C/j_0$ for $j_0\le g/2$. The loss is $2j_0$ in both sums.
This proves \eqref{eq:FE13}, including $p=1$.
\end{proof}

\begin{lemma}[First-order estimates]\label{lem:first-order}
Fix $M_0>0$. For all sufficiently large $h$, all $1\le p\le M_0\sqrt h$
and all degree tuples,
\[
 0\le1-c_{g,p}(\mathbf d)
 \le\min\{1,C_{M_0}D(D+p+1)/h\},\qquad
 |\mathsf C_p(h)-1|+|\mathsf G_p(h)-1|\le C_{M_0}(p+1)/h.
\]
\end{lemma}

\begin{proof}
Subtract the recursion at adjacent degrees, as in the proof of
\eqref{eq:TL23}, and bound its separating part by Lemma~\ref{lem:tails} at
the stability minima. For the profile, each increment costs
$C[(D+p+1)/h+(D+1)(p+2)^2/h^2]$; sum the path.
Dilaton sums the child bound against absolute $b_d$ moments, KdV uses
degree one and $T^{[1]}$, and the exact descent gives $p=1$.
These estimates concern the true quantities, in a range of marking counts
that grows with $h$.
\end{proof}

\subsection{Remainder bounds: proof of Proposition~\ref{prop:hierarchy}}\label{app:remainder-bounds}

\begin{proof}[Proof of Proposition~\ref{prop:hierarchy}]
The coefficient bounds are Lemma~\ref{lem:coefficient-bounds}, with
$S_j^{2j}\le(D+p+Mj)^{2j}$. It remains to prove the remainder bounds
\eqref{eq:FE7}, the path independence and the finiteness of the
coefficient rules.

For $L\ge1$ define the normalized remainders
\[
 \mathscr E_L=\sup\frac{\bigl|c_{g,p}(\mathbf d)-\sum_{j<L}e_{p,j}(\mathbf d)h^{-j}\bigr|}{S_L^{2L}h^{-L}},
 \qquad
 \mathscr C_L=\sup\frac{\bigl|\mathsf C_p(h)-\sum_{j<L}\mathsf C_{p,j}h^{-j}\bigr|}{S_L^{2L}h^{-L}},
\]
and $\mathscr H_L,\mathscr G_L$ in the same way for $\mathsf G_p$ with
$p\ge2$ and with $p=1$; the suprema are over real $Y\in[0,y_0]$, all
stable states with $h\ge1$, $p\ge1$, and all degree tuples, invalid
tuples having profile zero. We prove by induction on $L$, in the stage
order profile, cusp, genus, descent, that these numbers are finite and
at most $CB^L$; this is \eqref{eq:FE7}.

\textbf{Large weights.} Put $\delta=S_L^2/h$ and fix $\delta_0\in(0,1]$
with $B_0\delta_0<1/2$. The exact quantities are bounded: $0\le c\le1$,
$|\mathsf C_p(h)|\le C$ by \eqref{eq:TL16}, $\mathsf G_p(h)\le1$ for
$p\ge2$ by \eqref{eq:TL17}, and $\mathsf G_1(h)\le C$ by the case $p=1$ of \eqref{eq:TL27}.
If $\delta\ge\delta_0$, then Lemma~\ref{lem:coefficient-bounds} and $S_j^{2j}h^{-j}\le\delta^j$
($j\le L$) bound the normalized remainder by
\[
 \delta^{-L}\Bigl(C+\sum_{j<L}CB_0^{\,j}\delta^{j}\Bigr)
 \le C\delta_0^{-L}\Bigl(1+\sum_{j<L}(B_0\delta_0)^j\Bigr)\le2C\delta_0^{-L},
\]
since $\delta^{j-L}\le\delta_0^{j-L}$ for $j<L$. Invalid-degree parents
satisfy $D\ge h$, hence $\delta\ge h\ge1$, and lie in this region. In the remaining region
$\delta<\delta_0$ one has $p+ML\le S_L<\sqrt h$, so every marking count
that occurs, including those of children, is at most $p+L+2\le4\sqrt h$:
Lemmas~\ref{lem:tails} and~\ref{lem:first-order} apply with $M_0=4$. Also $h>M^2L^2/\delta_0$, and
shrinking $\delta_0$ once makes Lemma~\ref{lem:first-order} give $\mathsf C_p(h)\ge1/2$
wherever a cusp ratio is inverted. From now on $\delta<\delta_0$.

\textbf{Products of exact factors.} Every term of the exact adjacent
difference, of the exact dilaton identity, of the exact KdV identity and
of the exact descent is a product of exact factors: a prefactor
$A^r\prod(h-c_k)^{-1}$, cusp and genus ratios at shifted complexities,
and profiles of children. Each exact factor $F$ has, for every $q\ge1$,
a decomposition $F=T_{<q}F^{\mathrm{form}}+R_F(q)$ into its truncated
formal series and a remainder, and the formal rule at order $L$ is the
truncation of the product of the formal series. The product of exact
factors and the truncated formal product differ by the finite identity
\begin{equation}
 \prod_kF_k-T_{<q}\prod_kF_k^{\rm form}
 =R_{F_1}(q)\prod_{k>1}F_k
 +\sum_{j<q}f_{1,j}h^{-j}\Bigl(\prod_{k>1}F_k-T_{<q-j}\prod_{k>1}F_k^{\rm form}\Bigr),
 \label{eq:FE-product}
\end{equation}
applied again to the bracket on the right, and so on. Each resulting term
contains exactly one remainder $R_{F_k}(q_k)$, multiplied by known
coefficients of the factors before it and by exact factors after it. For
an inverse, with known formal inverse coefficients $d_j$,
\begin{equation}
 \mathsf C\,T_{<q}(1/\mathsf C)-1
 =\sum_{j<q}d_jh^{-j}\bigl(\mathsf C-T_{<q-j}\mathsf C\bigr),
 \label{eq:FE-inverse}
\end{equation}
so the remainder of $1/\mathsf C$ is $1/\mathsf C$ times a sum of known
coefficients against remainders of $\mathsf C$.

\textbf{Bounding one term.} Normalize by $S_L^{2L}h^{-L}$. A known
coefficient of order $j'$ of a factor with weight $S'$ is at most
$CB_0^{j'}(S')^{2j'}$ by Lemma~\ref{lem:coefficient-bounds}, where $S'\le S_L+\ell$, $S_L+d+1$,
$S_L+2$ or $S_L$ by (a) of Subsection~\ref{app:weights}, applied with $j'\le L$. A remainder
$R_{F_k}(q_k)$ of a factor at complexity $h-c$, $|c|\le L$, is at most
$\mathscr X_{q_k}(S')^{2q_k}(h-c)^{-q_k}$, where $\mathscr X$ is the error
sequence of its type, and
$(h-c)^{-q_k}\le h^{-q_k}\mathrm{e}^{2|c|q_k/h}\le\mathrm{e}^{2\delta_0/M^2}h^{-q_k}$
because $q_k\le L$ and $L^2/h\le\delta_0/M^2$. The exact factors after
the remainder are bounded by the constants above. The known shift tails
are bounded as follows. If the orders of the known coefficients and of
the prefactor sum to $a=L-q$ and the first omitted shift degree is $q$,
the elementary inequality
\[
 \sum_{n\ge q}\frac{(a)_n}{n!}x^n\le\frac{(a)_q}{q!}x^q(1-x)^{-a-q}
 \qquad(a\ge1,\ 0\le x<1),
\]
which follows from $(a)_{q+k}=(a)_q(a+q)_k$ and $(q+k)!\ge q!k!$, bounds
the tail by $\binom{L-1}{q}(c/h)^q(1-c/h)^{-L}$ with $c\le L$; after
division by $S_L^{2q}h^{-q}$ this is at most
$\mathrm{e}^{2\delta_0/M^2}/q!$, since $c(L-1)\le L^2\le S_L^2$
and $(1-c/h)^{-L}\le\mathrm e^{2cL/h}$ in the small-weight region.
The counts, kernel moments and path sums are
handled exactly as in the proof of Lemma~\ref{lem:coefficient-bounds}, with the same powers of $S_L$: a linear
term supplies $S_L^2$ against the prefactor order one, and a separating
term of loss $r$ supplies $S_L^{r+2}\le S_L^{2r}$ ($r\ge2$) in the profile
rule and $S_L^{r}\le S_L^{2r}$ ($r\ge1$) in the KdV rule.

\textbf{Forcing terms.} Three contributions do not involve the error
sequences. The exact topological types of loss $r>L$, which the formal
rule omits, have total mass bounded by \eqref{eq:FE13} with $L+1$ in
place of $L$; after the kernel sum and the coordinate path their profile
contribution is at most
$CD(D+1)(B(p+L)/h)^{L}\le CB^LS_L^{L+2}h^{-L}\le CB^LS_L^{2L}h^{-L}$ for
$L\ge2$, and the KdV contribution has no path or pair factor. The
difference between the exact kernel sum over $\ell\le m$ and the formal
sum over all $\ell$, and the macroscopic cutoffs in the proof of Lemma~\ref{lem:tails}, are
$O(\mathrm{e}^{-ch})$, and
$\mathrm{e}^{-ch}\le C^LL!\,h^{-L}\le C^LS_L^{2L}h^{-L}$. At $L=1$ the
normalized remainders are bounded by Lemma~\ref{lem:first-order}.

\textbf{The linear majorant system.} Let $F_0=C/(1-B_0t)$ be the majorant
of the normalized known coefficients from Lemma~\ref{lem:coefficient-bounds}, and
$V=C_{\rm ex}+F_0$, where $C_{\rm ex}\ge1$ bounds all exact factors and
$1/\mathsf C$ in the region $\delta<\delta_0$. Collecting the bounds
above, the error sequences $\mathscr E,\mathscr C,\mathscr H,\mathscr G$,
regarded as formal series in $t$, satisfy coefficientwise
\begin{equation}
\begin{aligned}
 \mathscr E&\preceq F+K_E\mathscr R,&
 \mathscr C&\preceq F+b\mathscr E,\\
 \mathscr H&\preceq F+b\mathscr E+K_H\mathscr R,&
 \mathscr G&\preceq F+b(\mathscr H+\mathscr C)+K_G\mathscr R,
 \qquad \mathscr R=\mathscr E+\mathscr C+\mathscr H+\mathscr G,
\end{aligned}
\label{eq:FE-error-system}
\end{equation}
with the zero-constant kernels
\[
 K_E=b\,\mathrm{e}^{bt}\Bigl[tV^2+\sum_{r\ge2}(2r+1)(bt)^rV^{2r}\Bigr],\quad
 K_H=b\,\mathrm{e}^{bt}\sum_{r\ge1}2r(bt)^rV^{2r},\quad
 K_G=b\bigl[\mathrm{e}^{bt}V^3-V(0)^3\bigr],
\]
and the forcing
\[
 F=\frac{bV^3\mathrm{e}^{bt}}{(1-btV^2)^3}+\frac{b}{1-\delta_0^{-1}t},
\]
for a constant $b$ depending only on the constants already fixed. The
factors $2r+1$ and $2r$ count the possible positions of the single
remainder among the factors of a product with $r$ cusp, at most $r$ genus
and at most one profile factor; the term $b\mathscr E$ in the second and
third lines is the current-order remainder of the profile, which enters
the dilaton and KdV rules without a prefactor $t$, and $b(\mathscr H+\mathscr C)$
plays the same role in the descent. The first summand of $F$ dominates
the shift tails, the omitted topology and the cutoff terms, each of which
is a product of at most $2r+1$ known or exact factors times $C^q/q!$ or
$C^LL!/S_L^{2L}$; the second summand dominates the region
$\delta\ge\delta_0$.

\textbf{Induction.} Substituting the first line of
\eqref{eq:FE-error-system} into the others gives
$\mathscr R\preceq A_0F+A_0(K_E+K_H+K_G)\mathscr R$ for a constant
$A_0$. Since $K_E+K_H+K_G$ has no constant term, the coefficient of
$t^L$ on the right involves $\mathscr R_j$ only for $j<L$. Let
$\mathscr R_*=A_0F/(1-A_0(K_E+K_H+K_G))$, a series with nonnegative
coefficients, analytic near $t=0$. If $\mathscr R_j\le[\mathscr R_*]_j$
for all $j<L$, then
$\mathscr R_L\le[A_0F+A_0(K_E+K_H+K_G)\mathscr R_*]_L=[\mathscr R_*]_L$
by monotonicity of series with nonnegative coefficients; the case $L=1$
is the base estimate. Cauchy's estimate for $\mathscr R_*$ gives
$\mathscr R_L\le CB^L$ and proves \eqref{eq:FE7}; no infinite
inverse-complexity series is evaluated. For each fixed degree tuple, two
coordinate paths give asymptotic expansions of the same ratio, and
uniqueness of asymptotic coefficients as $g\to\infty$ proves the asserted
path independence.

\textbf{Finiteness of the coefficient rules.} It remains to justify the finite coefficient rule and polynomial
bounds used above. At each fixed order, the degree lattice admits a
finite boundary stratification: each coordinate is either one of
finitely many low integers or lies above that threshold. The
coefficients are polynomial on every stratum. A merger has an
exceptional low child degree only when $d_i+d_j+\ell-1$ is bounded;
otherwise polynomial summation uses finitely many $w$ moments.
For a nonseparating pair, separate bounded $u$, bounded $v$, and the
interior. The interior sum is a Faulhaber polynomial in $u+v$;
overlap of the two boundary ranges adds finitely many exceptions.
In a retained split, dimension bounds $u$ and $\mathbf d_I$ by
$3j+s-2$, so the same argument applies to the large child.
Summing the coordinate path preserves this class. One admissible (not
optimal) choice of the low-degree thresholds is given by $b_0^{\rm deg}=0$ and
$b_K^{\rm deg}=2b_{K-1}^{\rm deg}+2K+4$.
The adjacent coefficient has degree at most $2K-1$, and its path sum
has degree at most $2K$, proving the claimed growth.
The single-degree specialization in dilaton is eventually polynomial;
its sum uses finitely many signed $b_d$ moments and low $b_d$.
Exact small amplitudes are finite polynomials in the $\theta_\ell$. These facts
reduce every displayed infinite sum in the coefficient algorithm to
finite kernel data. The common analytic disks and removable critical
normalization give analyticity of that data throughout $[0,y_0]$.
\end{proof}

\subsection{Determining the constants at the pure-\texorpdfstring{$\psi$}{psi} endpoint}\label{sec:endpoint-expansion}

To integrate the transport equation to all orders, we need the value of
every coefficient at the pure-$\psi$ endpoint, with constants controlled
as the order grows. This subsection supplies these values from the
Painlev\'e~I asymptotics and then proves Theorem~\ref{thm:full-expansion}.

Invert the cusp expansion and write, at each fixed order,
\begin{equation}
 \frac1{\mathsf C_p(h)}
 =1+\frac{\beta_p}h+\sum_{j=1}^K\frac{\gamma_{p,j}}{h^{j+1}}
      +O_{p,K}(h^{-K-2}).
 \label{eq:FE8}
\end{equation}
Here $\gamma_{p,1}=\gamma_p$. The exact endpoint identity following
\eqref{eq:ST12} gives $1/\mathsf C_p(h;y_0)=1-p/(5h)$.
Uniformity and uniqueness of successive coefficients imply
\begin{equation}
 \gamma_{p,j}(y_0)=0\quad(j\ge1),\qquad
 \gamma_{p,j}/(Ar)\ \hbox{is analytic at }y_0.
 \label{eq:FE9}
\end{equation}
Analyticity follows from the finite coefficient rules, not from remainder
estimates. Division by the simple zero of $Ar$ is a bounded operation
on a fixed smaller complex neighborhood: write a vanishing function
as $(Y-y_0)\int_0^1 f'(y_0+t(Y-y_0))\,\mathrm{d}t$ and apply Cauchy's bound
once on the larger neighborhood. No successive shrinkage in order occurs.

We need a quantitative version of the pure-$\psi$ endpoint input.
Use the conventions of Garoufalidis--Its--Kapaev--Mari\~no~\cite{GIKM}:
$-u''/6+u^2=Z$, zero-instanton coefficients $u_{g,0}$,
$\mathfrak a=8\sqrt3/5$, $\nu_g=2g-1/2$ and
$S_1=-i3^{1/4}/(2\sqrt\pi)$. Put
\[
 Z_g=\frac{\pi i}{S_1}
       \frac{\mathfrak a^{\nu_g}u_{g,0}}{\Gamma(\nu_g)}.
\]
The Borel representation of Costin--Costin--Huang~\cite[Theorem~2]{CCH}
makes the coefficient formula in \cite[Section~2.3]{GIKM} quantitative,
as follows. Set $\rho=6^{-1/5}\mathrm{e}^{i\pi/5}$;
$y(Z_c)=\rho^2u(\rho Z_c)$ satisfies $y''=6y^2+Z_c$.
Their Boutroux coordinate is $\xi=\mathfrak a Z^{5/4}$, and the
normalized remainder is $u(Z)/\sqrt Z-1+4/(25\xi^2)$. Their variable
$x$ is our $\xi$, and their unknown $h(x)$ is this normalized remainder.
The adjacent GIKM solutions $u_1,u_0$ have sectors
$(-2\pi/5,6\pi/5)$ and $(-6\pi/5,2\pi/5)$ respectively.
Their normalized remainders are $O(\xi^{-4})$ on $+i\infty$ and
$-i\infty$. Thus the respective lateral constants in the exact Borel
representation vanish: multiplying by $\mathrm{e}^\xi\xi^{1/2}$ gives zero
on those rays. This identifies the lateral Borel sums with the
solutions themselves, not merely with their formal expansions.

The Borel transform at its first positive singularity has the form
$S(1-s)^{-1/2}\alpha(s)+\beta(s)$, where $\alpha,\beta$ are analytic
near one and $\alpha(1)=1$~\cite[Proposition~3(i)]{CCH}. Keep its multiplier symbolic
and normalize the jump between the two solutions in the convention of~\cite{GIKM}:
\[
 \mathfrak J(\xi)=
 \frac{u_1(Z)-u_0(Z)}{-S_1\mathrm{e}^{-\xi}Z^{-1/8}},\qquad
 \mathfrak J(\xi)\longrightarrow1.
\]
Choose a fixed small $r>0$ and angles $\pm\phi$ so
$(1+r)\cos\phi>1+r/2$. Deform only the finite initial parts of the
two Laplace rays into a local Hankel contour about $[1,1+r]$ and
compact joining arcs on $|s|=1+r$. The infinite ray tails are kept
unchanged. The analytic continuation and weighted $L^1$ estimates
of~\cite[Lemma~17(c),(f)]{CCH} bound these outer parts by $C\mathrm{e}^{-(1+\varepsilon)\xi}$,
with constants independent of Taylor order.
The local analytic term has no jump. Cauchy's estimate on $|s-1|\le2r$
bounds the degree-$L-1$ remainder of $\alpha(1+t)$ by
$C(2r)^{-L}t^L$. Its normalized Laplace integral is bounded by
$CB^L\Gamma(L+1/2)\xi^{-L}$. Extending the retained polynomial
integrals from $r$ to infinity obeys the same bound; the outer term
$O(\sqrt\xi \mathrm{e}^{-\varepsilon\xi})$ is at most $CB^LL!\xi^{-L}$.
Consequently, for all $L\ge1$ and real $\xi\ge\xi_0$,
\[
 \left|\mathfrak J(\xi)-\sum_{j<L}v_j\xi^{-j}\right|
       \le CB^LL!\xi^{-L},\qquad v_j=\mathfrak a^ju_{j,1}.
\]
No numerical Stokes multiplier from the other normalization is used.

The first coefficient is $v_1=-1/8$, and we derive it here because it
determines the first correction at the endpoint. The difference
$u_1-u_0$ is exponentially small and satisfies the linearization
$(u_1-u_0)''=12u_0(u_1-u_0)$ of $-u''/6+u^2=Z$ up to a term quadratic in
$u_1-u_0$, which is $O(\mathrm{e}^{-2\xi})$ and does not affect the
coefficients $v_j$. Here $12u_0=12Z^{1/2}-Z^{-2}/4+O(Z^{-9/2})$, because
$u_{1,0}=-1/48$. Put $f=Z^{-1/8}\mathrm{e}^{-\xi}$. Since
$\xi'=2\sqrt3Z^{1/4}$ and $\xi''=(\sqrt3/2)Z^{-3/4}$,
\[
 \frac{f''}f=\Bigl(\frac1{8Z}+\xi'\Bigr)^2+\frac1{8Z^2}-\xi''
 =12Z^{1/2}+\frac9{64}Z^{-2};
\]
the terms of order $Z^{-3/4}$ cancel, which is why the prefactor is
$Z^{-1/8}$. Writing $u_1-u_0=-S_1f\mathfrak J$ and dividing the
linearized equation by $-S_1f$ gives
$(f''/f)\mathfrak J+2(f'/f)\mathfrak J'+\mathfrak J''=12u_0\mathfrak J$.
The local Laplace representation above gives the expansion of
$\mathfrak J$ in a sector about the positive ray; Cauchy's estimate
therefore permits differentiation of its remainder on a smaller sector.
With $\mathfrak J=1+v_1\xi^{-1}+O(\xi^{-2})$ one has
$\mathfrak J'=-v_1\xi'\xi^{-2}+O(Z^{-7/2})=-(2\sqrt3v_1/\mathfrak a^2)Z^{-9/4}+O(Z^{-7/2})$,
so $2(f'/f)\mathfrak J'=-2\xi'\mathfrak J'+O(Z^{-13/4})=(24v_1/\mathfrak a^2)Z^{-2}+O(Z^{-13/4})$,
while $\mathfrak J''=O(Z^{-13/4})$. Comparing the coefficients of
$Z^{-2}$ gives $9/64+24v_1/\mathfrak a^2=-1/4$, that is
$v_1=-25\mathfrak a^2/1536=-1/8$, since $\mathfrak a^2=192/25$.
Equivalently $u_{1,1}=-5/(64\sqrt3)$.

The exact coefficient contour in~\cite[Section~2.3]{GIKM} expresses
$u_{g,0}$ as a circle integral and deforms it onto the five rays
$\arg t=2\pi k/5$, $k=0,\pm1,\pm2$; the symmetry of the tritronqu\'ee
solutions reduces their sum to five times the positive ray
\cite[equations~(2.25) and~(2.29)]{GIKM}:
\[
 u_{g,0}=-\frac5{2\pi i}\int_\rho^\infty t^{5g-2}
        \bigl(u_1(t^2)-u_0(t^2)\bigr)\,\mathrm{d}t+r_g(\rho),
 \qquad |r_g(\rho)|\le C\rho^{5g},
\]
for a fixed $\rho$, large enough for~\cite{GIKM} and such that
$\xi_0=\mathfrak a\rho^{5/2}$ is at least the threshold in the jump
expansion above. Put $\xi=\mathfrak at^{5/2}$, so that
$t^{5g-9/4}\,\mathrm{d}t=\tfrac25\mathfrak a^{-\nu_g}\xi^{\nu_g-1}\,\mathrm{d}\xi$,
and write $u_1-u_0=-S_1\mathrm{e}^{-\xi}t^{-1/4}\mathfrak J(\xi)$. The
factor $5\cdot\tfrac25$ cancels the $2$ in $2\pi i$, and the definition
of $Z_g$ gives exactly
\[
 Z_g=\frac1{\Gamma(\nu_g)}\int_{\xi_0}^{\infty}
        \mathrm{e}^{-\xi}\xi^{\nu_g-1}\mathfrak J(\xi)\,\mathrm{d}\xi
       +O\!\left(\frac{D^g}{\Gamma(\nu_g)}\right);
\]
the error term is $(\pi i/S_1)\mathfrak a^{\nu_g}r_g(\rho)/\Gamma(\nu_g)$. Applying the jump remainder
with $L=K+1\le g/2$ gives
$CB^LL!\Gamma(\nu_g-L)/\Gamma(\nu_g)\le CB^LL!g^{-L}$.
The omitted low integrals and inner circle are
$O(\mathrm{e}^{-c g\log g})$ uniformly for $K\le c_0g$, with $c_0$ small.
For the low integrals, bound the finite sum by
$C\xi_0^{\nu_g}\sum_{j<L}B^jj!\xi_0^{-j}/\Gamma(\nu_g)$
and use Stirling's lower bound.
Thus the finite falling-factorial expansion of $Z_g$ has a Gevrey-1
remainder. Re-expansion in $g^{-1}$ preserves it:
\[
 \frac{\Gamma(\nu_g-j)}{\Gamma(\nu_g)}
 =(2g)^{-j}\prod_{a=1}^j\left(1-\frac{a+1/2}{2g}\right)^{-1}.
\]
Cauchy on radius $1/(D(K+1))$ for the product, followed by
$j!(K+1)^{K+1-j}\le \mathrm{e}^{K+1}(K+1)!$, bounds the total omitted
inverse-genus powers by $CB^{K+1}(K+1)!g^{-K-1}$.

The exact identity of Itzykson--Zuber~\cite[equation~(6.10)]{IZ} is
\[
 \langle\tau_2^{3g-3}\rangle_g
 =\frac{2^g(3g-3)!}{(5g-5)(5g-3)}\,\phi_g,\qquad g\ge2,
\]
where $\phi(t)=\sum_{g\ge0}\phi_g t^{(1-5g)/2}$, $\phi_0=-1$,
satisfies $\phi''/3+\phi^2=t$. The substitution
$\phi(t)=-2^{1/5}u(2^{-2/5}t)$ gives exactly
$\phi_g=-2^g u_{g,0}$.
Combined with the expansion of $Z_g$ above, this gives, in the normalization of \eqref{eq:H4} and for every fixed $K$,
\begin{equation}
 \pi C_g(2^{3g-3})=1+\sum_{j=1}^K d_jg^{-j}+O_K(g^{-K-1}).
 \label{eq:FE10}
\end{equation}
The leading constant is also independently given by
\cite[Theorem~1]{GYZ}. More precisely, put $h_0=2g-2$, $h=h_0+p$,
$a_*=2/15$, and define the exact normalized endpoint
\[
 E_{g,p}=\frac{\pi\overline U_{g,p}(0^p;y_0)}
 {\Gamma(h)4^{-g}a_*^{-h}\sqrt{a_*/(\pi h)}}.
\]
For $p\ge1$ this is the limit of $\pi\Phi_g^{(p)}/\mathcal J_{g,p}$;
at $p=0$ it is defined by the same exact pure-$\psi$ expression.
String and dilaton give
$E_{g,p+1}/E_{g,p}=(h_0+4p/5)(h_0+p)^{-1}
\sqrt{(h_0+p+1)/(h_0+p)}$. The base conversion gives exactly
$\overline U_{g,0}(y_0)=-108^gu_{g,0}/[27(5g-5)(5g-3)]$.
Substitution therefore proves
\begin{equation}
 \begin{split}
 E_{g,p}&=Z_g\,\frac{\Gamma(h_0+3/2)}{\Gamma(h_0)h_0^{3/2}}
             \frac{R_p(h_0)}{1+4/(5h_0)},\\
 R_p(h_0)&=\sqrt{1+p/h_0}\prod_{i=0}^{p-1}
                       \frac{1+4i/(5h_0)}{1+i/h_0}.
 \end{split}
 \label{eq:FE14}
\end{equation}
This identity is exact for $g\ge2,p\ge0$, with the empty product one.

For fixed $p$, the endpoint value and its logarithm therefore have
Gevrey-1 remainders through order $K\le c_pg$. It remains to treat the
Gamma factor; for fixed $\alpha>0$,
\[
 \log\frac{\Gamma(h_0+\alpha)}{\Gamma(h_0)h_0^\alpha}
 =\int_0^\infty \mathrm{e}^{-h_0t}
 \left[\alpha-\frac{1-\mathrm{e}^{-\alpha t}}{1-\mathrm{e}^{-t}}\right]\frac{\mathrm{d}t}{t}.
\]
The kernel is analytic at zero on a fixed disk and bounded outside a
small positive interval. Its global Taylor remainder is bounded by
$CB^Lt^{L-1}$, giving $CB^L(L-1)!h_0^{-L}$ after integration.
The finite insertion factors have convergent expansions at $h_0^{-1}=0$.
Finite shifts from $h_0$ to $h=h_0+p$ preserve the factorial error by the
same Cauchy estimate for rational products just used above.
For logarithms, exponentials and inverses, take the degree-$K$ polynomial
on radius $1/(D(K+1))$; its unit or zero-constant part stays bounded
away from the relevant singularity. Cauchy's tail and the bounded
real derivative of the composed function preserve the remainder bound.
This proves the stated endpoint bound, including its log coefficients
$a^*_{p,j}$, defined by $\log E_{g,p}\sim\sum_{j\ge1}a^*_{p,j}h^{-j}$.
At first order, \eqref{eq:FE14}, the value $v_1=-1/8$ derived above,
and $\nu_g-1=h_0+1/2$ give
\[
 a^*_{p,1}=-\frac18+\frac38-\frac45+\frac p2-\frac{p(p-1)}{10}
 =\frac{-2p^2+12p-11}{20}=t_p,
\]
the constant in \eqref{eq:ST13}; replacing $h_0$ by $h=h_0+p$ does not
change the first-order coefficient.

In particular the first coefficient in \eqref{eq:FE10} is $d_1=-2/9$.
Indeed, at $p=0$ the definition of $E_{g,0}$ and \eqref{eq:PF13} give
\[
 \pi C_g(2^{3g-3})=\frac{E_{g,0}}{N_{h_0}},\qquad
 N_{h_0}=\frac{\mathcal M_{g,0}(y_0)}
          {\Gamma(h_0)4^{-g}a_*^{-h_0}/\sqrt{2\pi h_0v}}
     =\sqrt{2\pi h_0v}\,\Pr(S_{h_0}=3h_0/2),\qquad v=\frac{15}4,
\]
where $S_{h_0}$ is the sum of $h_0$ independent copies of the geometric law
with generating function $(2/5)/(1-3w/5)$; this is the normalized
negative-binomial coefficient of Theorem~\ref{thm:first-correction} at
displacement zero, so $N_{h_0}=1-19/(180h_0)+O(h_0^{-2})$. On the other side,
$Z_g=1+v_1\Gamma(\nu_g-1)/\Gamma(\nu_g)+O(g^{-2})=1-1/(8h_0)+O(h_0^{-2})$,
because $\nu_g-1=h_0+1/2$ and $v_1=-1/8$;
$\Gamma(h_0+3/2)/(\Gamma(h_0)h_0^{3/2})=1+3/(8h_0)+O(h_0^{-2})$; and
$(1+4/(5h_0))^{-1}=1-4/(5h_0)+O(h_0^{-2})$. Hence, by \eqref{eq:FE14},
\[
 \pi C_g(2^{3g-3})
 =1+\frac1{h_0}\Bigl(-\frac18+\frac38-\frac45+\frac{19}{180}\Bigr)+O(h_0^{-2})
 =1-\frac4{9h_0}+O(h_0^{-2})
 =1-\frac2{9g}+O(g^{-2}),
\]
since $h_0=2g-2$. This value agrees with the expansion recorded by Guo,
Yang and Zagier~\cite[equation~(25)]{GYZ}, which rests on the same
Painlev\'e~I input; it is used in the proof of
Theorem~\ref{thm:first-correction}.

Define
\begin{equation}
 a_{p,j}(Y)=a^*_{p,j}-\int_Y^{y_0}\frac{\gamma_{p,j}(v)}{A(v)r(v)}\,\mathrm{d}v,
 \qquad
 \mathcal L_{p,j}(Y)=[t^j]\exp\!\left(\sum_{i=1}^j a_{p,i}(Y)t^i\right).
 \label{eq:FE11}
\end{equation}
Proposition~\ref{prop:hierarchy}, inversion and the bounded division
and integration above give common-domain bounds
$\|a_{p,j}\|+\|\mathcal L_{p,j}\|\le C_p\mathfrak B_p^j(j!)^2$.
For formal exponentiation this follows also by Cauchy on the finite
log polynomial at radius $1/(Dj^2)$. The first coefficient is
$\mathcal L_{p,1}=\mathcal L_p$. These are the coefficients in
Theorem~\ref{thm:full-expansion}.

\begin{proof}[Proof of Theorem~\ref{thm:full-expansion}]
We first expand the logarithmic derivative, then fix its integration
constant near the endpoint. After integration, exponentiation gives
the claimed expansion. The extra orders retained below compensate for
the error accumulated near the zero of $Ar$.

Put $F_g(Y)=\log(\pi\Phi_g^{(p)}/\mathcal J_{g,p})$.
For fixed $p$ the inverse cusp ratio is bounded, and the finite
inverse identity in Proposition~\ref{prop:hierarchy} transfers its
remainder to $1/\mathsf C_p$ with error
$CB^N(N!)^2h^{-N}$, for $N\le c\sqrt h$.
The exact dilaton and \eqref{eq:PF11} give
\[
 F'_g=\frac{h}{Ar}(1/\mathsf C_p-1-\beta_p/h).
\]
For $L=K+1$, use $N=L+2=K+3$, retaining $\gamma_{p,j}$ through
$j=L$. Its unknown derivative error is
$CB^{L+2}((L+2)!)^2h^{-L-1}/(Ar)$.
Choose $\delta_*=h^{-L-1}$.
The positive comparison \eqref{eq:PF15} and the exact first equality
of \eqref{eq:PF16} change $F_g$ from its endpoint value by at most
$Ch\delta_*=Ch^{-L}$ on this strip. The leading harmonic
approximation in the second equality of \eqref{eq:PF16} is not used.
The endpoint expansion proved above supplies the value at the matching point.
The bounded coefficient derivatives from \eqref{eq:FE9} show that the
whole retained coefficient sum changes by $O(\delta_*/h)$ there,
after choosing the square-root order constant small.

Integrate backwards from $y_0-\delta_*$. Since $Ar$ has a simple zero,
the unknown error costs at most
\[
 CB^{L+2}((L+2)!)^2h^{-L-1}[1+(L+1)\log h].
\]
Relative to $(L!)^2h^{-L}$, its extra factor is a fixed exponential
base times $(L+2)^4[1+(L+1)\log h]/h$, bounded by another fixed
exponential in $L$. Indeed $L\le c\sqrt h$ implies
$h\ge c^{-2}L^2$; the decrease of $\log h/h$ for $h\ge e$ bounds
this factor by $C(L+2)^3\log(L+2)\le C D^L$.
The bounded exceptional values of $h,L$ are absorbed into $C$.
The known coefficients integrate to
\eqref{eq:FE11}. Discarding the last term $a_{p,L}h^{-L}$ proves
the asserted log error, on both the strip and its complement.
Choose $c_p$ smaller and $h_p$ larger once so that the inverse-cusp
remainder estimate remains valid through order $K+3$. The endpoint range is linear in $h$.

Finally the finite degree-$K$ log polynomial is uniformly bounded on
$|t|\le1/(D(K+1)^2)$. Cauchy's exponential tail at $t=1/h$ is
$C[D(K+1)^2/h]^{K+1}\le C(\mathrm{e}^2D)^{K+1}((K+1)!)^2h^{-K-1}$.
At the real evaluation point the polynomial is $O(h^{-1})$ and the
function $F_g$ itself is bounded by Proposition~\ref{prop:prefactor}.
Thus the exponential is Lipschitz on the required bounded interval
and transfers the error bound for the logarithm to \eqref{eq:FE12}.
Choosing $K=\lfloor c'_p\sqrt h\rfloor$ after the final error base,
with $\mathfrak B_p(c'_p)^2<1$, gives the exponential upper error.
\end{proof}

\section{Coefficient extraction with inverse-genus errors}\label{app:extraction}

This appendix proves Theorems~\ref{thm:sharp-extraction}
and~\ref{thm:canonical-expansion}. Subsection~\ref{sec:quantitative-local}
makes the local limit theorem of Section~\ref{extracting-fixed-cusp-counts}
quantitative; its complex comparison and Fourier estimates are reused
afterwards. Subsections~\ref{app:local-complex}--\ref{app:small-indices}
prove Theorem~\ref{thm:sharp-extraction}, and
Subsections~\ref{app:phase-change}--\ref{app:canonical-proof} prove
Theorem~\ref{thm:canonical-expansion} and compute its first coefficient.

\subsection{A quantitative local limit}\label{sec:quantitative-local}

We first make the local limit theorem \eqref{eq:TL47} quantitative. The argument transfers a real-axis error to a complex neighborhood and applies Fourier inversion; the positive convolution factor controls the remaining part of the integral. This gives a power-rate bound. The complex comparison and Fourier estimates established in the proof will also be used in Subsection~\ref{app:local-complex}, where retaining a Poisson factor gives an inverse-genus error uniformly over all coefficient indices.

\begin{proposition}\label{prop:quantitative-local}
Fix $p\ge1$, and use $h,\mu,\sigma$ from \eqref{eq:TL37}. There are
constants $C_p,\kappa_p>0$ such that, with $M=\min(h,\mu)$,
\begin{equation}
 \sup_{k\in\mathbb Z}
 \left|\sigma\mathbb P_x(N=k)-
 \frac{\mathrm{e}^{-(k-\mu)^2/(2\sigma^2)}}{\sqrt{2\pi}}\right|
 \le C_p M^{-\kappa_p},
 \label{eq:QL1}
\end{equation}
whenever $M$ is sufficiently large, uniformly for $0<x<x_c$.
\end{proposition}

\begin{proof}
Write $Z=(N-\mu)/\sigma$ and $\mathcal L(z)=\mathbb E_x \mathrm{e}^{zZ}$.
The integral identity \eqref{eq:TL38} and the convergent endpoint
expansions \eqref{eq:TL36} give constants $c,C_p>0$ for which
\begin{equation}
 \left|\log\mathcal L(u)-\frac{u^2}{2}\right|
 \le C_p M^{-1/2}(|u|+|u|^3),\qquad |u|\le c\sqrt M.
 \label{eq:QL2}
\end{equation}
To check the uniformity, note that near zero $\Psi,\mathscr D\Psi$ and
$\mathscr D^2\Psi$ are comparable to $x$, so the model Taylor remainder
is $O(|u|^3/\sqrt\mu)$ and the integrated error is
$O(|u|x/\sqrt\mu)$. On an interior compact interval both are
$O_p(h^{-1/2}(|u|+|u|^3))$. Near criticality,
$\sigma\asymp\sqrt h/\varepsilon$; for $|u|\le c\sqrt h$ the tilted
distance to the critical point stays comparable to $\varepsilon$.
The bounds $\Psi=O(\varepsilon^{-1})$ and
$\mathscr D^2\Psi=O(\varepsilon^{-3})$ give the same estimate.
Shrinking $c$ ensures the real moment generating function is defined
on the whole interval. As before, the error in \eqref{eq:TL30} is never
differentiated.

Put $T=\sqrt{\log M/80}$ and $f(z)=\mathcal L(z)-\mathrm{e}^{z^2/2}$.
The function is analytic on $|z|\le4T$. Positivity gives
$|\mathcal L(z)|\le\mathcal L(\Re z)$, whence \eqref{eq:QL2} implies
\begin{equation}
 \sup_{|z|\le4T}|f(z)|\le2\mathrm{e}^{9T^2}=B,\qquad
 \sup_{u\in[-T,T]}|f(u)|
 \le C_p M^{-1/2}(1+T)^3\mathrm{e}^{T^2}=b.
 \label{eq:QL3}
\end{equation}
For sufficiently large $M$, $b\le B$. In the disk $|z|<4T$ slit along
$[-T,T]$, the harmonic measure of the slit is at least $1/2$ for
$|z|\le T$. To check this, rescale $T$ to one and use
$G(w)=\log|w+\sqrt{w^2-1}|$, with the branch at infinity.
It vanishes on the slit, is at most $\log(1+\sqrt2)$ on $|w|\le1$,
and is at least $\log(4+\sqrt{15})$ on $|w|=4$.
The maximum principle compares harmonic measure to
$1-G/\log(4+\sqrt{15})$, which is at least $1/2$ on the inner disk.
Applying the two-constants principle to $\log|f|$ therefore yields
\begin{equation}
 |f(it)|\le\sqrt{bB}
 \le C_p M^{-1/4}(1+T)^{3/2}\mathrm{e}^{5T^2},\qquad |t|\le T.
 \label{eq:QL4}
\end{equation}
Its integral over this interval is $O_p(M^{-1/8})$.

The Fourier estimates of Section~\ref{uniform-fourier-bounds-and-the-lattice-local-limit}
give constants $c_0,c_1,C>0$ such that
\begin{equation}
 \int_{T<|t|\le\pi\sigma}|\mathbb E_x \mathrm{e}^{itZ}|\,\mathrm{d}t
 \le C\mathrm{e}^{-c_0T^2}+C\mathrm{e}^{-c_0h},\qquad 1\le T\le c_1\sqrt M.
 \label{eq:QL5}
\end{equation}
Away from a fixed critical neighborhood this follows from
\eqref{eq:TL40}: the product of the zero and one probabilities of
$V(xz)/V(x)$, divided by $\mathscr D\Psi(x)$, is bounded below.
Near zero both numerator and denominator are comparable to $x$.
In the critical neighborhood \eqref{eq:TL43} becomes
$(1+c t^2/h)^{-h/8}$. For $|t|\le\sqrt h$ this is bounded by
$\mathrm{e}^{-c't^2}$. Beyond $\sqrt h$, split the exponent in half, bound
one half at $\sqrt h$, and integrate the other against a fixed
integrable power. The integral is $O(\mathrm{e}^{-c'h})$.
The outer arc is exponentially small uniformly in $\varepsilon$ by
\eqref{eq:TL46}. These observations prove \eqref{eq:QL5}.

Fourier inversion, \eqref{eq:QL4}, \eqref{eq:QL5}, and the Gaussian
Fourier tail prove \eqref{eq:QL1}. The three contributions are
$O_p(M^{-1/8})$, $O(M^{-c_0/80})+O(\mathrm{e}^{-c_0h})$, and
$O(\mathrm{e}^{-T^2/2})=O(M^{-1/160})$, so one may take
$\kappa_p=\min(c_0,1/2)/80$. This exponent is not intended to be sharp.
\end{proof}

For the finite set of marking counts $p=1,2,3$, choose a common exponent
$\kappa>0$. At the common saddle \eqref{eq:TL48}, put
$M=\min(2g-1,n-1)$. Repeating the coefficient comparisons in
\eqref{eq:TL49}--\eqref{eq:TL54} now gives
\begin{equation}
 B_{g,n}=\frac{\pi^2}{2r(Y)}(1+O(M^{-\kappa})),\qquad
 C_{g,n}=\frac{\pi^2(n-1)}{gx}(1+O(M^{-\kappa}))
 \label{eq:QL6}
\end{equation}
when $g,n\to\infty$. The center shift in \eqref{eq:TL52} contributes
only $O(g^{-1})$, since its Gaussian argument is $O(g^{-1/2})$.


The next three subsections prove Theorem~\ref{thm:sharp-extraction}. We
first compare the true and the model generating functions on small
complex neighborhoods of the real saddle point, and then treat large
and small coefficient indices separately.

\subsection{Local complex comparison}\label{app:local-complex}

To extract a coefficient by a contour integral we need the marked
expansion on complex neighborhoods of the real saddle, not only on the
real axis. We transfer a finite-order expansion to such neighborhoods and
compare the true and model Fourier integrals.

Choose a fixed finite order $K$, to be enlarged below, and define
\begin{equation}
 F_{h,[K]}(z)=\frac{\mathcal J_{g,p}(z)}\pi
       \left(1+\sum_{j=1}^K\frac{\mathcal L_{p,j}(z)}{h^j}\right).
 \label{eq:CE3}
\end{equation}
This is initially a local analytic function near the positive real
interval. We do not assume that it is analytic on a whole cusp-weight
disk. Theorem~\ref{thm:full-expansion} gives its real relative error
$O_{p,K}(h^{-K-1})$.

The convergent endpoint expansions and compactness give a fixed $c>0$
such that all functions in \eqref{eq:CE3} are analytic at
$z=x\exp(w/\sigma)$ for $|w|\le c\sqrt M$. Near zero this is an
ordinary inverse-function neighborhood. Near criticality the radius
in $z$ is proportional to $x_c-x$, and
$\sigma(x_c-x)/x\asymp\sqrt h$; the square-root inverse and the
analytic $\delta$ coefficients give the required branch. Each fixed
collection of $\mathcal L_{p,j}$ is bounded on a smaller such
neighborhood. The model derivatives obey
\begin{equation}
 \frac{h\mathscr D^mS}{\sigma^m}=O_m(M^{1-m/2})\quad(m\ge3),
 \qquad
 \frac{\mathscr D^m\log B_p}{\sigma^m}=O_{p,m}(M^{-m/2})\quad(m\ge1).
 \label{eq:CE4}
\end{equation}
At zero use $\mathscr D^mS=2x+O_m(x^2)$ and
$\mathscr D^m\log B_p=O_{p,m}(x)$. At criticality, with
$\varepsilon=1-x/x_c$, both numerator derivatives before their
$h$ factor are $O(\varepsilon^{-m})$, whereas
$\sigma\asymp\sqrt h/\varepsilon$.
These bounds also hold locally in the complex neighborhoods by the
same convergent expansions.

Put
\[
 \mathcal L_h(w)=\mathrm{e}^{-kw/\sigma}
   \frac{F_h(x\mathrm{e}^{w/\sigma})}{F_h(x)},\qquad
 \mathcal L_{h,[K]}(w)=\mathrm{e}^{-kw/\sigma}
   \frac{F_{h,[K]}(x\mathrm{e}^{w/\sigma})}{F_{h,[K]}(x)}.
\]
Taylor expansion gives the disk bound $C_{p,K}\mathrm{e}^{C|w|^2+C|w|}$
for the centered, normalized true generating function and finite analytic
model. For the true function, positivity bounds its modulus by its value
at $\Re w$, where the real expansion applies. No error term is
differentiated. The explicit model also satisfies
\begin{equation}
 \left|\frac{\mathcal J_{g,p}(x\mathrm{e}^{it/\sigma})}
                    {\mathcal J_{g,p}(x)}\right|
       \le C_p \mathrm{e}^{-c't^2},\qquad |t|\le c\sqrt M,
 \label{eq:CE5}
\end{equation}
for sufficiently large $M$. Indeed the real action has leading term
$-t^2/2$; its first omitted real term is $O(t^4/M)$, and the real
amplitude change is $O_p(t^2/M)$. Shrinking $c$ absorbs the former,
and growing $M$ absorbs the latter.

Take $T^2=A_0\log h$, with $A_0$ large enough that the Fourier
tail \eqref{eq:QL5} of the true distribution, the polynomial-Gaussian tails below, and the
sparse-model Fourier tail below are $O(h^{-3})$. These tail
constants depend on the fixed action $S$ and amplitude $B_p$, not
on the finite truncation order.
The central truncation order will first be fixed below; then choose $C_0$
so that $M\ge C_0\log h$ implies $4T\le c\sqrt M$ for its
analytic neighborhood.

\begin{lemma}[Local comparison]\label{lem:local-comparison}
There are a fixed order $K_c$ and a constant $C_0$ such that, uniformly
for $M\ge C_0\log h$,
\[
 \sigma\mathbb P_x(N=k)=\frac1{2\pi}\int_{-T}^{T}\mathcal L_{h,[K_c]}(it)\,\mathrm dt
 +O(h^{-2}),\qquad T^2=A_0\log h.
\]
\end{lemma}

\begin{proof}
The difference between the two centered, normalized functions has real-slit bound
$C h^{-K-1}\mathrm{e}^{CT^2}$ on $[-T,T]$ and disk bound $C\mathrm{e}^{C'T^2}$.
The harmonic-measure argument of \eqref{eq:QL4} then gives
\begin{equation}
 \sup_{|t|\le T}|\mathcal L_h(it)-\mathcal L_{h,[K]}(it)|
       \le C_{p,K}h^{-(K+1)/2+C''A_0}.
 \label{eq:CE6}
\end{equation}
The exponential constants here are independent of $K$: for each fixed
$K$, the modulus of its finite correction factor lies between $1/2$ and $2$ on the
chosen neighborhoods for all sufficiently large $h$. Only the threshold
and multiplicative constants depend on $K$. Fix the central order $K=K_c$ sufficiently large
that the integral of \eqref{eq:CE6} is $O(h^{-2})$, and then fix $C_0$.
Thus $\sigma\mathbb P_x(N=k)$ equals the centered finite-model integral
on $[-T,T]$, divided by $2\pi$, up to $O(h^{-2})$.
The normalization $F_{h,[K]}(x)/F_h(x)=1+O(h^{-K-1})$ is retained
when recovering coefficients. All orders chosen here are fixed
independently of $g,x,k$.
\end{proof}

\subsection{Cancellation of the Poisson singularity}

For indices with $M\ge h^{1/4}$ we compute the model coefficient by the
saddle-point method. The coefficients of its Gaussian expansion have poles
$\Psi^{-j}$ as the cusp weight tends to zero. These poles are exactly the
terms of the expansion of $\sqrt{2\pi k}\,\mathfrak p_k$, so they cancel after
division by the Poisson factor in~\eqref{eq:CE2}.

\begin{lemma}[Large indices]\label{lem:large-indices}
The estimate \eqref{eq:CE2} holds uniformly for integers $k\ge h^{1/4}$.
\end{lemma}

\begin{proof}
Here $M=\min(h,k)\ge h^{1/4}$. Put $b=\log B_p$. For a formal
variable $u$ define the model polynomials through the exponential
\begin{equation}
 \exp\!\left\{
 \sum_{m\ge3}u^{m-2}\frac{\mathscr D^{m-1}\Psi}{m!v^{m/2}}w^m
 +\sum_{m\ge1}u^m\frac{\mathscr D^mb}{m!v^{m/2}}w^m\right\}.
 \label{eq:CE7}
\end{equation}
Let $E_j(x)$ be the Gaussian Fourier integral of its $u^{2j}$
coefficient: replace $w^{2a}$ by $(-1)^a(2a-1)!!$ and odd powers
by zero. Each coefficient uses only finitely many derivatives.
By Lemma~\ref{lem:local-comparison} it suffices to evaluate the central
model integral. Expansion under the Gaussian bound \eqref{eq:CE5} gives
\begin{equation}
 \frac1{2\pi}\int_{-T}^T \mathrm{e}^{-ikt/\sigma}
       \frac{\mathcal J_{g,p}(x\mathrm{e}^{it/\sigma})}{\mathcal J_{g,p}(x)}\,\mathrm{d}t
 =\frac1{\sqrt{2\pi}}\left(1+\sum_{j=1}^3\frac{E_j(x)}{h^j}
                                  +O_p(M^{-4}+h^{-2})\right).
 \label{eq:CE8}
\end{equation}
To justify the stated remainder, retain weights through seven in
\eqref{eq:CE7}; a term of weight $q$ has parity $q$ in $w$, so odd
weights integrate to zero. The remaining terms are bounded by $M^{-4}$ times
a fixed polynomial in $|t|$ times $\mathrm{e}^{-c't^2}$. Their integral is
$O(M^{-4})$. The cutoff is $O(\sqrt{\log h})$, which is smaller than
any fixed positive power of $M$ in this range; the Taylor remainder
and exponential remainder bounds therefore apply uniformly. The
removed polynomial-Gaussian tails are $O(h^{-2})$ by the initial
choice of $A_0$. Replacing the finite model by $\mathcal J$ in its central
integral costs $O_p(h^{-1})$, by boundedness of its corrections and
\eqref{eq:CE5}.

The coefficients $E_j$ are bounded at criticality by
\eqref{eq:CE4}. Near zero, $v=\Psi+O(\Psi^2)$ and
$\mathscr D^m\Psi=\Psi+O_m(\Psi^2)$. In a monomial of weight $2j$
with $a$ action factors and $b_*$ amplitude factors, the total degree
in $w$ is $2j+2a$. Its leading power in $\Psi$ is therefore
$\Psi^{b_*-j}$. Every amplitude factor, or higher term in an action
coefficient, improves the pole by at least one. Consequently
\begin{equation}
 E_j(x)=s_j\Psi(x)^{-j}+O_{p,j}(\Psi(x)^{1-j})
       \quad(x\downarrow0),
 \label{eq:CE9}
\end{equation}
where $s_j$ is the coefficient of $k^{-j}$ in
$\sqrt{2\pi k}\,\mathfrak p_k$. Indeed the leading action terms in
\eqref{eq:CE7} are exactly those of a centered Poisson law.
Parity makes the remaining Laurent powers integral.
For example, direct Gaussian integration gives
\[
 E_1=-\frac{\mathscr D^2b+(\mathscr Db)^2}{2v}
      +\frac{(\mathscr Db)\mathscr D^2\Psi}{2v^2}
      +\frac{\mathscr D^3\Psi}{8v^2}
      -\frac{5(\mathscr D^2\Psi)^2}{24v^3}.
\]
The amplitude terms are needed because $k=h\Psi$ is a model center,
not an assertion about the exact finite-genus mean.

Divide \eqref{eq:CE8} by the central Poisson expansion
$\sqrt{2\pi k}\,\mathfrak p_k=1+\sum_{j=1}^3s_j/(h\Psi)^j+O(k^{-4})$.
The principal pole cancels at each order. The quotient coefficients
$R_j$ satisfy
\begin{equation}
 |R_j(x)|\le C_{p,j}(1+\Psi(x)^{1-j}),\qquad j=1,2,3.
 \label{eq:CE10}
\end{equation}
This follows inductively: a product of a Poisson pole of order $i$
and a previous quotient pole of order at most $j-i-1$ has order
at most $j-1$. The quotient is bounded at the critical end as well.
Since $k=h\Psi\ge1$, every $h^{-j}R_j$ is $O_p(h^{-1})$.
Also $M^{-4}\le h^{-1}$. Combining this with \eqref{eq:CE6}
and the bounded finite amplitude corrections proves
\eqref{eq:CE2} in this range.
\end{proof}

\subsection{Small coefficient indices}\label{app:small-indices}

For $k\le h^{1/4}$ the saddle is too close to the origin for the Gaussian
method. We expand the model coefficient directly around the Poisson
coefficient of $\mathrm e^{\mathfrak s_1hz}$.

\begin{lemma}[Small indices]\label{lem:small-indices}
The estimate \eqref{eq:CE2} holds uniformly for integers $0\le k\le h^{1/4}$.
\end{lemma}

\begin{proof}
Near zero write $S(z)=\mathfrak s_1z+\mathfrak s_2z^2+O(z^3)$, $\mathfrak s_1=2$, and
$B_p(z)=1+d_1z+O(z^2)$. The full model \eqref{eq:CE3} is analytic
on a fixed neighborhood of zero. For $k\le h^{1/4}$, coefficient
majorants give
\begin{equation}
 \frac{k![z^k]F_{h,[K]}(z)}
      {(\mathcal J_{g,p}(0)/\pi)(\mathfrak s_1h)^k}
 =1+\frac{(\mathfrak s_2/\mathfrak s_1^2)k(k-1)+(d_1/\mathfrak s_1)k+\mathcal L_{p,1}(0)}h
       +O_{p,K}\!\left(\frac{(k+1)^4}{h^2}\right).
 \label{eq:CE11}
\end{equation}
Factor $\mathrm{e}^{\mathfrak s_1hz}$ from the model. Its retained multiplier is
$1+h\mathfrak s_2z^2+d_1z+\mathcal L_{p,1}(0)/h$.
Analytic absolute-coefficient majorants bound the omitted terms by
convergent products starting with
$hz^3,h^2z^4,z^2,z/h,h^{-2}$. In convolution with the exponential,
a degree-$a$ term contributes the factor $(k)_a/(\mathfrak s_1h)^a$, bounded
by $(k/(\mathfrak s_1h))^a$. Evaluation of the majorants at that last argument
gives the remainder in \eqref{eq:CE11}; their exponent is
$O(k^2/h)=o(1)$. This avoids a spurious $\sqrt k$ loss from a
supremum Cauchy estimate.

At the saddle, direct expansion gives the same polynomial without
$\mathcal L_{p,1}(0)$ for the right-hand side of \eqref{eq:CE2},
relative to this Poisson baseline. Indeed,
$hS-k+k\log(k/(\mathfrak s_1hx))$ contributes $\mathfrak s_2k^2/(\mathfrak s_1^2h)$,
$\log B_p$ contributes $d_1k/(\mathfrak s_1h)$, and
$\tfrac12\log(k/(hv))$ contributes $-\mathfrak s_2k/(\mathfrak s_1^2h)$.
Their error is $O_p((k+1)^4/h^2)$. Thus model coefficient divided
by the proposed saddle expression is $1+O_p(h^{-1})$ for
$k\le h^{1/4}$.

For $C_0\log h\le k\le h^{1/4}$, comparison
\eqref{eq:CE6} applies. Here the entire circle $|z|=x$ is inside
the fixed zero-neighborhood. The convergent series for $S$ and
$|1-\cos(m\theta)|\le m^2(1-\cos\theta)$ give
\[
 \Re(S(x\mathrm{e}^{i\theta})-S(x))\le-cx(1-\cos\theta),
 \qquad |\theta|\le\pi.
\]
The model amplitude is bounded on this circle, and $hx\asymp k$,
so its standardized Fourier tail beyond $T$ is $O(h^{-2})$ by the
initial choice of $A_0$. Complete the central model integral to its full
coefficient integral and apply \eqref{eq:CE11}. This proves
\eqref{eq:CE2} in the middle sparse range. The scaled model central mass is bounded below by a
positive constant here, by \eqref{eq:CE11} and the saddle comparison;
therefore the additive central-integral errors remain relative errors.

For $0\le k\le C_0\log h$ use a different contour and, if needed,
a different fixed truncation order $K=K_\ell$, chosen after $C_0$. Put
$\rho=(k+1)/(2h)$ and consider
$(F_h(\rho z)-F_{h,[K]}(\rho z))/(\mathcal J_{g,p}(0)/\pi)$
in the disk $|z-1|<16$ slit along $[0,2]$.
Both functions are analytic there for all large $h$; only a fixed
neighborhood of zero is used. Positivity, the real expansion, and
the explicit model give disk bound $C\mathrm{e}^{C'(k+1)}$ and slit bound
$C_{p,K}h^{-K-1}\mathrm{e}^{C'(k+1)}$.
On $|z|=1$ the slit harmonic measure is greater than $1/2$:
translate $w=z-1$, so $|w|\le2$, and compare the Green function
in the proof of \eqref{eq:QL4} on radii two and sixteen.
Its lower bound is
\[
 1-\frac{\log(2+\sqrt5)}{\log(16+\sqrt{255})}>\frac12.
\]
Thus the unit-circle error is at most
$C_{p,K}h^{-(K+1)/2}\mathrm{e}^{C''(k+1)}=O(h^{-2})$ upon fixing $K$
sufficiently large. Cauchy extraction changes the coefficient,
relative to the baseline of \eqref{eq:CE11}, by at most
\[
 Ch^{-2}\frac{k!}{(2h\rho)^k}
       =Ch^{-2}\frac{k!}{(k+1)^k}\le Ch^{-2}.
\]
Equation~\eqref{eq:CE11} and its saddle comparison finish this
range, including $k=0$. The two finite orders $K_c,K_\ell$ are fixed before $h$ tends to
infinity, and each is used only in its own range. Enlarging the
low-mean order does not change the already fixed central analytic
neighborhood or $C_0$.
\end{proof}

\begin{proof}[Proof of Theorem~\ref{thm:sharp-extraction}]
Lemmas~\ref{lem:large-indices} and~\ref{lem:small-indices} cover all
$k\ge0$; taking common constants proves the uniform assertion.
\end{proof}

\subsection{The analytic phase change at zero}\label{app:phase-change}

To obtain an expansion that is uniform down to $k=0$, we change variables
so that the phase becomes exactly that of a Poisson coefficient, following
Temme~\cite{Temme}.

\begin{lemma}[Analytic phase change]\label{lem:phase-change}
Let $S$ be analytic near zero with $S(0)=0$ and $S'(0)=s>0$, and let
$x=x(\nu)$ be the solution of $xS'(x)=\nu$ near zero, so that
$x=\nu/s+O(\nu^2)$. There is a jointly analytic map $z=Z(u,\nu)$ on a fixed
polydisk, with $Z(0,\nu)=0$, such that
\begin{equation}
 S(z)-\nu\log z=u-\nu\log u+\mathcal H(\nu),\qquad
 \mathcal H(\nu)=S(x)-\nu+\nu\log(\nu/x),\qquad Z(\nu,\nu)=x.
 \label{eq:CG5}
\end{equation}
If $H$ is analytic near zero and
$f(u,\nu)=H(Z(u,\nu))\,uZ_u(u,\nu)/Z(u,\nu)$, then $f$ is jointly analytic,
including at $u=0$, and for all sufficiently small $\nu=k/h$,
\begin{equation}
 [z^k]\mathrm{e}^{hS(z)}H(z)
       =\mathrm{e}^{h\mathcal H(\nu)}[u^k]\mathrm{e}^{hu}f(u,\nu).
 \label{eq:CG7}
\end{equation}
\end{lemma}

\begin{proof}
The logarithmic equality in \eqref{eq:CG5} uses the nonvanishing
coordinate quotient \(Z/u\). We construct $Z$ directly. Put $c=\nu/x$, $d=z-x$ and
$\phi=S(z)-cz+\nu\log c-\mathcal H$.
Then $\phi(x,\nu)=\phi_z(x,\nu)=0$, and its double-zero quotient is
$\int_0^1(1-t)S''(x+td)\,\mathrm{d}t$, jointly analytic.
Seek $u=z(c+d\eta)$. The phase equation becomes
\begin{equation}
 \frac{\phi}{d^2}-\eta+
 \nu\sum_{m\ge2}\frac{(-1)^{m-1}}m c^{-m}d^{m-2}\eta^m=0.
 \label{eq:CG6}
\end{equation}
At $z=\nu=0$ it is $S''(0)/2-\eta=0$, with derivative $-1$ in
$\eta$. The analytic implicit-function theorem gives $\eta(z,\nu)$.
Since $u_z(0,0)=s$, the inverse $Z$ is jointly analytic on a fixed
polydisk. The construction gives $u(x,\nu)=\nu$ and $u(z,0)=S(z)$.
Only local analyticity of $S$ is required.

The function $f$ is jointly analytic because $Z/u$ is analytic and
nonvanishing. For \eqref{eq:CG7}, write the coefficient as a residue and
change variables $z=Z(u,\nu)$; the circle $|u|=\nu$ lies in the fixed
analytic neighborhood for all sufficiently small $\nu$, and its image
surrounds zero once.
\end{proof}

At the saddle,
twice differentiating \eqref{eq:CG5} gives
$f(\nu,\nu)=H(x)\sqrt{\nu/v}$, with the positive branch.
For $H=B_p$ the normalization on the right is exactly
$\mathrm{e}^{h\mathcal H}h^k B_p(x)\sqrt{\nu/v}/k!$,
which equals the $x$-dependent part of \eqref{eq:CG1}.

\subsection{Poisson central polynomials with a uniform remainder}

In the new variable the coefficient is a Poisson coefficient of an
analytic amplitude. Expanding the amplitude about $u=k/h$ expresses it
through the central polynomials $T_m(k)$ below; we bound the remainder
uniformly in $k$.

Define the normalized coefficient operator by
\[
 \mathscr P_{h,k}f=k!h^{-k}[u^k]\mathrm{e}^{hu}f(u,\nu),\qquad \nu=k/h.
\]
Its central polynomial moments are
\begin{equation}
 \mathscr P_{h,k}(u-\nu)^m=h^{-m}T_m(k),\qquad
 \sum_{m\ge0}T_m(k)\frac{t^m}{m!}=\mathrm{e}^{-kt}(1+t)^k.
 \label{eq:CG8}
\end{equation}
Thus $T_0=1$, $T_1=0$; for $m\ge2$, $T_m$ has degree at most
$\lfloor m/2\rfloor$ and zero constant term. Write
$T_m(k)=\sum_{j=\lceil m/2\rceil}^{m-1}t_{m,j}k^{m-j}$.

\begin{lemma}[Poisson central expansion]\label{lem:poisson-central}
Let $f(u,\nu)$ be jointly analytic and bounded by one on a fixed polydisk.
There are $\nu_0,C>0$ such that, uniformly for $0\le\nu=k/h\le\nu_0$ and
$0\le K\le h/2$,
\begin{equation}
 \begin{split}
 \mathscr P_{h,k}f&=\sum_{j=0}^Kh^{-j}\mathscr T_j[f](\nu)
       +O\!\left(C^{K+1}(K+1)!h^{-K-1}\right),\\
 \mathscr T_0[f]&=f(\nu,\nu),\qquad
 \mathscr T_j[f]=\sum_{m=j+1}^{2j}
 t_{m,j}\nu^{m-j}\frac{\partial_u^mf(\nu,\nu)}{m!}\quad(j\ge1).
 \end{split}
 \label{eq:CG10}
\end{equation}
Moreover $\|\mathscr T_j[f]\|\le C^{j+1}j!$ on a smaller complex disk in
$\nu$, each $\mathscr T_j[f]$ with $j\ge1$ vanishes at $\nu=0$, and
$|\mathscr P_{h,k}f|\le C$.
\end{lemma}

\begin{proof}
Choose $\nu_0>0$ small and a fixed Cauchy radius $r$ about every
$u=\nu$, with $2\nu_0<r/2$. The Taylor remainder through degree
$2K+1$ on $|u|=\nu$ is at most
$2r^{-2K-2}|u-\nu|^{2K+2}$.
For $k\ge1$, normalized coefficient extraction therefore bounds it by
\begin{equation}
 \begin{split}
 &Cr^{-2K-2}\mathfrak p_k^{-1}\nu^{2K+2}
   \int_{-\pi}^{\pi}|\theta|^{2K+2}\mathrm{e}^{-ck\theta^2}\,\mathrm{d}\theta\\
 &\hspace{20pt}\le C^{K+1}(K+1)!\nu^{K+1}h^{-K-1}
       \le C_1^{K+1}(K+1)!h^{-K-1}.
 \end{split}
 \label{eq:CG9}
\end{equation}
Here $\mathfrak p_k^{-1}\le C\sqrt{k+1}$ and
$1-\cos\theta\ge c\theta^2$ on $[-\pi,\pi]$.
The Gaussian integral is an upper bound at every order, even when
$K$ exceeds $k$. At $k=0$ the operator is evaluation at zero and
this remainder vanishes.

The retained moment terms above inverse order $K$ must also be bounded.
The generating function in \eqref{eq:CG8} gives
\[
 |[k^a]T_m(k)|\le\frac{m!}{a!}
 [t^m]\left(\sum_{r\ge2}\frac{t^r}{r}\right)^a
 \le2^m\frac{m!}{a!},\qquad 1\le a\le m/2.
\]
For $\ell=m-a=K+1+d>K$, the retained range has
$0\le d\le K-1$, $m\le2K+1$, and
\[
 \frac{m!}{(m-\ell)!}\le\frac{(2K+1)!}{(K-d)!}
       \le C^K(K+1)!K^d.
\]
The sum of $(K/h)^d$ is bounded when $K\le h/2$; the polynomial
number of terms is absorbed into $C^K$. At $K=0$ this sum is empty.
This proves \eqref{eq:CG10}. At order $j$, $m!/(m-j)!\le(2j)!/j!\le4^jj!$.
Cauchy's bound thus gives
$\|\mathscr T_j[f]\|\le C^{j+1}j!$ on one smaller complex
$\nu$ disk. Each positive-order coefficient vanishes at zero.
The same contour estimate with no vanishing factor gives
$|\mathscr P_{h,k}f|\le C\|f\|$.
\end{proof}


Apply Lemmas~\ref{lem:phase-change} and~\ref{lem:poisson-central} to
$f_s(u,\nu)=B_p(Z(u,\nu))\mathcal L_{p,s}(Z(u,\nu))uZ_u/Z$.
Put $f_0^*(\nu)=B_p(x)\sqrt{\nu/v}$, which is bounded away from
zero on a fixed small complex disk, with $f_0^*(0)=1$.
The extracted coefficients there are
\begin{equation}
 \mathcal A_{p,j}(Y)=\frac1{f_0^*(\nu)}
       \sum_{s=0}^j\mathscr T_{j-s}[f_s](\nu),\qquad \nu=\Psi(X(Y)).
 \label{eq:CG11}
\end{equation}
This proves their zero-end analyticity and endpoint values.
The common norms of Theorem~\ref{thm:full-expansion} give
$\|f_s\|\le C_p\mathfrak B_p^s(s!)^2$ on one polydisk. Applying
\eqref{eq:CG10} at order $K-s$ uses
$(s!)^2(K-s+1)!\le((K+1)!)^2$.
For a finite marked model of order $J=\lfloor\kappa\sqrt h\rfloor>K$,
the bounded operator controls all unused amplitudes by their first
geometric tail term, since $\mathfrak B_p(s+1)^2/h<1/2$ through $J$.
Thus both model errors and coefficient norms have the required Gevrey
bound. The functions \eqref{eq:CG11} agree with \eqref{eq:CG3} where
both are defined; this is proved at the end of the proof of
Theorem~\ref{thm:canonical-expansion}, from the two expansions on a
common range of densities.

\Needspace{7\baselineskip}
\subsection{Proof of the uniform coefficient expansion}\label{app:canonical-proof}

\begin{proof}[Proof of Theorem~\ref{thm:canonical-expansion}]
We first replace the true generating function by one finite marked
expansion, with exponentially small error after coefficient extraction.
For small cusp density the Poisson calculation above applies. For the
remaining densities we use uniform Gaussian coordinates and divide by
the exact Poisson normalization. The last step absorbs the exponential
comparison error into the stated remainder at every allowed order.

Write $F_h=\Phi_g^{(p)}$ and $J_h=\mathcal J_{g,p}/\pi$.
Choose a sufficiently small fixed $\kappa>0$, set
$J=\lfloor\kappa\sqrt h\rfloor$, and use the single finite model
\[
 M_h(z)=J_h(z)P_J(z),\qquad
 P_J=\sum_{s=0}^J\mathcal L_{p,s}h^{-s}.
\]
Theorem~\ref{thm:full-expansion} gives
$|F_h/M_h-1|\le C\mathrm{e}^{-a_0\sqrt h}$ on the entire real interval.
Its common coefficient norm also gives $\|P_J-1\|\le C/h$ on one
complex parameter neighborhood, since the consecutive majorant ratio
is $\mathfrak B_p(s+1)^2/h<1/2$. All choices below use this same model.
The constants in the two coefficient ranges can be chosen in the order
$d$, then $\eta$, then $\tau$; the small-disk argument below specifies
the first two choices.

We first compare the true coefficients with those of this model. Put
$m_* =\min(h,k)$ and, for $k\ge1$,
\[
 \mathscr F_x(w)=\mathrm{e}^{-kw/\sigma}
       \frac{F_h(x\mathrm{e}^{w/\sigma})}{F_h(x)},\qquad
 \mathscr M_x(w)=\mathrm{e}^{-kw/\sigma}
       \frac{M_h(x\mathrm{e}^{w/\sigma})}{M_h(x)}.
\]
There is a fixed $c>0$ on whose disks $|w|\le c\sqrt{m_*}$
the complex inverse parameter is defined and stays in the common
coefficient neighborhood. At criticality,
$\varepsilon(x\mathrm{e}^{w/\sigma})/\varepsilon(x)=1+O(w/\sqrt h)$;
choosing $c$ small fixes its square-root branch and keeps the tilted
real radius below $x_c$. Near zero use ordinary analytic inversion,
and use compactness between the endpoints.
The centered model action has zero first derivative and bounded second
derivative: its second derivative is
$v(x\mathrm{e}^{w/\sigma})/v(x)$. The convergent endpoint expansions bound
this ratio and the normalized logarithmic amplitude derivative.
Consequently
\[
 |\mathscr M_x(w)|\le C\mathrm{e}^{C|w|^2+C|w|}.
\]
Positivity bounds $|\mathscr F_x(w)|$ by $\mathscr F_x(\Re w)$.
On the real diameter the exact identity
\[
 \mathscr F_x(u)=\mathscr M_x(u)
 \frac{1+e_h(x\mathrm{e}^{u/\sigma})}{1+e_h(x)},\qquad
 e_h=F_h/M_h-1,
\]
therefore gives the same disk bound and an additional factor
$\mathrm{e}^{-a_0\sqrt h}$ for their difference on that diameter.

For a fixed small $\eta>0$, take $k\ge\eta\sqrt h$ and
$T=\tau h^{1/4}$. Choose $\tau>0$ after $\eta,a_0$, so that
$4T\le c\sqrt{m_*}$, $T\le c_1\sqrt{m_*}$ in \eqref{eq:QL5},
and $T<\pi\sigma$. The scale-invariant harmonic measure used in
\eqref{eq:QL4} gives
\[
 \sup_{|t|\le T}|\mathscr F_x(it)-\mathscr M_x(it)|
 \le C\exp\{-a_0\sqrt h/2+C T^2+C T\}
 \le C\mathrm{e}^{-a_1\sqrt h}.
\]
Its integral has the same bound after decreasing $a_1$.
The integrated tail \eqref{eq:QL5} of the true distribution is also exponentially small
in $\sqrt h$; the outer-arc estimate introduces no factor $\sigma$.
The model central integral is bounded by \eqref{eq:CE5} and the
bounded $P_J$ ratio. Restoring the real normalization thus shows
\begin{equation}
 \left|[z^k]F_h-
 \frac{x^{-k}}{2\pi\sigma}\int_{-T}^T
             \mathrm{e}^{-ikt/\sigma}M_h(x\mathrm{e}^{it/\sigma})\,\mathrm{d}t\right|
 \le C\mathcal Q_{g,p,k}\mathrm{e}^{-a_2\sqrt h}.
 \label{eq:CG14}
\end{equation}
Indeed the error before division is at most
$CF_h(x)x^{-k}\sigma^{-1}\mathrm{e}^{-a_2\sqrt h}$, and
$\mathfrak p_k\sqrt k$ is bounded above and below for integer $k\ge1$.
The real matching factor $F_h(x)/M_h(x)=1+O(\mathrm e^{-a_2\sqrt h})$
has been absorbed into the error.

For the remaining $k$, put $r=d/\sqrt h$.
On $|z|\le4r$, positivity and the near-zero model give
$|(F_h-M_h)/J_h(0)|\le C\mathrm{e}^{Chr}$; on $[0,r]$ this improves to
$C\mathrm{e}^{-a_0\sqrt h+Chr}$. In the disk $|z-r/2|<2r$ slit along
$[0,r]$, the slit harmonic measure is greater than $1/4$ on
$|z|\le r/4$. To see this, rescale to $|w|<4$ with slit $[-1,1]$.
On the target $|w|\le3/2$ its Green function is less than $\log4$,
and on $|w|=4$ it is greater than $\log7$; comparison gives
$1-\log4/\log7>1/4$.
Choose $d$ small, then $\eta$ small, so the two-constants bound is
$C\mathrm{e}^{-a_3\sqrt h}$ and the saddle for $1\le k<\eta\sqrt h$
lies in $|z|\le r/4$. Cauchy's bound on $|z|=x$ is on the
$\mathcal Q$ scale with a cost at most $C\sqrt{k+1}$:
$v/\Psi$ is bounded and $J_h(x)/J_h(0)\ge c$ near zero.
This polynomial cost is absorbed into a smaller exponential.
At $k=0$ use the marked estimate at zero directly.

Choose a switching density $\nu_*>0$ strictly inside the fixed
Poisson neighborhood. For $\Psi(x)\le\nu_*$, the full model circle
is analytic and
\[
 \Re(S(x\mathrm{e}^{i\theta})-S(x))\le-cx(1-\cos\theta).
\]
This follows from $S(z)=2z+\sum_{m\ge2}\mathfrak s_mz^m$ and
$\sum_{m\ge2}m^2|\mathfrak s_m|x^m=O(x^2)$, using
$1-\cos(m\theta)\le m^2(1-\cos\theta)$.
The amplitudes are bounded, and $hv\asymp k$, so the standardized
model tail is Gaussian on the whole circle. For $k\ge\eta\sqrt h$
it completes the integral in \eqref{eq:CG14} to $[z^k]M_h$ at
exponential precision. Together with the small disk, this compares the
true coefficients with those of the full model for every
$\Psi(x)\le\nu_*$.
The preceding Poisson argument now gives \eqref{eq:CG4}, with one
additional fixed error $C\mathrm{e}^{-a\sqrt h}$, throughout this range.

\textbf{Uniform coordinates near criticality.} For the other range we justify the growing-order Gaussian expansion
uniformly through the critical point. In $\delta=y_0-Y$, write the
convergent local factors as
\[
 x_c-X(Y)=\delta^2\alpha(\delta),\quad
 \mathcal Z(Y)=\delta^{5/2}\zeta(\delta),\quad
 B_p(X(Y))=\delta^{p/2}\mathfrak b_p(\delta),\quad
 v(X(Y))=\delta^{-4}v_0(\delta),
\]
where all four factors are analytic and nonzero at zero.
For $z=X(Y)\mathrm{e}^{iu/\sqrt v}$ the inverse has
$\delta_z=\delta q(\delta,u)$, jointly analytic and close to one:
divide its equation $x_c-X(y_0-\delta q)=x_c-z$ by $\delta^2$
and apply the implicit-function theorem at $q=1,u=\delta=0$.
Its $q$ derivative is $2\alpha(0)\ne0$.
Thus $S(z)-S(x)$ and $B_p(z)/B_p(x)$ are analytic joint families,
expressed through $\log q$, $\zeta(\delta q)/\zeta(\delta)$ and
$q^{p/2}\mathfrak b_p(\delta q)/\mathfrak b_p(\delta)$.
Also $\Psi/\sqrt v$ extends analytically. Compactness gives common
neighborhoods for $\Psi\ge\nu_*/2$.

Set
\[
 \psi_Y(u)=S(x\mathrm{e}^{iu/\sqrt v})-S(x)-i\Psi u/\sqrt v,
 \qquad W_Y(u)=u\sqrt{-2\psi_Y(u)/u^2}.
\]
The removable quotient is one at zero; $W'_Y(0)=1$.
On one fixed disk $W_Y$ has a jointly analytic inverse $U_Y$,
with $\psi_Y(U_Y(w))=-w^2/2$.
Put $u_0=T/\sqrt h=\tau h^{-1/4}$.
The image endpoints are $\pm u_0+O(u_0^2)$. Straight connectors
to $\pm u_0$ have length $O(u_0^2)$ and
$\Re(w^2)\ge u_0^2/2$. A finite contour deformation therefore
costs at most $C\sqrt h\,u_0^2\mathrm{e}^{-hu_0^2/4}$.
After factoring $J_h(x)x^{-k}/\sigma$, the central integral becomes
\[
 \frac{\sqrt h}{2\pi}\int_{-u_0}^{u_0}\mathrm{e}^{-hw^2/2}
       \sum_{s=0}^J\widetilde f_s(w,Y)h^{-s}\,\mathrm{d}w
       +O(\mathrm{e}^{-a\sqrt h}),
\]
where
$\widetilde f_s=[B_p(z(U_Y(w)))/B_p(x)]
 \mathcal L_{p,s}(z(U_Y(w)))U'_Y(w)$ has common norm
$C_p\mathfrak B_p^s(s!)^2$ and $\widetilde f_0(0,Y)=1$.

\textbf{Gaussian expansion.} For $s\le K$, Taylor expand through degree $2(K-s)+1$ on that
fixed disk. The absolute Gaussian moment bounds the remainder by
$C^{K-s+1}(K-s+1)!h^{-K+s-1}\|\widetilde f_s\|$.
Odd terms vanish. To extend the polynomial moments to the whole line,
use, for $|w|\ge u_0$,
\[
 \mathrm{e}^{-hw^2/2}\le \mathrm{e}^{-hu_0^2/4}\mathrm{e}^{-hw^2/4}.
\]
The remaining moment majorants $C^j j!h^{-j}$ have bounded sum
for $j\le K\le J=O(\sqrt h)$. All polynomial tails are therefore
$O(\mathrm{e}^{-a\sqrt h})$, even at growing Taylor degree.
The unused $s>K$ amplitudes again form a geometric tail.
The factorial convolution used above yields a numerator remainder
$CB^{K+1}((K+1)!)^2h^{-K-1}+C\mathrm{e}^{-a\sqrt h}$.
After factoring $1/\sqrt{2\pi}$, its coefficients have the same
common-domain bounds and agree with \eqref{eq:CG2}. Indeed, for fixed
$Y$ the model integral is defined for every real $h$, and both
families are the coefficients of its asymptotic expansion as
$h\to\infty$ through real values: \eqref{eq:CG2} before and the present
ones after the analytic change of variable $u=U_Y(w)$. Such expansion
coefficients are unique.

\textbf{The exact Poisson normalization.} The divisor in \eqref{eq:CG1} is exact. Put
$R(k)=\sqrt{2\pi k}\,\mathfrak p_k$. Its logarithm satisfies
\[
 \log R(k)=-\int_0^\infty \mathrm{e}^{-kt}
 \left(\frac1{\mathrm{e}^t-1}-\frac1t+\frac12\right)\frac{\mathrm{d}t}{t}.
\]
The kernel has a removable origin, analytic radius $2\pi$, and is
bounded on the positive axis. Its Taylor remainder through degree
$L-2$ is bounded globally there by $CD^Lt^{L-1}$: use Cauchy near
zero and bound each retained power outside a fixed small interval.
Integration gives $CD^L(L-1)!k^{-L}$.
For exponentiation and inversion, the degree-$K$ log polynomial is
bounded on $|t|\le1/(D_0(K+1))$. Cauchy's tail at $t=1/k$, for
$K\le c k$, is at most
$C[D_0(K+1)/k]^{K+1}\le C(\mathrm{e}D_0)^{K+1}(K+1)!k^{-K-1}$.
Since the exponential is bounded on the relevant real interval, the
remainder for $\log R$ transfers to $R$ and to $R^{-1}$, which
therefore have Gevrey-1 remainders. Here $k\ge(\nu_*/2)h$ and $\Psi^{-1}$ is
bounded analytic through criticality. By \eqref{eq:CG1},
$\mathcal Q_{g,p,k}/(\mathcal J_{g,p}(x)x^{-k}/(\pi\sigma))
=R(k)/\sqrt{2\pi}$. The two factors $1/\sqrt{2\pi}$ cancel.
Dividing the normalized numerator by $R(k)$ consequently preserves
the Gevrey-2 error and gives exactly
\eqref{eq:CG3}.

\textbf{Identification on the overlap.} For $\nu_*/2\le\Psi\le\nu_*$
both expansions hold; write $\mathcal A^{\mathrm P}_{p,j}$ for the
Poisson coefficients \eqref{eq:CG11} and $\mathcal A^{\mathrm G}_{p,j}$
for the Gaussian coefficients \eqref{eq:CG3}. Fixed-$Y$ uniqueness is
not available directly, since $k=h\Psi(X(Y))$ must be an integer, so we
argue by density. Suppose the two families agree below order $j$.
Subtracting the expansions at order $K=j$ gives
$\mathcal A^{\mathrm P}_{p,j}(Y_{h,k})-\mathcal A^{\mathrm G}_{p,j}(Y_{h,k})=O(h^{-1})$,
uniformly over all admissible $(h,k)$ with $Y_{h,k}$ in the overlap,
where $h\Psi(X(Y_{h,k}))=k$. Given $Y$ in the open overlap, take
$k=\lfloor h\Psi(X(Y))\rfloor$; since $\Psi\circ X$ is continuous and
strictly increasing, $Y_{h,k}\to Y$ as $h\to\infty$. Both coefficients
are continuous, so they agree at $Y$. By induction the two families
coincide on the open overlap. Both are analytic on
$0<\Psi(X(Y))<\nu_*$, so the identity theorem extends the equality to
this whole interval. Hence \eqref{eq:CG11} is the analytic continuation
of \eqref{eq:CG3} to $Y=0$, and the coefficients are analytic on one
complex neighborhood of $[0,y_0]$.

\textbf{Absorbing the residual error.} It remains to absorb the fixed exponential error at every allowed
order. After fixing the model, neighborhoods and cutoff, choose
$c_p<\kappa/2$ so small that $2c_p<1$ and
$4c_p\log(\mathrm{e}/(2c_p))<a$. For $L=K+1\le2c_p\sqrt h$ and $\mathfrak B_p\ge1$,
\[
 \mathfrak B_p^L(L!)^2h^{-L}
 \ge \mathrm{e}^{-2L\log(\mathrm{e}\sqrt h/L)}
 \ge \mathrm{e}^{-4c_p\log(\mathrm{e}/(2c_p))\sqrt h}.
\]
This bounds the fixed residual by the stated remainder simultaneously
for every $K\le c_p\sqrt h$, including zero.
Finally choose $c'_p<c_p$ after the final $\mathfrak B_p$, with
$\mathfrak B_p(c'_p)^2<1$. The upper bound
$\mathfrak B_p^L(L!)^2h^{-L}\le(\mathfrak B_pL^2/h)^L$ proves the exponential
approximation at $K=\lfloor c'_p\sqrt h\rfloor$.
\end{proof}

The first coefficient after extraction is
\begin{equation}
 \begin{split}
 \mathcal A_{p,1}={}&\mathcal L_{p,1}
 -\frac{\mathscr D^2b+(\mathscr Db)^2}{2v}
 +\frac{(\mathscr Db)\mathscr D^2\Psi}{2v^2}
 +\frac{\mathscr D^3\Psi}{8v^2}
 -\frac{5(\mathscr D^2\Psi)^2}{24v^3}+\frac1{12\Psi},
 \qquad b=\log B_p.
 \end{split}
 \label{eq:CG12}
\end{equation}
Its apparent pole at zero is removable by \eqref{eq:CG11}.
In particular, retaining this term improves the relative remainder
to $O_p(h^{-2})$ over every coefficient index.

At the critical endpoint its value is
\begin{equation}
 \mathcal A_{p,1}(y_0)=-\frac{p^2}{8}+\frac p2-\frac{37}{60}.
 \label{eq:CG13}
\end{equation}
Indeed, writing $\varepsilon=1-x/x_c$, the leading terms are
$S=-(5/4)\log\varepsilon+O(1)$ and
$\log B_p=(p/4)\log\varepsilon+O(1)$.
Their differentiated convergent expansions give
$E_1(y_0)=-(p^2+4p)/40-1/15$ in \eqref{eq:CG12}.
Combine this with $\mathcal L_{p,1}(y_0)=t_p$ from
\eqref{eq:ST13} and $\Psi^{-1}\to0$.


\section{The rate for adjacent coefficient ratios}
\label{sec:quantitative-coefficients}

This appendix proves Theorem~\ref{thm:quantitative-coefficient}, which
contains the rate in Theorem~\ref{thm:coefficient}. We quantify the three
degree ranges treated in Section~\ref{uniform-one-boundary-coefficient-ratios}.
The logarithmic comparison and concavity give the rate away from the
lowest degrees, including bounded positive codimension at the top.
Bracket difference estimates supply the remaining range. Balancing the
bounds gives $O(g^{-1/8})$ uniformly; we do not claim that this exponent
is optimal. The coefficients and the exact density are those of
Theorem~\ref{thm:coefficient}.

\begin{theorem}[Quantitative adjacent ratios]\label{thm:quantitative-coefficient}
There is a constant $C$ such that, for all sufficiently large $g$,
\begin{equation}
 \sup_{0\le k<3g-2}
 \left|\frac{a_{g,k}/a_{g,k+1}}{4\pi^2/t_{g,k}^2}-1\right|
 \le Cg^{-1/8}.
 \label{eq:QC1}
\end{equation}
For every fixed $\varepsilon\in(0,3)$ there is $C_\varepsilon$ such that
\begin{equation}
 \sup_{\varepsilon g\le k<3g-2}
 \left|\frac{a_{g,k}/a_{g,k+1}}{4\pi^2/t_{g,k}^2}-1\right|
 \le C_\varepsilon g^{-1/4}.
 \label{eq:QC2}
\end{equation}
In the second supremum $k$ is an integer. Both bounds include bounded
positive codimension from the top coefficient.
\end{theorem}

\subsection{The logarithmic model up to its endpoints}

We first make the logarithmic comparison of
Section~\ref{uniform-one-boundary-coefficient-ratios} quantitative, up to
both ends of the degree range.

Use $H,\widetilde A_m,\mu$ from \eqref{eq:C4} and \eqref{eq:C10}.
Let $\zeta(u)$ be the inverse of $\mu(z)$ on $0\le u\le3/2$,
with $\zeta(0)=0$. Since $\mu(z)=z/5+O(z^2)$ at zero and
$\mu'(z)>0$ throughout $[0,\pi^2/2]$, the function
\[
 \chi(u)=u/\zeta(u),\qquad \chi(0)=1/5,
\]
is positive and analytic in a neighborhood of the closed real interval.
Write $f$ for the function in \eqref{eq:C14}, extended continuously to
$[0,3]$, and set
\begin{equation}
 \mathfrak q(\alpha)=4\alpha\chi((3-\alpha)/2),\qquad
 f'(\alpha)=-\log\mathfrak q(\alpha),\qquad
 |f''(\alpha)|\le C(1+\alpha^{-1})\quad(0<\alpha\le3).
 \label{eq:QC3}
\end{equation}
These follow from \eqref{eq:C15}; in particular
$\mathfrak q(3)=12/5$ and $f(3)=0$.

We need the endpoint-uniform version of the logarithmic saddle estimate
\eqref{eq:C12}.

\begin{lemma}[Endpoint-uniform saddle estimate]\label{lem:uniform-saddle}
For integers $0\le m\le3h/2$, let $z=\zeta(m/h)$, interpreting the
$m=0$ expressions by continuity. Then, uniformly in $m$,
\begin{equation}
 \log \widetilde A_m(h)=-h\log H(z)-m\log z+O(\log(h+1)).
 \label{eq:QC4}
\end{equation}
\end{lemma}

\begin{proof}
Tilt $H^{-1}$ at $z$ and let
$S_h$ be the sum of $h$ independent variables with probability
generating function $H(z)/H(zw)$. Exactly,
\[
 \widetilde A_m(h)=H(z)^{-h}z^{-m}\Pr(S_h=m).
\]
There is $c>0$ such that $\Pr(S_h=m)\ge c/\sqrt{m+1}$ uniformly.
We give the details at the degenerate endpoint $m/h\to0$.
When $m\to\infty$ and $m/h\to0$, one has $z\asymp m/h$,
$\operatorname{Var}(S_h)\asymp m$, and each fixed cumulant of the
sum is $O(m)$. The probabilities at zero and one of one summand have
product comparable to $z$. The Fourier bound \eqref{eq:TL40} therefore
gives $\exp(-cm\theta^2)$ on $[-\pi,\pi]$, and the centered Taylor
expansion gives the local Gaussian limit at $m$, with central mass
asymptotic to $(2\pi\operatorname{Var}(S_h))^{-1/2}$.
When $m/h$ stays away from zero, the same argument uses a compact
tilted family with variance comparable to $h$ and positive zero and
one probabilities. If $m$ stays bounded, pass to a fixed integer $m$;
then $z=5m/h+O_m(h^{-2})$ and
$\widetilde A_m(h)=(h/5)^m/m!\,(1+O_m(h^{-1}))$, so the central mass tends to
$\mathrm{e}^{-m}m^m/m!>0$. At $m=0$ it is one.
These cases and a subsequence contradiction prove the asserted uniform
lower bound. The upper bound is one, proving \eqref{eq:QC4}.
\end{proof}

Put
\[
 G_g(k)=\log\frac{J_k}{J_Ng^{2(N-k)}},\qquad N=3g-2.
\]

\begin{lemma}[Logarithmic comparison up to the endpoints]\label{lem:log-comparison}
There is a constant $C$ such that
\begin{equation}
 \max_{0\le k\le N}|G_g(k)-g f(k/g)|\le C\sqrt g.
 \label{eq:QC5}
\end{equation}
\end{lemma}

\begin{proof}
This follows from the model \eqref{eq:C4}, the comparison \eqref{eq:C5} and
Lemma~\ref{lem:uniform-saddle}. The passage to $f$ is uniform, even at
$m=0$. To see this, use the
exact model for the intersection number,
$3^h\Gamma(h)4^{-g}m!\widetilde A_m(h)/(2k+1)!!$ from Section~\ref{uniform-one-boundary-coefficient-ratios}
and apply Stirling's bound with logarithmic error to its factorials,
using $0\log0=0$. Replacing $h=2g-1$ by $2g$ and
$m=3g-2-k$ by $3g-k$ changes its saddle action by $O(\log(g+1))$:
the two partial derivatives of the action are $-\log H(z)$ and
$-\log z$, the latter having only an integrable logarithmic
singularity at $m=0$. The powers of $g$ cancel as in
\eqref{eq:C13}--\eqref{eq:C14}. Subtracting the top model value gives
$g f(k/g)+O(\log(g+1))$ uniformly. The two logarithmic errors between
the true and the model values are at most $C\sqrt g$ by \eqref{eq:C5}, including the top
coefficient, which proves \eqref{eq:QC5}.
\end{proof}

\subsection{Quantitative secants and the top endpoint}

\begin{proof}[Proof of Theorem~\ref{thm:quantitative-coefficient}]
The sequence $G_g$ is concave by \eqref{eq:C7}. Write
$q_k=J_k/J_{k+1}$. For an integer $\ell\ge1$ with
$\ell\le k/2$ and $k+\ell\le N$, its secants imply
\[
 -\frac{G_g(k)-G_g(k-\ell)}\ell
 \le\log(q_k/g^2)
 \le-\frac{G_g(k+\ell)-G_g(k)}\ell.
\]
Equations \eqref{eq:QC3} and \eqref{eq:QC5} yield
\begin{equation}
 \left|\log\frac{q_k}{g^2\mathfrak q(k/g)}\right|
 \le C\left(\frac{\sqrt g}\ell+\frac\ell k\right).
 \label{eq:QC6}
\end{equation}
If $k+\ell>N$, retain the lower secant and use
$q_k\le q_{N-1}$ for the upper bound. The exact top anchor
\eqref{eq:C18} gives $q_{N-1}/g^2=12/5+O(g^{-1})$.
Since $\log\mathfrak q$ is Lipschitz near three and
$3-k/g\le(\ell+2)/g$, this gives the same bound
\eqref{eq:QC6}, with a changed constant.

Take $\ell=\lfloor g^{5/8}\rfloor$. For $k\ge g^{3/4}$ it is
at most $k/2$ for large $g$, and \eqref{eq:QC6} is $O(g^{-1/8})$.
For $k\ge\varepsilon g$, instead take
$\ell=\lfloor g^{3/4}\rfloor$ to obtain $O_\varepsilon(g^{-1/4})$.

These estimates for the intersection numbers translate into the
statement of the theorem without losing the bounded-codimension range.
Indeed, put
$m=N-k$, $h=2g-1$, and $z_{g,k}=t_{g,k}^2/2$. By \eqref{eq:C11},
$\mu(z_{g,k})=m/h$. Therefore the value of $q_k=J_k/J_{k+1}$ that
corresponds, through \eqref{eq:C3}, to the proposed equivalent
$4\pi^2/t_{g,k}^2$ of $a_{g,k}/a_{g,k+1}$ is
\begin{equation}
 \frac{(2k+3)m}{z_{g,k}}
 =(2k+3)h\chi(m/h)
 =g^2\mathfrak q(k/g)\left(1+O(k^{-1}+g^{-1})\right).
 \label{eq:QC7}
\end{equation}
The last error is uniform for $k\ge1$: the two arguments of $\chi$
differ by $O(g^{-1})$, while $\chi$ is positive with bounded
logarithmic derivative on the closed interval. This calculation
retains $m$ exactly, including $m=1$. By \eqref{eq:C3}, the two
secant estimates prove \eqref{eq:QC1} for $k\ge g^{3/4}$ and all
of \eqref{eq:QC2}.

\subsection{The small-degree range}

It remains to consider $k\le g^{3/4}$. Put
$\rho_k=b_{g,k}/b_{g,k+1}\ge1$ and
$\omega_k=(2k+3)(N-k)$ as before. Equation~\eqref{eq:C25} gives
$\rho_k=1+O(g^{-1/2})$ for $k\le\delta_0\sqrt g$ with a fixed
small $\delta_0>0$. For the remaining indices we claim
\begin{equation}
 0\le\rho_k-1\le C\left(
       \frac{\log(g+1)}{\sqrt k}+\sqrt{\frac kg}+\frac1k\right)
 \quad(\delta_0\sqrt g\le k\le g^{3/4}).
 \label{eq:QC8}
\end{equation}
First $\rho_k<2$ uniformly in this range for large $g$. Otherwise
monotonicity of $q_j$ and $\omega_j/\omega_k\le1+\delta$ propagate a fixed
ratio gap over a block of length comparable to $k$. Its product
would force $\log(V_{g,1}/b_{g,j})\ge c k$ at an index $j\asymp k$,
contradicting \eqref{eq:C23}, whose upper bound is $o(k)$ uniformly
on this range.

Now put $\eta=\rho_k-1\in[0,1]$ and $\delta=\eta/6$.
For $k\le j\le\lfloor(1+\delta)k\rfloor$, the same monotonicity
and the exact factor $\omega_j$ give
$\rho_j\ge(1+\eta)/(1+\delta)\ge1+\eta/2$.
If $\eta k\ge12$, this block has at least $\eta k/12$ indices.
Summing logarithms over the block and using \eqref{eq:C23} at its
last index plus one gives
\[
 c\eta^2 k\le C\log^2(g+1)+Ck^2/g.
\]
All these indices lie below the fixed $c_0g$ range of \eqref{eq:C23}
for large $g$. If $\eta k<12$, then $\eta<12/k$ directly.
This proves \eqref{eq:QC8}. Its right side is $O(g^{-1/8})$ on
the displayed range. Thus $\rho_k=1+O(g^{-1/8})$ uniformly for
$0\le k\le g^{3/4}$.

Finally the saddle equation at $t=\pi$ gives
$\pi^2/t_{g,k}^2=1+O((k+1)/g)$ uniformly on this small-degree
range. Since $a_{g,k}/a_{g,k+1}=4\rho_k$, the desired
relative error is $O(g^{-1/8})$ there as well. Together with the
secant range this proves both assertions.
\end{proof}

The argument is a quantitative version of the one used for the leading
equivalent, based on the same logarithmic comparison and concavity. Liu--Xu's fixed-codimension expansion
has finer errors at each fixed codimension~\cite{LX}; the bounds here
hold with a single constant across all valid degrees.

\section{Proofs for the tight length spectrum}\label{app:tight}

This appendix proves Theorem~\ref{thm:tight-spectrum} and
Corollary~\ref{cor:tight-boltzmann}, following the outline in
Subsection~\ref{subsec:tight-spectrum}. The scale $\alpha$ and the
intensity $\nu$ are those of \eqref{eq:tight-scale}.

\subsection{Tight volume generating functions}\label{subsec:tight-gf}

Put $P=2\pi^2$. Let $T_{g,p,k}(\mathbf L)$ be the tight volume with
$p$ protected geodesic boundaries and $k$ cusps that may be filled,
in the conventions of~\cite{Budd,BZ,BL}. For stable $(g,p)$ define
\[
 \mathscr T_{g,p}(x;\mathbf L)
 =\sum_{k\ge0}\frac{T_{g,p,k}(\mathbf L)}{P^{3g-3+p+k}}
                 \frac{x^k}{k!},\qquad F_{g,p}=\Phi_g^{(p)}.
\]
Unmarked functions are used only for $g\ge2$. Set
$\varrho(Y)=\sqrt{r(Y)/P}$ and $S(t)=\sinh(t)/t$, with $S(0)=1$.

\begin{lemma}[Tight generating functions]\label{lem:tight-gf}
Let $(g,p)$ be stable and $0\le Y<y_0$. The tight generating function is
expressed through the shifted amplitudes by the exact identity
\begin{equation}
 \mathscr T_{g,p}(X(Y);\mathbf L)
 =\sum_{\mathbf d}U^Y_{g,p}(\mathbf d)
       \prod_{i=1}^p\frac{L_i^{2d_i}}{P^{d_i}(2d_i+1)!},
 \label{eq:tight-insertion}
\end{equation}
which fixes both the normalization and the meaning of the protected
boundaries. Its ratio to the ordinary marked generating function satisfies
\begin{equation}
 \begin{split}
 1\le\frac{\mathscr T_{g,p}(X(Y);\mathbf L)}{F_{g,p}(X(Y))}
   &\le\prod_{i=1}^p S(\varrho(Y)L_i),\\
 \frac{\mathscr T_{g,p}(X(Y);\mathbf L)}{F_{g,p}(X(Y))}
   &=\prod_{i=1}^p S(\varrho(Y)L_i)
        \bigl(1+O_{p,B}(g^{-1})\bigr)
       \quad\text{if }\max_i\varrho(Y)L_i\le B,
 \end{split}
 \label{eq:tight-real}
\end{equation}
where the inequality holds at every stable genus. Finally, for each fixed
physical length vector,
\begin{equation}
 \mathscr T_{g,p}(x;\mathbf L)=V(x)^{2g-2+p}R_{g,p,\mathbf L}(x),
 \qquad [x^j]R_{g,p,\mathbf L}(x)\ge0,
 \label{eq:tight-positive}
\end{equation}
and the series converge for $0\le x<x_c$.
\end{lemma}

\begin{proof}
The identity \eqref{eq:tight-insertion} follows from~\cite[Proposition~12]{BZ}. Indeed, their root and
moments are $R=Y/P$ and $M_j=P^jX^{(j+1)}(Y)$. With $q=P/2$,
their shifted times are $q^{k-1}t'_k$, where the $t'_k$ are the
times in \eqref{eq:H8}. The dimension constraint gives the factor
$q^{3g-3+p-|\mathbf d|}$. Combining their external insertion
$2L^{2d}/(4^d d!)$ with the free-energy factor $2^{3g-3}$ and
the double factorial in \eqref{eq:H6} gives
\eqref{eq:tight-insertion}. This also covers genus zero with $p\ge3$.

Equations~\eqref{eq:TL15} and~\eqref{eq:TL23} give \eqref{eq:tight-real}. For
the asymptotic estimate,
the factorial series with the extra weight
$|\mathbf d|(|\mathbf d|+1)$ is uniformly summable on the indicated
box. Thus the lengths may grow as $r(Y)^{-1/2}$.

For each fixed physical length vector, the proof of
\eqref{eq:TL33} applies to every amplitude in
\eqref{eq:tight-insertion}, which gives \eqref{eq:tight-positive}. The
finite sum in \eqref{eq:tight-insertion} and \eqref{eq:tight-real} justify
convergence throughout $0\le x<x_c$.
\end{proof}

\subsection{Prescribing the cusp count}\label{subsec:tight-count}

We record the coefficient argument for the cut surfaces that occur in
the tight integration formula. Fix a number $k$ of cuts and put
$h_0=2g-2$. When $n\to\infty$, choose $x$ by
\[
 (h_0+1)\Psi(x)=n-1,\qquad
 \sigma^2=(h_0+1)\mathscr D\Psi(x),\qquad x=X(Y),
\]
and write $\alpha_x=2\varrho(Y)$. If $n$ stays bounded, we pass to a
subsequence on which $n$ is fixed and take instead $x=c/h_0$ with $c>0$.

\begin{lemma}[Connected cut configurations]\label{lem:tight-connected}
Fix $k$ and a compact box of scaled lengths $\mathbf t\in[0,\infty)^{2k}$.
With $x$ chosen as above, uniformly for $\mathbf t$ in the box,
\begin{equation}
 \frac{[z^n]\mathscr T_{g-k,2k}(z;\mathbf t/\alpha_x)}
      {[z^n]\Phi_g(z)}
 =(4r(Y))^k\prod_{i=1}^{2k}S(t_i/2)(1+o(1)).
 \label{eq:tight-connected}
\end{equation}
Moreover $\alpha_x/\alpha(n/g)\to1$.
\end{lemma}

\begin{proof}
Suppose first that $n\to\infty$. The denominator
$F_{g,1}=\Phi_g'$ has its local limit at $n-1$.
For the connected-complement numerator
$G(x)=\mathscr T_{g-k,2k}(x;\mathbf t/\alpha_x)$, the physical
lengths are held fixed when its argument is varied. Its model center
is $h_0\Psi(x)$, and
\[
 \frac{n-h_0\Psi(x)}\sigma
 =\frac{1+\Psi(x)}\sigma\longrightarrow0.
\]
The endpoint expansions give
\[
 \sup_{0<x<x_c}
 \frac{|\mathscr D\log r(Y(x))|}{\sqrt{\mathscr D\Psi(x)}}<\infty.
\]
Consequently a real tilt $x\mathrm e^{u/\sigma}$ changes $r$ by a
relative $O_u(g^{-1/2})$, uniformly up to either endpoint.
Equation~\eqref{eq:tight-real} therefore transfers the real Gaussian
limit from $F_{g-k,2k}$ to $G$. The positive factor
\eqref{eq:tight-positive} transfers the Fourier estimates
\eqref{eq:TL40}--\eqref{eq:TL46}. The two local probabilities at
$n$ and $n-1$ are asymptotic to $(\sqrt{2\pi}\sigma)^{-1}$.

We use the derivative only in the denominator: adding a protected cusp
to a tight surface would change which homotopies are allowed.
For a generating function $J$ with nonnegative coefficients write
$\mathbb P_x^J(N=m)=[z^m]J(z)x^m/J(x)$. The exact identity is
\begin{equation}
 \frac{[z^n]G(z)}{[z^n]\Phi_g(z)}
 =\frac n x\frac{G(x)}{\Phi_g'(x)}
       \frac{\mathbb P_x^G(N=n)}
            {\mathbb P_x^{\Phi_g'}(N=n-1)}.
 \label{eq:tight-coefficient}
\end{equation}
Equation~\eqref{eq:TL24} at $p=0$ shows that
$n\Phi_g(x)/(x\Phi_g'(x))\to1$. Iterating
\eqref{eq:TL31} now yields \eqref{eq:tight-connected}.

For bounded $n$, with $x=c/h_0$, equation~\eqref{eq:TL24} gives
\eqref{eq:TL30} also at $p=0$. Integrating that estimate between
$x$ and $xz$, and using $\Psi(x)=2x+O(x^2)$, proves convergence of
the probability generating functions to $\exp(2c(z-1))$.
The same holds for $G$ by \eqref{eq:tight-real}. The probabilities at
the fixed index $n$ thus converge to the same positive Poisson
probability. This proves \eqref{eq:tight-connected} in this case too.

Finally $\alpha_x/\alpha(n/g)\to1$. At finite density this follows
by continuity; at infinite density use $r(Y_\lambda)\sim
(3\sqrt5j_{0,1}/2)\lambda^{-1/2}$. This relative comparison is
needed because both scales tend to zero at infinite density.
\end{proof}

\subsection{Separating configurations}\label{subsec:tight-sep}

\begin{lemma}[Separating configurations]\label{lem:tight-separating}
Fix $k$ and $L$. Uniformly in $n$, with lengths scaled by $\alpha_x$,
\begin{equation}
 \mathbb E^{\mathrm{WP}}_{g,n}
 \bigl[\text{ordered separating tight $k$-tuples with scaled
 lengths at most $L$}\bigr]=o(1),
 \label{eq:tight-separating-count}
\end{equation}
where only simple, pairwise disjoint tuples are counted, and separating
means that their union disconnects the filled surface.
\end{lemma}

\begin{proof}
For a stable connected cut graph with $k$ edges and $Q$ vertices, let
$(g_i,p_i)$ be the component types after filling the cusps. Then
\[
 p_i\ge1,\quad h_i=2g_i-2+p_i\ge1,\quad
 \sum_i h_i=h_0,\quad \sum_i p_i=2k,\quad
 \sum_i g_i=g-k+Q-1.
\]
For fixed positive $p$, Proposition~\ref{prop:prefactor} implies the
two-sided bound
\begin{equation}
 F_{g,p}(X(Y))\asymp_p
 4^{-g}\mathcal Z(Y)^{-h}r(Y)^{p/2}f(h),
 \qquad f(h)=\frac{\Gamma(h)}{\sqrt h},\quad h=2g-2+p.
 \label{eq:tight-gamma}
\end{equation}
The constants are uniform in $Y$. Bounded stable genera are included
by positivity and continuity of the normalized finite amplitudes on
$[0,y_0]$. For $p=0$, $g\ge2$, the same bound follows from $p=1$
and \eqref{eq:TL24}.

The identities for the component types cancel all powers of
$\mathcal Z$ in the product-to-parent ratio. They leave
\[
 r(Y)^{-k}\frac{\prod_i F_{g_i,p_i}(X(Y))}{\Phi_g(X(Y))}
 \le C_k\frac{\prod_i f(h_i)}{f(h_0)}.
\]
To sum over genera, observe that
$\sum_{a=1}^{\ell-1}f(a)f(\ell-a)\le C f(\ell-1)$.
Indeed,
\[
 \frac{\Gamma(a)\Gamma(\ell-a)}{\Gamma(\ell)}
 =\frac{1}{(\ell-1)\binom{\ell-2}{a-1}},
\]
the reciprocal binomial coefficients have uniformly bounded sum,
and $\sqrt{\ell}/\sqrt{a(\ell-a)}$ is bounded. Iteration gives
$f^{*Q}(\ell)\le C_Q f(\ell-Q+1)$. Since the marking allocations and
gluings have bounded multiplicity for fixed $k$, it follows that
\begin{equation}
 r(Y)^{-k}\sum_{\substack{\text{stable cut types}\\Q\text{ components}}}
 \frac{\prod_i F_{g_i,p_i}(X(Y))}{\Phi_g(X(Y))}
 \le C_k g^{-(Q-1)},\qquad 2\le Q\le k+1.
 \label{eq:tight-separating-real}
\end{equation}

For scaled lengths in a fixed compact box,
\eqref{eq:tight-real} bounds each boundary factor uniformly, including
on small-genus components. Every product in this positive sum has
the factor $V^{h_0}$ by \eqref{eq:tight-positive}. Its largest tilted
point probability is at most $C/\sigma$, by Lemma~\ref{lem:point-mass}. Applying \eqref{eq:tight-coefficient}, against the
denominator probability at least $c/\sigma$, therefore transfers
\eqref{eq:tight-separating-real} to coefficients. For bounded $n$,
use $x=c/h_0$ and the positive limiting denominator probability;
the numerator probability is at most one.

In the tight integration formula~\cite[Proposition~3.2]{BL}, the
physical normalization costs $P^{-k}$ and the length Jacobian is
$\alpha_x^{-2k}=(P/(4r))^k$. These cancel the remaining $r^k$.
The symmetry factors are at most one. This proves
\eqref{eq:tight-separating-count}.
\end{proof}

\subsection{Proof of Theorem~\ref{thm:tight-spectrum}}\label{subsec:tight-proof}

\begin{proof}[Proof of Theorem~\ref{thm:tight-spectrum}]
For fixed disjoint intervals in $(0,\infty)$, count ordered tuples
of simple, pairwise disjoint tight curves in those intervals.
The connected-complement symmetry is $2^{-k}$ in the tight integration
formula. Lemma~\ref{lem:tight-connected}, with each cut length
appearing on two boundary components, gives the limiting integrand
\[
 \prod_{i=1}^k\frac12 t_i S(t_i/2)^2\,\mathrm dt_i
 =\prod_{i=1}^k\frac{\cosh t_i-1}{t_i}\,\mathrm dt_i.
\]
Lemma~\ref{lem:tight-separating} removes all other cut types.
Thus the disjoint-tuple moments have the Poisson limits of every
fixed order.

Here the disjoint-tuple moments must be related to the actual count
moments: a small probability of intersections alone would not suffice.
Choose $L$ above all interval endpoints, and let $A_L$ be the event that
scaled lengths at most $L$ include a
self-intersecting tight curve or two intersecting tight curves. Write
$\widehat N_{\mathbf r}$ for the disjoint-tuple count in the prescribed
intervals, with multiplicities $\mathbf r=(r_1,\ldots,r_j)$.
We use the deterministic constructions in the proofs of
\cite[Propositions~5.2--5.3]{BL}, rather than their probabilistic
conclusions in the near-critical model. These constructions apply after
rescaling by any positive length scale.

Fix $\varepsilon>0$ below the interval endpoints. On $A_L$, outside
the event of a tight systole shorter than $\varepsilon$, the construction
extends a counted tuple to a separating tight multicurve. Its number of
components and scaled lengths are bounded in terms of
$\varepsilon,L,\mathbf r$, independently of $g,n$. If there is a shorter
tight systole, either it adds a disjoint curve in $[0,\varepsilon]$, or
the tuple again extends to a separating multicurve. The original ordered
tuple is retained in these extensions. Thus the separating terms vanish
by \eqref{eq:tight-separating-count}, and the mixed disjoint-tuple limit gives
\[
 \limsup_{g\to\infty}\mathbb E[
   \widehat N_{\mathbf r}\mathbf1_{A_L}]
 \le \nu([0,\varepsilon])\prod_{i=1}^j\nu(I_i)^{r_i},
\]
where $I_i$ are the prescribed intervals. Letting $\varepsilon\downarrow0$
removes the right side. The same construction without a marked tuple
gives $\mathbb P(A_L)\to0$. Conditional on $A_L^c$, factorial counts
are exactly disjoint-tuple counts. Their limiting moments therefore
remain unchanged, and the moment-determinate Poisson law gives joint
count convergence.

The first-moment bound for simple tight curves in $[0,\varepsilon]$
tends to $\nu([0,\varepsilon])\to0$. A shortest tight curve is
simple, so the same bound controls the event of any tight curve in
this window. This extends the convergence to $[0,\infty)$ and gives
tightness on compact intervals. Finally the disjoint-tuple moments
bound the number of curves in a fixed window in probability.
For any fixed upper bound on that number,
\eqref{eq:tight-separating-count} rules out every disconnected
complement. Letting the upper bound increase proves the last claim.
\end{proof}

\subsection{Proof of Corollary~\ref{cor:tight-boltzmann}}\label{subsec:tight-boltzmann-proof}

\begin{proof}[Proof of Corollary~\ref{cor:tight-boltzmann}]
Write $x=P\mu_g$ and $h=2g-2$. The case $x=0$ has $N=0$.
Otherwise \eqref{eq:TL24} at $p=0$ gives
\[
 \mathbb E N=h\Psi(x)+O(\Psi(x)).
\]
If $h\Psi(x)$ stays bounded, $N$ is tight and both $N/g$ and
$2\Psi(x)$ tend to zero. If $h\Psi(x)\to\infty$, integrating
\eqref{eq:TL30} at $p=0$ gives the same real central limit theorem
as \eqref{eq:TL39}, hence
\[
 N=h\Psi(x)+O_{\mathbb P}\!\left(\sqrt{h\mathscr D\Psi(x)}\right).
\]
The endpoint estimates \eqref{eq:TL36} imply
$\mathscr D\Psi/\Psi^2\le C(1+1/\Psi)$. It follows that
$N/(h\Psi(x))\to1$ in probability in this second case.
Continuity of $\alpha$ at finite density and its
$\lambda^{-1/4}$ asymptotic at infinity therefore give, in both cases,
\[
 \frac{\alpha(N/g)}{\alpha(2\Psi(x))}
 =\frac{\alpha(N/g)}{a_g}\longrightarrow1
 \quad\text{in probability}.
\]
Subsequences cover arbitrary choices of $x$.

For each fixed nonnegative compactly supported continuous test function,
the conditional Laplace-functional convergence in
Theorem~\ref{thm:tight-spectrum} is uniform in $n$: a sequence violating
uniformity would contradict that theorem. Averaging over $N$ gives the
Poisson limit with random normalization $\alpha(N/g)$.
The displayed scale comparison and continuity of dilation of point
measures replace this by $a_g$. For $\mu_g\to\mu<\mu_{\rm cr}$, a change of
scale gives the stated unscaled intensity. At the critical endpoint,
$r\sim3(y_0-Y)$ and $x_c-x\sim(-X''(y_0))(y_0-Y)^2/2$ show that
$a_g$ is asymptotic to the scale in~\cite[Theorem~1.2]{BL}.
\end{proof}

\section{Proofs for the critical length spectrum}\label{app:critical}

This appendix proves Theorem~\ref{thm:critical-cluster}, following the
outline in Subsection~\ref{sec:critical-cluster}. Throughout,
$n/\sqrt g\to c\in(0,\infty)$, and $P_s$, $K_s$, $\mu_c$ and $\nu$ are as
in \eqref{eq:cluster-intensities}.

\subsection{Ordinary cuts and cusp allocations}

Short multicurves are organized by their cut graphs: the vertices are the
components of the complement, labeled by their genera, and the edges are
the curves. Lemma~\ref{lem:critical-cuts} bounds the expected number of
multicurves with a given cut graph. It shows that, in the limit, only
nonseparating curves and boundaries of two-cusp pairs of pants
contribute; the later subsections compute these contributions jointly.

Fix a connected cut graph with $k\ge1$ edges and $Q$ vertices;
loops and multiple edges are allowed. At vertex $i$ let $a_i$ be the
genus and $p_i\ge1$ its degree. Then
\[
 \sum_i a_i=g-k+Q-1,\qquad \sum_i p_i=2k.
\]
Partition the vertices into $u$ of type $(0,1)$, $v$ of type $(0,2)$,
and $w$ with $h_i=2a_i-2+p_i\ge1$. Put $h_0=2g-2$. For large
$g$, $w\ge1$, and the latter complexities sum to $H=h_0+u$.
A component of type $(0,1)$ needs at least two cusps, and one of
type $(0,2)$ at least one.

\begin{lemma}[Critical cut estimate]\label{lem:critical-cuts}
Fix $k$ and a bound on the cut lengths. At $n/\sqrt g\to c>0$,
the expected number of cut multicurves of the specified graph and
vertex partition, summed over all admissible genera and cusp allocations,
is $O(g^{1-w-v/2})$. If a specified $(0,1)$ component has at least
three cusps, the bound gains a factor $g^{-1/2}$.
\end{lemma}

\begin{proof}
Write $h_n=2g-2+n$ for the complexity of the cusped surface. The volume bounds in \cite[Lemmas~18--19]{SW} give
\[
 V_{a,b}\le C\frac{(4\pi^2)^{d-1}\Gamma(d)}{\sqrt d},
 \quad d=2a-2+b\ge1,\qquad
 V_{g,n}\ge c_0\frac{(4\pi^2)^{h_n-1}\Gamma(h_n)}{\sqrt{h_n}}.
\]
The first bound includes $d=1$ after enlarging $C$ for $(0,3)$ and
$(1,1)$; the second uses $n=O(\sqrt g)$. If vertex $i$ receives
$q_i$ cusps, put $d_i=h_i+q_i$. Then $\sum_i d_i=h_n$ and
$\sqrt{h_n}/\prod_i\sqrt{d_i}\le\sqrt Q$. Mirzakhani integration,
the bounded-boundary estimate \cite[Lemma~3.2]{Mir}, and the
multinomial allocation of cusp labels therefore bound the expected
cut count $E_{g,n}$ by
\[
 E_{g,n}\le\frac{C_{k,c}\,n!}{\Gamma(h_n)}
       \sum_{\substack{\text{admissible genera}\\q_1+\cdots+q_Q=n}}
       \prod_i\frac{\Gamma(h_i+q_i)}{q_i!}.
\]
Here and below the constants may also depend on the length bound.
The powers of $4\pi^2$ leave only the fixed factor $(4\pi^2)^{1-Q}$.

The cusp sums at the two unstable vertex types are bounded using
\begin{equation}
 \phi_{0,1}(z)=\sum_{q\ge2}\frac{z^q}{q(q-1)}=\frac{z^2}{2}+O(z^3),
 \qquad \phi_{0,2}(z)=\sum_{q\ge1}\frac{z^q}{q}=z+O(z^2).
 \label{eq:cluster-unstable}
\end{equation}
At a stable vertex their generating function is
$\Gamma(h)(1-z)^{-h}$. Moreover
\[
 \sum_{\substack{h_1+\cdots+h_w=H\\h_i\ge1}}
       \prod_i\Gamma(h_i)\le C_w\Gamma(H-w+1).
\]
For $w=2$ this follows by dividing by $\Gamma(H-1)$ and summing
$\binom{H-2}{h_1-1}^{-1}$; iteration gives the general bound.
Dropping genus parity constraints only increases the sum.
Let $a_m=[z^m]\phi_{0,1}(z)^u\phi_{0,2}(z)^v$. Expanding $(1-z)^{-H}$ gives
\[
 E_{g,n}\le C_{k,c}\frac{\Gamma(H-w+1)}{\Gamma(H)}
       \sum_{m\le n}a_m(n)_m\frac{\Gamma(h_n+u-m)}{\Gamma(h_n)}.
\]
Uniformly for $m\le n=O(\sqrt{h_n})$, the finite-product formula for
the Gamma ratio gives
$\Gamma(h_n+u-m)/\Gamma(h_n)\le C_{k,c}h_n^{u-m}$: the logarithmic
error is $O((m+1)^2/h_n)$. Since $(n)_m\le n^m$, positivity yields
\begin{equation}
 E_{g,n}\le C_{k,c}h_n^{u-w+1}\phi_{0,1}(n/h_n)^u\phi_{0,2}(n/h_n)^v
          =O(g^{1-w-v/2}).
 \label{eq:critical-graph-bound}
\end{equation}
Requiring three cusps at a specified $(0,1)$ vertex replaces one
$\phi_{0,1}(z)$ by $\phi_{0,1}(z)-z^2/2$ and gains $n/h_n=O(g^{-1/2})$.
\end{proof}

Only $w=1$, $v=0$ can survive: the graph then has one large-genus
vertex, loops at that vertex, and two-cusp pants attached as leaves.
All other graphs, and all leaves with more than two cusps, vanish.

\subsection{Excluding complex short subsurfaces}\label{subsec:critical-exclusion}

We need expectation bounds, not only a bound on the probability of a
bad configuration, so we give the multiplicity argument explicitly.

\begin{lemma}[Bounded primitive count]\label{lem:primitive-count}
For fixed $K,L$, a finite-area hyperbolic surface $S$ with geodesic
boundary and $0<-\chi(S)\le K$ has at most $C(K,L)$ primitive
unoriented geodesics of length at most $L$, uniformly in its metric.
\end{lemma}

\begin{proof}
Double $S$ along its boundaries. The double has area at
most $4\pi K$. Fix a Margulis constant $\varepsilon>0$. A primitive
geodesic wholly in a thin component is its core; there are boundedly
many such cores. Every other geodesic meets the thick part. A maximal
$\varepsilon/4$-separated net there has bounded cardinality by area
packing. At each net point, orbit points of a lift are uniformly
separated, so hyperbolic ball packing bounds the number of based
loops of length at most $L+\varepsilon$. Connecting a short geodesic
to a net point proves the assertion. Primitivity is preserved upon
inclusion in the double or in an ambient surface: a root would have
the same geodesic image and would already be a closed curve in $S$.
\end{proof}

\begin{lemma}[Complex short configurations]\label{lem:critical-complex}
Fix $r$ and $L$. As $g\to\infty$ with $n/\sqrt g\to c\in(0,\infty)$,
\begin{equation}
 \mathbb E[\text{intersecting ordered simple $r$-tuples of lengths
 at most $L$}]=O_{r,L,c}(g^{-1/2}),
 \label{eq:critical-intersections}
\end{equation}
\begin{equation}
 \mathbb E[\text{nonsimple primitive curves of length at most $L$
 not filling two-cusp pants}]=O_{L,c}(g^{-1/2}),
 \label{eq:critical-bad-curves}
\end{equation}
and the expected number of simple separating curves of length at most $L$
that are not waists of two-cusp pairs of pants tends to zero.
\end{lemma}

\begin{proof}
For fixed $r,L$, a tuple of $r$ distinct simple geodesics of length
at most $L$ has bounded total intersection number. We use the local
segment argument in the proof of \cite[Proposition~4.5]{MP}; it
applies equally in the presence of cusps. Thicken each nonannular
component of their union, fill complementary disks and peripheral
once-punctured disks, and tighten the essential boundaries. The
resulting pieces have bounded topology and total boundary length at
most $2rL$. Keep one ambient geodesic for each parallel boundary class.
Empty annuli between parallel copies collapse; their two incidences
are the two sides of a single edge, possibly a loop. An annulus with
cusps does not collapse and is a $(0,2)$ component counted among the
$v$ vertices.
Isolated tuple curves are retained as additional cut edges, unless
they are already one of the tightened boundaries. Minimal intersection
ensures that these boundaries are disjoint from the original tuple.

There are finitely many resulting graph shapes and markings at fixed
$r,L$. For a fixed cut multicurve, Lemma~\ref{lem:primitive-count}
bounds the choices of original curves in each marked piece. Hence the
tuple count is bounded pointwise by a constant times a finite sum of
marked cut-multicurve counts. An intersecting simple tuple cannot fill
two-cusp pants. Its marked component either has $h_i\ge1$, forcing
a second stable vertex to carry the unbounded genus, or has type
$(0,2)$, or has type $(0,1)$ with at least three cusps. The bounds in
Lemma~\ref{lem:critical-cuts} give \eqref{eq:critical-intersections}.

For a single nonsimple primitive curve of length at most $L$, use
the filling construction of \cite[Definition~4.10 and Lemma~4.11]{Hide}.
The filled piece has boundary length at most $2L$ and area at most
$4L$, so its topology is bounded. The same recording retains formal
boundary sides that tighten to the same ambient curve as sides of a
cut edge. Lemma~\ref{lem:primitive-count} controls the number of curves
filling each piece. The preceding classification therefore gives
\eqref{eq:critical-bad-curves}.
A filled piece cannot be the whole ambient surface for large $g$,
since its topology is bounded. Finally, the one-edge case of
Lemma~\ref{lem:critical-cuts} shows that simple separating curves
of bounded length other than two-cusp pants waists vanish in expectation.
\end{proof}

\subsection{Joint nonseparating and pants counts}

Let $\mathcal N_g$ count simple nonseparating lengths and let
$\mathcal Q_g$ count two-cusp pants waist lengths.

\begin{lemma}[Joint Poisson limit]\label{lem:critical-skeleton}
As $g\to\infty$ with $n/\sqrt g\to c\in(0,\infty)$,
\begin{equation}
 (\mathcal N_g,\mathcal Q_g)\Longrightarrow
 (\mathcal N,\mathcal Q_c)
 \label{eq:critical-skeleton}
\end{equation}
jointly in the vague topology on $[0,\infty)$, where $\mathcal N$ and
$\mathcal Q_c$ are independent Poisson point processes of intensities $\nu$
and $\mu_c$.
\end{lemma}

\begin{proof}
In the mixed
factorial moments of $(\mathcal N_g,\mathcal Q_g)$, \eqref{eq:critical-intersections} removes intersecting
tuples. Lemma~\ref{lem:critical-cuts} removes every disjoint cut type
except the graph with a single large component, $k$ nonseparating
loops and $r$ two-cusp leaves. For prescribed bounded intervals for
the $k+r$ lengths, its exact contribution is
\begin{equation}
 \frac{(n)_{2r}}{2^{r+k}V_{g,n}}
 \int V_{g-k,n-r+2k}(\mathbf s,\mathbf t,\mathbf t,0^{n-2r})
        \prod_{j=1}^r s_j\,\mathrm ds_j
        \prod_{i=1}^k t_i\,\mathrm dt_i.
 \label{eq:critical-mixed-moment}
\end{equation}
Here $(n)_{2r}$ is the falling factorial. The factor $2^{-r}$ counts
unordered pairs of cusp labels at the ordered pants leaves, as in
\cite[Lemma~28]{SW}; $2^{-k}$ exchanges the sides of each
nonseparating cut, as in \cite[Proposition~4.2]{MP}. There is no
factor $r!$ or $k!$, since both tuples are ordered.

The volume ratio is a product of $r$ cusp ratios and $k$ genus ratios.
The cusp ratios satisfy $V_{g,j+1}/V_{g,j}\sim8\pi^2g$ for $j\le n$ by
\cite[Lemma~5.1(ii)]{MZ}, since $n=o(g)$.
For the genus ratios, \cite[Lemma~5.1(iii)]{MZ} gives only an upper bound,
and the two-sided limit at $n\asymp\sqrt g$ is the zero-density case
$\mathcal B(0)=1$ of Theorem~\ref{thm:volume}. Together,
\[
 \frac{V_{g-k,n-r+2k}}{V_{g,n}}\sim(8\pi^2g)^{-r}.
\]
Proposition~\ref{prop:bounded-boundaries}, applied at zero limiting
density, where $\mathfrak b_0(L)=\sinh(L/2)/(L/2)$, supplies, uniformly
on the fixed length box, the boundary factor
$\prod_j S(s_j/2)\prod_i S(t_i/2)^2$, with $S(t)=\sinh(t)/t$; compare
Shen--Wu's bounded-boundary estimate~\cite[Proposition~25]{SW}, which
they attribute essentially to \cite[Proposition~3.1]{MP}
and~\cite[Lemma~20]{NWX}.
Consequently the limiting factors in \eqref{eq:critical-mixed-moment}
are
\[
 \frac12 tS(t/2)^2\,\mathrm dt=\nu(\mathrm dt),\qquad
 \frac{c^2}{16\pi^2}sS(s/2)\,\mathrm ds=\mu_c(\mathrm ds).
\]
All mixed factorial moments on disjoint bounded intervals converge to
those of independent Poisson counts. Higher fixed moments give uniform
integrability, and the multivariate Poisson law is moment determinate.
The first-moment bounds near zero are $O(\varepsilon^2)$. This proves the
lemma.
\end{proof}

\subsection{Proof of Theorem~\ref{thm:critical-cluster}}

\begin{proof}[Completion of the proof of Theorem~\ref{thm:critical-cluster}]
Fix $L<\infty$. By Lemma~\ref{lem:critical-complex}, with probability
tending to one every nonsimple primitive curve of length at most $L$
fills two-cusp pants with waist at most $2L$. By
\eqref{eq:critical-intersections}, the pants waists up to $2L$ are
pairwise disjoint with probability tending to one. The pants interiors
are then disjoint as well. Other simple separating curves up to $L$
are absent with probability tending to one. On this event, the full
primitive counting measure restricted to $[0,L]$ is exactly
\[
 \left.\left(\mathcal N_g+
       \sum_{s\in\mathcal Q_g\cap[0,2L]}K_s\right)\right|_{[0,L]}.
\]
Every interior primitive geodesic of $P_s$ is nonsimple and fills it.
The filling boundary bound implies $s\le2\ell$ for such a geodesic
of length $\ell$. Thus pants with waist greater than $2L$ cannot
contribute, and the limiting sum is locally finite.

For $\varepsilon>0$, the family $P_s$, $\varepsilon\le s\le2L$,
is compact in its fixed-topology moduli space. Lengths of marked
classes vary continuously, and only finitely many such classes have
length at most $L$ somewhere on this compact family. For the latter
assertion, use a finite cover by uniformly bilipschitz markings,
with standard neighborhoods at the two cusps, and the finiteness of
each center metric's length spectrum below a fixed bound. Hence
$s\mapsto\langle f,K_s\rangle$ is continuous away from zero for
continuous compactly supported $f$. Apply \eqref{eq:critical-skeleton}
with waist cutoff $\varepsilon$, choosing $L$ larger than the support
of $f$. The probability of removing any waist is at most
$\mathbb E\mathcal Q_g([0,\varepsilon])$, whose limit is
$\mu_c([0,\varepsilon])=O(\varepsilon^2)$. Letting
$\varepsilon\downarrow0$ proves \eqref{eq:critical-cluster}.
The independent Poisson Laplace functionals give
\eqref{eq:cluster-laplace}.
\end{proof}

\section*{Acknowledgements}
We thank Yang Shen for his valuable comments on the exposition and the discussion of prior work. His suggestions prompted us to give a clearer account of the literature on Weil--Petersson volume asymptotics, the geometry of moduli spaces, and the spectra of random hyperbolic surfaces, and to clarify the relationship between our spectral applications and the earlier results and methods of Shen and Wu.

We thank Hao Xu for helpful correspondence concerning the work of Liu and Xu and its relationship to the present paper. In particular, his comments prompted us to distinguish more clearly the uniformity in the spectral-curve parameter at a fixed number of markings from the uniformity in the total cusp count established here, and to explain that the determination of Zograf's constant in our proof does not require the latter uniformity. This exchange helped us improve the comparison of the two deformation arguments and the description of their respective scopes.

We acknowledge Gewu Intelligence Lab for providing the collaborative research environment in which this project was developed.

\section*{Statements and Declarations}
\noindent\textit{Funding.} The authors received no funding for this work.

\noindent\textit{Competing interests.} The authors have no relevant financial or non-financial interests to disclose.

\noindent\textit{Data availability.} Data sharing is not applicable to this article, as no datasets were generated or analysed.

\section*{Disclosure of automated assistance}
OpenAI GPT-6.1 Sol, OpenAI GPT-6 Astra, Anthropic Claude Fable 5.1, Anthropic Claude Opus 5.5, and Google Gemini 4 Argon contributed to the mathematical development, checking, and preparation of this paper. Their combined work included developing proof arguments and uniform asymptotic estimates, checking identities and references, comparing prior results, and revising the exposition.

\begin{thebibliography}{99}
\bibitem{AggAbelian} Amol Aggarwal, \emph{Large genus asymptotics for volumes of strata of abelian differentials}, with an appendix by Anton Zorich, J. Amer. Math. Soc. \textbf{33} (2020), no.~4, 941--989. \url{https://doi.org/10.1090/jams/947}; arXiv:1804.05431v3.
\bibitem{Agg} Amol Aggarwal, \emph{Large genus asymptotics for intersection numbers and principal strata volumes of quadratic differentials}, Invent. Math. \textbf{226} (2021), no.~3, 897--1010. \url{https://doi.org/10.1007/s00222-021-01059-9}.
\bibitem{AMNO} Mohamad Alameddine, Olivier Marchal, Nikita Nikolaev, and Nicolas Orantin, \emph{The WKB Riemann--Hilbert problem for the Painlev\'e~I system}, preprint (2026), \href{https://arxiv.org/abs/2610.20123v1}{arXiv:2610.20123v1}.
\bibitem{AM} Nalini Anantharaman and Laura Monk, \emph{A high-genus asymptotic expansion of Weil--Petersson volume polynomials}, J. Math. Phys. \textbf{63} (2022), no.~4, 043502. \url{https://doi.org/10.1063/5.0039385}.
\bibitem{AMPartOne} Nalini Anantharaman and Laura Monk, \emph{Friedman--Ramanujan functions in random hyperbolic geometry and application to spectral gaps}, preprint (2023), \href{https://arxiv.org/abs/2304.02678v1}{arXiv:2304.02678v1}.
\bibitem{AMGap} Nalini Anantharaman and Laura Monk, \emph{Friedman--Ramanujan functions in random hyperbolic geometry and application to spectral gaps II}, preprint (2025), \href{https://arxiv.org/abs/2502.12268v2}{arXiv:2502.12268v2}.
\bibitem{AMTwoNinths} Nalini Anantharaman and Laura Monk, \emph{Typical hyperbolic surfaces have a spectral gap greater than $2/9-\epsilon$}, Int. Math. Res. Not. IMRN \textbf{2026} (2026), no.~15, article~rnag117. \url{https://doi.org/10.1093/imrn/rnag117}; arXiv:2604.09792v1.
\bibitem{AC} Enrico Arbarello and Maurizio Cornalba, \emph{Combinatorial and algebro-geometric cohomology classes on the moduli spaces of curves}, J. Algebraic Geom. \textbf{5} (1996), no.~4, 705--749; arXiv:alg-geom/9406008.
\bibitem{BEG} Ga\"etan Borot, Bertrand Eynard, and Alessandro Giacchetto, \emph{The factorial growth of topological recursion}, Lett. Math. Phys. \textbf{115} (2025), no.~3, article~62. \url{https://doi.org/10.1007/s11005-025-01950-z}.
\bibitem{BFJ} S\'ebastien Boucksom, Charles Favre, and Mattias Jonsson, \emph{Differentiability of volumes of divisors and a problem of Teissier}, J. Algebraic Geom. \textbf{18} (2009), no.~2, 279--308. \url{https://doi.org/10.1090/S1056-3911-08-00490-6}; arXiv:math/0608260.
\bibitem{Budd} Timothy Budd, \emph{Irreducible metric maps and Weil--Petersson volumes}, Comm. Math. Phys. \textbf{394} (2022), no.~2, 887--917. \url{https://doi.org/10.1007/s00220-022-04418-6}; arXiv:2012.11318v2.
\bibitem{BC} Timothy Budd and Nicolas Curien, \emph{Random punctured hyperbolic surfaces \& the Brownian sphere}, preprint (2025), \href{https://arxiv.org/abs/2508.18792v1}{arXiv:2508.18792v1}.
\bibitem{BL} Timothy Budd and Tanguy Lions, \emph{The tight length spectrum of large-genus random hyperbolic surfaces with many cusps}, preprint (2025), revised 2026, \href{https://arxiv.org/abs/2506.02611v3}{arXiv:2506.02611v3}.
\bibitem{BZ} Timothy Budd and Bart Zonneveld, \emph{Topological recursion of the Weil--Petersson volumes of hyperbolic surfaces with tight boundaries}, J. Math. Phys. \textbf{65} (2024), no.~9, 092302. \url{https://doi.org/10.1063/5.0192711}.
\bibitem{Buser} Peter Buser, \emph{A note on the isoperimetric constant}, Ann. Sci. \'Ecole Norm. Sup. (4) \textbf{15} (1982), no.~2, 213--230. \url{https://doi.org/10.24033/asens.1426}.
\bibitem{CMSZ} Dawei Chen, Martin M\"oller, Adrien Sauvaget, and Don Zagier, \emph{Masur--Veech volumes and intersection theory on moduli spaces of Abelian differentials}, Invent. Math. \textbf{222} (2020), no.~1, 283--373. \url{https://doi.org/10.1007/s00222-020-00969-4}; arXiv:1901.01785.
\bibitem{ChenMullane} Dawei Chen and Scott Mullane, \emph{A tale of two volumes of moduli spaces: Weil--Petersson and Masur--Veech}, preprint (2026), \href{https://arxiv.org/abs/2603.07318v1}{arXiv:2603.07318v1}.
\bibitem{CCH} Ovidiu Costin, Rodica D. Costin, and Min Huang, \emph{Tronqu\'ee solutions of the Painlev\'e equation $P_{\rm I}$}, Constr. Approx. \textbf{41} (2015), no.~3, 467--494. \url{https://doi.org/10.1007/s00365-015-9287-1}; arXiv:1310.5330v2.
\bibitem{DGZZBounds} Vincent Delecroix, \'Elise Goujard, Peter Zograf, and Anton Zorich, \emph{Uniform lower bound for intersection numbers of $\psi$-classes}, SIGMA \textbf{16} (2020), article~086, 13~pp. \url{https://doi.org/10.3842/SIGMA.2020.086}.
\bibitem{DGZZ} Vincent Delecroix, \'Elise Goujard, Peter Zograf, and Anton Zorich, \emph{Masur--Veech volumes, frequencies of simple closed geodesics, and intersection numbers of moduli spaces of curves}, Duke Math. J. \textbf{170} (2021), no.~12, 2633--2718. \url{https://doi.org/10.1215/00127094-2021-0054}. Original preprint (2019): \href{https://arxiv.org/abs/1908.08611v1}{arXiv:1908.08611v1}.
\bibitem{DVV} Robbert Dijkgraaf, Herman Verlinde, and Erik Verlinde, \emph{Loop equations and Virasoro constraints in non-perturbative two-dimensional quantum gravity}, Nuclear Phys. B \textbf{348} (1991), no.~3, 435--456. \url{https://doi.org/10.1016/0550-3213(91)90199-8}.
\bibitem{ET} Lorenz Eberhardt and Gustavo J. Turiaci, \emph{2D dilaton gravity and the Weil--Petersson volumes with conical defects}, Comm. Math. Phys. \textbf{405} (2024), no.~4, article~103. \url{https://doi.org/10.1007/s00220-024-04964-1}; arXiv:2304.14948v2.
\bibitem{EZ} Alex Eskin and Anton Zorich, \emph{Volumes of strata of Abelian differentials and Siegel--Veech constants in large genera}, Arnold Math. J. \textbf{1} (2015), no.~4, 481--488. \url{https://doi.org/10.1007/s40598-015-0028-0}; arXiv:1507.05296.
\bibitem{EGGGL} Bertrand Eynard, Elba Garcia-Failde, Alessandro Giacchetto, Paolo Gregori, and Danilo Lewa\'nski, \emph{Resurgent large genus asymptotics of intersection numbers}, Forum Math. Pi \textbf{14} (2026), article~e21. \url{https://doi.org/10.1017/fmp.2026.10037}; arXiv:2309.03143v1.
\bibitem{EGGLS} Bertrand Eynard, Elba Garcia-Failde, Paolo Gregori, Danilo Lewa\'nski, and Ricardo Schiappa, \emph{Resurgent asymptotics of Jackiw--Teitelboim gravity and the nonperturbative topological recursion}, Ann. Henri Poincar\'e \textbf{25} (2024), no.~9, 4121--4193. \url{https://doi.org/10.1007/s00023-023-01412-z}.
\bibitem{EO} Bertrand Eynard and Nicolas Orantin, \emph{Weil--Petersson volume of moduli spaces, Mirzakhani's recursion and matrix models}, preprint (2007), \href{https://arxiv.org/abs/0705.3600v1}{arXiv:0705.3600v1}.
\bibitem{GIKM} Stavros Garoufalidis, Alexander Its, Andrei Kapaev, and Marcos Mari\~no, \emph{Asymptotics of the instantons of Painlev\'e I}, Int. Math. Res. Not. IMRN \textbf{2012} (2012), no.~3, 561--606. \url{https://doi.org/10.1093/imrn/rnr029}; arXiv:1002.3634v3.
\bibitem{GPRS} Luca Griguolo, Jacopo Papalini, Lorenzo Russo, and Domenico Seminara, \emph{Asymptotics of Weil--Petersson volumes and two-dimensional quantum gravities}, SciPost Phys. \textbf{17} (2024), no.~6, 156. \url{https://doi.org/10.21468/SciPostPhys.17.6.156}.
\bibitem{Gru} Samuel Grushevsky, \emph{An explicit upper bound for Weil--Petersson volumes of the moduli spaces of punctured Riemann surfaces}, Math. Ann. \textbf{321} (2001), no.~1, 1--13. \url{https://doi.org/10.1007/PL00004496}.
\bibitem{Guillermit} Marthe Guillermit, \emph{A refinement of the asymptotic expansion of Weil--Petersson volumes}, preprint (2026), \href{https://arxiv.org/abs/2607.20341v1}{arXiv:2607.20341v1}.
\bibitem{GY} Jindong Guo and Di Yang, \emph{On the large genus asymptotics of psi-class intersection numbers}, Math. Ann. \textbf{388} (2024), no.~1, 61--97. \url{https://doi.org/10.1007/s00208-022-02505-6}; arXiv:2110.06774v3.
\bibitem{GYZ} Jindong Guo, Di Yang, and Don Zagier, \emph{On uniform large genus asymptotics of Witten's intersection numbers}, preprint (2026), \href{https://arxiv.org/abs/2603.15233v1}{arXiv:2603.15233v1}.
\bibitem{GHS} Qi Guo, Bobo Hua, and Yang Shen, \emph{Spectral gaps for noncompact hyperbolic surfaces with linearly many cusps}, preprint (2025), revised 2026, \href{https://arxiv.org/abs/2507.16794v3}{arXiv:2507.16794v3}.
\bibitem{HMO} Yasuyuki Hatsuda, Takaki Matsumoto, and Kazumi Okuyama, \emph{Non-perturbative effects in JT gravity from KdV equations}, J. High Energy Phys. \textbf{2026} (2026), no.~3, 056. \url{https://doi.org/10.1007/JHEP03(2026)056}; arXiv:2505.16433v1.
\bibitem{HSWX} Yuxin He, Yang Shen, Yunhui Wu, and Yuhao Xue, \emph{Non-simple systoles on random hyperbolic surfaces for large genus}, preprint (2023), revised 2025, \href{https://arxiv.org/abs/2308.16447v2}{arXiv:2308.16447v2}.
\bibitem{HWX} Yuxin He, Yunhui Wu, and Yuhao Xue, \emph{Uniform spectral gaps for random hyperbolic surfaces with not many cusps}, preprint (2026), \href{https://arxiv.org/abs/2602.08352v1}{arXiv:2602.08352v1}.
\bibitem{Hide} Will Hide, \emph{Spectral gap for Weil--Petersson random surfaces with cusps}, Int. Math. Res. Not. IMRN \textbf{2023} (2023), no.~20, 17411--17460. \url{https://doi.org/10.1093/imrn/rnac293}.
\bibitem{HMT} Will Hide, Davide Macera, and Joe Thomas, \emph{Spectral gap with polynomial rate for Weil--Petersson random surfaces}, preprint (2025), \href{https://arxiv.org/abs/2508.14874v2}{arXiv:2508.14874v2}.
\bibitem{HideMagee} Will Hide and Michael Magee, \emph{Near optimal spectral gaps for hyperbolic surfaces}, Ann. of Math. (2) \textbf{198} (2023), no.~2, 791--824. \url{https://doi.org/10.4007/annals.2023.198.2.6}; arXiv:2107.05292v3.
\bibitem{HTLargeN} Will Hide and Joe Thomas, \emph{Large-$n$ asymptotics for Weil--Petersson volumes of moduli spaces of bordered hyperbolic surfaces}, Comm. Math. Phys. \textbf{406} (2025), no.~9, article~203. \url{https://doi.org/10.1007/s00220-025-05369-4}.
\bibitem{HTSphere} Will Hide and Joe Thomas, \emph{Short geodesics and small eigenvalues on random hyperbolic punctured spheres}, Comment. Math. Helv. \textbf{100} (2025), no.~3, 463--506. \url{https://doi.org/10.4171/CMH/588}; arXiv:2209.15568v3.
\bibitem{HTMany} Will Hide and Joe Thomas, \emph{Small eigenvalues of hyperbolic surfaces with many cusps}, preprint (2024), revised July 2026, \href{https://arxiv.org/abs/2410.06093v2}{arXiv:2410.06093v2}.
\bibitem{IZ} Claude Itzykson and Jean-Bernard Zuber, \emph{Combinatorics of the modular group II: the Kontsevich integrals}, Internat. J. Modern Phys. A \textbf{7} (1992), no.~23, 5661--5705. \url{https://doi.org/10.1142/S0217751X92002581}.
\bibitem{Johnson} Clifford V. Johnson, \emph{A boundary-addition identity for Weil--Petersson volumes and their supersymmetric generalizations}, preprint (2026), \href{https://arxiv.org/abs/2609.21030v1}{arXiv:2609.21030v1}.
\bibitem{JR} Clifford V. Johnson and Jo\~ao Rodrigues, \emph{Non-perturbative data for Weil--Petersson volumes and intersection numbers using ordinary differential equations}, J. High Energy Phys. \textbf{2026} (2026), no.~8, 083. \url{https://doi.org/10.1007/JHEP08(2026)083}; arXiv:2601.03351v2.
\bibitem{JK} Nalini Joshi and Alexander V. Kitaev, \emph{On Boutroux's tritronqu\'ee solutions of the first Painlev\'e equation}, Stud. Appl. Math. \textbf{107} (2001), no.~3, 253--291. \url{https://doi.org/10.1111/1467-9590.00187}.
\bibitem{Kap} Andrei A. Kapaev, \emph{Quasi-linear Stokes phenomenon for the Painlev\'e first equation}, J. Phys. A \textbf{37} (2004), no.~46, 11149--11167. \url{https://doi.org/10.1088/0305-4470/37/46/005}.
\bibitem{KMZ} Ralph M. Kaufmann, Yuri I. Manin, and Don Zagier, \emph{Higher Weil--Petersson volumes of moduli spaces of stable $n$-pointed curves}, Comm. Math. Phys. \textbf{181} (1996), no.~3, 763--787. \url{https://doi.org/10.1007/BF02101297}.
\bibitem{Keel} Sean Keel, \emph{Basepoint freeness for nef and big line bundles in positive characteristic}, Ann. of Math. (2) \textbf{149} (1999), no.~1, 253--286. \url{https://doi.org/10.2307/121025}; arXiv:math/9901149.
\bibitem{Kim} Yusuke Kimura, \emph{JT gravity and the asymptotic Weil--Petersson volume}, Phys. Lett. B \textbf{811} (2020), article~135989. \url{https://doi.org/10.1016/j.physletb.2020.135989}; arXiv:2008.04141v2.
\bibitem{Kontsevich} Maxim Kontsevich, \emph{Intersection theory on the moduli space of curves and the matrix Airy function}, Comm. Math. Phys. \textbf{147} (1992), no.~1, 1--23. \url{https://doi.org/10.1007/BF02099526}.
\bibitem{Laz} Robert Lazarsfeld, \emph{Positivity in Algebraic Geometry I}, Ergeb. Math. Grenzgeb. (3) \textbf{48}, Springer, Berlin, 2004. \url{https://doi.org/10.1007/978-3-642-18808-4}.
\bibitem{LWZenodo} Zhiyi Li and Qihang Wang, \emph{Uniform asymptotics of Weil--Petersson volumes and their coefficients}, preprint, version~1.0, Zenodo, 16 October 2026. \url{https://doi.org/10.5281/zenodo.23414478}.
\bibitem{LipWright} Michael Lipnowski and Alex Wright, \emph{Towards optimal spectral gaps in large genus}, Ann. Probab. \textbf{52} (2024), no.~2, 545--575. \url{https://doi.org/10.1214/23-AOP1657}.
\bibitem{LX} Kefeng Liu and Hao Xu, \emph{A remark on Mirzakhani's asymptotic formulae}, Asian J. Math. \textbf{18} (2014), no.~1, 29--52. \url{https://doi.org/10.4310/AJM.2014.v18.n1.a2}; arXiv:1103.5136.
\bibitem{LXConstant} Kefeng Liu and Hao Xu, \emph{Uniform asymptotics and the Weil--Petersson constant}, preprint, 24 October 2026, \href{https://arxiv.org/abs/2610.33752v1}{arXiv:2610.33752v1}.
\bibitem{MNP} Michael Magee, Fr\'ed\'eric Naud, and Doron Puder, \emph{A random cover of a compact hyperbolic surface has relative spectral gap $\frac{3}{16}-\varepsilon$}, Geom. Funct. Anal. \textbf{32} (2022), no.~3, 595--661. \url{https://doi.org/10.1007/s00039-022-00602-x}.
\bibitem{ManZ} Yuri I. Manin and Peter Zograf, \emph{Invertible cohomological field theories and Weil--Petersson volumes}, Ann. Inst. Fourier (Grenoble) \textbf{50} (2000), no.~2, 519--535. \url{https://doi.org/10.5802/aif.1764}.
\bibitem{MT} Henry Maxfield and Gustavo J. Turiaci, \emph{The path integral of 3D gravity near extremality; or, JT gravity with defects as a matrix integral}, J. High Energy Phys. \textbf{2021} (2021), no.~1, article~118. \url{https://doi.org/10.1007/JHEP01(2021)118}; arXiv:2006.11317v2.
\bibitem{McShane} Greg McShane, \emph{Simple geodesics and a series constant over Teichmuller space}, Invent. Math. \textbf{132} (1998), no.~3, 607--632. \url{https://doi.org/10.1007/s002220050235}.
\bibitem{MirInv} Maryam Mirzakhani, \emph{Simple geodesics and Weil--Petersson volumes of moduli spaces of bordered Riemann surfaces}, Invent. Math. \textbf{167} (2007), no.~1, 179--222. \url{https://doi.org/10.1007/s00222-006-0013-2}.
\bibitem{MirJAMS} Maryam Mirzakhani, \emph{Weil--Petersson volumes and intersection theory on the moduli space of curves}, J. Amer. Math. Soc. \textbf{20} (2007), no.~1, 1--23. \url{https://doi.org/10.1090/S0894-0347-06-00526-1}.
\bibitem{MirCount} Maryam Mirzakhani, \emph{Growth of the number of simple closed geodesics on hyperbolic surfaces}, Ann. of Math. (2) \textbf{168} (2008), no.~1, 97--125. \url{https://doi.org/10.4007/annals.2008.168.97}.
\bibitem{Mir} Maryam Mirzakhani, \emph{Growth of Weil--Petersson volumes and random hyperbolic surface of large genus}, J. Differential Geom. \textbf{94} (2013), no.~2, 267--300. \url{https://doi.org/10.4310/jdg/1367438650}.
\bibitem{MP} Maryam Mirzakhani and Bram Petri, \emph{Lengths of closed geodesics on random surfaces of large genus}, Comment. Math. Helv. \textbf{94} (2019), no.~4, 869--889. \url{https://doi.org/10.4171/CMH/477}; arXiv:1710.09727v2.
\bibitem{MZ} Maryam Mirzakhani and Peter Zograf, \emph{Towards large genus asymptotics of intersection numbers on moduli spaces of curves}, Geom. Funct. Anal. \textbf{25} (2015), no.~4, 1258--1289. \url{https://doi.org/10.1007/s00039-015-0336-5}.
\bibitem{Monk} Laura Monk, \emph{Benjamini--Schramm convergence and spectra of random hyperbolic surfaces of high genus}, Anal. PDE \textbf{15} (2022), no.~3, 727--752. \url{https://doi.org/10.2140/apde.2022.15.727}.
\bibitem{MonkNaud} Laura Monk and Fr\'ed\'eric Naud, \emph{Spectral gaps on large hyperbolic surfaces}, in \emph{Proceedings of the International Congress of Mathematicians 2026}, Vol.~3: Invited Lectures (Sections 1--4), SIAM, Philadelphia, PA, 2026, 434--453. \url{https://doi.org/10.1137/25m1805527}.
\bibitem{MS} Motohico Mulase and Brad Safnuk, \emph{Mirzakhani's recursion relations, Virasoro constraints and the KdV hierarchy}, Indian J. Math. \textbf{50} (2008), no.~1, 189--218; \href{https://arxiv.org/abs/math/0601194v3}{arXiv:math/0601194v3}.
\bibitem{NWX} Xin Nie, Yunhui Wu, and Yuhao Xue, \emph{Large genus asymptotics for lengths of separating closed geodesics on random surfaces}, J. Topol. \textbf{16} (2023), no.~1, 106--175. \url{https://doi.org/10.1112/topo.12276}.
\bibitem{OS} Kazumi Okuyama and Kazuhiro Sakai, \emph{JT gravity, KdV equations and macroscopic loop operators}, J. High Energy Phys. \textbf{2020} (2020), no.~1, 156. \url{https://doi.org/10.1007/JHEP01(2020)156}; arXiv:1911.01659v2.
\bibitem{Penner} R. C. Penner, \emph{Weil--Petersson volumes}, J. Differential Geom. \textbf{35} (1992), no.~3, 559--608. \url{https://doi.org/10.4310/jdg/1214448257}.
\bibitem{Petri} Bram Petri, \emph{Bass notes of random hyperbolic surfaces of large genus}, S\'eminaire Bourbaki survey, preprint (2026), \href{https://arxiv.org/abs/2607.06331v1}{arXiv:2607.06331v1}.
\bibitem{RZ} Kasra Rafi and Anton Zorich, \emph{Asymptotics of Masur--Veech volume of strata of quadratic differentials with many simple poles}, preprint (2026), \href{https://arxiv.org/abs/2610.28963v1}{arXiv:2610.28963v1}.
\bibitem{Ren} Irene Ren, \emph{Mirzakhani's frequencies of simple closed geodesics on hyperbolic surfaces in large genus and with many cusps}, Geom. Dedicata \textbf{219} (2025), no.~1, article~14. \url{https://doi.org/10.1007/s10711-024-00972-6}.
\bibitem{Rivin} Igor Rivin, \emph{Geodesics with one self-intersection, and other stories}, Adv. Math. \textbf{231} (2012), no.~5, 2391--2412. \url{https://doi.org/10.1016/j.aim.2012.07.018}; arXiv:0901.2543v4.
\bibitem{SSS} Phil Saad, Stephen H. Shenker, and Douglas Stanford, \emph{JT gravity as a matrix integral}, preprint (2019), \href{https://arxiv.org/abs/1903.11115v1}{arXiv:1903.11115v1}.
\bibitem{ST} Georg Schumacher and Stefano Trapani, \emph{Estimates of Weil--Petersson volumes via effective divisors}, Comm. Math. Phys. \textbf{222} (2001), no.~1, 1--7. \url{https://doi.org/10.1007/s002200100493}.
\bibitem{SW} Yang Shen and Yunhui Wu, \emph{Arbitrarily small spectral gaps for random hyperbolic surfaces with many cusps}, Comm. Math. Phys. \textbf{406} (2025), no.~8, article~190. \url{https://doi.org/10.1007/s00220-025-05366-7}; arXiv:2203.15681v4.
\bibitem{Temme} Nico M. Temme, \emph{Laplace type integrals: transformation to standard form and uniform asymptotic expansions}, Quart. Appl. Math. \textbf{43} (1985), no.~1, 103--123. \url{https://doi.org/10.1090/qam/782260}.
\bibitem{Witten} Edward Witten, \emph{Two-dimensional gravity and intersection theory on moduli space}, Surveys in Differential Geometry \textbf{1} (1991), 243--310. \url{https://doi.org/10.4310/SDG.1990.v1.n1.a5}.
\bibitem{WittenDef} Edward Witten, \emph{Matrix models and deformations of JT gravity}, Proc. R. Soc. A \textbf{476} (2020), no.~2244, 20200582. \url{https://doi.org/10.1098/rspa.2020.0582}; arXiv:2006.13414v1.
\bibitem{Wolpert} Scott A. Wolpert, \emph{On the Weil--Petersson geometry of the moduli space of curves}, Amer. J. Math. \textbf{107} (1985), no.~4, 969--997. \url{https://doi.org/10.2307/2374363}.
\bibitem{Wright} Alex Wright, \emph{A tour through Mirzakhani's work on moduli spaces of Riemann surfaces}, Bull. Amer. Math. Soc. (N.S.) \textbf{57} (2020), no.~3, 359--408. \url{https://doi.org/10.1090/bull/1687}.
\bibitem{WXPrime} Yunhui Wu and Yuhao Xue, \emph{Prime geodesic theorem and closed geodesics for large genus}, J. Eur. Math. Soc. (2025), published online first. \url{https://doi.org/10.4171/JEMS/1653}.
\bibitem{WX} Yunhui Wu and Yuhao Xue, \emph{Random hyperbolic surfaces of large genus have first eigenvalues greater than $\frac{3}{16}-\epsilon$}, Geom. Funct. Anal. \textbf{32} (2022), no.~2, 340--410. \url{https://doi.org/10.1007/s00039-022-00595-7}.
\bibitem{XZ} Dingkang Xiong and Jian Zhou, \emph{Eynard--Orantin topological recursion, moduli-space Virasoro constraints, and semisimple CohFTs}, preprint (2026), \href{https://arxiv.org/abs/2610.24980v1}{arXiv:2610.24980v1}.
\bibitem{ZogSpectral} Peter G. Zograf, \emph{Small eigenvalues of automorphic Laplacians in spaces of cusp forms}, Zap. Nauchn. Sem. Leningrad. Otdel. Mat. Inst. Steklov. (LOMI) \textbf{134} (1984), 157--168. \url{https://www.mathnet.ru/eng/znsl4746}.
\bibitem{Zog93} Peter Zograf, \emph{The Weil--Petersson volume of the moduli space of punctured spheres}, in \emph{Mapping Class Groups and Moduli Spaces of Riemann Surfaces (G\"ottingen, 1991/Seattle, WA, 1991)}, Contemp. Math. \textbf{150}, Amer. Math. Soc., Providence, RI, 1993, 367--372. \url{https://doi.org/10.1090/conm/150/01300}.
\bibitem{Zog98} Peter Zograf, \emph{Weil--Petersson volumes of moduli spaces of curves and the genus expansion in two dimensional gravity}, preprint (1998), \href{https://arxiv.org/abs/math/9811026}{arXiv:math/9811026}.
\bibitem{Zog} Peter Zograf, \emph{On the large genus asymptotics of Weil--Petersson volumes}, preprint (2008), \href{https://arxiv.org/abs/0812.0544}{arXiv:0812.0544}.
\end{thebibliography}

\end{document}
